\xchapter{Dedication}




TO THE RIGHT HONOURABLE


LORD BLACHFORD, K.C.M.G., P.C.,


IN LOVING REMEMBRANCE


OF OLD DAYS


OF PLEASANT INTIMATE COMPANIONSHIP,


FROM HIS AFFECTIONATE FRIEND,



JOHN H. NEWMAN.




\xchapter{Advertisement}




BOTH these Essays were written when the
author was Fellow of Oriel College, Oxford.


The former of them, on the Miracles of Scripture, was written in
1825-26 for the ``Encyclopædia Metropolitana,'' being the
sequel to a Life of Apollonius Tyanæus.


The latter, on the Miracles of the first age of Christianity, was
written in 1842-43, as a Preface to a Translation of a portion of
Fleury's Ecclesiastical History.


In the first of the two, the Miracles of Scripture are regarded as
mainly addressed to religious inquirers, of an evidential nature, the
instruments of conversion, and the subjects of an inspired record. In
the second, the Ecclesiastical Miracles are regarded as addressed to
Christians, the rewards of faith, and the matter of devotion, varying
in their character from simple providences to distinct innovations
upon physical order, and coming to us by tradition or in legend,
trustworthy or not, as it may happen in the particular case.

These distinct views of miraculous agency, thus contrasted, involve
no inconsistency with each other; but it must be owned that, in the
Essay upon the Scripture Miracles, the Author goes beyond both the
needs and the claims of his argument, when, in order to show their
special dignity and beauty, he depreciates the purpose and value of
the Miracles of Church History. To meet this undue disparagement, in
his first Essay, of facts which have their definite place in the
Divine Dispensation, he points out, in his second, the essential
resemblance which exists between many of the Miracles of Scripture and
those of later times; and it is with the same drift that, in this
Edition, a few remarks at the foot of the page have been added in
brackets.


With the exception of these bracketed additions in both Essays, and
of a Memorandum at the end of the volume, the alterations made,
whether in text or notes, are simply of a literary character. As to
the latter, no verification has been made of the references which they
contain, much pains having been bestowed on them, as it is believed,
in the original Edition.
\emph{June} 29, 1870.

\xchapter{Essay I. The Miracles of Scripture Compared with those reported elsewhere, as regards Their Nature, Credibility, and Evidence}

\xsection{Introduction. On the Miracles of Scripture}

I PROPOSE
to attempt an extended comparison between the Miracles of Scripture
and those elsewhere related, as regards their nature, credibility, and
evidence. I shall divide my observations under the following heads:—

§ 1. On the Idea and Scope of a Miracle.

§ 2. On the antecedent Credibility of a Miracle,
considered as a Divine Interposition.

§ 3. On the Criterion of a Miracle, considered
as a Divine Interposition.

§ 4. On the direct Evidence for the Christian
Miracles.

\xsection{Section 1. On the Idea and Scope of a Miracle}

A MIRACLE
may be considered as an event inconsistent with the constitution of
nature, that is, with the established course of things in which it is
found. Or, again, an event in a given system which cannot be referred
to any law, or accounted for by the operation of any principle, in
that system. It does not necessarily imply a violation of nature, as
some have supposed,—merely the interposition of an external cause,
which, we shall hereafter show, can be no other than the agency of the
Deity. And the effect produced is that of unusual or increased action
in the parts of the system.

It is then a relative term, not only as it
presupposes an assemblage of laws from which it is a deviation, but
also as it has reference to some one particular system; for the same
event which is anomalous in one, may be quite regular when observed in
connexion with another. The Miracles of Scripture, for instance, are
irregularities in the economy of nature, but with a moral end;
forming one instance out of many, of the providence of God, that is,
an instance of occurrences in the natural world with a final cause.
Thus, while they are exceptions to the laws of one system, they may
coincide with those of another. They profess to be the evidence of a
Revelation, the criterion of a divine message. To consider them as
mere exceptions to physical order, is to take a very incomplete view
of them. It is to degrade them from the station which they hold in the
plans and provisions of the Divine Mind, and to strip them of their
real use and dignity; for as naked and isolated facts they do but
deform an harmonious system.

From this account of a Miracle, it is evident
that it may often be difficult exactly to draw the line between
uncommon and strictly miraculous events. Thus the production of ice
might have seemed at first sight miraculous to the Siamese; for it was
a phenomenon referable to none of those laws of nature which are in
ordinary action in tropical climates. Such, again, might magnetic
attraction appear, in ages familiar only with the attraction of
gravity \footnote{Campbell, On Miracles, Part i. Sec. 2.}. On the other
hand, the extraordinary works of Moses or St. Paul appear miraculous,
even when referred to those simple and elementary principles of nature
which the widest experience has confirmed. As far as this affects the
discrimination of supernatural facts, it will be considered in its
proper place; meanwhile let it suffice to state, that those events
only are connected with our present subject which have no assignable
second cause or antecedent, and which, on that account, are from the
nature of the case referred to the immediate agency of the Deity.

A Revelation, that is, a direct message from God
to man, itself bears in some degree a miraculous character; inasmuch
as it supposes the Deity actually to present Himself before His
creatures, and to interpose in the affairs of life in a way above the
reach of those settled arrangements of nature, to the existence of
which universal experience bears witness. And as a Revelation itself,
so again the evidences of a Revelation may all more or less be
considered miraculous. Prophecy is an evidence only so far as
foreseeing future events is above the known powers of the human mind,
or miraculous. In like manner, if the rapid extension of Christianity
be urged in favour of its divine origin, it is because such extension,
under such circumstances, is supposed to be inconsistent with the
known principles and capacity of human nature. And the pure morality
of the Gospel, as taught by illiterate fishermen of Galilee, is an
evidence, in proportion as the phenomenon disagrees with the
conclusions of general experience, which leads us to believe that a
high state of mental cultivation is ordinarily requisite for the
production of such moral teachers. It might even be said that,
strictly speaking, no evidence of a Revelation is conceivable which
does not partake of the character of a Miracle; since nothing but a
display of power over the existing system of things can attest the
immediate presence of Him by whom it was originally established; or,
again, because no event which results entirely from the ordinary
operation of nature can be the criterion of one that is extraordinary
\footnote{Hence it is that in the Scripture accounts of Revelations to the
Prophets, etc., a sensible Miracle is so often asked and given; as if
the vision itself, which was the medium of the Revelation, was not a
sufficient evidence of it, as being perhaps resolvable into the
ordinary powers of an excited imagination; \emph{e.g}., Judg. vi.
36-40, etc.}.

In the present argument I confine myself to the
consideration of Miracles commonly so called; such events, that is,
for the most part, as are inconsistent with the constitution of the
physical world.

Miracles, thus defined, hold a very prominent
place in the evidence of the Jewish and Christian Revelations. They
are the most striking and conclusive evidence; because, the laws of
matter being better understood than those to which mind is conformed,
the transgression of them is more easily recognised. They are the most
simple and obvious; because, whereas the freedom of the human will
resists the imposition of undeviating laws, the material creation, on
the contrary, being strictly subjected to the regulation of its Maker, looks to Him alone for a change in its constitution. Yet
Miracles are but a branch of the evidences, and other branches have
their respective advantages. Prophecy, as has been often observed, is
a growing evidence, and appeals more forcibly than Miracles to those
who are acquainted with the Miracles only through testimony. A
philosophical mind will perhaps be most strongly affected by the fact
of the very existence of the Jewish polity, or of the revolution
effected by Christianity. While the beautiful moral teaching and
evident honesty of the New Testament writers is the most persuasive
argument to the unlearned but single-hearted inquirer. Nor must it be
forgotten that the evidences of Revelation are cumulative, that they
gain strength from each other; and that, in consequence, the argument
from Miracles is immensely stronger when viewed in conjunction with
the rest, than when considered separately, as in an inquiry of the
present nature.

As the relative force of the separate evidences
is different under different circumstances, so again has one class of
Miracles more or less weight than another, according to the accidental
change of times, places, and persons addressed. As our knowledge of
the system of nature, and of the circumstances of the particular case
varies, so of course varies our conviction. Walking on the sea, for
instance, or giving sight to one born blind, would to us perhaps be a
Miracle even more astonishing than it was to the Jews; the laws of
nature being at the present day better understood than formerly, and
the fables concerning magical power being no longer credited. On the
other hand, stilling the wind and waves with a word may by all but
eye-witnesses be set down to accident or exaggeration without the
possibility of a full confutation; yet to eye-witnesses it would carry
with it an overpowering evidence of supernatural agency by the voice
and manner that accompanied the command, the violence of the wind at
the moment, the instantaneous effect produced, and other
circumstances, the force of which a narrative cannot fully convey. The
same remark applies to the Miracle of changing water into wine, to the
cure of demoniacal possessions, and of diseases generally. From a
variety of causes, then, it happens that Miracles which produced a
rational conviction at the time when they took place, have ever since
proved rather an objection to Revelation than an evidence for it, and
have depended on the rest for support; while others, which once were
of a dubious and perplexing character, have in succeeding ages come
forward in its defence. It is by a process similar to this that the
anomalous nature of the Mosaic polity, which might once be an obstacle
to its reception, is now justly alleged in proof of the very Miracles
by which it was then supported \footnote{See Sumner's ``Records of Creation,'' Vol. i.}.
It is important to keep this remark in view, as it is no uncommon
practice with those who are ill-affected to the cause of Revealed
Religion to dwell upon such Miracles as at the present day rather
require than contribute evidence, as if they formed a part of the
present proof on which it rests its pretensions \footnote{See Hume, On Miracles: ``Let us examine those Miracles related in
Scripture, and, \emph{not to lose ourselves in too wide a field}, let
us confine ourselves to such as we find in the \emph{Pentateuch}, etc.
It gives an account of the state of the world and of human nature
entirely different from the present; of our fall from that state; of
the age of man extended to near a thousand years,'' etc. See Berkeley's
``Minute Philosopher,'' Dial. vi. Sec. 30.}.

In the foregoing remarks, the being of an
intelligent Maker has been throughout assumed; and, indeed, if the
peculiar object of a Miracle be to evidence a message from God, it is
plain that it implies the admission of the fundamental truth, and
demands assent to another beyond it. His particular interference it
directly proves, while it only reminds of His existence. It professes
to be the signature of God to a message delivered by human
instruments; and therefore supposes that signature in some degree
already known, from His ordinary works. It appeals to that moral sense
and that experience of human affairs which already bear witness to His
ordinary presence. Considered by itself, it is at most but the token
of a superhuman being. Hence, though an additional instance, it is not
a distinct species of evidence for a Creator from that contained
in the general marks of order and design in the universe. A proof
drawn from an interruption in the course of nature is in the same line
of argument as one deduced from the existence of that course, and in
point of cogency is inferior to it. Were a being who had experience
only of a chaotic world suddenly introduced into this orderly system
of things, he would have an infinitely more powerful argument for the
existence of a designing Mind, than a mere interruption of that system
can afford. A Miracle is no argument to one who is deliberately, and
on principle, an atheist.

Yet, though not abstractedly the more convincing,
it is often so in effect, as being of a more striking and imposing
character. The mind, habituated to the regularity of nature, is
blunted to the overwhelming evidence it conveys; whereas by a Miracle
it may be roused to reflection, till mere conviction of a superhuman
being becomes the first step towards the acknowledgment of a Supreme
Power. While, moreover, it surveys nature as a whole, it is not
capacious enough to embrace its bearings, and to comprehend what it
implies. In miraculous displays of power the field of view is
narrowed; a detached portion of the divine operations is taken as an
instance, and the final cause is distinctly pointed out. A Miracle,
besides, is more striking, inasmuch as it displays the Deity in
action; evidence of which is not supplied in the system of
nature. It may then accidentally bring conviction of an intelligent
Creator; for it voluntarily proffers a testimony which we have
ourselves to extort from the ordinary course of things, and forces
upon the attention a truth which otherwise is not discovered, except
upon examination.

And as it affords a more striking evidence of a
Creator than that conveyed in the order and established laws of the
Universe, still more so does it of a Moral Governor. For, while nature
attests the being of God more distinctly than it does His moral
government, a miraculous event, on the contrary, bears more directly
on the fact of His moral government, of which it is an immediate
instance, while it only implies His existence. Hence, besides
banishing ideas of Fate and Necessity, Miracles have a tendency to
rouse conscience, to awaken to a sense of responsibility, to remind of
duty, and to direct the attention to those marks of divine government
already contained in the ordinary course of events \footnote{Farmer, On Miracles, Chap. i. Sec. 2.}.

Hitherto, however, I have spoken of solitary
Miracles; a system of miraculous interpositions, conducted with
reference to a final cause, supplies a still more beautiful and
convincing argument for the moral government of God.

\xsection{Section 2. On the Antecedent Credibility of a Miracle, Considered as a Divine Interposition}

IN
proof of miraculous occurrences, we must have recourse to the same
kind of evidence as that by which we determine the truth of historical
accounts in general. For though Miracles, in consequence of their
extraordinary nature, challenge a fuller and more accurate
investigation, still they do not admit an investigation conducted on
different principles,—Testimony being the main assignable medium of
proof for past events of any kind. And this being indisputable, it is
almost equally so that the Christian Miracles are attested by evidence
even stronger than can be produced for any of those historical facts
which we most firmly believe. This has been felt by unbelievers who
have been, in consequence, led to deny the admissibility of even the
strongest testimony, if offered in behalf of miraculous events, and
thus to get rid of the only means by which they can be proved to have
taken place. It has accordingly been asserted, that all events
inconsistent with the course of nature bear in their very front such
strong and decisive marks of falsehood and absurdity, that it is
needless to examine the evidence adduced for them \footnote{\emph{I.e}., it is pretended to try \emph{past}
events on the principles used in conjecturing \emph{future}; viz., on
antecedent probability and examples. (Whately's Treatise on
Rhetoric.) See Leland's ``Supplement to View of Deistical
Writers,'' Let. 3.}. ``Where men are heated by zeal and enthusiasm,'' says Hume, with
a distant but evident allusion to the Christian Miracles, ``there is
no degree of human testimony so strong as may not be procured for the
greatest absurdity; and those who will be so silly as to examine the
affair by that medium, and seek particular flaws in the testimony, are
almost sure to be confounded.'' \footnote{Essays, Vol. ii. Note 1.} Of
these antecedent objections, which are supposed to decide the
question, the most popular is founded on the frequent occurrence of
wonderful tales in every age and country—generally, too, connected
with Religion; and since the more we are in a situation to examine
these accounts, the more fabulous they are proved to be, there would
certainly be hence a fair presumption against the Scripture narrative,
did it resemble them in its circumstances and proposed object. A more
refined argument is that advanced by Hume, in the first part of his \emph{Essay
On Miracles}, in which it is maintained against the credibility of
a Miracle, that it is more probable that the testimony should be
false than that the Miracle should be true.

This latter objection has been so ably met by
various writers, that, though prior in the order of the argument to
the former, it need not be considered here. It derives its force from
the assumption, that a Miracle is strictly a causeless phenomenon, a
self-originating violation of nature; and is solved by referring the
event to divine agency, a principle which (it cannot be denied) has
originated works indicative of power at least as great as any Miracle
requires. An adequate cause being thus found for the production of a
Miracle, the objection vanishes, as far as the mere question of power
is concerned; and it remains to be considered whether the anomalous
fact be of such a character as to admit of being referred to the
Supreme Being. For if it cannot with propriety be referred to Him, it
remains as improbable as if no such agent were known to exist. At this
point, then, I propose taking up the argument; and by examining what
Miracles are in their nature and circumstances referable to Divine
agency, I shall be providing a reply to the former of the objections
just noticed, in which the alleged similarity of all miraculous
narratives one to another, is made a reason for a common rejection of
all.

In examining what Miracles may properly be
ascribed to the Deity, Hume supplies us with an observation so just, when taken in its full extent, that I shall make it the
groundwork of the inquiry on which I am entering. As the Deity, he
says, discovers Himself to us by His works, we have no rational
grounds for ascribing to Him attributes or actions dissimilar from
those which His works convey. It follows, then, that in discriminating
between those Miracles which can and those which cannot be ascribed to
God, we must be guided by the information with which experience
furnishes us concerning His wisdom, goodness, and other attributes.
Since a Miracle is an act out of the known track of Divine agency, as
regards the physical system, it is almost indispensable to show its
consistency with the Divine agency, at least, in some other point of
view if, that is, it is recognised as the work of the same power. Now,
I contend that this reasonable demand is satisfied in the Jewish and
Christian Scriptures, in which we find a narrative of Miracles
altogether answering in their character and circumstances to those
general ideas which the ordinary course of Divine Providence enables
us to form concerning the attributes and actions of God.

While writers expatiate so largely on the laws of
nature, they altogether forget the existence of a moral system: a
system which, though but partially understood, and but general in its
appointments as acting upon free agents, is as intelligible in its
laws and provisions as the material world. Connected with this moral government, we find certain instincts of mind; such as
conscience, a sense of responsibility, and an approbation of virtue;
an innate desire of knowledge, and an almost universal feeling of the
necessity of religious observances; while, in fact, Virtue is, on the
whole, rewarded, and Vice punished. And though we meet with many and
striking anomalies, yet it is evident they are but anomalies, and
possibly but in appearance so, and with reference to our partial
information.

These two systems, the Physical and the Moral,
sometimes act in union, and sometimes in opposition to each other; and
as the order of nature certainly does in many cases interfere with the
operation of moral laws (as, for instance, when good men die
prematurely, or the gifts of nature are lavished on the bad), there is
nothing to shock probability in the idea that a great moral object
should be effected by an interruption of physical order. \footnote{See Butler's
``Analogy,'' Part i. Chap.
iii. \textbf{[}This footnote appears on
page 17 without a reference on the page—\textbf{NR]}} But, further than this, however physical laws may embarrass the
operation of the moral system, still on the whole they are subservient
to it; contributing, as is evident, to the welfare and convenience of
man, providing for his mental gratification as well as animal
enjoyment, sometimes even supplying correctives to his moral
disorders. If, then, the economy of nature has so constant a reference
to an ulterior plan, a Miracle is a deviation from the
subordinate for the sake of the superior system, and is very far
indeed from improbable, when a great moral end cannot be effected
except at the expense of physical regularity. Nor can it be fairly
said to argue an imperfection in the Divine plans, that this
interference should be necessary. For we must view the system of
Providence as a whole; which is not more imperfect because of the
mutual action of its parts, than a machine, the separate wheels of
which effect each other's movements.

Now the Miracles of the Jewish and Christian
Religions must be considered as immediate effects of Divine Power
beyond the action of nature, for an important moral end; and are in
consequence accounted for by producing, not a physical, but a final
cause \footnote{Divine Legation, Book ix. Chap. v. Vince, On
Miracles, Sermon i.}. We are not left to contemplate
the bare anomalies, and from the mere necessity of the case to refer
them to the supposed agency of the Deity. The power of displaying them
is, according to the Scripture narrative, intrusted to certain
individuals, who stand forward as their interpreters, giving them a
voice and language, and a dignity demanding our regard; who set them
forth as evidences of the greatest of moral ends, a Revelation from
God,—as instruments in His hand of affecting a direct intercourse
between Himself and His creatures, which otherwise could not have
been effected,—as vouchers for the truth of a message which they
deliver \footnote{As, for instance, Exod. iv. 1-9, 29-31; vii.
9, 17; Numb. xvi. 3, 28, 29; Deut. iv. 36-40; xviii. 21, 22; Josh.
iii. 7-13; 1 Sam. x. 1-7; xii. 16-19; 1 Kings xiii. 3; xvii. 24; xviii.
36-39; 2 Kings i. 6, 10; v. 15; xx. 8-11; Jer. xxviii. 15-17; Ezek.
xxxiii. 33; Matt. x. 1-20; xi. 3-5, 20-24; Mark xvi. 15-20; Luke i.
18-20; ii. 11, 12; v. 24; vii. 15, 16; ix. 2; x. 9; John ii. 22; iii.
2; v. 36, 37; ix. 33; x. 24-38; xi. 15, 41, 42; xiii. 19; xiv. 10, 11,
29; xvi. 4; xx. 30-31; Acts i. 8; ii. 22, 33; iii. 15, 16; iv. 33; v.
32; viii. 6; x. 38; xiii. 8-12; xiv. 3; Rom. xv. 18, 19; 1 Cor. ii.
4, 5; 2 Cor. xii. 12; Heb. ii. 3, 4; Rev. xix. 10.}. This is plain and
intelligible; there is an easy connection between the miraculous
nature of their works and the truth of their words; the fact of their
superhuman power is a reasonable ground for belief in their superhuman
knowledge. Considering, then, our instinctive sense of duty and moral
obligation, yet the weak sanction which reason gives to the practice
of virtue, and withal the uncertainty of the mind when advancing
beyond the first elements of right and wrong; considering, moreover,
the feeling which wise men have entertained of the need of some
heavenly guide to instruct and confirm them in goodness, and that
unextinguishable desire for a Divine message which has led men in all
ages to acquiesce even in pretended revelations, rather than forego
the consolation thus afforded them; and again, the possibility (to say
the least) of our being destined for a future state of being, the
nature and circumstances of which it may concern us much to know,
though from nature we know nothing; considering, lastly, our
experience of a watchful and merciful Providence, and the
impracticability already noticed of a Revelation without a Miracle, it
is hardly too much to affirm that the moral system points to an
interference with the course of nature, and that Miracles wrought in
evidence of a Divine communication, instead of being antecedently
improbable, are, when directly attested, entitled to a respectful and
impartial consideration.

When the various antecedent objections which
ingenious men have urged against Miracles are brought together, they
will be found nearly all to arise from forgetfulness of the existence
of moral laws \footnote{Vince, On Miracles, Sermon
i.}. In their zeal to perfect
the laws of matter they most unphilosophically overlook a more sublime
system, which contains disclosures not only of the Being but of the
Will of God. Thus, Hume, in a passage above referred to, observes, ``Though the Being to whom the Miracle is ascribed be Almighty, it
does not, upon that account, become a whit more probable, since it is
impossible for us to know the attributes or actions of such a Being,
otherwise than from the experience which we have of His productions in
the usual course of nature. This still reduces us to past observation,
and obliges us to compare the instances of the violation of truth in
the testimony of men with those of the violation of the laws of
nature by Miracles, in order to judge which of them is most likely and
probable.'' Here the moral government of God, with the course of
which the Miracle entirely accords, is altogether kept out of sight.
With a like heedlessness of the moral character of a Miracle, another
writer, notorious for his irreligion \footnote{Voltaire.},
objects that it argues mutability in the Deity, and implies that the
physical system was not created good, as needing improvement. And a
recent author adopts a similarly partial and inconclusive mode of
reasoning, when he confuses the Christian Miracles with fables of
apparitions and witches, and would examine them on the strict
principle of those legal forms which from their secular object go far
to exclude all religious discussion of the question \footnote{Bentham, Preuves Judiciares, Liv. viii.}. Such reasoners seem to suppose, that when the agency of the Deity
is introduced to account for Miracles, it is the illogical
introduction of an unknown cause, a reference to a mere name, the
offspring, perhaps, of popular superstition; or, if more than a name,
to a cause that can be known only by means of the physical creation;
and hence they consider Religion as founded in the mere weakness or
eccentricity of the intellect, not in actual intimations of a Divine
government as contained in the moral world. From an apparent
impatience of investigating a system which is but partially
revealed, they esteem the laws of the material system alone worthy the
notice of a scientific mind; and rid themselves of the annoyance which
the importunity of a claim to miraculous power occasions them, by
discarding all the circumstances which fix its antecedent probability,
all in which one Miracle differs from another, the professed author,
object, design, character, and human instruments.

When this partial procedure is resisted, the \emph{à
priori} objections of sceptical writers at once lose their force.
Facts are only so far improbable as they fall under no general rule;
whereas it is as parts of an existing system that the Miracles of
Scripture demand our attention, as resulting from known attributes of
God, and corresponding to the ordinary arrangements of His providence.
Even as detached events they might excite a rational awe towards the
mysterious Author of nature. But they are presented to us, not as
unconnected and unmeaning occurrences, but as holding a place in an
extensive plan of Divine government, completing the moral system,
connecting Man and his Maker, and introducing him to the means of
securing his happiness in another and eternal state of being. That
such is the professed object of the body of Christian Miracles, can
hardly be denied. In the earlier Religion it was substantially the
same, though, from the preparatory nature of the Dispensation a less enlarged view was given of the Divine counsels. The express
purpose of the Jewish Miracles is to confirm the natural evidence of
one God, the Creator of all things, to display His attributes and will
with distinctness and authority, and to enforce the obligation of
religious observances, and the sinfulness of idolatrous worship \footnote{Exod. iii.; xiv.; xx. 22, 23; xxxiv. 6-17;
Deut. iv. 32-40; Josh. ii. 10, 11; iv. 23, 24; 1 Sam. v. 3, 4; xii.
18; 2 Sam. vii. 23; 1 Kings viii. 59, 60; xviii. 36, 37; xx. 28; 2
Kings xix. 15-19, 35; 2 Chron. xx. 29; Isaiah vi. 1-5; xix. 1; xliii.
10-12.}. Whether we turn to the earlier or later ages of Judaism, in the
plagues of Egypt, in the parting of Jordan, and the arresting of the
sun's course by Joshua, in the harvest thunder at the prayer of
Samuel, in the rending of the altar at Bethel, in Elijah's sacrifice
on Mount Carmel, and in the cure of Naaman by Elisha, we recognise
this one grand object throughout. Not even in the earliest ages of the
Scripture history are Miracles wrought at random, or causelessly, or
to amuse the fancy, or for the sake of mere display; nor prodigally,
for the mere conviction of individuals, but for the most part on a
grand scale, in the face of the world, to supply whole nations with
evidence concerning the Deity. Nor are they strewn confusedly over the
face of the history, being with few exceptions reducible to three
eras; the formation of the Hebrew Church and polity, the reformation
in the times of the idolatrous Kings of Israel, and the promulgation of the Gospel. Let it be observed, moreover, that the
power of working them, instead of being assumed by any classes of men
indiscriminately, is described as a prerogative of the occasional
Prophets, to the exclusion of the Priests and Kings; a circumstance
which, not to mention its remarkable contrast to the natural course of
an imposture, is deserving attention from its consistency with the
leading design of Miracles already specified. For the respective
claims of the Kings and Priests were already ascertained, when once
the sacred office was limited to the family of Aaron, and the regal
power to David and his descendants; whereas extraordinary messengers,
as Moses, Samuel, and Elijah, needed some supernatural display of
power to authenticate their pretensions. In corroboration of this
remark I might observe upon the unembarrassed manner of the Prophets
in the exercise of their professed gift; their disdain of argument or
persuasion, and the confidence with which they appeal to those before
whom they are said to have worked their Miracles.

These and similar observations do more than
invest the separate Miracles with a dignity worthy of the Supreme
Being; they show the coincidence of them all in one common and
consistent object. As parts of a system, the Miracles recommend and
attest each other, evidencing not only general wisdom, but a digested
and extended plan. And while this appearance of design connects
them with the acknowledged works of a Creator, who is in the natural
world chiefly known to us by the presence of final causes, so, again,
a plan conducted as this was, through a series of ages, evinces not
the varying will of successive individuals, but the steady and
sustaining purpose of one Sovereign Mind. And this remark especially
applies to the coincidence of views observable between the Old and New
Testament; the latter of which, though written after a long interval
of silence, the breaking up of the former system, a revolution in
religious discipline, and the introduction of Oriental tenets into the
popular Theology, still unhesitatingly takes up and maintains the
ancient principles of miraculous interposition.

An additional recommendation of the Scripture
Miracles is their appositeness to the times and places in which they
were wrought; as, for instance, in the case of the plagues of Egypt,
which, it has been shown \footnote{See Bryant.}, were directed against the
prevalent superstitions of that country. Their originality, beauty,
and immediate utility, are further properties falling in with our
conceptions of Divine agency. In their general character we discover
nothing indecorous, light, or ridiculous; they are grave, simple,
unambiguous, majestic. Many of them, especially those of the later
Dispensation, are remarkable for their benevolent and merciful
character; others are useful for a variety of subordinate purposes, as
a pledge of the certainty of particular promises, or as comforting
good men, or as edifying the Church. Nor must we overlook the moral
instruction conveyed in many, particularly in those ascribed to
Christ, the spiritual interpretation which they will often bear, and
the exemplification which they afford of particular doctrines \footnote{Jones, On the Figurative Language of
Scripture, Lecture x. Farmer, On Miracles, Chap. iii. Sec. 6, 2.}.

Accepting, then, what may be called Hume's
canon, that \emph{no work can be reasonably ascribed to the agency of
God, which is altogether different from those ordinary works from
which our knowledge of Him is originally obtained}, I have shown
that the Miracles of Scripture, far from being exceptionable on that
account, are strongly recommended by their coincidence with what we
know from nature of His Providence and Moral Attributes. That there
are some few among them in which this coincidence cannot be traced, it
is not necessary to deny. As a whole they bear a determinate and
consistent character, being great and extraordinary means for
attaining a great, momentous, and extraordinary object.

I shall not, however, dismiss this criterion of
the antecedent probability of a Miracle with which Hume has furnished
us, without showing that it is more or less detrimental to the
pretensions of all professed Miracles but those of the Jewish and
Christian Revelations; in other words, that none else are likely to
have occurred, because none else can with any probability be referred
to the agency of the Deity, the only known cause of miraculous
interposition. We exclude then

\xsubsection{1. Those which are not even referred by the workers of them to Divine Agency}

Such are the extraordinary works attributed by
some to Zoroaster; and, again, to Pythagoras, Empedocles, Apollonius,
and others of their School; which only claim to be the result of their
superior wisdom, and were quite independent of a Supreme Being \footnote{See, in contrast, Gen. xl. 8; xli. 16; Dan.
ii. 27-30, 47; Acts iii. 12-16; xiv. 11-18; a contrast sustained, as
these passages show, for 1500 years.}. Such are the supposed effects of witchcraft or of magical charms,
which profess to originate with Spirits and Demons; for, as these
agents, supposing them to exist, did not make the world, there is
every reason for thinking they cannot of themselves alter its
arrangements \footnote{Sometimes charms are represented as having an
inherent virtue, independent of invisible agents, as in the account
given by Josephus of Eleazar's drawing out a devil through the
nostrils of a patient by means of a ring, which contained in it a drug
prescribed by Solomon. Josephus, Antiq. viii. 2, Sec. 5. See Acts
viii. 19.}. And those, as in some accounts of apparitions, which are silent respecting their origin, and are
referred to God from the mere necessity of the case.

\xsubsection{2. Those which are unworthy of an All-wise Author}

1. As, for example, the Miracles of Simon Magus,
who pretended he could assume the appearance of a serpent, exhibit
himself with two faces, and transform himself into whatever shape he
pleased \footnote{Lavington, Enthusiasm of Meth. and Papists
comp. Part iii. Sec. 43.}. Such are most of the Miracles recorded in the
apocryphal accounts of Christ \footnote{Jones, On the Canon, Part iii.}; \emph{e.g}., the sudden
ceasing of all kinds of motion at His birth, birds stopping in the
midst of their flight, men at table with their hands to their mouths,
yet unable to eat, etc.; His changing, when a child, His playmates
into kids, and animating clay figures of beasts and birds; the
practice attributed to Him of appearing to His disciples sometimes as
a youth, sometimes as an old man, sometimes as a child, sometimes
large, sometimes less, sometimes so tall as to reach the Heavens; and
the obeisance paid Him by the military standards when He was brought
before Pilate. Of the same cast is the story of His picture presented
by Nicodemus to Gamaliel, which, when pierced by the Jews, gave forth
blood and water.

2. Under this head of exception fall many of the
Miracles related by the Fathers \footnote{Middleton, Free Enquiry.}; \emph{e.g}., that of the
consecrated bread changing into a live coal in the hands of a woman,
who came to the Lord's supper after offering incense to an idol; of
the dove issuing from the body of Polycarp at his martyrdom; of the
petrifaction of a fowl dressed by a person under a vow of abstinence;
of the exorcism of the demoniac camel; of the stones shedding tears at
the barbarity of the persecutions; of inundations rising up to the
roofs of churches without entering the open doors; and of pieces of
gold, as fresh as from the mint, dropt from heaven into the laps of
the Italian Monks \footnote{[Vide, however, Essay ii., \emph{infra}, n.
48-50, 54, 58, etc.]}.

3. Of the same character are the Miracles of the
Romish Breviary \footnote{[Vide ibid.]}; as the prostration of wild beasts before
the martyrs they were about to devour; the miraculous uniting of two
chains with which St. Peter had been at different times bound; and the
burial of Paul the Hermit by lions.

4. Such again are the Rabbinical Miracles, as
that of the flies killed by lightning for settling on a rabbi's
paper. And the Miracles ascribed by some to Mahomet, as that the trees
went out to meet him, the stones saluted him, and a camel complained
to him \footnote{The offensiveness of these, and many others
above instanced, consists in attributing \emph{moral} feelings to
inanimate or irrational beings.}. The exorcism in the Book of Tobit must here be
mentioned, in which the Evil Spirit who is in love with Sara is driven
away by the smell of certain perfumes \footnote{[So the Protestant version.] It seems to have
been a common notion that possessed persons were beloved by the Spirit
possessing them. See Philostr. iv. 25. Gospel of the Infancy,
xiv.-xvi., xxxiii. Justin Martyr, Apol. p. 113, Ed. Thirlb. We find
nothing of this kind in the account of Scripture demoniacs.}.

5. Hence the Scripture accounts of Eve's
temptation by the serpent; of the speaking of Balaam's ass; of Jonah
and the whale; and of the devils sent into the herd of swine, are by
themselves more or less improbable, being unequal in dignity to the
rest. They are then supported by the system in which they are found,
as being a few out of a multitude, and therefore but exceptions (and,
as we suppose, but apparent exceptions) to the general rule. In some
of them, too, a further purpose is discernible, which of itself
reconciles us to the strangeness of their first appearance, and
suggests the possibility of similar reasons, though unknown, being
assigned in explanation of the rest. As the Miracle of the swine, the
object of which may have been to prove to us the reality of demoniacal
possessions \footnote{Divine Legation, Book ix. Chap. v.}.

6. Miracles of mere power, even when connected
with some ultimate object, are often improbable for the same
general reason, viz., as unworthy of an All-wise Author. Such as that
ascribed to Zoroaster \footnote{Brucker, Vol. i. p. 147.}, of suffering melted brass to be
poured upon his breast without injury to himself. Unless indeed their
immediate design be to exemplify the greatness of God, as in the
descent of fire from heaven upon Elijah's sacrifice, and in
Christ's walking on the sea \footnote{Power over the elements conveyed the most
striking proof of Christ's mission from the God of nature, who in
the Old Testament is frequently characterized as ruling the sea,
winds, etc. Psalm lxv. 7; lxxvii. 19; Job xxxviii. 11, etc. It is
said, that a drawing of feet upon the water was the hieroglyphic for
impossibility. Christ moreover designed, it appears, to make trial of
His disciples' faith by this miracle. See Matt. xiv. 28-31; Mark vi.
52. We read of the power to ``move mountains,'' but evidently as a
proverbial expression. The transfiguration, if it need be mentioned,
has a \emph{doctrinal} sense, and seems besides to have been intended
to lead the minds of the Apostles to the consideration of the
Spiritual Kingdom. One of Satan's temptations was to induce our Lord
to work a Miracle of mere power. Matt. iv. 6, 7. See Acts x. 38, for
the general character of the Miracles.}, which evidently possess a
dignity fitting them to be works of the Supreme Being. The propriety
indeed of the Christian Miracles, contrasted with the want of decorum
observable in those elsewhere related, forms a most striking evidence
of their divinity.

7. Here, too, ambiguous Miracles find a place, it
being antecedently improbable that the Almighty should rest the credit
of His Revelation upon events which but obscurely implied His
immediate presence.

8. And, for the same reason, those are in some
measure improbable which are professed by different Religions; because
from a Divine Agent may be expected distinct and peculiar specimens of
divine agency. Hence the claims to supernatural power in the primitive
Church are in general questionable \footnote{[Vide Essay ii. \emph{infra}, n. 81, etc.]}, as resting upon the
exorcism of evil spirits, and the cure of diseases; works, not only
less satisfactory than others, as evidence of a miraculous
interposition, but suspicious, from the circumstance that they were
exhibited also by Jews and Gentiles of the same age \footnote{Middleton, Stillingfleet, Orig. Sacr. ii. 9,
Sec. 1.}. In the
plagues of Egypt and Elijah's sacrifice, which seem to be of this
class, there is a direct contest between two parties; and the object
of the divine messenger is to show his own superiority in the very
point in which his adversaries try their powers. Our Saviour's use
of the clay in restoring sight has been accounted for on a similar
principle, such external means being in repute among the Heathen in
their pretended cures.

\xsubsection{3. Those which have no professed Object}

1. Hence a suspicion is thrown on all miracles
ascribed by the Apocryphal Gospels to Christ in His infancy; for,
being prior to His preaching, they seem to attest no doctrine,
and are but distantly connected with any object.

2. Those again on which an object seems to be
forced. Hence many harmonizing in one plan arrest the attention more
powerfully than a detached and solitary miracle, as converging to one
point, and pressing upon our notice the end for which they are
wrought. This remark, as far as it goes, is prejudicial to the miracle
wrought (as it is said) in Hunneric's persecution, long after the
real age of Miracles was past; when the Athanasian confessors are
reported to have retained the power of speech after the loss of their
tongues.

3. Those, too, must be viewed with suspicion
which are disjoined from human instruments, and are made the vehicle
of no message \footnote{Farmer, On Miracles, Chap. v.}; since, according to our foregoing view,
Miracles are only then divested of their \emph{à priori}
improbability when furthering some great moral end, such as
authenticating a divine communication. It is an objection then to
those ascribed to relics generally, and in particular to those
attributed to the tomb of the Abbé Paris, that they are left to tell
their own story, and are but distantly connected with any object
whatever. As it is, again, to many tales of apparitions, that they do
not admit of a meaning, and consequently demand at most only an otiose
assent, as Paley terms it. Hence there is a difficulty in the
narrative contained in the first verses of John v.; because we cannot
reduce the account of the descent of the Angel into the water to give
it a healing power under any known arrangement of the divine economy.
We receive it, then, on the general credit of the Revelation of which
it forms part \footnote{The verse containing the account of the Angel
is wanting in many MSS. of authority, and is marked as suspicious by
Griesbach. The mineral spring of Bethesda is mentioned by Eusebius as
celebrated even in his day.}.

4. For the same reason, \emph{viz}., the want of a
declared object, a prejudice is excited when the professed worker is
silent, or diffident as to his own power; since our general experience
of Providence leads us to suppose that miraculous powers will not be
committed to an individual who is not also prepared for his office by
secret inspiration. This speaks strongly against the cures ascribed by
Tacitus to Vespasian, and would be an objection to our crediting the
prediction uttered by Caiaphas, if separated from its context, or
prominently brought forward to rest an argument upon. It is in general
a characteristic of the Scripture system, that Miracles and
inspiration go together \footnote{Douglas, Criterion. Warburton, Sermon on
Resurrection.}.

5. With a view to specify the object distinctly,
some have required that the Miracle should be wrought after the
delivery of the message \footnote{Fleetwood, Farmer, and others.}. A message delivered an indefinite
time after the Miracle, while it cannot but excite attention from
the general reputation of the messenger for an extraordinary gift, is
not so expressly stamped with divine authority, as when it is ushered
in by his claiming, and followed by his displaying, supernatural
powers. For if a Miracle, once wrought, ever after sanctions the
doctrines taught by the person exhibiting it, it must be attended by
the gift of infallibility,—a sustained miracle, which is
inconsistent with that frugality in the application of power which is
observable in the general course of Providence \footnote{The idea is accordingly discountenanced,
Matt. vii. 22, 23. Heb. vi. 4-6. Gal. ii. 11-14.}. On the other
hand, when an unambiguous Miracle having been first distinctly
announced, is wrought with the professed object of sanctioning a
message from God, it conveys an irresistible evidence of its divine
origin. Accident is thus excluded, and the final cause indissolubly
connected with the supernatural event. I may remark that the Miracles
of Scripture were generally wrought on this plan \footnote{St. Mark ends his Gospel by saying, that the
Apostles ``went forth and preached everywhere, the Lord working
with them, and \emph{confirming the word by signs} FOLLOWING,''
chap. xvi. 20. See also Exodus iv. 29, 30; 1 Kings xiii. 2, 3; 2 Kings
xx. 8-11; Acts xiv. 3, etc.}. In
conformity to which we find moreover that the Apostles, etc., could
not work miracles when they pleased \footnote{\emph{E.g}., Acts xx. 22, 23; Phil. ii. 27; 2
Tim. iv. 20. In the Book of Acts we have not a few instances of the
Apostles acting under the immediate direction of the Holy Spirit. The
gift of tongues is an exception to the general remark, as we know it
was abused; but this from its nature was, when once given, possessed
as an ordinary talent, and needed no fresh divine influence for
subsequent exercise of it. It may besides be viewed as a medium of
conveying the message, as well as being the seal of its divinity, and
as such needed not in every instance to be marked out as a
supernatural gift. Miracles in Scripture are not done by wholesale, \emph{i.e}.,
indiscriminately and at once, without the particular will and act of
the individual; the contrary was the case with the cures at the tomb
of the Abbé Paris. Acts xix. 11, 12, perhaps forms an exception; but
the Miracles there mentioned are expressly said to be \emph{special},
and were intended to put particular honour on the Apostle; Cf. Luke
vi. 19; viii. 46, which seem to illustrate John iii. 34. [But vide
Essay ii., n. 83-85.]}; a circumstance more
consistent with our ideas of the Divine government, and
connecting the extraordinary acts more clearly with specific objects,
than if the supernatural gifts were unlimited and irrevocable.

6. Lastly, under this head I may notice professed
miracles which, as those attributed to Apollonius, may be separated
from a narrative without detriment to it. The prodigies of Livy, for
instance, form no part in the action of the history, which is equally
intelligible without them \footnote{\emph{E.g}., he says,
``ADJICIUNT
\emph{miracula huic pugnæ},'' ii. 7.}. The miraculous events of the
Pentateuch, on the contrary, or of the Gospels and Acts, though of
course they may be rejected together with the rest of the narrative,
can be rejected in no other way; since they form its substance and
groundwork, and, like the figure of Phidias on Minerva's
shield, cannot be erased without spoiling the entire composition \footnote{Whereas other extraordinary accounts are like
the statue of the Goddess herself, which could readily be taken to
pieces, and resolved into its constituent parts, the precious metal
and the stone. For the Jewish Miracles, see Graves, On the Pentateuch,
Part i. It has been observed that the \emph{discourses} of Christ so
constantly grow out of His \emph{Miracles}, that we can hardly admit
the former without admitting the latter also. But His \emph{discourses}
form His \emph{character}, which is by no means an obvious or easy one
to imagine, had it never existed.}.

\xsubsection{4. Those which are exceptionable as regards their Object}

1. If the professed object be trifling and
unimportant; as in many related by the Fathers, \emph{e.g}.,
Tertullian's account of the vision of an Angel to prescribe to a
female the exact length and measure of her veil, or the divine
admonition which Cyprian professes to have received to mix water with
wine in the Eucharist, in order to render it efficacious \footnote{Middleton, Free Inquiry. [No question
relative to the Eucharistic rite can be unimportant.]}.
Among these would be reckoned the directions given to Moses relative
to the furnishing of the Tabernacle, and other regulations of the
ceremonial law, were not further and important objects thereby
effected; such as, separating the Israelites from the surrounding
nations, impressing upon them the doctrine of a particular Providence,
prefiguring future events, etc.

2. Miracles wrought for the gratification of mere
\emph{curiosity} are referable to this head of objection. Hence the
triumphant invitations which some of the Fathers make to their heathen
opponents to attend their exorcisms excite an unpleasant feeling in
the mind, as degrading a solemn spectacle into a mere popular
exhibition.

3. Those, again, which have a \emph{political or
party object}, as the cures ascribed to Vespasian, or as those
attributed to the tomb of the Abbé Paris, and the Eclectic prodigies,
all which, viewed in their best light, tend to the mere aggrandizement
of a particular Sect, and have little or no reference to the good of
Mankind at large. It tells in favour of the Christian Miracles, that
the Apostles, generally speaking, were not enabled to work them for
their own personal convenience, to avoid danger, escape suffering, or
save life. St. Paul's preservation from the effects of the viper's
bite on the Isle of Melita is a solitary exception to this remark, no
mention being made of his availing himself of this Miracle to convert
the natives to the Christian faith \footnote{Rev. J. Blanco White, Against Catholicism,
Let. 6. The Breviary Miracles form a striking contrast to the
Christian in this point. [Not surely on the point of their benefiting
the worker.]}.

4. For a similar reason, those bear a less
appearance of probability which are wrought for the conviction of
individuals. I have already noticed the contrary character of the
Scripture Miracles in this respect; for instance, St. Paul's
miraculous conversion did not end with itself, but was followed by
momentous and inestimable consequences \footnote{Acts
xxvi. 16.}. Again, Miracles
attended the conversion of the Æthiopian Eunuch, Cornelius, and
Sergius Paulus; but these were heads and firstfruits of different
classes of men who were in time to be brought into the Church \footnote{Ibid. viii. 26, 39; x. 3, etc.; xiii. 12.
These three classes are mentioned together in prophecy. Isa. lvi. 4-8.}.

5. Miracles with a bad or vicious object are
laden with an extreme antecedent improbability; for they cannot at all
be referred to the only known cause of supernatural power, the agency
of God. Such are most of the fables concerning the heathen deities;
not a few of the professed Miracles of the primitive Church, which are
wrought to sanction doctrines opposed not only to Scriptural truth,
but to the light of nature \footnote{\emph{E.g}., to establish Monachism, etc. [Monachism
is not unnatural, unless we are prepared to maintain that an unnatural
state of life has the sanction of our Lord and St. Paul.]}; and some related in the
Apocryphal Gospels, especially Christ's inflicting death upon a
schoolmaster who threatened to strike Him, and on a boy who happened
to run violently against Him \footnote{Jones, On the Canon, Part iii.}. Here must be noticed several
passages in Scripture, in which a miraculous gift seems at first sight
to be exercised to gratify revengeful feelings, and which are,
therefore, received on the credit of the system \footnote{Gen. ix. 24-27; Judges xvi. 28-30; 2 Kings
ii. 24; 2 Chron. xxiv. 22.}.

6. Unnecessary Miracles are improbable; as those
wrought for an object attainable without an exertion, or with less
exertion, of extraordinary power \footnote{It does not follow, because all Miracles are
equally \emph{easy} to an Almighty Author, that all are equally \emph{probable};
for, as has been often remarked, a frugality in the application of
power is observable throughout His works.}. Of this kind, we contend,
would be the writing of the Gospel on the skies, which some
unbelievers have proposed as but an adequate attestation to a
Revelation; for, supposing the recorded fact of their once occurring
be sufficient for a rational conviction, a perpetual Miracle becomes
superfluous \footnote{Dr. Graves observes, of the miraculous agency
in the age of Moses and Joshua, that ``God continued it \emph{only so
long as} was indispensably necessary to introduce and settle the
Jewish nation in the land of their inheritance, and establish this
dispensation so as to answer the purposes of the divine economy. After
this, He gradually withdrew His supernatural assistance; He left the
nation collectively and individually to act according to their own
choice,'' etc.—Lectures on the Pentateuch, Part iii. Lecture 2.}. Such, again, would be the preservation of the
text of Scripture in its verbal correctness, which many have supposed
necessary for its infallibility as a standard of Truth.

7. The same antecedent objection presses on
Miracles wrought in attestation of truths already known. We do not,
for instance, require a Miracle to convince us the Sun shines, or
that Vice is blameable. The Socinian scheme is in a great measure
chargeable with bringing the Miracles of the Gospel under this
censure: for it prunes away the Christian system till little is left
for the Miracles to attest. On this ground an objection has been taken
to the Miracle wrought in favour of the Athanasians in Hunneric's
persecution, as above mentioned; inasmuch as it merely professes to
authorize a comment on the sacred text, \emph{i.e}., to sanction a
truth which is not new, unless Scripture be obscure \footnote{See Maclaine's Note on the subject,
Mosheim, Eccl. Hist. Cent. v. Part ii. Chap. v. [Vide Essay ii., n.
220, etc.]}. Here,
too, may be noticed Miracles wrought in evidence of doctrines already
established; such as those of the Papists, who seem desirous of
answering the unbeliever's demand for a perpetual Miracle. Popish
Miracles, as has often been observed, occur in Popish countries, where
they are least wanted; whereas, if real, they would be invaluable
among Protestants \footnote{[This is answered \emph{infra}, Essay ii., n.
97, etc.]}. Hence the primitive Miracles become
suspicious, in proportion as we find Christianity established, not
only from the increasing facility of fraud, but moreover from the
apparent needlessness of the extraordinary display. And hence,
admitting the Miracles of Christ and His followers, future Miracles
with the same end are somewhat improbable. For enough have been
wrought to attest the doctrine; and attention, when once excited
by supernatural means, may be kept alive by a standing Ministry, just
as inspiration is supplied by human learning.

8. I proceed to notice inconsistency in the
objects proposed, as creating a just prejudice against the validity of
miraculous pretensions. This applies to the claims of the Romish
Church, in which Miracles are wrought by hostile sects in support of
discordant tenets \footnote{Douglas, Criterion, p. 105, Note (8vo. Edit.
1807).}. It constitutes some objection to the bulk
of the Miracles of the primitive Church, when viewed as a continuation
of the original gift, that they differ so much in manner, design, and
attendant circumstances, from those recorded in Scripture \footnote{[All this is answered \emph{infra}, Essay
ii., n. 96, 101.]}. ``We see,'' says Middleton (in the ages subsequent to the Christian
era) ``a dispensation of things ascribed to God, quite different from
that which we meet with in the New Testament. For in those days the
power of working Miracles was committed to none but the Apostles, and
to a few of the most eminent of the other disciples, who were
particularly commissioned to propagate the Gospel and preside in the
Church of Christ. But, upon the pretended revival of the same powers
in the following Ages, we find the administration of them committed,
not to those who were intrusted with the government of the Church, not
to the successors of the Apostles, to the Bishops, the Martyrs,
nor to the principal champions of the Christian cause; but to \emph{boys},
to \emph{women}, and, above all, \emph{to private and obscure laymen},
not only of an inferior but sometimes also of a \emph{bad character}
\footnote{Scripture sometimes attributes miraculous
gifts to men of bad character; but we have no reason for supposing
such could work miracles at pleasure (see Numb. xxii. 18; xxiii. 3, 8,
12, 20; xxiv. 10-13), or attest any doctrine but that which Christ and
His Apostles taught; nor is our faith grounded upon their preaching.
Moreover, their power may have been given them for some further
purpose; for though to attest a divine message be the primary object
of Miracles, it need not be the only object. ``It would be highly
ridiculous,'' says Mr. Penrose in his recent work on Miracles, ``to
erect a steam engine for the mere purpose of opening and shutting a
valve; but the engine being erected is very wisely employed both for
this and for many other purposes, which, comparatively speaking, are
of very little significance.'' [This applies to ecclesiastical
miracles.]}.

9. Hence, to avoid the charge of inconsistency in
the respective objects of the Jewish and Christian Miracles, it is
incumbent upon believers in them to show that the difference between
the two systems is a difference in appearance only, and that Christ
came not to destroy but to fulfil the Law. Here, as far as its
antecedent appearance is concerned, the Miracle said to have occurred
on Julian's attempt to rebuild the Jewish Temple is seen to great
advantage. The object was great, the time critical, its consequences
harmonize very happily with the economy of the Mosaic Dispensation,
and the general spirit of the Prophetical writings, and the fact
itself has some correspondence with the prodigies which preceded the
final destruction of Jerusalem \footnote{See Warburton's Julian.}.

10. Again, Miracles which do not tend to the
accomplishment of their proposed end are open to objection; and those
which have not effected what they had in view. Hence some kind of
argument might be derived against the Christian Miracles, were they
not accompanied by a prediction of their temporary failure in
effecting their object; or, to speak more correctly, were it not their
proposed object gradually to spread the doctrines which they
authenticate \footnote{See Parables in Matt. xiii. 3, 24, 31, 33,
47; xxiv. 12; Acts xx. 29, 30; 2 Thess. ii. 3; 2 Tim. iii. 1-5, etc.}. There is nothing, however, to break the force
of this objection when directed against the Miracles ascribed to the
Abbé Paris; since the Jansenist interest, instead of being advanced
in consequence of them, soon after lost ground, and was ultimately
ruined \footnote{Paley, Evidences, Part i. Prop. 2.}.

11. These Miracles are also suspicious, as having
been stopped by human authority; it being improbable that a Divine
Agent should permit any such interference with His plan. The same
objection applies to the professed gift of exorcising demoniacs in the
primitive Church; which was gradually lost after the decree of the
Council of Laodicea confined the exercise of it to such as were
licensed by the Bishop \footnote{It had hitherto been in the hands of the
meaner sort of the Christian laity. After that time, ``few or none of
the clergy, nor indeed of the laity, were any longer able to cast out
devils; so that the old Christian exorcism or prayer for the
energumens in the church began soon after to be omitted as useless.''
Whiston, in Middleton. [Vid. Essay ii., n. 59.]}. And lastly, to the supernatural
character of Prince Hohenlohe's cures, which were stopped at Bamberg
by an order from authority, that ``none should be wrought except in
the presence of Magistrates and medical practitioners.'' \footnote{Bentham, Preuves Judiciaires, Liv. viii.
Chap. x. [This fact requires testimony stronger than Bentham's.
However, as to the Abbé Paris, the epigram is well known,

\begin{quote}

``De par le roi, defense à Dieu,

De faire miracles en ce lieu.'']

\end{quote}}

These are the most obvious objections which may
be fairly made to the antecedent probability of miraculous narratives.
It will be observed, however, that none of them go so far as to
deprive testimony for them of the privilege of being heard. Even where
the nature of the facts related forbids us to refer the Miracle to
divine agency, as when it is wrought to establish some immoral
principle, still it is not more than extremely improbable and to be
viewed with strong suspicion. Christians at least must acknowledge
that the \emph{à priori} view which Reason takes would in some cases
lead to an erroneous conclusion. A Miracle, for instance, ascribed to
an Evil Spirit is, prior to the information of Scripture,
improbable; and if it stood on its own merits would require very
strong testimony to establish it, as being referred to an unknown
cause. Yet, on the authority of Scripture, we admit the occasional
interference of agents short of divine with the course of nature.
This, however, only shows that these \emph{à priori} tests are not
decisive. Yet if we cannot always ascertain what Miracles are
improbable, at least we can determine what are not so; moreover, it
will still be true that the more objections lie against any professed
Miracle, the greater suspicion justly attaches to it, and the less
important is the fact, even if capable of proof.

On the other hand, even when the external
appearance is altogether in favour of the Miracle, it must be
recollected, nothing is thereby proved concerning the fact of its
occurrence. We have done no more than recommend to notice the
evidence, whatever it may be, which is offered in its behalf. Even,
then, could Miracles be found with as strong an antecedent case as
those of Scripture, still direct testimony must be produced to
substantiate their claims on our belief. At the same time, since there
are none such, a fair prepossession is indirectly created in favour of
the latter, over and above their intrinsic claims on our attention.

Some few indeed of the Scripture Miracles are
open to exception; and have accordingly been noticed in the course of the above remarks as by themselves improbable. These,
however, are seldom such in more than one respect; whereas the other
Miracles which came before us were open to several or all of the
specified objections at the same time. And, further, as they are but a
few in the midst of an overpowering majority pointing consistently to
one grand object, they must not be torn from their moral context, but,
on the credit of the rest, they must be considered but apparent
exceptions to the rule. It is obvious that a large system must consist
of various parts of unequal utility and excellence; and to expect each
particular occurrence to be complete in itself, is as unreasonable as
to require the parts of some complicated machine, separately taken, to
be all equally finished and fit for display \footnote{In thus refusing to admit the existence of \emph{real}
exceptions to the general rule, in spite of \emph{appearances}, we are
not exposing ourselves to that charge of excessive systematizing which
may justly be brought against those who, with Hume, reject the very
notion of a Miracle, as implying an interruption of physical
regularity. For the Revelation which we admit, on the authority of the
general system of Miracles, imparts such accurate and extended
information concerning the attributes of God, over and about the
partial and imperfect view of them which the world affords, as
precludes the supposition of \emph{any} work of His being evil or
useless. Whereas there is no \emph{voice} in the mere analogy of
nature which expressly denies the possibility of real exceptions to
its general course.}.

Let these remarks suffice on the question of the
antecedent probability or improbability of a miraculous narrative. Enough, it may be hoped, has been said to separate the
Miracles of Scripture from those elsewhere related, and to invest them
with an importance exciting in an unprejudiced mind a just interest in
their behalf, and a candid attention to the historical testimony on
which they rest; inasmuch as they are ascribed to an adequate cause,
recommended by an intrinsic dignity, and connected with an important
object, while all others are more or less unaccountable, unmeaning,
extravagant, and useless. And thus, \emph{viz}., on the ground of this utter
dissimilarity between the Miracles of Scripture and those reported
elsewhere, we are enabled to account for the incredulity with which
believers in Revelation listen to any extraordinary account at the
present day; and which sometimes is urged against them as inconsistent
with their assent to the former. It is because they admit the
Scripture Miracles. Belief in these has pre-occupied their minds, and
created a fair presumption against those of a different class;—the
prospect of a recurrence of supernatural agency being in some measure
discountenanced by the Revelation already given; and again, the
weakness and insipidity, the want of system and connexion, the
deficiency in the evidence, and the transient repute of marvellous
stories ever since, creating a strong and just prejudice against those
similar accounts which now from time to time are noised abroad.

\xsection{Section 3. On the Criterion of a Miracle, considered as a Divine Interposition}

IT has
sometimes been asked, whether Miracles are a sufficient evidence of
the interposition of the Deity? under the idea that other causes,
besides divine agency, might be assigned for their production. This is
obviously the reverse objection to that I have as yet considered,
which was founded on the assumption that they could be referred to no
known cause whatever. After showing, then, that the Scripture Miracles
may be ascribed to the Supreme Being, I proceed to show that they
cannot reasonably be ascribed to those other causes which have been
sometimes assigned for them, for instance, to unknown laws of nature,
or to the secret agency of Spirits.

1. Now it is evidently unphilosophical to
attribute them to the power of invisible Beings, short of God;
because, independently of Scripture (the truth of which, of course,
must not be assumed in this question), we have no evidence of the
existence of such beings. Nature attests, indeed, the being of a
God, but not of a race of intelligent creatures between Him and Man.
In assigning a Miracle, therefore, to the influence of Spirits, an
hypothetical cause is introduced merely to remove a difficulty. And
even did analogy lead us to admit their possible existence, yet it
would tend rather to disprove than to prove their power over the
visible creation. They may be confined to their own province, and
though superior to Man, still may be unable to do many things which he
can effect; just as Man in turn is superior to birds and fishes,
without having, in consequence, the power of flying or of inhabiting
the water \footnote{Campbell, On Miracles, Part ii. Sec. 3.
Farmer, Ch. ii. Sec. 1.}.

Still it may be necessary to show that on our own
principles we are not open to any charge of inconsistency. That is, it
has been questioned, whether, in admitting the existence and power of
Spirits on the authority of Revelation, we are not in danger of
invalidating the evidence upon which that authority rests. For the
cogency of the argument from Miracles depends on the assumption, that
interruptions in the course of nature must ultimately proceed from
God; which is not true, if they may be effected by other beings
without His sanction. And it must be conceded that, explicit as
Scripture is in considering Miracles as signs of divine agency, it
still does seem to give created Spirits some power of working
them; and even, in its most literal sense, intimates the possibility of
their working them in opposition to the true doctrine \footnote{Deut. xiii. 1-3; Matt. xxiv. 24; 2 Thess. ii. 9-11.}. With a view of meeting this difficulty, some writers have
attempted to make a distinction between great and small, many and few
Miracles; and have thus inadvertently destroyed the intelligibility of
any, as the \emph{criterion} of a divine interposition \footnote{More or less, Sherlock, Clarke, Locke, and others.}. Others, by referring to the nature of the doctrine attested,
in order to determine the author of the Miracle, have exposed
themselves to the plausible charge of adducing, first, the Miracle to
attest the divinity of the doctrine, and then, the doctrine to prove
the divinity of the Miracle \footnote{Prideaux, Clarke, Chandler, etc., seem hardly to have guarded
sufficiently against the charge here noticed. There is an appearance
of doing honour to the Christian doctrines in representing them as \emph{intrinsically}
credible, which leads many into supporting opinions which, carried to
their full extent (as they were by Middleton), supersede the need of
Miracles altogether. It must be recollected, too, that they who are
allowed to praise have the privilege of finding fault, and may reject,
according to their \emph{à priori} notions, as well as receive.
Doubtless the divinity of a clearly immoral doctrine could not be
evidenced by Miracles; for our belief in the moral attributes of God
is much stronger than our conviction of the negative proposition, that
none but He can interfere with the system of nature. But there is
always the danger of extending this admission beyond its proper
limits, of supposing ourselves adequate judges of the \emph{tendency}
of doctrines, and because unassisted Reason informs us what is moral
and immoral in our own case, of attempting to decide on the abstract
morality of actions; for many have rejected the miraculous narrative
of the Pentateuch, from an unfounded and an unwarrantable opinion,
that the means employed in settling the Jews in Canaan were in
themselves immoral. These remarks are in nowise inconsistent with
using (as was done in a former section) our actual knowledge of God's
attributes, obtained from a survey of nature and human affairs, in
determining the probability of certain professed Miracles having
proceeded from Him. It is one thing to infer from the experience of
life, another to imagine the character of God from the gratuitous
conceptions of our own minds. From experience we gain but general and
imperfect ideas of wisdom, goodness, etc., enough (that is) to bear
witness to a Revelation when given, not enough to supersede it. On the
contrary, our speculations concerning the Divine Attributes and
designs, professing, as they do, to decide on the truth of revealed
doctrines, in fact go to supersede the necessity of a Revelation
altogether.}.
Others, on the contrary, have thought themselves obliged to deny
the power of Spirits altogether, and to explain away the Scripture
accounts of demoniacal possessions, and the narrative of our Lord's
Temptation \footnote{Especially Farmer.}. Without,
however, having recourse to any of these dangerous modes of answering
the objection, it may be sufficient to reply, that since, agreeably to
the antecedent sentiment of reason, God has adopted Miracles as the
seal of a divine message, we believe He will never suffer them to be
so counterfeited as to deceive the humble inquirer. Thus the
information given by Scripture in nowise undoes the original
conclusions of Reason; for it anticipates the objection which itself
furnishes, and by revealing the express intention of God in miraculous
displays, guarantees to us that He will allow no interference of
created power to embarrass the proof thence resulting, of His special
interposition \footnote{Fleetwood, On Miracles, Disc. ii. p. 201. Van Mildert's Boyle
Lectures, Serm. xxi.}. It is
unnecessary to say more on this subject; and questions concerning the
existence, nature, and limits of spiritual agency will find their
place when Christians are engaged in settling among themselves the
doctrines of Scripture. We take it, therefore, for granted, as an
obvious and almost undeniable principle, that real Miracles, \emph{i.e}.,
interruptions in the course of nature, cannot reasonably be referred
to any power but divine, because it is natural to refer an alteration
in the system to its original author, and because Reason does not
inform us of any other being but God exterior to nature; and lastly,
because in the particular case of the Scripture Miracles, the workers
of them confirm our previous judgment by expressly attributing them to
Him.

2. A more subtle question remains, respecting the
possible existence of causes in nature, to us unknown, by the supposed
operation of which the apparent anomalies may be reconciled to the
ordinary laws of the system. It has already been admitted that some
difficulty will at times attend the discrimination of miraculous from
merely uncommon events; and it must be borne in mind that in this, as
in all questions from which demonstration is excluded, it is
impossible, from the nature of the case, absolutely to disprove any,
even the wildest, hypothesis which may be framed. It may freely be
granted, moreover, that some of the Scripture Miracles, if they stood
alone, might reasonably be referred to natural principles of which we
were ignorant, or resolved into some happy combination of accidental
circumstances. For our purpose, it is quite sufficient if there be a
considerable number which no sober judgment would attempt to deprive
of their supernatural character by any supposition of our ignorance of
natural laws, or of exaggeration in the narrative. Raising the dead
and giving sight to the blind by a word, feeding a multitude with the
casual provisions which one individual among them had with him,
healing persons at a distance, and walking on the water, are facts,
even separately taken, far beyond the conceivable effects of artifice
or accident; and much more so when they meet together in one and the
same history. And here Hume's argument from general experience is in
point, which at least proves that the ordinary powers of nature are
unequal to the production of works of this kind. It becomes, then, a
balance of opposite probabilities, whether gratuitously to suppose a
multitude of perfectly unknown causes, and these, moreover, meeting in
one and the same history, or to have recourse to one, and that a known
power, then miraculously exerted for an extraordinary and worthy
object. We may safely say no sound reasoner will hesitate on which
alternative to decide. While, then, a fair proportion of the Scripture
Miracles are indisputably deserving of their name, but a weak
objection can be derived from the case of the few which, owing to
accidental circumstances, bear at the present day less decisive marks
of supernatural agency. For, be it remembered (and it is a strong
confirmatory proof that the Jewish and Christian Miracles are really
what they profess to be) that though the miraculous character of some
of them is more doubtful in one age than in another, yet the progress
of Science has made no approximation to a general explication of them
on natural principles. While discoveries in Optics and Chemistry have
accounted for a host of apparent miracles, they hardly touch upon
those of the Jewish and Christian systems. Here is no phantasmagoria
to be detected, no analysis or synthesis of substances, ignitions,
explosions, and other customary resources of the juggler's art \footnote{See Farmer, Ch. i. Sec. 3.}. But, as before, we shall best be able to estimate their
character in this respect by contrasting them with other occurrences
which have sometimes been considered miraculous. Thus, too, a second
line of difference will be drawn between them and the mass of rival
prodigies, whether religious or otherwise, to which they are often
compared.

A Miracle, then, as far as it is an evidence of
Divine interposition, being an ascertained anomaly in an established
system, or an event without assignable physical cause, those facts, of
course, have no title to the name—

\xsubsection{1. Which may be referred to misstatements in the testimony}

1. Such are many of the prodigies of the Heathen
Mythology and History, which have been satisfactorily traced to an
exaggeration of natural events. For instance, the fables of the
Cyclops, Centaurs, of the annual transformation of a Scythian nation
into wolves, as related by Herodotus, etc. Or natural facts
allegorized, as in the fable of Scylla and Charybdis. Or where the
fact may be explained by supplying a probable omission; as we should
account for a story of a man sailing in the air by supposing a balloon
described \footnote{Bentham, Preuves Judiciaires, Liv. viii. Ch. x.}.

2. Or where the Miracle is but verbal, as the
poetical prodigy of thunder without clouds, which is little better
than a play upon words; for, supposing it to occur, it would not be
called thunder. Or as when Herodotus speaks of wool growing on trees;
for, even were it in substance the same as wool, it could not be
called so without a contradiction in terms.

3. Or where the Miracle is one simply of degree,
for then exaggeration is more easily conceivable;—thus many supposed
visions may have been but natural dreams.

4. Or where it depends on the combination of a
multitude of distinct circumstances, each of which is necessary for
the proof of its supernatural character, and where, as in fine
experiments, a small mistake is of vast consequence. As those which
depend on a coincidence of time, which it is difficult for any person
to have ascertained. For instance, the exclamation which Apollonius is
said to have uttered concerning the assassination of Domitian at the
time of its taking place; and, again, the alleged fact of his
appearing at Puteoli on the same morning in which he was tried at
Rome. Such, too, in some degree, is the professed revelation made to
St. Basil, who is said to have been miraculously informed of the death
of the Emperor Julian at the very moment that it took place \footnote{Middleton, Free Inquiry.}. Here we may instance many stories of apparitions; as the
popular one concerning the appearance of a man to the club which he
used to frequent at the moment after his death, who was afterwards
discovered to have escaped from his nurses in a fit of delirium
shortly before it took place, and actually to have joined his friends.
We may add the case related to M. Bonnet, of a woman who pretended to
know what was passing at a given time at any part of the globe,
and who was detected by the simple expedient of accurately marking the
time, and comparing her account with the fact \footnote{Bentham, Preuves Judiciaires, Liv. viii. Ch. x.}. In the same class must be reckoned not a few of the answers
of the Heathen Oracles, if it be worth while to allude to them; as
that which informed Crœsus of his occupation at a certain time agreed
upon. In the Gospel, the nobleman's son begins to amend at the very
time that Christ speaks the word; but this circumstance does not
constitute, it merely increases the Miracle. The argument from
Prophecy is, in this point of view, somewhat deficient in simplicity
and clearness, as implying the decision of many previous questions:
such, for instance, as to the existence of the professed prediction
before the event, the interval between the prediction and its
accomplishment, the completeness of its accomplishment, etc. Hence
Prophecy affords a more learned and less popular proof of Divine
interposition than physical Miracles, and, except in cases where it
contributes a very strong evidence, is commonly of inferior cogency.

\xsubsection{2. Those which, from suspicious circumstances attending them, may not unfairly be referred to an unknown cause}

1. As those which take place in departments of
nature little understood; for instance, Miracles of
Electricity.—Again, an assemblage of Miracles confined to one line
of extraordinary exertion in some measure suggests the idea of a cause
short of divine. For while their repetition looks like the profession,
their similarity argues a want, of power. This remark is
disadvantageous to the Miracles of the primitive Church, which
consisted almost entirely of exorcisms and cures \footnote{[Vide, however, \emph{infra}, Essay ii., n. 82, etc.]}; to the Pythagorean, which were principally Miracles of
sagacity; and, again, to those occurring at the tomb of the Abbé
Paris, which were limited to cures, and cures, too, of particular
diseases. While the Miracles of Scripture are frugally dispensed as
regards their object and seasons, they are carefully varied in their
nature; like the work of One who is not wasteful of His riches, yet
can be munificent when occasion calls for it.

2. Here we may notice tentative Miracles, as
Paley terms them; that is, where out of many trials only some succeed;
for inequality of success seems to imply accident, in other words, the
combination of unknown physical causes. Such are the cures of scrofula
by the King's touch, and those effected in the Heathen Temples \footnote{Stillingfleet, Orig. Sacr. Book ii. Ch. x. Sec. 9.}; and, again, those at the tomb of the Abbé Paris, there being
but eight or nine well-authenticated cures out of the multitude of
trials that were made \footnote{Douglas, Criterion, p. 133.}.
One of the peculiarities of the cures ascribed to Christ is their
invariable success \footnote{Ibid. p. 260, cites the following texts: Matt. iv. 23, 24; viii. 16;
ix. 35; xii. 15; xiv. 12; Luke iv. 40; vi. 19.}.

3. Here, for a second reason, diffidence in the
agent casts suspicion on the reality of professed Miracles; for at
least we have the sanction of his own opinion for supposing them to be
the effect of accident or unknown causes.

4. Temporary Miracles also, as many of the
Jansenist and other extraordinary cures \footnote{Douglas, Criterion, p. 190. Middleton, Free Inquiry, iv. Sec. 3.}, may be similarly accounted for; for, if ordinary causes can
undo, it is not improbable they may be able originally to effect. The
restoration of Lazarus and the others was a restoration to their
former condition, which was mortal; their subsequent dissolution,
then, in the course of nature, does not interfere with the
completeness of the previous Miracle.

5. The Jansenist cures are also unsatisfactory,
as being gradual, and, for the same reason, the professed liquefaction
of St. Januarius's blood; a progressive effect being a characteristic,
as it seems, of the operations of nature. Hence those Miracles are
most perspicuous which are wrought at the word of command; as those of
Christ and His Apostles. For this as well as other reasons, incomplete
Miracles, as imperfect cures, are no evidence of supernatural
agency; and here, again, we have to instance the cures effected at the
tomb of the Abbé Paris.

6. Again, the use of means is suspicious; for a
Miracle may almost be defined to be an event without means. Hence,
however miraculous the production of ice might appear to the Siamese,
considered abstractedly, they would hardly so account it in an actual
experiment, when they saw the preparation of nitre, etc., which in
that climate must have been used for the purpose. In the case of the
Steam-vessel or the Balloon, which, it has been sometimes said, would
appear miraculous to persons unacquainted with Science, the chemical
and mechanical apparatus employed could not fail to rouse suspicion in
intelligent minds. Hence professed Miracles are open to suspicion, if
confined to one spot; as were the Jansenist cures. For they thereby
became connected with a necessary condition, which is all we
understand by a means: for instance, such may often be imputed to a
confederacy, which (as is evident) can from its nature seldom shift
the scene of action. ``The Cock-lane ghost could only knock and scratch
in one place.'' \footnote{Hey's Lectures, Book i. Ch. xvi. Sec. 10.} The
Apostles, on the contrary, are represented as dispersed about, and
working Miracles in various parts of the world \footnote{Douglas, Criterion, p. 337.}. These remarks are, of course, inapplicable in a case
where the apparent means are known to be inadequate, and are not
constantly used; as our Lord's occasional application of clay to the
eyes, which, while it proves that He did not need such
instrumentality, conveys also an intimation that all the efficacy of
means is derived from His appointment.

\xsubsection{3. Those which may be referred to the supposed operation of a cause known to exist}

1. Professed Miracles of knowledge or mental
ability are often unsatisfactory for this reason; being in many cases
referable to the ordinary powers of the intellect. Of this kind is the
boasted elegance of the style of the Koran, alleged by Mahomet in
evidence of his divine mission. Hence most of the Miracles of
Apollonius, consisting, as they do, in knowing the thoughts of others,
and predicting the common events of life, are no criterion of a
supernatural gift; it being only under certain circumstances that such
power can clearly be discriminated from the natural exercise of
acuteness and sagacity. Accordingly, though a knowledge of the hearts
of men is claimed by Christ, it seems to be claimed rather with a view
to prove to Christians the doctrine of His Divine Nature than to
attest to the world His authority as a messenger from God. Again, St.
Paul's prediction of shipwreck on his voyage to Rome was intended to
prevent it; and so was the prediction of Agabus concerning the
same Apostle's approaching perils at Jerusalem \footnote{Acts xxi. 10-14; xxvii. 10, 21.}.

2. For a second reason, then, the argument from
Prophecy is a less simple and striking proof of divine agency than a
display of Miracles; it being impossible, in all cases, to show that
the things foretold were certainly beyond the ordinary faculties of
the mind to have discovered. Yet when this is shown, Prophecy is one
of the most powerful of conceivable evidences; strict foreknowledge
being a faculty not only above the powers, but even above the
comprehension of the human mind.

3. And much more fairly may apparent Miracles be
attributed to the supposed operation of an existing physical cause,
when they are parallel to its known effects; as chemical,
meteorological, etc., phenomena. For though the cause may not,
perhaps, appear in the particular case, yet it is known to have acted
in others similar to it. For this reason, no stress can be laid on
accounts of luminous crosses in the air, human shadows in the clouds,
appearances of men and horses on hills, and spectres when they are
speechless, as is commonly the case, ordinary causes being assignable
in all of these; or, again, on the pretended liquefaction of the blood
of St. Januarius, or on the exorcism of demoniacs, which is the most
frequent Miracle in the Primitive Church.

4. The remark applies, moreover, to cases of
healing, so far as they are not instantaneous, complete, etc.:
conditions which exclude the supposition of natural means being
employed, and which are strictly fulfilled in the Gospel narrative.

5. Again, some cures are known as possible
effects of an excited imagination; particularly when the disease
arises from obstruction and other disorders of the blood and spirits,
as the cures which took place at the tomb of the Abbé Paris \footnote{Douglas, Criterion, p. 172.}.

6. We should be required to add those cases of
healing in Scripture where the faith of the petitioners was a
necessary condition of the cure, were not these comparatively few, and
some of them such as no imagination could have effected (for instance,
the restoration of sight), and some wrought on persons absent; and
were not faith often required, not of the patient, but of the relative
or friend who brought him to be healed \footnote{Mark x. 51, 52; Matt. viii. 5-13. See Douglas, Criterion, p. 258.
``Where persons petitioned themselves for a cure, a declaration of
their faith was often required, that none might be encouraged to try
experiments out of curiosity, in a manner which would have been very
indecent, and have tended to many bad consequences.'' Doddridge on Acts
ix. 34.}.

7. The force of imagination may also be alleged
to account for the supposed visions and voices which some enthusiasts
have believed they saw and heard; for instance, the trances of Montanus
and his followers, the visions related by some of the Fathers,
and those of the Romish saints \footnote{[The visions of Catholic saints were granted to them, as is said in
the next sentence about Scripture visions, ``for purposes distinct from
that of evidencing the doctrine.'']}; lastly, Mahomet's pretended night-journey to heaven: all
which, granting the sincerity of the reporters, may not unreasonably
be referred to the effects of disease or of an excited imagination.

8. Such, it is obvious, might be some of the
Scripture Miracles; for instance, the various appearances of Angels to
individuals, the vision of St. Paul when he was transported to the
third heaven, etc., which accordingly were wrought, as Scripture
professes, for purposes distinct from that of evidencing the doctrine,
viz., in order to become the medium of a revelation, or to confirm
faith, etc. In other cases, however, the supposition of imagination is
excluded by the vision having been witnessed by more than one person,
as the Transfiguration; or by its correspondence with distinct visions
seen by others, as in the circumstances which attended the conversion
of Cornelius; or by its connection with a permanent Miracle, as the
appearance of Christ to St. Paul in his conversion, is connected with
his blindness in consequence, which remained three days \footnote{Paley's Evidences, Part i. Prop. 2.}.

9. Much more inconclusive are those which are
actually attended by a physical cause known or suspected to be
adequate to their production. Some of those who were cured at the
tomb of the Abbé Paris were at the time making use of the usual
remedies; the person whose inflamed eye was relieved was, during his
attendance at the sepulchre, under the care of an eminent oculist;
another was cured of a lameness in the knee by the mere effort to
kneel at the tomb \footnote{Douglas, Criterion, pp. 143, 184, Note.}.
Arnobius challenges the Heathens to produce one of the pretended
miracles of their gods performed without the application of some
prescription \footnote{Stillingfleet, Book ii. Ch. x. Sec. 9.}.

10. Again, Hilarion's cures of wounds, as
mentioned by Jerome, were accompanied by the application of
consecrated oil \footnote{Middleton, Free Inquiry, iv. Sec. 2.}. The
Apostles indeed made use of oil in some of their cures \footnote{Mark vi. 13.}, but they more frequently healed without a medium of any kind.
A similar objection might be urged against the narrative of Hezekiah's
recovery from sickness, both on account of the application of the
figs, and the slowness of the cure, were it anywhere stated to have
been miraculous \footnote{2 Kings xx. 4-7.}.
Again, the dividing of the Red Sea, accompanied as it was by a strong
east wind, would not have been clearly miraculous, had it not been
effected at the word of Moses.

11. Much suspicion, too, is (as some think) cast
upon the miraculous nature of the fire, etc., which put a stop to
Julian's attempt to rebuild the Temple at Jerusalem, by the
possibility of referring it to the operation of chemical causes.

12. Lastly, answers to prayer, however
providential, are not miraculous; for in granting them, God acts by
means of, not out of, His usual system, making the ordinary course of
things subservient to a gracious purpose. Such events, then, instead
of evidencing the Divine approbation to a certain cause, must be
proved from the goodness of the cause to be what they are interpreted
to be. Yet by supposed answers to prayer, appeals to Heaven, pretended
judgments, etc., enthusiasts in most ages have wished to sanction
their claims to divine inspiration. By similar means the pretensions
of the Romish hierarchy have been supported \footnote{[But not ultimately founded and rested upon them, as has been the way
with enthusiasts.]}.

Here we close our remarks on the criterion of a
Miracle; which, it has been seen, is no one definite peculiarity,
applicable to all cases, but the combined force of a number of varying
circumstances determining our judgment in each particular instance. It
might even be said, that a determinate criterion is almost
inconceivable. For when once settled, it might appear, as was
above remarked, to be merely the physical antecedent of the
extraordinary fact; while, on the other hand, from the direction thus
given to the ingenuity of impostors, it would soon itself need a
criterion to distinguish it from its imitations. Certain it is, that
the great variety of circumstances under which the Christian Miracles
were wrought, furnishes an evidence for their divine origin, in
addition to that derived from their publicity, clearness, number,
instantaneous production, and completeness.

The exorcism of demoniacs, however, has already
been noticed as being, perhaps, in every case deficient in the proof
of its miraculous nature. Accordingly, this class of Miracles seems
not to have been intended as a primary evidence of a divine mission,
but to be addressed to those who already admitted the existence of
evil spirits, in proof of the power of Christ and His followers over
them \footnote{See Div. Leg. Book ix. Ch. v. Hence the exercise of this gift seems
almost to have been confined to Palestine. At Philippi St. Paul casts
out a spirit of divination in self-defence (Acts xvi. 16-18). In the
transaction related Acts xix. 11-17, Jews are principally concerned.}. To us, then, it
is rather a doctrine than an evidence, manifesting our Lord's power,
as other doctrines instance His mercy.

With regard to the argument from Prophecy, which
some have been disposed to abandon on account of the number of
conditions necessary for the proof of its supernatural character, it
should be remembered, that inability to fix the exact boundary of
natural sagacity is no objection to such prophecies as are undeniably
beyond it; and that the mere inconclusiveness of some of those in
Scripture, as proofs of Divine Prescience, has no positive force
against others contained in it, which furnishes a full, lasting, and,
in many cases, growing evidence of its inspiration \footnote{Some unbelievers have urged the irrelevancy of St. Matthew's citations
from the Old Testament Prophecies in illustration of the events of
Christ's life, \emph{e.g}. ch. ii. 15. It must be recollected,
however, that what is evidence in one age is often not so in another.
That certain of the texts adduced by the Evangelist furnish at the
present day no proof of Divine Prescience, is very true; but unless
some kind of argument could have been drawn from them at the time the
Gospel was written, from traditional interpretations of their sense,
we can scarcely account for St. Matthew's introducing them. The
question is, has there been a loss of what was evidence formerly, (as
is often the case,) or did St. Matthew bring forward as a prophetical
evidence what was manifestly not so, as if to hurt the effect of those
other passages, as ch. xxvii. 35, which have every appearance of being
real predictions? It has been observed, that Prophecy in general must
be obscure, in order that the events spoken of may not be understood
before their accomplishment.}.

\xsection{Section 4. On the Direct Evidence for the Christian Miracles}

IMPORTANT
as are the inquiries which I have hitherto prosecuted, it is obvious
that they do not lead to any positive conclusion, whether certain
miraculous accounts are true or not. However necessary a direct
anomaly in the course of nature may be to rouse attention, and an
important final cause to excite interest and reverence, still the
quality of the testimony on which the accounts rest can alone
determine our belief in them. The preliminary points, however, have
been principally dwelt upon, because objections founded on them form
the strong ground of unbelievers, who seem in some degree to allow the
strength of the direct evidence for the Scripture Miracles. Again, an
examination of the direct evidence is less necessary here, because,
though antecedent questions have not been neglected by Christian
writers \footnote{Especially by Vince, in his valuable Treatise
on the Christian Miracles; and Hey, in his Lectures.}, yet the evidence
itself, as might be expected, has chiefly engaged their attention
\footnote{As of Paley, Lyttelton, Leslie, etc.}. Without entering, then,
into a minute consideration of the facts and arguments on which the
credibility of the Sacred History rests, I proceed to contrast its
evidence generally with that produced for other miraculous narratives;
and thus to complete a comparison which has been already instituted,
as regards the antecedent probability and the criterion of Miracles.

For the present, then, I forego the advantage
which the Scripture Miracles have gained in the preceding Sections
over all professed facts of a similar nature. In reality, indeed, the
very same evidence which would suffice to prove the former, might be
inadequate when offered in behalf of those of the Eclectic School or
the Romish Church. For the Miracles of Scripture, and no other, are
unexceptionable, and worthy of a Divine Agent; and Bishop Butler has
clearly shown, that, in a practical question, as the divinity of a
professed Revelation must be considered, even the weakest reasons are
decisive when not counteracted by any opposite arguments \footnote{The only fair objection that can be made to this statement is, that it
is antecedently improbable that the Almighty should work Miracles with
a view to general conviction, without furnishing strong evidence that
they really occurred. This was noticed above, when the antecedent
probability of Miracles was discussed. That it is unsatisfactory to
decide on scanty evidence is no objection, as in other most important
practical questions we are constantly obliged to make up our minds and
determine our course of action on insufficient evidence.}. Whatever evidence, then, is offered for them is entirely
available to the proof of their actual occurrence; whereas evidence
for the truth of other similar accounts, supposing it to exist, would
be first employed in overcoming the objections which attach to them
all from their very character, circumstances, or object. If, however,
it can be shown that the Miracles of Scripture as far surpass all
others in their direct evidence, as they excel them in their \emph{à
priori} probability, a much stronger case will be made out in their
favour, and an additional line of distinction drawn between them and
others.

The credibility of testimony arises from the
belief we entertain of the character and competency of the witnesses;
and this is true, not only in the case of Miracles, but when facts of
any kind are examined into. It is obvious that we should be induced to
distrust the most natural and plausible statement when made by a
person whom we suspected of a wish to deceive, or of relating facts
which he had no sufficient means of knowing. Or if we credited his
narrative, we should do so, not from dependence on the reporter, but
from its intrinsic likelihood, or from circumstantial evidence. In the
case of ordinary facts, therefore, we think it needless, as indeed it
would be endless, to inquire rigidly into the credibility of the testimony by which they are conveyed to us, because they in a manner
speak for themselves. When, however, the information is unexpected, or
extraordinary, or improbable, our only means of determining its truth
is by considering the credit due to the witnesses; and then, of
course, we exercise that right of scrutiny which we before indeed
possessed, but did not think it worth while to claim. A Miracle, then,
calls for no distinct species of testimony from that offered for other
events, but for a testimony strong in proportion to the improbability
of the particular fact attested; and it is as impossible to draw any
line, or to determine how much is required, as to define the quantity
and quality of evidence necessary to prove the occurrence of an
earthquake, or the appearance of any meteoric phenomenon. Everything
depends on those attendant circumstances, of which I have already
spoken,—the object of the Miracle, the occasion, manner, and human
agent employed. If, for instance, a Miracle were said to be wrought
for an immoral object, then of course the fact would rest on the
credibility of the testimony alone, and would challenge the most rigid
examination. Again, if the object be highly interesting to us, as that
professed by the Scripture Miracles, we shall naturally be careful in
our inquiry, from an anxious fear of being biassed. But in any case
the testimony cannot turn out to be more than that of competent and
honest men; and an inquiry must not be prosecuted under the idea
of finding something beyond this, but to obtain proofs of this.

And since the existence of competency and honesty
may be established in various ways, it follows that the credibility of
a given story may be proved by distinct considerations, each of which,
separately taken, might be sufficient for the purpose. It is obvious,
moreover, as indeed is implied by the very nature of moral evidence,
that the proof of its credibility may be weaker or stronger, and yet
in both cases be a proof; and hence, that no limit can be put to the
conceivable accumulation of evidence in its behalf. Provided, then,
the existing evidence be sufficient to produce a rational conviction,
it is nothing to the purpose to urge, as has sometimes been alleged
against the Scripture Miracles, that the extraordinary facts might
have been proved by different or more over-powering evidence. It has
been said, for instance, that no testimony can fairly be trusted which
has not passed the ordeal of a legal examination. Yet, calculated as
that mode of examination undoubtedly is to elicit truth, surely truth
may be elicited by other ways also. Independent and circumstantial
writers may confirm a fact as satisfactorily as witnesses in Court.
They may be questioned and cross-questioned, and, moreover, brought up
for re-examination in any succeeding age; whereas, however great may
be the talents and experience of the men who conduct the legal
investigation, yet when they have once closed it, and given in their
verdict, we believe upon their credit, and we have no means of
examining for ourselves. To say, however, that this kind of evidence
might have been added to the other, in the case of the Christian Miracles
\footnote{Some of our Saviour's Miracles, however, \emph{were} subjected to
judicial examination. (See John v. and ix.) In v. 16, the measures of
the Pharisees are described by the technical word [\emph{edi\={o}kon}].}, is merely to assert that
the proof of the credibility of Scripture might have been stronger
than it is; which I have already allowed it might have been, without
assignable limit.

The credibility, then, of Testimony depending on
the evidence of honesty and competency in those who give it, it is
prejudicial, first, to their character for honesty—

1. If desire of gain, power, or other temporal
advantage may be imputed to them. This would detract materially from
the authority of Philostratus, even supposing him to have been in a
situation for ascertaining the truth of his own narrative, as he
professes to write his account of Apollonius at the instance of his
patroness, the Empress Julia, who is known to have favoured the
Eclectic cause. Again, the account of the Miracle performed on the
door-keeper at the cathedral at Saragossa, on which Hume insists,
rests principally upon the credit of the Canons, whose interest was
concerned in its establishment. This remark, indeed, obviously applies
to the Romish Miracles generally \footnote{[The Miracles of Catholic Saints as little benefited their workers as
the Miracles of the Apostles.]}. The Christian Miracles, on the contrary, were attested by the
Apostles, not only without the prospect of assignable worldly
advantage, but with the certainty and after the experience of actual
suffering.

2. When there is room for suspecting party spirit
or rivalry, as in the miraculous biographies of the Eclectic
philosophers; in those of Loyola and other saints of the rival orders
in the Romish Church; and in the present Mahometan accounts of the
Miracles of Mahomet, which, not to mention other objections to them,
are composed with an evident design of rivalling those of Christ \footnote{See Professor Lee's Persian Tracts, pp. 446, 447.}.

3. Again, a tale once told may be persisted in
from shame of retracting, after the motives which first gave rise to
it have ceased to act, even at the risk of suffering. This remark
cannot apply to the case of the Apostles, until some reason is
assigned for their getting up their miraculous story in the first
instance. If necessary, however, it could be brought with force
against any argument drawn from the perseverance of the witnesses
for the cures professedly wrought by Vespasian, ``postquam nullum
mendacio pretium;'' for, as they did not suffer for persisting in their
story, had they retracted, they would have gratuitously confessed
their own want of principle.

4. A previous character for falsehood is almost
fatal to the credibility of a witness of an extraordinary narrative;
for instance, the notorious insincerity and frauds of the Church of
Rome in other things are in themselves enough to throw a strong
suspicion on its testimony to its own Miracles \footnote{[There \emph{have} been frauds among Catholics, and for gain, as among
Protestants, whether churchmen or dissenters, or among antiquarians,
or transcribers of MSS., or picture-dealers, or horse-dealers; for the
``Net gathers of every kind;'' but that does not prove the Church to be
fraudulent, unless geological or chemical frauds are slurs upon the
character of the British Association.]}. The primitive Church is in some degree open to a charge of a
similar nature \footnote{Hey, Lectures, book i. ch. xii. sec. 15.}. Or an
intimacy with suspicious characters; for instance, Prince Hohenlohe's
connection with the Romish Church, and that of Philostratus with the
Eclectics, since both the Eclectic and Romish Schools have
countenanced the practice of what are called pious frauds \footnote{Hey, Lectures, book i. ch. xii. sec. 15.}.

5. Inconsistencies or prevarications in the
testimony, marks of unfairness, exaggeration, suppression of
particulars, etc. Of all these, Philostratus stands convicted, whose
memoir forms a remarkable contrast to the artless and candid
narratives of the Evangelists. The Books of the New Testament,
containing as they do separate accounts of the same transactions,
admit of a minute cross-examination, which terminates so decidedly in
favour of their fidelity, as to recommend them highly on the score of
honesty, even independently of the known sufferings of the writers.

6. Lastly, objection may be taken to witnesses
who have the opportunity of being dishonest; as those who write at a
distance from the time and place of the professed Miracle, or without
mentioning particulars, etc. But on these points I shall speak
immediately in a different connection.

Secondly, witnesses must be not only honest, but
competent also; that is, such as have ascertained the facts which they
attest, or who report after examination. Here then I notice—

1. Deficiency of examination implied in the
circumstances of the case. As when it is first published in an age or
country remote from the professed time and scene of action; for in
that case room is given to suspect failure of memory, imperfect
information, etc., whereas to write in the presence of those who know
the circumstances of the transactions is an appeal which increases the
force of the testimony by associating them in it. Accounts, however,
whether miraculous or otherwise, possess very little intrinsic
authority, when written so far from the time or place of the
transactions recorded, as the biographies of Pythagoras, Apollonius,
Gregory Thaumaturgus, Mahomet, Loyola, or Xavier \footnote{Paley, Evidences, Part i. Prop. 2.}. The opposite circumstances of the Christian Testimony have
often been pointed out. Here we may particularly notice the
providential dispersion of the Jews over the Roman Empire before the
age of Christ; by which means the Apostles' testimony was given in
heathen countries, as well as in Palestine, in the face of those who
had both the will and the power to contradict it, if incorrect.

While the testimony of contemporaries is
necessary to guarantee the truth of ordinary history, Miracles require
the testimony of eye-witnesses. For ordinary events are believed in
part from their being natural, but testimony being the main support of
a miraculous narrative, must in that case be the best of its kind.
Again, we may require the testimony to be circumstantial in reference
to dates, places, persons, etc.; for the absence of these seems to
imply an imperfect knowledge, and at least gives less opportunity of
inquiry to those who wish to ascertain its fidelity \footnote{The vagueness of the accounts of miraculous interpositions related by
the Fathers is pointed out by Middleton. (Free Inquiry, ii. p. 22.) [\emph{Vide
infra}, Essay ii., n. 137, 138.]}.

Miracles which are not lasting do not admit of
adequate examination; as visions, extraordinary voices, etc. The cure
of diseases, on the other hand, is a permanent evidence of a
divine interposition; particularly such cures of bodily imperfections
as are undeniably miraculous in their nature, as well as permanent; to
these, then, our Lord especially appeals in evidence of His divine
mission \footnote{Matt. xi. 5.}. Lastly,
statements are unsatisfactory in which the miracle is described as
wrought before a very \emph{few}; for room is allowed for suspecting
mistake, or an understanding between the witnesses. Or, on the other
hand, those wrought \emph{in a confused crowd}; such are many standing
miracles of the Romanists, which are exhibited with the accompaniment
of imposing pageants, or on a stage, or at a distance, or in the midst
of candles and incense \footnote{[Candles and incense are commonly used in the daytime; and our Lord
wrought many of His miracles in a throng which was pressing upon Him.]}.
Our Saviour, on the contrary, bids the lepers He had cleansed \emph{show}
themselves to the Priests, and make the customary offering as a \emph{memorial}
of their cures \footnote{Luke v. 14; xvii. 14.}. And
when He appeared to the Apostles after His Resurrection, He allowed
them to examine His hands and feet \footnote{Luke xxiv. 39, 40.}. Those of the Scripture Miracles which were wrought before
few, or in a crowd, were permanent; as cures \footnote{Mark viii. 22-26.}, and the raising of Jairus's daughter; or were of so vast a
nature, that a crowd could not prevent the witnesses from
ascertaining the fact, as the standing still of the Sun at the word of
Joshua.

2. Deficiency of examination implied in the
character, etc., of the witnesses: (1) for instance, if there be any
suspicion of their derangement, or if there be an evident defect in
those bodily or mental faculties which are necessary for examining the
Miracle, as when the intellect or senses are impaired. Number in the
witnesses refutes charges of this nature; for it is not conceivable
that many should be deranged or mistaken at once, and in the same way.

(2) Enthusiasm, ignorance, and habitual
credulity, are defects which no number of witnesses removes. The
Jansenist Miracles took place in the most ignorant and superstitious
district of Paris \footnote{The Faubourg St. Marcel. Less.}.
Alexander Pseudo-mantis practised his arts among the Paphlagonians, a
barbarous people. Popish Miracles and the juggles of the Heathen
Priests have been most successful in times of ignorance \footnote{[Might not the same insinuations be thrown out against the miracles of
Elisha? On the other hand, was the age of St. Ambrose and St.
Augustine ignorant? or that of St. Philip Neri?]}.

Yet, while we reasonably object to gross
ignorance or besotted credulity in witnesses for a miraculous story,
we must guard against the opposite extreme of requiring the testimony
of men of science and general knowledge. Men of philosophical minds
are often too fond of inquiring into the causes and mutual
dependence
of events, of arranging, theorizing, and refining, to be accurate and
straightforward in their account of extraordinary occurrences. Instead
of giving a plain statement of facts, they are insensibly led to
correct the evidence of their senses with a view to account for the
strange phenomenon; as Chinese painters, who, instead of drawing in
perspective, give lights and shadows their supposed meaning, and
depict the prospect as they think it should be, not as it is \footnote{It is well known, that those persons are accounted the best
transcribers of MSS. who are ignorant of the language transcribed; the
habit of \emph{correcting} being almost involuntary in men of letters.}. As Miracles differ from other events only when considered
relatively to a general system, it is obvious that the same persons
are competent to attest miraculous facts who are suitable witnesses of
corresponding natural ones. If a peasant's testimony be admitted to
the phenomenon of meteoric stones, he may evidence the fact of an
unusual and unaccountable darkness. A physician's certificate is not
needed to assure us of the illness of a friend; nor is it necessary
for attesting the simple fact that he has instantaneously recovered.
It is important to bear this in mind, for some writers argue as if
there were something intrinsically defective in the testimony given by
ignorant persons to miraculous occurrences \footnote{Hume, On Miracles, Part ii. Reason 1.}. To say that unlearned persons are not judges of the fact
of a miraculous event, is only so far true as all testimony is
fallible and liable to be distorted by predjudice. Every one, not only
superstitious persons, is apt to interpret facts in his own way; if
the superstitious see too many prodigies, men of science may see too
few. The facility with which the Japanese ascribed the ascent of a
balloon, which they witnessed at St. Petersburgh, to the powers of
magic, (a circumstance which has been sometimes urged against the
admission of unlearned testimony \footnote{Bentham, Preuves Judiciaires, Liv. viii. Ch. ii.},) is only the conduct of theorists accounting for a novel
phenomenon on the principles of their own system.

It may be said, that ignorance prevents a witness
from discriminating between natural and supernatural events, and thus
weakens the authority of his judgment concerning the miraculous nature
of a fact. It is true; but if the fact be recorded, we may judge for
ourselves on that point. Yet it may be safely said, that not even
before persons in the lowest state of ignorance could any great
variety of professed Miracles be displayed without their
distinguishing rightly, on the whole, between the effects of nature
and those of a power exterior to it; though in particular instances
they doubtless might be mistaken. Much more would this be the case
with the lower ranks of a civilized people. Practical intelligence is
insensibly diffused from class to class; if the upper ranks are
educated, numbers besides them, without any formal and systematic
knowledge, almost instinctively discriminate between natural and
supernatural events. Here science has little advantage over common
sense; a peasant is quite as certain that a resurrection from the dead
is miraculous as the most able physiologist \footnote{It has been observed, that more suitable witnesses could not be
selected of the fact of a miraculous draught of fishes than the
fishermen of the lake wherein it took place.}.

The original witnesses of our Saviour's Miracles
were very far from a dull or ignorant race. The inhabitants of a
maritime and border country, as Galilee was, engaged, moreover, in
commerce, composed of natives of various countries, and, therefore,
from the nature of the case, acquainted with more than one language,
have necessarily their intellects sharpened and their minds
considerably enlarged, and are of all men least disposed to acquiesce
in marvellous tales \footnote{See Less, Opuscul.}.
Such a people must have examined before they suffered themselves to be
excited in the degree which the Evangelists describe \footnote{If, on the other hand, we would see with how unmoved an unconcern men
receive accounts of miracles, when they believe them to be events of
every-day occurrence, we may turn to the conduct of the African
Christians in the Age of Austin, whom that Father in vain endeavoured
to interest in miraculous stories of relics, etc., by formal accounts
and certificates of the cures wrought by them. (See Middleton, p.
138.) The stir, then, which the miracles of Christ made in Galilee
implies, that they were not received with an indolent belief. It must
be noticed, moreover, in opposition to the statement of some
unbelievers, that great numbers of the Jews were converted (Acts ii.
41; iv. 4; v. 13, 14; vi. 7; ix. 35; xv. 5; xxi. 20). On this subject,
see Jenkin, On the Christian Religion, Vol. ii. Ch. xxxii.}. But even supposing those among them who were in
consequence convinced of the divine mission of Christ, were of a more
superstitious turn of mind than the rest, still this is not sufficient
to account for their conviction. For superstition, while it might
facilitate the bare admission of miraculous events, would at the same
time weaken their practical influence. Miracles ceasing to be
accounted strange, would cease to be striking also. Whereas the
conviction wrought in the minds of these men was no bare and indolent
assent to facts which they might have thought antecedently probable or
not improbable, but a conversion in principles and mode of life, and a
consequent sacrifice of all that nature holds dear, to which none
would submit except after the fullest examination of the authority
enjoining it. If additional evidence be required, appeal may be made
to the multitude of Gentiles in Greece and Asia, in whose principles
and mode of living belief in the Miracles made a change even more
striking and complete than was effected in the case of the Jews. In a
word, then, the conversion which Christ and His Apostles effected
invalidates the charge of blind credulity in the witnesses; the
practical nature of the belief wrought in them proving that it was
founded on an examination of the Miracles.

(3) Again, it weakens the authority of the
witnesses, if their belief can be shown to have been promoted by the
influence of superiors; for then they virtually cease to be themselves
witnesses, and report the facts on the authority (as it were) of their
patrons. It is observable, that the national conversions of the Middle
Ages generally began with the princes, and descended to their
subjects; those of the Apostolic Age obviously proceeded in the
reverse order \footnote{Mosheim, Eccl. Hist. Cent. vi. viii. ix.}.

(4) It is almost fatal to the validity of the
testimony, if the miracle which is attested coincides with a previous
system, or supports a cause already embraced by the witnesses. Men are
always ready to believe what flatters their own opinions, and of all
prepossessions those of Religion are the strongest. There is so much
in the principle of all Religion that is true and good, so much
conformable to the best feelings of our nature, which perceives itself
to be weak and guilty, and looks out for an unseen and superior being
for guidance and support; and the particular worship in which each
individual is brought up is so familiarized to him by habit, so
endeared to his affections by the associations of place and the
recollections of past years, so connected too with the ordinary
transactions and most interesting events of life, that even should
that form be irrational and degrading, still it will in most cases
preserve a strong influence over his mind, and dispose him to
credit upon slight examination any arguments adduced in its defence.
Hence an account of Miracles in confirmation of their own Religion
will always be favourably received by men whose creed has already led
them to expect such interpositions of superior beings. This
consideration invalidates at once the testimony commonly offered for
Pagan and Popish Miracles, and in no small degree that for the
Miracles of the primitive Church \footnote{[Vide Essay ii. n. 36-45. Ecclesiastical Miracles are mainly the
rewards of faith; not, strictly speaking, evidence.]}. The professed cures of Vespasian were performed in honour of
Serapis in the midst of his worshippers; and the people of Saragossa,
who attested the Miracle wrought in the case of the door-keeper of the
Cathedral, had previous faith in the virtues of holy oil \footnote{It has been noticed as a suspicious circumstance in the testimony to
the reported miracle wrought in the case of the Confessors in the
persecution of the Arian Hunneric, that Victor Vitensis, one of the
principal witnesses, though writing in Africa, where it professedly
took place, and where the individuals thus distinguished were then
living, yet refers only to one of them, who was then living at the
Athanasian Court at Constantinople, and held in particular honour by
Zeno and the Empress.—``If any one doubt the fact, let him go to
Constantinople.'' See the whole evidence in Milner's Church History,
Cent. v. Ch. xi.; who, however, strongly defends the miracle. Gibbon
pretends to do the same, with a view to provide a rival to the Gospel
Miracles.}.

Here the evidence for the Scripture Miracles is
unique. In other cases the previous system has supported the
Miracles, but here the Miracles introduced and upheld the system. The
Christian Miracles in particular \footnote{Not to mention those of Moses and Elijah.} were received on their own merits; and the admission of them
became the turning-point in the creed and life of the witnesses, which
thenceforth took a new and altogether different direction. But,
moreover, as if their own belief in them were not enough, the Apostles
went out of their way to debar any one from the Christian Church who
did not believe them as well as themselves \footnote{Campbell, On Miracles, Part ii. Sec. 1.}. Not content that men should be converted on any ground, they
fearlessly challenged refutation, by excluding from their fellowship
of suffering any who did not formally assent, as a necessary condition
of admittance and a first article of faith, to one of the most
stupendous of all the miracles, their Master's Resurrection from the
dead;—a procedure this, which at once evinces their own unqualified
conviction of the fact, and associates, too, all their converts with
them as believers in a miracle contemporary with themselves. Nor is
this all; a religious creed necessarily prejudices the mind against
admitting the miracles of hostile sects, in the very same proportion
in which it leads it to acquiesce in such as support its own dogmas \footnote{Ibid. Part i. Sec. 4.}. The Christian Miracles, then, have the strongest of
conceivable attestations, in the conversion of many who at first
were prejudiced against them, and in the extorted confession of
enemies, who by the embarrassment which the admission occasioned them,
at least showed that they had not made it till after a full and
accurate investigation of the extraordinary facts.

(5) It has been sometimes objected, that the
minds of the first converts might be wrought upon by the doctrine of a
future state which the Apostles preached, and be thus persuaded to
admit the miracles without a rigorous examination \footnote{Gibbon, particularly Ch. xv.}. But, as Paley well replies, evidence of the truth of the
promise would still be necessary; especially as men rather demand than
dispense with proof when some great and unexpected good is reported to
them. Yet it is more than doubtful whether the promise of a future
life would excite this interest; for the desire of immortality, though
a natural, is no permanent or powerful feeling, and furnishes no
principle of action. Most men, even in a Christian country, are too
well satisfied with this world to look forward to another with any
great and settled anxiety. Supposing immortality to be a good, it is
one too distant to warm or influence them. Much less are they disposed
to sacrifice present comfort, and strip themselves of former opinions
and habits, for the mere contingency of future bliss. The hope of
another life, grateful as it is under affliction, will not induce
a man to rush into affliction for the sake of it. The inconvenience of
a severe complaint is not out-balanced by the pleasure of a remedy. On
the other hand, though we know that gratuitous declarations of coming
judgments and divine wrath may for a time frighten weak minds, they
will neither have effect upon strong ones, nor produce a permanent and
consistent effect upon any. Persons who are thus wrought upon in the
present day believe the denunciations because they are in Scripture,
not Christianity because Scripture contains them. The authority of
Revealed Religion is taken for granted both by the preacher and his
hearers. On the whole, then, it seems inconceivable that the promise
or threat of a future life should have supplied the place of previous
belief in Christianity, or have led the witnesses to admit the
Miracles on a slight examination.

(6) Lastly, love of the marvellous, of novelty,
etc., may be mentioned as a principle influencing the mind to
acquiesce in professed miracles without full examination. Yet such
feelings are more adapted to exaggerate and circulate a story than to
invent it. We can trace their influence very clearly in the instances
of Apollonius and the Abbé Paris, both of whom had excited attention
by their eccentricities, before they gained reputation for
extraordinary power \footnote{See the Author's memoir of Apollonius.—Of the Abbé, Mosheim says, ``Diem
vix obierat, voluntariis cruciatibus et pœnis exhaustus, mirabilis
iste homo, quum immensa hominum multitudo ad ejus corpus conflueret;
quorum alii pedes ejus osculabantur, alii partem capillorum
abscindebant, quam sancti loco pignoris ad mala quævis averruncanda
servarent, alii libros et lintea quæ attulerant, cadaveri admovebant
quod virtute quadam divina plenum esse putabant. Et statim vis illa
mirifica, quâ omne, quod in terrâ hâc reliquit, præditum esse
fertur, apparebat,'' etc. Inquisit. in verit. Miraculor. F. de Paris,
Sec. 1.}.
Such principles, moreover, are not in general practical, and have
little power to sustain the mind under continued opposition and
suffering \footnote{Paley, Evidences, Part i. Prop. 2.}.

———————

These are some of the obvious points which will
come into consideration in deciding upon the authority of testimony
offered for miracles; and they enable us at once to discriminate the
Christian story from all others which have been set up against it.
With a view of simplifying the argument, the evidence for the Jewish
miracles has been left out of the question \footnote{The truth of the Mosaic narrative is proved from the genuineness of
the Pentateuch, as written to contemporaries and eye-witnesses of the
miracles; from the predictions contained in the Pentateuch; from the
very existence of the Jewish system (Sumner's Records); and from the
declarations of the New Testament writers. The miracles of Elijah and
Elisha are proved to us by the authority of the Books in which they
are related, and by means of the New Testament.}; because, though strong and satisfactory, it is not at the
present day so directly conclusive as that on which the Christian
rest. Nor is it necessary, I conceive, to bring evidence for more than
a fair proportion of the Miracles; supposing, that is, those which
remain unproved are shown to be similar to them, and indissolubly
connected with the same system. It may be even said, that if the
single fact of the Resurrection be established, quite enough will have
been proved for believing all the Miracles of Scripture.

Of course, however, the argument becomes far
stronger when it is shown that there \emph{is} evidence for the great bulk of
the miracles, though not equally strong for some as for others; and
that the Jewish, sanctioned as they are by the New Testament, may also
be established on distinct and peculiar grounds. Nor let it be
forgotten, that the Christian story itself is supported, over and
above the evidence that might fairly be required for it, by several
bodies of testimony quite independent of each other \footnote{The fact of the Christian miracles may be proved, first, by the
sufferings and consistent story of the original witnesses; secondly,
from the actual conversion of large bodies of men in the age in which
they are said to have been wrought; thirdly, from the institution, at
the time, of a day commemorative of the Resurrection, which has been
observed ever since; fourthly, by collateral considerations, such as
the tacit assent given to the miracles by the adversaries of
Christianity, the Eclectic imitations of them, and the pretensions to
miraculous power in the primitive Church. These are distinct
arguments; no one of them absolutely presupposes the genuineness of
the Scripture narrative, though the force of the whole is much
increased when it is proved.}. By separate processes of reasoning it may be shown, that
if Christianity was established without miracles, it was, to say the
least, an altogether singular and unique event in the history of
mankind; and the extreme improbability of so many distinct and
striking peculiarities uniting, as it were, by chance in one and the
same case, raises the proof of its divine origin to a moral certainty.
In short, it is only by being made unnatural that the Christian
narrative can be deprived of a supernatural character; and we may
safely affirm that the strongest evidence we possess for the most
certain facts of other history, is weak compared to that on which we
believe that the first preachers of the Gospel were gifted with
miraculous powers.

And thus a case is established so strong, that
even were there an antecedent improbability in the facts attested, in
most judgments it would be sufficient to overcome it. On the contrary,
we have already shown their intrinsic character to be exactly such as
our previous knowledge of the attributes and government of the
Almighty would lead us to expect in works ascribed to Him. Their
grandeur, beauty, and consistency; the clear and unequivocal marks
they bear of superhuman agency; the importance and desirableness of
the object they propose to effect, are in correspondence with the
variety and force of the evidence itself.

Such, then, is
the contrast they present to all other professed miracles, from those
of Apollonius downwards— which have been shown, more or less,
to be improbable from the circumstances of the case, inconclusive when
considered as marks of divine interference, and quite destitute of
good evidence for their having really occurred.

Lastly, it must
be observed, that the proof derived from interruptions in the course
of nature, though a principal, is yet but one out of many proofs on
which the cause of Revealed Religion rests; and that even supposing
(for the sake of argument) it were altogether inconclusive at the
present day, still the other evidences \footnote{Such as the system of doctrine, marks of design, gradual disclosure of
unknown truths, etc., connecting together the whole Bible as the work
of one mind:—Prophecy:—the character of Christ:—the morality of
the Gospel:—the wisdom of its doctrines, displaying at once
knowledge of the human heart, and skill in engaging its affections,
cet.}, as they are called, would be fully equal to prove to us the
divine origin of Christianity.

\xchapter{Essay II. The Miracles of Early Ecclesiastical History, Compared with those of Scripture, as regards Their Nature, Credibility, and Evidence}

\xsection{Chapter 1. Introduction}

1.
SACRED History is distinguished
from Profane by the nature of the facts which enter into its
composition, and which are not always such as occur in the ordinary
course of things, but are extraordinary and divine. Miracles are its
characteristic, whether it be viewed as biblical or ecclesiastical: as
the history of a reign or dynasty more or less approximates to
biography, as the history of a wandering tribe passes into romance or
poetry, as a constitutional history borders on a philosophical
dissertation, so the history of Religion is necessarily of a
theological cast, and is occupied with the supernatural. It is a
record of ``the kingdom of heaven,'' a manifestation of the Hand of
God; and, ``the temple of God being opened,'' and ``the ark of His
testament,'' there are ``lightnings and voices,'' the momentary yet
recurring tokens of that conflict between good and evil, which is
waging in the world of spirits from age to age. This supernatural
agency, as far as it is really revealed to us, is from its very nature
the most important of the characteristics of sacred history, and the
mere rumour of its manifestation excites interest in consequence of
the certainty of its existence. But since the miraculous statements
which are presented to us are often not mere rumours or surmises, but
in fact essential to the narrative, it is plain that to treat any such
series of events, (for instance, the history of the Jews, or of the
rise of Christianity, or of the Catholic Church,) without taking them
into account, is to profess to write the annals of a reign, yet to be
silent about the monarch,—to overlook, as it were, his personal
character and professed principles, his indirect influence and
immediate acts.

2.
Among the subjects, then, which the history of the early centuries of
Christianity brings before us, and which are apt more or less to
startle those who with modern ideas commence the study of Church
History generally, (such as the monastic rule, the honour paid to
celibacy, and the belief in the power of the keys,) it seems right to
bestow attention in the first place on the supernatural narratives
which occur in the course of it, and of which various specimens will
be found in any portion of it which a reader takes in hand. It will
naturally suggest itself to him to form some judgment upon them,
and a perplexity, perhaps a painful perplexity, may ensue from the
difficulty of doing so. This being the case, it is inconsiderate and
almost wanton to bring such subjects before him, without making at
least the attempt to assist him in disposing of them. Accordingly, the
following remarks have been written in discharge of a sort of duty
which a work of Ecclesiastical History involves \footnote{[The occasion of this Essay was the publication
of a portion of Fleury's Ecclesiastical History in English.]},—not indeed without a deep sense of the arduousness of such an essay,
or of the incompleteness and other great defects of its execution, but
at the same time, as the writer is bound to add, without any apology
at all for discussing in his own way a subject which demands
discussion, and which, if any other, is an open question in the
English Church, and has only during the last century been viewed in a
light which he believes to be both false in itself, and dangerous
altogether to Revealed Religion.

3.
It may be advisable to state in the commencement the conclusions to
which the remarks which follow will be found to tend; they are such as
these:—that Ecclesiastical Miracles, that is, Miracles posterior to
the Apostolic age, are on the whole different in object, character,
and evidence, from those of Scripture on the whole, so that the one
series or family ought never to be confounded with the other; yet that the former are not therefore at once to be rejected; that
there was no Age of Miracles, after which miracles ceased; that there
have been at all times true miracles and false miracles, true accounts
and false accounts; that no authoritative guide is supplied to us for
drawing the line between the two; that some of the miracles reported
were true miracles; that we cannot be certain how many were not true;
and that under these circumstances the decision in particular cases is
left to each individual, according to his opportunities of judging.

\xchapter{Chapter 2. On the Antecedent Probability of the Ecclesiastical Miracles}

4.
A FACT is properly called ``improbable,'' only when it has some quality or circumstance
attached to it which operates to the disadvantage of evidence adduced
in its behalf. We can scarcely avoid forming an opinion for or against
any statement which meets us; we feel well-disposed towards some
accounts or reports, averse from others, sometimes on no reason
whatever beyond our accidental frame of mind at the moment, sometimes
because the facts averred flatter or thwart our wishes, coincide or
interfere with the view of things familiar to us, please or startle
our imagination, or on other grounds equally vague and untrustworthy.
Such anticipations about facts are as little blameable as the fancies
which spontaneously rise in the mind about a person's stature and
appearance before seeing him; and, like such fancies, they are
dissipated at once when the real state of the case is in any way
ascertained. They are simply notional; and form no presumption
in reason, for or against the facts, or the evidence of the facts, to
which they relate.

5.
An antecedent improbability, then, in certain facts, to be really
such, must avail to prejudice the evidence which is offered in their
behalf, and must be of a nature to diminish or destroy its force. Thus
it is improbable, in the highest degree, that our friend should have
done an act of fraud or injustice; and improbable again, but in a
slight degree, that our next-door neighbour should have been highly
promoted, or that he should have died suddenly. We do not acquiesce in
any evidence whatever that comes to hand even for the latter
occurrence, and in none but the very best for the former. Again, there
is a general improbability attaching to the notion that the members of
certain sects or of certain political parties should commit themselves
to this or that cast of opinions, or line of conduct; and, on the
other hand, though there is no general improbability that individuals
of the poorest class should make large fortunes, yet a strong
probability may lie against certain given persons of that class in
particular.

6.
Now it may be asserted that there is no presumption whatever against
miracles generally in the ages after the Apostles, though there may be
and is a certain antecedent improbability in this or that particular
miracle.

There is no presumption against Ecclesiastical
Miracles generally, because inspiration has stood the brunt of any
such antecedent objection, whatever it be worth, by its own
supernatural histories, and in establishing their certainty in fact,
has disproved their impossibility in the abstract. If miracles are
antecedently improbable, it is either from want of a cause to which
they may be referred, or of experience of similar events in other
times and places. What neither has been before, nor can be attributed
to an existing cause, is not to be expected, or is improbable. But
Ecclesiastical Miracles are occurrences not without a parallel; for
they follow upon Apostolic Miracles, and they are referable to the
Author of the Apostolic as an All-sufficient Cause. Whatever be the
regularity and stability of nature, interference with it can be,
because it has been; there is One who both has power over His own
work, and who before now has not been unwilling to exercise it. In
this point of view, then, Ecclesiastical Miracles are more
advantageously circumstanced than those of Scripture.

7.
What has happened once, may happen again; the force of the presumption
against Miracles lies in the opinion entertained of the inviolability
of nature, to which the Creator seems to ``have given a law which
shall not be broken.'' When once that law is shown to be but general,
not necessary, and (if the word may be used) when its \emph{prestige}
is once destroyed, there is nothing to shock the imagination in
a miraculous interference twice or thrice, as well as once. What never
has yet happened is improbable in a sense quite distinct from that in
which a thing is improbable which has before now happened; the
improbability of the latter class of facts may be greater or less, it
may be very great; but whatever the strength of the improbability, it
is different in kind from the improbability attaching to such as admit
of being called impossible by those who reject them.

8.
It may be urged in reply, that the precedent of Scripture is no
special recommendation of Ecclesiastical Miracles; for the abstract
argument against miracles, as such, has little or no force, as soon as
the mere doctrine of a Creator and Supreme Governor is admitted, and
even prior to any reference to inspired history; that there is no
question among religious men of the existence of a Cause \emph{adequate}
to the production of miracles anywhere or at any period; the question
rather is whether He \emph{will} work them; whether the \emph{Ecclesiastical}
Miracles themselves, being what and when they were, are probable, not
whether there is a general presumption against them all simply \emph{as
miracles}; on the other hand, that while the Scripture Miracles
avail little as a precedent for subsequent miracles, \emph{as miracles},
for no precedent is wanted, they do actually tend to discredit them, \emph{as
being subsequent}, for from the nature of the case irregularities
can be but rarely allowed in any system. It is at first sight
not to be expected that the Author of nature should interrupt His own
harmonious order at all, though He is powerful to do so; and therefore
the fact of His having done so once makes it only less probable that
He will do so again. Moreover, if any recurrence of miraculous action
is to be anticipated, it is the recurrence of a similar action, not a
manifestation of power, ever so different from it; whereas the
miracles of the ages subsequent to the Apostles are on the whole so
very unlike those of which we read in Scripture, in their object,
circumstances, nature, and evidence, as even to be disproved by the
very contrast. This is what may be objected.

9.
Now as far as this representation involves the discussion of the
special character and circumstances of the Ecclesiastical Miracles, it
will come under consideration in the next Chapter; here we are only
engaged with the abstract question, whether the fact that miracles
have once occurred, and that under certain circumstances and with
certain characteristics, does or does not prejudice a proof when
offered, of their having occurred again, and that under other
circumstances
and with other characteristics.

10.
On this point many writers have expressed opinions which it is
difficult to justify. Thus Bishop Warburton, in the course of some
excellent remarks on the Christian miracles, is led to propose a
certain test of true miracles, founded on their professed \emph{object},
and suggests that this will furnish us with means of drawing the line
of supernatural agency in the early Church. ``If [the \emph{final cause}],''
he says, ``be so \emph{important} as to make the miracle \emph{necessary
to the ends of the dispensation}, this is all that can be
reasonably required to entitle it to our belief;'' so far he is
vindicating the Apostolic Miracles, and his reasoning is
unexceptionable; but he adds in a note, ``Here, by the way, let me
observe, that what is now said gives that \emph{criterion} which Dr.
Middleton and his opponents, in a late controversy concerning
miracles, demanded of one another, and which yet both parties, for
some reasons or other, declined to give; namely, some certain mark to
enable men to distinguish, for all the purposes of religion, between
true and certain miracles, and those which were false or
doubtful.'' \footnote{Pp. 239-241, Edit. 4.} He begins
by saying that miracles which subserve a certain object deserve our
consideration, he ends by saying that those which do not subserve it
do not deserve our consideration, and he makes himself the judge
whether they subserve it or not.

11.
Bishop Douglas, too, after observing that the miracles of the second
and third centuries have a character less clearly supernatural and an
evidence less cogent than those of the New Testament, and that the
fourth and fifth are ``ages of credulity and superstition,''
and the miracles which belong to them are ``wild and ridiculous,''
proceeds to lay down a \emph{decisive criterion} between true miracles
and their counterfeits, and this criterion he considers to be the gift
of inspiration in their professed workers. ``Though it may be a
matter more of curiosity than of use, to endeavour to determine the
exact time when miraculous powers were withdrawn from the Church, yet \emph{I
think that it may be determined with some degree of exactness}. The
various opinions of learned Protestants, who have extended them at all
after the Apostles, show how much they have been at a loss with regard
to this, which has been urged by Papists with an air of triumph, as
if, Protestants not being able to agree when the age of miracles was
closed, this were an argument of its not being closed as yet. If there
be anything in this objection, though perhaps there is not, \emph{I think
I have it in my power to obviate it}, by fixing upon a period,
beyond which we may be \emph{certain} that miraculous powers did not
subsist.'' Then he refers to his argument in favour of the New
Testament miracles, that ``what we know of the attributes of the
Deity, and of the usual methods of His government, \emph{inclines us to
believe} that miracles will never be performed by the agency and
instrumentality of men, but when these men are set apart and chosen by
God to be His ambassadors, as it were, to the world, to deliver some
message or to preach some doctrine as a law from heaven; and in
this case their being vested with a power of working miracles is the
best credential of the divinity of their mission.'' So far, as
Warburton, this author keeps within bounds; but next he proceeds, as
Warburton also, to extend his argument from a defence of what is true
to a test of what is false. ``If we set out with this as a \emph{principle},
then shall we easily determine when it was that miracles ceased to be
performed by Christians; for we shall be led to conclude that the age
of Christian miracles must have ceased with the age of Christian
inspiration. So long as Heaven thought proper to set apart any
particular set of men to be the authorized preachers of the new
religion revealed to mankind, so long, may we rest satisfied,
miraculous powers were continued. But whenever this purpose was
answered, and inspiration ceased to be any longer necessary, by the
complete publication of the Gospel, then would the miraculous powers, \emph{whose
end was to prove the truth of inspiration}, be \emph{of course}
withdrawn.'' \footnote{Pp. 239-241, Edit. 4.}

12.
Here he determines \emph{à priori} in the most positive manner the ``end'' or object of miracles in the designs of Providence. That it
is very natural and quite consistent with humility to form antecedent
notions of what is likely and what not likely, as in other matters, so
as regards the Divine dealings with us, has been implied above; but it
is neither reverent nor philosophical in a writer to ``think he
has it in his power'' to dispense with good evidence in behalf of
what professes to be a work of God, by means of a summary criterion of
his own framing. His very mode of speech, as well as his procedure,
reminds us of Hume, who in like manner, when engaged in invalidating
the evidence for all miracles whatever, observes that ``nothing is so
convenient as a \emph{decisive} argument,'' (such as Archbishop
Tillotson's against the Real Presence,) ``which must at least
silence the most arrogant bigotry and superstition, and free us from
their impertinent solicitations,'' and then ``\emph{flatters himself
that he has discovered} an argument of a like nature, which, if
just, will, with the wise and learned, be an everlasting check to all
kinds of superstitious delusion, and, consequently, will be useful as
long as the world endures.''

13.
It is observable that in another place Douglas had said, that ``though we may be certain that God will never reverse the course of
nature but for important ends, (the course of nature being the plan of
government laid down by Himself,) Infinite Wisdom may see ends highly
worthy of a miraculous interposition, the importance of \emph{which may
lie hid front our shallow comprehension}. Were, therefore, the
miracles, about the credibility of which we now dispute, events
brought about by invisible agency, though our being able to discover
an important end served by a miracle would be no weak additional
motive to our believing it; yet our \emph{not} being able to discover
any such end \emph{could be no motive to induce us to reject it}, if
the testimony produced to confirm it be unexceptionable.'' \footnote{Page 217.} The author is here speaking of the miracles of the Old and New
Testaments, which he believes; and, like a religious man, he feels,
contrariwise to Hume, that it is not ``convenient,'' but dangerous,
to allow of an antecedent test, which, for what he knows, and before
he is aware, may be applied in disproof of one or other instance of
those gracious manifestations. But it is far otherwise when he comes
to speak of Ecclesiastical Miracles, which he begins with disbelieving
without much regard to their evidence, and is engaged, not in
examining or confuting, but in burdening with some test or criterion
which may avail, in Hume's words, ``to silence bigotry and
superstition, and to free us from their impertinent solicitations.''
He acts towards the miracles of the Church, as Hume towards the
miracles of Scripture.

14.
And surely with less reason than Hume, from a consideration already
suggested; because, in being a believer in the miracles of Scripture,
he deprives himself of that strong antecedent ground against all
miracles whatever, both Scriptural and Ecclesiastical, on which Hume
took his stand. Allowing, as he is obliged to allow, that the
ecclesiastical miracles are possible, because the Scripture
miracles are true, he rejects ecclesiastical miracles as not
subserving the object which he arbitrarily assigns for miracles under
the Gospel, while he protects the miracles of Scripture by the
cautious proviso, that ``Infinite Wisdom may see ends'' for an
interposition, ``the importance of which may lie hid from our shallow
comprehension.'' Yet it is a fairer argument against miraculous
agency in a particular instance, before it is known in any case to
have been employed, that its object is apparently unimportant, than
after such agency has once been manifested. What has been introduced
for greater ends may, when once introduced, be made subservient to
secondary ones. Parallel cases are of daily occurrence in matters of
this world; and if it is allowable, as it is generally understood to
be, to argue from final causes in behalf of the being of a God—that
is, to apply the analogy of a human framer and work to the relation
subsisting between the physical world and a Creator—surely it is
allowable also to illustrate the course of Divine Providence and
Governance by the methods and procedures of human agents. Now, nothing
is more common in scientific and social arrangements than that works
begun for one purpose should, in the course of operation, be made
subservient, as a matter of course, to lesser ones. A mechanical
contrivance or a political organization is continued for secondary
objects, when the primary has been attained; and thus miracles
begun either for Warburton's object or Douglas's may be continued
for others, ``the importance of which,'' in the language of the
latter, ``may lie hid from our shallow comprehension.''

15.
Hume judges of professedly Divine acts by \emph{experience}; Bishops
Warburton and Douglas by the \emph{probable objects} which a Divine
Agent must pursue. Both parties draw extravagant conclusions, and that
unphilosophically; but surely we know much less of the designs and
purposes of Divine Providence, on which Warburton and Douglas insist,
than we know of that physical course of things on which Hume takes his
stand. Facts actually come before us; the All-wise Mind is hidden from
us. We have a right to form anticipations about facts; we may not,
except very reverently and humbly, attempt to trace, and we dare not
prescribe, the rules on which Providence conducts the government of
the world. The Apostle warns us, ``Who hath known the mind of the
Lord? and who hath been His counsellor?'' And surely, a fresh or
additional object in the course of Providence presents a less
startling difficulty to the mind than an interposition in the laws of
nature. If we conquer our indisposition towards the news of such an
interposition by reflecting on the Sovereignty of the Creator, let us
not be religious by halves, let us submit our imaginations to the full
idea of that inscrutable Sovereignty, nor presume to confine it
within bounds narrower than are prescribed by His own attributes.

16.
This, then, is the proper answer to the objection urged against the
post-apostolic miracles, on the ground that the first occurrence of
miracles does in itself discredit their recurrence, and that the
miracles subsequent to those of Scripture differ, in fact, from the
Scripture miracles in their objects and circumstances. The ordinary
Providence of God is conducted upon a \emph{system}; and as even the
act of creation is now contemplated by some philosophers as possibly
subject to law, so it is more probable than not that there is also a
law of supernatural manifestations. And thus the occurrence of
miracles is rather a presumption for than against their recurrence;
such events being not isolated acts, but the indications of the
presence of an agency. And again, since every system consists of parts
varying in importance and value, so also as regards a dispensation of
miracles, ``God hath set every one of them in the body as it hath
pleased Him;'' and even ``those members which seem to be more
feeble'' and less ``comely'' are ``necessary,'' and are sustained
by their fellowship with the more honourable.

17.
It may be added that Scripture, as in Mark xvi. 17, 18, certainly does
give a \emph{primâ facie} countenance to the idea that miracles are a
privilege accorded to true believers, and that where is faith, there
will be the manifested signs of its invisible Author. Hence it
was the opinion of Grotius \footnote{On Mark xvi. 17, Grotius avows his belief in the continuance of a
miraculous agency down to this day. He illustrates that text from St.
Justin, St. Irenæus, Origen, Tertullian, Minucius Felix, and
Lactantius, as regards the power of exorcism, and refers to the acts
of Victor of Cilicia in the Martyrology of Ado, and to the history of
Sabinus, Bishop of Canusium, in Greg. Turon., for instances of
miraculous protection against poison. As to missions, he asserts that
the presence of miraculous agency is even a test whether the doctrine
preached is Christ's. ``Si quis etiam nunc gentibus Christi ignaris,
(illis enim proprie miracula inserviunt, 1 Cor. xiv. 22), ita ut ipse
annunciari voluit, annunciat, promissionis vim duraturam arbitror.
Sunt enim [\emph{ametamel\={e}ta tou theou d\={o}ra}]. Sed nos,
cujus rei culpa est in nostra ignaviâ aut diffidentiâ, id solemus in
Deum rejicere.'' Elsewhere he professes his belief in the miracle
wrought upon the Confessors under Hunneric, who spoke after their
tongues were cut out; and in the ordeals of hot iron in the middle
ages (De Verit. i. 17); and in the miracles wrought at the tombs of the
Martyrs. Ibid. iii. 7, fin. Vide also De Antichr. p. 502, col. 2.},
who is here quoted from his connection with English Theology, and of
Barrow, Dodwell, and others, that miracles are at least to be expected
as attendants on the labours of Missionaries. Now this Scripture
intimation, whether fainter or stronger, does, as far as it goes, add
to the presumption in favour of the miracles of ecclesiastical
history, by authoritatively assigning them a place in the scheme of
Christianity. But this subject, as well as others touched upon in this
Chapter, will more distinctly come into review in those which follow.

\xchapter{Chapter 3. On the Internal Character of the Ecclesiastical Miracles}

18.
THE miracles wrought in times
subsequent to the Apostles are of a very different character, viewed
as a whole, from those of Scripture viewed as a whole; so much so,
that some writers have not scrupled to say that, if they really took
place, they must be considered as forming another dispensation \footnote{Vide Middleton's Inquiry, p. 24. et alib.
Campbell on Miracles, p. 121.}; and at least they are in some sense supplementary to the
Apostolic. This will be evident both on a survey of some of them, and
by referring to the language used by the Fathers of the Church
concerning them.

1.

19.
The Scripture miracles are for the most part evidence of a Divine
Revelation, and that for the sake of those who have not yet been
instructed in it, and in order to the instruction of multitudes: but
the miracles which follow have sometimes no discoverable or
direct object, or but a slight object; they happen for the sake of
individuals, and of those who are already Christians, or for purposes
already effected, as far as we can judge, by the miracles of
Scripture. The Scripture miracles are wrought by persons consciously
exercising under Divine guidance a power committed to them for
definite ends, professing to be immediate messengers from heaven, and
to be evidencing their mission by their miracles: whereas
Ecclesiastical miracles are not so much wrought as displayed, being
effected by Divine Power without any visible media of operation at
all, or by inanimate or material media, as relics and shrines, or by
instruments who did not know at the time what they were effecting, or,
if they were hoping and praying for such supernatural blessing, at
least did not know when they were to be used as instruments, when not.
The miracles of Scripture are, as a whole, grave, simple, and
majestic: those of Ecclesiastical History often partake of what may
not unfitly be called a romantic character, and of that wildness and
inequality which enters into the notion of romance. The miracles of
Scripture are undeniably beyond nature: those of Ecclesiastical
History are often scarcely more than extraordinary accidents or
coincidences, or events which seem to betray exaggerations or errors
in the statement. The miracles of Scripture are definite and whole transactions, drawn out and carried through from first to last,
with beginning and ending, clear, complete, and compact in the
narrative, separated from extraneous matter, and consigned to
authentic statements: whereas the Ecclesiastical, for the most part,
are not contained in any authoritative form or original document; at
best they need to be extracted from merely historical works, and often
are only floating rumours, popular traditions, vague, various,
inconsistent in detail, tales which only \emph{happen} to have
survived, or which in the course of years obtained a permanent place
in local usages or in particular rites or on certain spots, recorded
at a distance from the time and country when and where they profess to
have occurred, and brought into shape only by the juxta-position and
comparison of distinct informations. Moreover, in Ecclesiastical
History true and false miracles are mixed: whereas in Scripture
inspiration has selected the true to the exclusion of all others.

2.

20.
The peculiarity of these miracles, as far as their nature and
character are concerned, which is the subject immediately before us at
present, will be best understood by an enumeration of some of them,
taken almost at random, in the order in which they occur in the
authors who report them.

The Life of St. Gregory of Neocæsarea in Pontus
(A.D. 250), is written by
his namesake of Nyssa, who lived about 120 years after him, and who,
being a native and inhabitant of the same country, wrote from the
traditions extant in it. He is called Thaumaturgus, from the
miraculous gift ascribed to him, and it is not unimportant to observe
that he was the original Apostle of the heathen among whom he was
placed. He found at first but seventeen Christians in his diocese, and
he was the instrument of converting the whole population both of town
and country. St. Basil (A.D.
370), whose see was in the neighbourhood, states this circumstance,
and adds, ``Great is the admiration which still attends on him among
the people of that country, and his memory resides in the Churches new
and ever fresh, impaired by no length of time. And therefore no usage,
no word, no mystic rite of any sort, have they added to the Church
beyond those which he left. Hence many of their observances seem
imperfect, on account of the ancient manner in which they are
conducted. For his successors in the government of the Churches did
not endure the introduction of anything which has been brought into
use since his date.'' \footnote{De Spir. S. 74.}

21.
St. Gregory of Nyssa tells us that, when he was first coming into his
heathen and idolatrous diocese, being overtaken by night and rain, he
was obliged, with his companions, to seek refuge in a temple
which was famous for its oracles. On entering he invoked the name of
Christ, and made the sign of the cross, and continued till morning in
prayer and psalmody, as was his custom. He then went forward, but was
pursued by the Priest of the temple, who threatened to bring him
before the magistrates, as having driven the evil spirit from the
building, who was unable to return. Gregory tore off a small portion
of the book he had with him, and wrote on it the words, ``Satan, enter.''
The Priest, on returning, finding that the permission took effect as
well as the former prohibition, came to him a second time, and asked
to be instructed about that God who had such power over the demons.
Gregory unfolded to him the mystery of the Incarnation; and the pagan,
stumbling at it, asked to see a miracle. Nyssen, who has spoken all
along as relating the popular account, now says that he has to relate
what is ``of all the most incredible.'' A stone of great size lay before
them; the Priest asked that it might be made to move by Gregory's
faith, and Gregory wrought the miracle. This was followed by the
Priest's conversion, but not as an isolated event; for, on his entry
into the city, all the inhabitants went out to meet him, and enough
were converted on the first day by his preaching to form a church. In
no long time he was in a condition to call upon his flock to build a
place of worship, the first public Christian edifice on record;
which remained to Nyssen's time, in spite of the serious earthquakes
which had visited the city.

22.
St. Gregory's fame extended into the neighbouring districts, and
secular causes were brought for his determination. Among those who
came to him were two brothers, who had come into their father's large
property, and litigated about the possession of a lake which formed
part of it. When his efforts to accommodate their difference failed,
and the disputants, being strong in adherents and dependants, were
even proceeding to decide the matter by force of arms, Gregory the day
before the engagement betook himself to the lake, and passed the night
there in prayer. The lake was dried up, and in Nyssen's time its bed
was covered with woods, pasture and corn land, and dwellings. Another
miracle is attributed to him of a similar character. A large and
violent stream, which was fed by the mountains of Armenia, from time
to time broke through the mounds which were erected along its course
in the flat country, and flooded the whole plain. The inhabitants, who
were heathen, having heard the fame of Gregory's miracles, made
application to him for relief. He journeyed on foot to the place, and
stationed himself at the very opening which the stream had made in the
mound. Then invoking Christ, he took his staff, and fixed it in the
mud; and then returned home. The staff budded, grew, and became a
tree, and the stream never passed it henceforth: since it was
planted by Gregory at the very time when the mound had burst, and was
appealed to by the inhabitants \footnote{[\emph{Mechri tou nun tois hepich\={o}riois
theama ginetai to phuton kai di\={e}g\={e}ma … onoma de mechri
tou nun hesti t\={o} dendr\={o}i h\={e} bact\={e}ria, mn\={e}mosunon
t\={e}s Gr\={e}goriou charitos kai duname\={o}s, tois ench\={o}riois
hen panti t\={o} chron\={o} s\={o}zomenon}]. T. ii. pp.
991, 992.}, who were converted in
consequence, and was still living in Nyssen's time, it became a sort
of monument of the miracle. On one of his journeys two Jews attempted
to deceive him; the one lay down as if dead, and the other pretended
to lament him, and asked alms of Gregory for a shroud. Gregory threw
his garment upon him, and walked on. His companion called on him to
rise, but found him really dead. One day when he was preaching, a boy
cried out that some one else was standing by Gregory, and speaking
instead of him; at the end of the discourse Gregory observed to the
bystanders that the boy was possessed, and taking off the covering
which was on his own shoulders, breathed on it, and cast it on the
youth, who forthwith showed all the usual symptoms of demoniacs. He
then put his hand on him, and his agitation ceased, and his delusion
with it.

23.
Now, concerning these and similar accounts, it is obvious to remark,
on the one hand, that the alleged miracles were wrought in order to
the conversion of idolaters; on the other hand, when we read of stones
changing their places, rivers restrained, and lakes dried up,
and, at the same time, of buildings remaining in spite of earthquakes,
we are reminded, as in the case of the Scripture miracle upon the
cities of the plain, that a volcanic country is in question, in which
such phenomena are to a great extent coincident with the course of
nature. It may be added, that the biographer not only is frequent in
the phrases ``it is said,'' ``it is still reported,'' but he assigns
as a reason for not relating more of St. Gregory's miracles, that he
may be taxing the belief of his readers more than is fitting, and he
throughout writes in a tone of apology as well as of panegyric.

24.
Next, let us turn to St. Athanasius's biographical notice of St.
Antony, who began the solitary life A.D.
270. Athanasius knew him personally, and writes whatever he was able
to learn from himself; for ``I followed him,'' he says, ``no small time,
and poured water upon his hands;'' and he adds, that ``everywhere he has
had an anxious regard to truth.'' The following are some of the
supernatural or extraordinary portions of his narrative. He relates
that the enemy of souls appeared to Antony, first like a woman, then
like a black child, when he confessed himself to be the spirit of
lewdness, and to have been vanquished by the young hermit. Afterwards,
when he was passing the night in the tombs, he was attacked by evil
spirits, and so severely stricken that he lay speechless till a friend
found him next day \footnote{Eusebius relates of one Natalis, a Confessor
of the end of the second century, that he fell into the heresy of
Theodotus, a sort of Unitarianism, and was warned by our Lord in
visions. On neglecting these, he was severely scourged by angels all
through the night. Hist. v. 28. Vide Hieron. adv. Rufin. p. 414.}. When he was on his first journey
into the desert, a large plate of silver lay in his way; he
soliloquized thus, ``Whence this in the desert? This is no beaten path,
no track of travellers; it is too large to be dropped without being
missed; or if dropped, it would have been sought after and found, for
there is no one else to take it. This is a snare of the devil; thou
shalt not, O devil, hinder thus my earnest purpose; unto perdition be
it with thee!'' As he spoke, the plate vanished. He exhorted his
friends not to fear the evil spirits: ``They conjure up phantoms to
terrify cowards; but sign yourselves with the cross, and go forth in
confidence.'' ``Once there appeared to me,'' he says, on another
occasion, ``a spirit very tall, with a great show, and presumed to say,
'I am the Power of God,' and 'I am Providence; what favour shall I do
thee?' But I the rather spit upon him, naming the Christ, and essayed
to strike him, and I think I did; and straightway this great personage
vanished with all his spirits at Christ's name. Once he came, the
crafty one, when I was fasting, and as a Monk, with the appearance of
loaves, and bade me eat: 'Eat, and have over thy many pains;
thou too art a man, and art like to be sick;' I, perceiving his craft,
rose up to pray. He could not bear it, but vanished through the door,
like smoke. Listen to another thing, and that securely and fearlessly;
and trust me, for I lie not. One time some one knocked at my door in
the monastery; I went out, and saw a person tall and high. 'Who art
thou?' says I; he answers, 'I am Satan.' Then I asked, 'Why art thou
here?' He says, 'Why do the Monks, and all other Christians, so
unjustly blame me? Why do they curse me hourly?' 'Why troublest thou
them?' I rejoin. He, 'I trouble them not; they harass themselves; I
have become weak. I have no place left, no weapon, no city. Christians
are now everywhere; at last even the desert is filled with Monks. Let
them attend to themselves, and not curse me, when they should not.'
Then I said to him, admiring the grace of the Lord, 'A true word
against thy will, who art ever a liar, and never speakest truth; for
Christ hath come and made thee weak, and overthrown thee and stripped
thee.' At the Saviour's name he vanished; it burned him, and he could
not bear it.''

25.
Once, when travelling to some brethren across the desert, water failed
them; they sat down in despair, and let the camel wander. Antony knelt
down and spread out his hands in prayer, when a spring of water burst
from the place where he was praying. A person came to him, who
was afflicted with madness or epilepsy, and begged his prayers; he
prayed for him, and then said, ``Go, and be healed.'' The man refusing
to go, Antony said, ``If thou remainest here, thou canst not be healed;
but go to Egypt, and thy cure shall be wrought in thee.'' He believed,
went, and was cured as soon as he got sight of Egypt. At another time
he was made aware that two brothers were overtaken in the desert by
want of water; that one was dead, and the other dying; he sent two
Monks, who buried the one and restored the other. Once, on entering a
vessel, he complained of a most loathsome stench; the boatmen said
that there was fish in it, but without satisfying Antony, when
suddenly a cry was heard from a youth on board, who was possessed by a
spirit. Antony used the name of our Lord, and the sick person was
restored. St. Athanasius relates a similar instance of Antony's power,
which took place in his presence. When the old man left Alexandria,
whither he had gone to assist the Church against the Arians,
Athanasius accompanied him as far as the gate. A woman cried after
him, ``Stop, thou man of God; my daughter is miserably troubled by a
spirit.'' Athanasius besought him too, and he turned round. The girl,
in a fit, lay on the ground; but on Antony praying, and naming the
name of Christ, she rose restored. It should be observed, that
Alexandria was at this time still in a great measure a heathen
city. Athanasius says that, while Antony was there, as many became
Christians in a few days as were commonly converted in the course of
the year. This fact is important, not only as showing us the purpose
which his miracles answered, but as informing us by implication that
pretensions such as Antony's were not of every day's occurrence then,
but arrested attention and curiosity at the time.

26.
We have a similar proof of the comparative rareness of such miraculous
power in St. Jerome's Life of Hilarion. When the latter visited
Sicily, one of his disciples, who was seeking him, heard in Greece
from a Jew that ``a Prophet of the Christians had appeared in Sicily,
and was doing so many miracles and signs, that men thought him one of
the old Saints.'' Hilarion was the first solitary in Palestine, and a
disciple of St. Antony. St. Jerome enumerates various miracles which
were wrought by him, such as his giving sight to a woman who had been
ten years blind, restoring a paralytic, procuring rain by his prayers,
healing the bites of serpents with consecrated oil, curing a dropsy,
curbing the violence of the sea upon a shore, exorcising the
possessed, and among these a camel which had killed many persons in
its fury. When he was solemnly buried, ten months after his death, his
Monk's dress was quite whole upon him, and his body was entire as if
he had been alive, and sent forth a most exquisite fragrance.

27.
Sulpicius gives us an account of his master St. Martin's miracles,
which encountered much incredulity when he first published it. ``I am
shocked to say what I lately heard,'' says his friend to him in his
Dialogues; ``but an unhappy man has asserted that you tell many lies in
your book.'' As St. Martin was the Apostle of Gaul, the purpose
effected by his miracles is equally clear and sufficient, as in the
instance of Thaumaturgus; and they are even more extraordinary and
startling than his. Sulpicius in his Dialogues solemnly appeals to our
Lord that he has stated nothing but what he saw himself, or knew, if
not on St. Martin's own word, at least on sure testimony. He also
appeals to living witnesses. The following are instances taken from
the first of his two works.

28.
Before Martin was a Bishop, while he was near St. Hilary at Poictiers,
a certain Catechumen, who lived in his monastery, died of a fever, in
Martin's absence, without baptism. On his return, the Saint went by
himself into the cell where the body lay, threw himself upon it,
prayed, and then raising himself with his eyes fixed on it, patiently
waited his restoration, which took place before the end of two hours.
The man, thus miraculously brought to life, lived many years, and was
known to Sulpicius, though not till after the miracle. At the same
period of his life he also restored a servant in a family, who had
hung himself, and in the same way. Near Tours, which was his
See, a certain spot was commonly considered to be the tomb of Martyrs,
and former Bishops had placed an altar there. No name or time was
known, and Martin found reason to suspect that the tradition was
unfounded. For a while he remained undecided, as being afraid of
encouraging either superstition or irreverence; at length he went to
the tomb, and prayed to Christ to be told who was buried there, and
what his character. On this a dismal shade appeared, who, on being
commanded to speak, confessed that he was a robber who had been
executed for his crimes, and was in punishment. Martin's attendants
heard the voice, but saw nothing. Once, when he was on a journey, he
saw at a distance a heathen funeral procession, and mistook it for
some idolatrous ceremonial, the country people of Gaul being in the
practice of carrying their gods about their fields. He made the sign
of the cross, and bade them stop and set down the body; this they were
constrained to do. When he discovered their real business, he suffered
them to proceed. At another time, on his giving orders for cutting
down a pine to which idolatrous honour was paid, a heathen said, ``If
thou hast confidence in thy God, let us hew the tree, and do thou
receive it as it falls; if thy Lord is with thee, thou wilt escape
harm.'' Martin accepted the condition, and when the tree was falling
upon him, made the sign of the cross; the tree reeled round and
fell on the other side. This miracle converted the vast multitude who
were spectators of it \footnote{Sulpicius adds, ``Et vere ante Martinum
pauci admodum, imo pæne nulli, in illis regionibus Christi nomen
receperant; quod adeò virtutibus illius exemploque convaluit, ut jam
ibi nullus locus sit, qui non aut ecclesiis frequentissimis aut
monasteriis sit repletus.'' V. Mart. 10.}. About the same time, when he had set
on fire a heathen temple, the flames spread to a house which joined
it. Martin mounted on to the roof of the building that was in peril,
and by his presence warned off the fire, and obliged it to confine
itself to the work intended for it. At Paris a leper was stationed at
the gate of the city; Martin went up and kissed and blessed him, and
his leprosy disappeared.

29.
St. Augustine, again, enumerates at the end of his \emph{De civitate Dei},
certain miracles which he himself had witnessed, or had on good
authority, such as these. An actor of the town of Curulis was cured of
the paralysis in the act of baptism; this Augustine knew, on what he
considered the best authority. A person known to Augustine, who had
received earth from the Holy Sepulchre, asked him and another Bishop
to place it in some oratory for the profit of worshippers. They did
so, and a country youth, who was paralytic, hearing of it, asked to be
carried to the spot. After praying there, he found himself recovered,
and walked home. By the relics of St. Stephen one man was cured
of a fistula, another of the stone, another of the gout; a child who
had been crushed to death by a wheel was restored to life; also a nun,
by means of a garment which had been taken to his shrine and thrown
over her corpse; and another female by the same means; and another by
the oil used at the shrine; and a dead infant who was brought to it.
In less than two years even the formal statements given in of miracles
wrought at St. Stephen's shrine at Hippo were almost seventy.

3.

30.
These miracles are recorded by writers of the fourth century, though
they belong, in one case wholly, in another partially, to the history
of the third. When we turn to earlier writers, we find similar
assertions of the presence of a miraculous agency in the Church, and
its manifestations have the same general character. Exorcisms, cures,
visions, are the chief miracles of the fourth century; and they are
equally so of the second and third, so that the former have a natural
claim to be considered the continuation of the latter. But there are
these very important differences between the two,—that the accounts
in the fourth century are much more in detail than those of the second
and third, which are commonly vague and general; and next, that in the
second and third those kinds of miraculous operations which are the
most decisive proofs of a supernatural presence are but
sparingly or scarcely mentioned.

31.
Middleton's enumeration of these primitive miracles, which on the
whole may be considered to be correct, is as follows: ``The power of
raising the dead, of healing the sick, of casting out devils, of
prophesying, of seeing visions, of discovering the secrets of men, of
expounding the Scriptures, of speaking with tongues.'' \footnote{Page 72.} Of
these the only two which are in their nature distinctly miraculous are
the first and last; and for both of these we depend mainly on the
testimony of St. Irenæus, who lived immediately after the Apostolical
Fathers, that is, close upon the period when even modern writers are
disposed to allow that miracles were wrought in the Church. Douglas
observes, ``If we except the testimonies of Papias and Irenæus, who
speak of raising the dead, … I can find no instances of miracles
mentioned by the Fathers before the fourth century, as what were
performed by Christians in their times, but the cures of diseases,
particularly the cures of demoniacs, by exorcising them; which last
indeed seems to be their favourite standing miracle, and the only one
which I find (after having turned over their writings carefully, and
with a view to this point,) they challenged their adversaries to come
and see them perform.'' \footnote{Page 232.}

32.
It must be observed, however, that though certain occurrences
are in their character more miraculous than others, yet that a miracle
of \emph{degree} may, in the particular case, be quite as clearly
beyond the ordinary course of nature. Imagination can cure the sick in
certain cases, in certain cases it cannot; and we shall have a very
imperfect view of the alleged miracles of the second and third
centuries, if, instead of patiently contemplating the instances
recorded, in their circumstances and details, we content ourselves
with their abstract character, and suffer a definition to stand in
place of examination. Thus if we take St. Cyprian's description of the
demoniacs, in which he is far from solitary \footnote{For ancient testimonies to the power of
exorcism, vid. Middleton, pp. 80-90. Douglas's Criterion, p. 232, Note
1. Farmer, On Miracles, pp. 241, 242. Whitby's Preface to Epp. § 10.}, we shall find
that while it is quite open to accuse him and others of misstatement,
we cannot accept his description as it stands, without acknowledging
that the conflict between the powers of heaven and the evil spirit was
then visibly proceeding as in the time of Christ and His Apostles.
``O would you listen to them,'' he says to the heathen Demetrian, ``and
see them, when they are adjured and tormented by us with spiritual
lashes, hurled with words of torture out of bodies they have
possessed, when shrieking and groaning at a human voice, and beneath a
power divine laid under lash and stripe, they confess the
judgment to come. You will find that we are entreated of them whom you
entreat, feared by them whom you fear, and whom you adore. Surely
thus, at least, will you be brought to confusion in these your errors,
when you behold and hear your gods at once, upon our questioning,
betraying what they are, and unable, even in your presence, to conceal
their tricks and deceptions.'' \footnote{Treat. viii. 8. Oxford tr.} Again, ``You may see them
by our voice, and through the operation of the unseen Majesty, lashed
with stripes, and scorched with fire; stretched out under the increase
of their multiplying penalty, shrieking, groaning, entreating,
confessing from whence they came, and when they depart, even in the
hearing of their own worshippers; and either leaping out suddenly, or
gradually vanishing, as faith in the sufferer aids, or grace in the
curer conspires.'' \footnote{Treat. ii. 4. Oxford tr.} Passages equally strong might be cited
from writers of the same period.

33.
And there are other occurrences of a distinctly miraculous character
in the earlier centuries, which come under none of Middleton's or
Douglas's classes, but which ought not to be overlooked. For instance,
a fragrance issued from St. Polycarp when burning at the stake, and on
his being pierced with a sword a dove flew out. Narcissus, Bishop of
Jerusalem, about the end of the second century, when oil failed for
the lamps on the vigil of Easter, sent persons to draw water
instead; which, on his praying over it, was changed into oil.
Eusebius, who relates this miracle, says that small quantities of the
oil were preserved even to his time. St. Cyprian speaks of a person
who had lapsed in persecution, attempting to communicate; when on
opening the arca, or receptacle in which the consecrated Bread was
reserved, fire burst out from it and prevented her. Another, on
attending at church with the same purpose, found that he had received
from the priest nothing but a cinder.

34.
Lastly, in this review of the miracles belonging to the early Church,
it will be right to include certain isolated ones which have an
historical character, and are accordingly more celebrated than the
rest. Such is the miracle of the thundering Legion, that is, of the
rain accorded to the prayers of Christian soldiers in the army of
Marcus Antoninus, when they were perishing by thirst; the appearance
of a Cross in the sky to Constantine's army, with the inscription, ``In
hoc signo vinces;'' the sudden death of Arius, close upon his proposed
re-admission into the Church, at the prayers of Alexander of
Constantinople; the discovery of the Cross, the multiplication of its
wood, and the miracles wrought by it; the fire bursting forth from the
foundations of the Jewish temple, which hindered its rebuilding; the
restoration of the blind man on the discovery of the relics of St.
Gervasius and St. Protasius; and the power of speech granted to
the African confessors who had lost their tongues in the Vandal
persecution \footnote{For other ancient testimonies to the
ecclesiastical miracles, vid. Dodwell. Dissert. in Irenæum. ii.
41-60. Middleton's Inquiry, pp. 2-19. Brook's Defens. Miracl. Eccl.
pp. 16-22. Mr. Isaac Taylor's Anc. Christ. part 7.}. These and other such shall be considered
separately, before I conclude.

35.
Imperfect as is this survey of the miracles ascribed to the ages later
than the Apostolic, it is quite sufficient for the purpose for which
it has been made; viz., to show that those miracles are on the whole
very different in their character and attendant circumstances from the
Gospel miracles, which certainly are very far from preparing us for
them, or rather at first sight indispose us for their reception \footnote{On the difference between the miracles of
Scripture and of Ecclesiastical History, vid. Douglas's Crit. pp.
221-237. Paley's Evidences, Part i. Prop. 2. Middl. pp. 21-26, 91-96,
etc. Bishop Blomfield's Sermons, note on p. 82. Dodwell attempts to
draw a line between the Ante-Nicene and the later miracles, in favour
of the former (Dissert. in Iren. ii. 62-66), as regards testimony,
nature, instrument, and object.}.

4.

36.
And in the next place this important circumstance must be considered,
which is as clear as it is decisive, that the Fathers speak of
miracles as having in one sense ceased with the Apostolic
period;—that is to say, whereas they sometimes speak of miracles as
existing in their own times, still they say also that \emph{Apostolic}
miracles, or miracles \emph{like} the Apostles', whether in their
object, cogency, impressiveness, or character, were no longer of
occurrence in the Church; an interpretation which they themselves in
some passages give to their own testimony. ``Argue not,'' says St.
Chrysostom, ``because miracles do not happen now, that they did not
happen then ... In those times they were profitable, and now they are
not.'' He proceeds to say that, in spite of this difference, the mode
of conviction was substantially the same. ``We persuade not by
philosophical reasonings, but from Divine Scripture, and we recommend
what we say by the miracles then done. And then, too, they persuaded
not by miracles only, but by discussion.'' And presently he adds, ``The
more evident and constraining are the things which happen, the less
room there is for faith.'' \footnote{Hom. in 1 Cor. vi. 2 and 3.} And again in another passage, ``Why
are there not those now who raise the dead and perform cures? I will
not say why not; rather, why are there not those now who despise the
present life? why serve we God for hire? When, however, nature was
weak, when faith had to be planted, then there were many such; but now
He wills not that we should hang on these miracles, but be ready for
death.'' \footnote{Hom. 8, in Col. § 5.}

37.
In like manner St. Augustine introduces his catalogue of
contemporary miracles, which has been partly given above, by stating
and allowing the objection that miracles were not then as they had
been. ``Why, say they, do not these miracles take place now, which, as
you preach to us, took place once? I might answer that they were
necessary before the world believed, that it might believe.'' \footnote{De Civ. Dei, xxii. 8, § 1.}
He then goes on to say that miracles were wrought in his time, only
they were not so public and well-attested as the miracles of the
Gospel.

38.
St. Ambrose, on the discovery of the bodies of the two Martyrs, uses
the language of surprise; which is quite in accordance with the
feelings which the miracles of Antony and Hilarion seem to have roused
in Alexandria and in Sicily. ``You know, you yourselves saw, that many
were cleansed from evil spirits; very many, on touching with their
hands the garment of the Saints, were delivered from the infirmities
which oppressed them. The \emph{miracles of the old time} are come \emph{again},
when by the advent of the Lord Jesus a fuller grace was shed upon the
earth.'' Under a similar feeling \footnote{Ep. i. 22, § 9. The same feeling of
reverence for times past must be taken partly to account for the
expressions [\emph{ichn\={e}}] and [\emph{hypoleleiptai}] in Origen,
Eusebius, etc., below note a \textbf{[}Note
21\textbf{]}.} he speaks of the two
corpses, which happened to be of large size, as ``miræ magnitudinis,
ut prisca ætas ferebat.'' \footnote{Ibid. § 2.}

39.
And Isidore of Pelusium, after observing that in the Apostles holiness
of life and power of miracles went together, adds, ``Now, too, if the
life of teachers rivalled the Apostolic bearing, perhaps miracles
would take place; though if they did not, such life would suffice for
the enlightening of those who beheld it.'' \footnote{Ep. iv. 80.}

40.
The doctrine, thus witnessed by the great writers of the end of the
fourth century, is supported by as clear a testimony two centuries
before and two centuries after. Pope Gregory, at the end of the sixth,
in commenting on the text, ``And these signs shall follow those that
believe,'' says, ``Is it so, my brethren, that, because ye do not these
signs, ye do not believe? On the contrary, they were necessary in the
beginning of the Church: for, that faith might grow, it required
miracles to cherish it withal; just as when we plant shrubs, we water
them till they seem to thrive in the ground, and as soon as they are
well rooted we cease our irrigation. This is what Paul teaches,
'Tongues are a sign, not for those who believe, but for those who
believe not;' and there is something yet to be said of these signs and
powers of a more recondite nature. For Holy Church doth spiritually
every day, what she then did through the Apostles, corporally. For
when the Priests by the grace of exorcism lay hands on believers, and
forbid evil spirits to inhabit their minds, what do they but
cast out devils? And any believers soever who henceforth abandon the
secular words of the old life, and utter holy mysteries, and rehearse,
as best they can, the praise and power of their Maker, what do they
but speak with new tongues? Moreover, while by their good exhortations
they remove evil from the hearts of others, they are taking up
serpents, etc.; ... which miracles are the greater, because they are
the more spiritual: the greater, because they are the means of
raising, not bodies, but souls; these signs, then, dearest brethren,
by God's aid, ye do if ye will.'' \footnote{In Evang. ii. 29.} And St. Clement of
Alexandria, at the end of the second century: ``If it was imputed to
Abraham for righteousness on his believing, and we are the seed of
Abraham, we too must believe by hearing. For Israelites we are, who
are obedient, not through signs \footnote{Strom. ii. 6, p. 444. So Mr. Osburn, (Errors
Apost. Fathers, p. 12,) and I think rightly. The Bishop of Lincoln,
however, observes, ``I find \emph{only one} passage in the writings of Clement
which has any bearing on the question of the existence of miraculous
powers in the Church;'' and proceeds to refer to the Extracts from the
writings of Theodotus. Kaye's Clement, p. 468. The Bishop argues, in
his work upon Tertullian, that miracles had then ceased, from a
passage in the De Pudicitiâ, in which, after saying that the Apostles
had spiritual powers peculiar to themselves, Tertullian adds, ``Nam et
mortuos suscitaverunt, quod Deus solus; et debiles redintegraverunt,
quod nemo nisi Christus; immo et plagas inflixerunt, quod voluit
Christus.'' c. 21.}, but through hearing.'' \footnote{The following passages will be found to
testify to the same general fact, that the special miraculous powers
possessed by the Apostles did not continue in the Church after them.
Eusebius says that, according to St. Irenæus, instances of miraculous
powers, [\emph{en ekkl\={e}siais tisin hypoleleipto}], Hist. v. 7.
[\emph{ichn\={e}}], of the miracles still remain, says Origen
contra Cels. i. 2, fin. [\emph{ichn\={e}, kai tina ge meizona}]. \emph{Ibid}.
ii. 8. [\emph{ichn\={e} par holigois}]. \emph{Ibid}. vii. 8, fin.
In two of these passages the gift is connected with holiness of life,
a doctrine which Dodwell denies to have existed till the middle ages,
Dissert. in Iren. ii. 64, though he is aware of these passages. [\emph{oude
ichnos upoleleiptai}], Chrysost. de Sacerd. iv. 3, fin. [\emph{oi de
nun pantes omou}] cannot do as much as St. Paul's handkerchiefs. \emph{Ibid}.
iv. 6. He implies that the dead were not raised in his day. ``If God
saw that the raising of the dead would profit the living, He would not
have omitted it.'' De Lazar. iv. 3. ``Where is the Holy Spirit now? a
man may ask; for then it was appropriate to speak of Him, when
miracles took place, and the dead were raised, and all lepers were
cleansed; but now,'' etc. De Sanct. Pent. i. 3. He adds that now we
have the sanctifying gifts instead. So, again, ``The Apostles indeed
enjoyed the grace of God in abundance; but if \emph{we} were bid raise
the dead, or open the eyes of the blind, or cleanse lepers, or
straighten the lame, or cast out devils, and heal the like disorders,''
etc. Ad Demetr. i. 8. ``When the knowledge of Him as yet was not spread
abroad, then miracles used to take place; but now there is no need of
that teaching, the facts themselves proclaiming and manifesting the
Lord.'' In Psalm cxlii. 5. Vid. also Inscript. Act. ii. 3. Speaking of
the miracles in the wilderness, he says, ``In our case also, when we
came out of error, many wonders were displayed; but after that they
stopped, when religion was planted everywhere. And if subsequently
they happened [to the Jews], they were few and scattered, as when the
sun stood, etc., and this too has appeared in our case;'' and then he
goes on to mention the ``fiery eruption at the temple,'' etc., in Matth.
Hom. iv. i. And ibid. Hom. xxxii. 7, after mentioning the Apostolic
miracles of cleansing lepers, exorcising spirits, and raising the
dead, he says, ``This is the greatest proof of your nobleness and love,
to believe God without pledges; for this is one reason, among others,
why God ceased miracles ... Seek not miracles, then, but health of
soul.'' And then he contrasts with visible miracles the ``greater'' ones
of beneficence, self-command, etc., to the end of the Homily. And in
Joan. ``Now, too, there are those who seek and say, Why are there not
miracles now? If thou art faithful as behoveth, and love Christ as
thou shouldest, miracles thou needest not.'' Hom. xxiv. 1. Elsewhere,
after speaking of the gift of the Spirit dwelling in us, he adds, ``Not
that we may raise the dead, nor cleanse lepers, but that we show forth
the greatest miracle of all, charity,'' in Rom. Hom. viii. 7. After
quoting the text, ``We are changed into the same image from glory to
glory,'' he adds, ``This was shown more manifestly when the gifts of
miracles were in operation; but even now it is not difficult to
discern it when a man has believing eyes,'' etc., in 2 Cor. Hom. vii.
5. In like manner, St. Augustine, after mentioning the Apostolic
miracles, ``Sanati languidi, mundati leprosi, incessus claudis, cæcis
visus, surclis auditus est redditus,'' and the changing of water into
wine, the multiplication of the loaves, etc., continues, ``Cur, inquis,
ista modò non fiunt? quia non moverent, nisi mira essent: at si
solita essent, mira non essent.'' De Util. cred. 16. He adds, in his
Retractations, ``Hoc dixi, quia non tanta, nec omnia modo, non quia
nulla fiunt etiam modo.'' Again, ``Cum Ecclesia Catholica per totum
orbem diffusa atque fundata sit, nec miracula illa in nostra tempora
durare permissa sunt, ne animus semper visibilia quæreret,'' etc. De
Ver. Rel. 25. He adds, in his Retractations, ``Non sic accipiendum est
quod dixi, ut nunc in Christi nomine fieri miracula nulla credantur.
Nam ego ipse, quando istum ipsum librum scripsi, ad Mediolanensium
corpora Martyrum in eâdem civitate cæcum illuminatum fuisse jam
noveram,'' etc. Vid. also Pope Greg. Mor. xxvii. 18.}

5.

41.
What the distinctions are between the Apostolic and the later
miracles, which allow of the Fathers saying in a true sense that
miracles ceased with the first age, has in many ways appeared from
what has already come before us. For instance, it has appeared that
the Ecclesiastical Miracles were but locally known, or were done in
private; or were so like occurrences which are not miraculous as to
give rise to doubt and perplexity, at the time or afterwards, as
to their real character; or they were so unlike the Scripture
Miracles, so strange and startling in their nature and circumstances,
as to need support and sanction rather themselves than to supply it to
Christianity; or they were difficult from their drift, or their
instruments or agents, or the doctrine connected with them. In a word,
they are not primarily and directly evidence of Revelation, though
they may become so accidentally, or to certain persons, or in
the way of confirmation. That they are not the direct evidence of
revealed truth, is fully granted by St. Augustine in the following
striking passage from one of his works against the Donatists:—

42. ``Let
him prove that we must hold to the Church in Africa only, to the loss
of the nations, or again that we must restore and complete it in all
nations from Africa; and prove it, not by saying 'It is true, because
I say it,' or 'because my associate says it,' or 'my
associates,' or 'these our Bishops,' 'Clerks,' or 'people;' or 'it is
true because Donatus, or Pontius, or any one else, did these or those
marvellous acts,' or 'because men pray at the shrines of our dead
brethren, and are heard,' or 'because this or that happens there,' or
'because this brother of ours,' or 'that our sister,' 'saw such and
such a vision when he was awake,' or 'dreamed such and such a vision
when he was asleep.' Put away what are either the fictions of men who
lie, or the wonders of spirits who deceive. For either what is
reported is not true, or, if among heretics wonders happen, we have
still greater cause for caution, inasmuch as our Lord, after declaring
that certain deceivers were to be, who should work some miracles, and
deceive thereby, were it possible, even the elect, added an earnest
charge, in the words, 'Behold, I have told you before.' Whence also
the Apostle warns us that 'the Spirit speaketh expressly, in the
latter times some shall depart from the faith, giving heed to seducing
spirits, and doctrines of devils.' Moreover, if any one is heard who
prays at the shrines of heretics, what he receives, whether good or
bad, is consequent not upon the merit of the place, but upon the merit
of his own earnest desire. For 'the Spirit of the Lord,' as it is
written, 'hath filled the whole world,' and 'the ear of His zeal
heareth all things.' And many are heard by God in anger; of whom saith
the Apostle, 'God gave them up to the desires of their own
hearts.' And to many God in favour gives not what they wish, that He
may give what is profitable ... Read we not that some were heard by
the Lord God Himself in the high places of Judah, which high places
notwithstanding were so displeasing to Him, that the kings who
overthrew them not were blamed, and those who overthrew them were
praised? Thus it appears that the state of heart of the suppliant is
of more avail than the place of supplicating.

43. ``Concerning
deceitful visions, they should read what Scripture says, that 'Satan
himself transforms himself into an angel of light,' and that 'dreams
have deceived many.' And they should listen, too, to what the Pagans
relate, as regards their temples and gods, of wonders either in deed
or vision; and yet 'the gods of the heathen are but devils, but it is
the Lord that made the heavens.' Therefore many are heard and in many
ways, not only Catholic Christians, but Pagans and Jews and heretics,
involved in various errors and superstitions; but they are heard
either by seducing spirits, (who do nothing, however, but by God's
permission, judging in a sublime and ineffable way what is to be
bestowed upon each;) or by God Himself, whether for the punishment of
their wickedness, or for the solace of their misery, or as a warning
to them to seek eternal salvation. But salvation itself and life
eternal no one attains, unless he hath Christ the Head. Nor can any
one have Christ the Head, who is not in His body, which is the
Church; which, as the Head Himself, we are bound to discern in holy
canonical Scripture, not to seek in the various rumours of men, and
opinions, and acts, and sayings, and sights.

44.
``Let no one therefore object such facts who is prepared to answer
me; for I too am far from claiming credit for my position, that the
communion of Donatus is not the Church of Christ, on the ground that
certain bishops in it are convicted, in records ecclesiastical, and
municipal, and judicial, of burning the sacred books, ... or that the
Circumcelliones have committed so much evil, or that some of them cast
themselves down precipices, or throw themselves into the fire,
... or that at their sepulchres herds of strollers, men and
women, in a state of drunkenness and abandonment, bury themselves in
wine day and night, or pollute themselves with deeds of profligacy.
Let all this be considered merely as their chaff, without prejudice to
the Church, if they themselves are really holding to the Church. But
whether this be so, let them prove only from canonical Scripture; just
as we do not claim to be recognized as in the Church of Christ,
because the body to which we hold has been graced by Optatus of
Milevis or Ambrose of Milan, or other innumerable Bishops of our
communion, or because it is set forth in the Councils of our
colleagues, or because through the whole world in holy places, which
are frequented by our communion, so great marvels take place,
whether answers to prayer, or cures; so that the bodies of Martyrs,
which had lain concealed so many years, (as they may hear from many if
they do but ask,) were revealed to Ambrose, and in presence of those
bodies a man long blind and perfectly well known to the citizens of
Milan recovered his eyes and sight; or because one man has seen a
vision, or because another has been taken up in spirit, and heard
either that he should not join, or that he should leave, the party of
Donatus. All such things which happen in the Catholic Church, are to
be approved because they are in the Catholic Church; not she
manifested to be Catholic, because these things happen in her.'' \footnote{De Unit. Eccl. 49, 50.}

6.

45.
So far St. Augustine; it being granted, however, that the object of
Ecclesiastical Miracles is not, strictly speaking, that of evidencing
Christianity, still they may have other uses, known or unknown,
besides that of being the argumentative basis of revealed truth; and
therefore it does not at once destroy the credibility of such
miraculous narratives, vouched to us on good authority, that they have
no assignable object, or an object different from those which are
specified in Scripture, as was observed in the foregoing Chapter.

46.
Here we are immediately considering the \emph{internal character} of
the miracles later than the Apostolic period: and what real prejudice
ought to attach to them from the dissimilarity or even contrariety of
many of them to the Scripture Miracles will be best ascertained by
betaking ourselves to the argument from Analogy, and attempting to
measure these occurrences by such rules and suggestions as the works
of God, brought before us whether in the visible creation or in
Scripture, may be found to supply. And first of the natural world as
it meets our senses:—

47. ``All
the works of the Lord are exceeding good,'' says the son of Sirach; ``a
man need not to say, What is this? Wherefore is that? for He hath made
all things for their uses.'' Yet an exuberance and variety, a seeming
profusion and disorder, a neglect of severe exactness in the
prosecution of its objects, and of delicate adjustment in the details
of its system, are characteristics of the world both physical and
moral, and characteristics of Scripture also; but still the Wise Man
assures us, that the purposes of the Creator are not forgotten by Him,
or missed because they are hidden, or the work faulty because it is
subordinate or incomplete. All things are not equally good in
themselves, because they are diverse, yet everything is good in its
place. ``All the works of the Lord are good, and He will give every
needful thing in due season. So that a man cannot say, This is
worse than that; for in time they shall all be well approved.'' \footnote{Ecclus. xxxix. 16-35.} To persons who have not commonly the opportunity of witnessing for
themselves this great variety of the Divine works, there is something
very strange and startling,—it may even be said, unsettling—in the
first view of nature as it is. To take, for instance, the case of
animal nature, let us consider the effect produced upon the mind on
seeing for the first time the many tribes of the animal world, as we
find them brought together for the purposes of science or exhibition
in our own country. We are accustomed, indeed, to see wild beasts more
or less from our youth, or at least to read of them; but even with
this partial preparation, many persons will be moved in a very
singular way on going for the first time, or after some interval, to a
menagerie. They have been accustomed insensibly to identify the
wonder-working Hand of God with the specimens of its exercise which
they see about them; the forms of tame and domestic animals, which are
necessary for us, and which surround us, are familiar to them, and
they learn to take these as a sort of rule on which to frame their
ideas of the animated works of the Creator generally. When an eye thus
habituated to certain forms, colours, motions, and habits in the
inferior animals, is suddenly brought into the full assemblage of
those mysterious beings, with which it has pleased Almighty
Wisdom to people the earth, a sort of dizziness comes over it, from
the impossibility of our reducing all at once the multitude of new
ideas poured in upon us to the centre of view habitual to us; the mind
loses its balance, and it is not too much to say, that in some cases
it even falls into a sort of scepticism. Nature seems to be too
powerful and various, or at least too strange, to be the work of God,
according to that Image which our imbecility has set up within us for
the Infinite and Eternal, and as we have framed to ourselves our
contracted notions of His attributes and acts; and if we do not submit
ourselves in awe to His great mysteriousness, and chasten our hearts
and keep silence, we shall be in danger of losing our belief in His
presence and providence altogether.

48.
We have hitherto known enough of Him for our personal guidance, but we
have not understood that only thus much has been the extent of our
knowledge of Him. Religion we know to be a grave and solemn subject,
and some few vague ideas of greatness, sublimity, and majesty, have
constituted for us our whole image of Him whom the Seraphim adore. And
then we are suddenly brought into the vast family of His works, hardly
one of which is a specimen of those particular and human ideas with
which we have identified the Ineffable. First, the endless number of
wild animals, their independence of man, and uselessness to him;
then their exhaustless variety; then their strangeness in shape,
colour, size, motions, and countenance; not to enlarge on the still
more mysterious phenomena of their natural propensities and passions;
all these things throng upon us, and are in danger of overpowering us,
tempting us to view the Physical Cause of all as disconnected from the
Moral, and that, from the impression borne in upon us, that nothing we
see in this vast assemblage is \emph{religious} in our sense of the
word ``religious.'' We see full evidence there of an Author,—of power,
wisdom, goodness; but not of a Principle or Agent correlative to our
religious ideas. But without pushing this remark to an extreme point,
or dwelling on it further than our present purpose requires, let two
qualities of the works of nature be observed before leaving the
subject, which (whatever explanation is to be given of them, and
certainly some explanation is not beyond even our limited powers) are
at first sight very perplexing. One is that principle of \emph{deformity},
whether hideousness or mere homeliness, which exists in the animal
world; and the other (if the word may be used with due soberness) is
the \emph{ludicrous};—that is, judging of things, as we are here
judging of them, by their impression upon our minds.

49.
It is obvious to apply what has been said to the case of the miracles
of the Church, as compared with those in Scripture. Scripture is to us
a garden of Eden, and its creations are beautiful as well as ``very
good;'' but when we pass from the Apostolic to the following ages, it
is as if we left the choicest valleys of the earth, the quietest and
most harmonious scenery, and the most cultivated soil, for the
luxuriant wildernesses of Africa or Asia, the natural home or kingdom
of brute nature, uninfluenced by man. Or rather, it is a great
injustice to the times of the Church, to represent the contrast as so
vast a one; and Adam might much more justly have been startled at the
various forms of life which were brought before him to be named, than
we may rationally presume to decide that certain alleged miracles in
the Church are not really such, on the ground that they are unlike
those to which our eyes have been accustomed in Scripture. There is
far greater difference between the appearance of a horse or an eagle
and a monkey, or a lion and a mouse, as they meet our eye, than
between even the most august of the Divine manifestations in Scripture
and the meanest and most fanciful of those legends which we are
accustomed without further examination to cast aside. Such contrary
properties, or rather such impressions of them on our minds, may be
the necessary consequence of Divine Agency moving on a system, and not
by isolated acts; or the necessary consequence of its deigning to work
with or through the eccentricities, the weaknesses, nay, the
wilfulness, of the human mind. As, then, birds are different from
beasts, as tropical plants differ from the productions of the
north, as one scene is severely beautiful, and another rich or
romantic, as the excellence of colours is incommensurate with
excellence of form, as pleasures of sight have nothing in common with
pleasures of scent, except that they are pleasures; so also in the
case of those works and productions which are above or beside the
ordinary course of nature, in spite of their variety, ``to every thing
there is a season, and a time to every purpose under the heaven,'' and ``He
hath made every thing beautiful in His time.'' And, as one description
of miracles may be necessary for evidence, viz., such as are at once
majestic and undeniable, so for those other and manifold objects which
the economy of the Gospel kingdom may involve, a more hidden and
intricate path, a more complex exhibition, a more exuberant method, a
more versatile rule, may be essential; and it may be as shallow a
philosophy to reject them merely because they are not such as we
should have expected from God's hand, or as we find in Scripture, as
to judge of universal nature by the standard of our own home, or
again, with the ancient heretics, to refuse to admit that the Creator
of the physical world is also the Father of our Lord Jesus Christ.

50.
Nay, it may even be urged that the variety of nature is antecedently a
reason for expecting variety in a supernatural agency, if such be
introduced; or, again, (as has been already observed,) if such
agency is conducted on a system, it must even necessarily involve
diversity and inequality in its separate parts; and, granting it was
intended to continue after the Apostolic age, the want of uniformity
between the miracles first wrought and those which followed, as far as
it is found, might have been almost foretold without the gift of
prophecy in that age, or at least may be fully vindicated in
this,—nay, even the inferiority of the Ecclesiastical Miracles to
the Apostolic; for, if Divine Wisdom had determined, as is not
difficult to believe, that the wonderful works which illuminate the
history of the first days of the Church should be the best and
highest, what was left to subsequent times, by the very terms of the
proposition, but miracles which are but second best, which must
necessarily have belonged to some other and independent system if they
too were the best, and which admit of belonging to the same system for
the very reason that they are not the best?

7.

51.
So much, then, on the general correspondence between the works of
nature, on the one hand, and the Miracles of sacred history, whether
Biblical or Ecclesiastical, viewed as one whole, on the other. And
while the physical system bears such an analogy to the supernatural
system, viewed in its Biblical and Ecclesiastical portions
together, as forms a strong argument in defence of the supernatural,
it is, on the other hand, so far unlike the Biblical portion of that
supernatural, when that portion is taken by itself, as to protect the
portion not Biblical from objections drawn from any differences
observable between it and the portion which is Biblical. If it be true
that the Ecclesiastical Miracles are in some sense an innovation upon
the idea of the Divine Economy, as impressed upon us by the Miracles
of Scripture, it is at least equally true that the Scripture Miracles
also innovate upon the impressions which are made upon us by the order
and the laws of the natural world; and as we reconcile our
imagination, nevertheless, to such deviation from the course of nature
in the Economy of Revelation, so surely we may bear without impatience
or perplexity that the subsequent history of Revelation should in turn
diverge from the path in which it originally commenced \footnote{This is Middleton's ground in the following
passage, with which should be compared the passages from Hume in the
text: ``The present question concerning the reality of the miraculous
powers of the primitive Church depends on the joint credibility of the
facts, pretended to have been produced by those powers, and of the
witnesses who attest them. If either part be infirm, their credit must
sink in proportion; and if the facts especially be incredible, must \emph{of
course} fall to the ground, \emph{because} no force of testimony
can alter the nature of things. The credibility of facts lies open to
the trial of our reason and senses, but the credibility of witnesses
depends on a variety of principles wholly concealed from us; and
though in many cases it may reasonably be presumed, yet in none can it
certainly be known. For it is common with men, out of crafty and
selfish views, to dissemble and deceive; or out of weakness and
credulity to embrace and defend with zeal what the craft of others had
imposed upon them; but plain facts cannot delude us—cannot speak any
other language, or give any other information but what flows from
nature and truth. The testimony, therefore, of facts, as it is offered
to our senses in this wonderful fabric and constitution of worldly
things, may properly be called the \emph{testimony of God Himself}, as
it carries with it the surest instruction in all cases, and to all
nations, which in the ordinary course of His providence He has thought
fit to appoint for the guidance of human life.'' Pp. ix. x.

Again, ``Our first care should be to inform
ourselves of the proper nature and condition of those miraculous
powers, ... as they are represented to us in the history of the
Gospel; for till we have learned from those sacred records what they
really were, for what purposes granted, and in what manner exerted by
the Apostles and first possessors of them, we cannot form a proper
judgment on those evidences which are brought either to confirm or
confute their continuance in the Church, and must dispute consequently
at random, as chance or prejudice may prompt us, about things unknown
to us.'' P. xi.

Again, ``The whole which the wit of man can
possibly discover either of the ways or will of the Creator, must be
acquired … not by imagining vainly within ourselves what may be
proper or improper for Him to do, but by looking abroad and
contemplating \emph{what He has actually} done; and attending
seriously to that revelation which He made of Himself from the
beginning, and placed continually before our eyes, in the wonderful
works and beautiful fabric of this visible world.'' P. xxii.}.

52.
Hume argues against miracles generally, ``Though the Being to whom
the miracle is ascribed be, in this case, Almighty, it does not, upon
that account, become a whit more probable; since it is impossible for
us to know the attributes or actions of such a Being, otherwise than
from the experience which we have of His productions in the usual
course of nature.'' \footnote{Essay on Miracles, Part ii. circ. fin.} And elsewhere he says, ``The Deity
is known to us only by His productions ... As the universe shows
wisdom and goodness, we infer wisdom and goodness. As it shows a
particular degree of these perfections, we infer a particular degree
of them, precisely adapted to the effect which we examine. But further
attributes, or further degrees of the same attributes, we can never be
authorized to infer or suppose, by any rules of just reasoning.'' \footnote{Essay on Providence.} And in a note he adds, ``In general, it may, I think, be
established as a maxim, that where any cause is known only by its
particular effects, it must be impossible to infer any new effects
from that cause ... To say that the new effects proceed only from a
continuation of the same energy, which is already known from the first
effects, will not remove the difficulty. For even granting this to be
the case, (which can seldom be supposed,) the very continuation and
exertion of a like energy, (for it is impossible it can be absolutely
the same,) I say, this exertion of a like energy, in a different
period of space and time, is a very arbitrary supposition, and what
there cannot possibly be any traces of in the effects, from
which all our knowledge of the cause is originally derived. Let the
inferred cause be exactly proportioned, as it should be, to the known
effect; and it is impossible that it can possess any qualities, from
which new or different effects can be inferred.''

53.
This is not the place to analyze a paradox which is sufficiently
refuted by the common sense of a religious mind; but the point which
concerns us to consider is, whether persons who, not merely question,
but prejudge the Ecclesiastical Miracles on the ground of their want
of resemblance, whatever that be, to those contained in
Scripture,—as if the Almighty could not do in the Christian Church
anything but what He had already done at the time of its foundation,
or under the Mosaic Covenant,—whether such reasoners are not siding
with the sceptic who in the above passages denies that the First Cause
can act supernaturally at all, because in nature He can but act
naturally, and whether it is not a happy inconsistency by which they
continue to believe the Scripture record, while they reject the
records of the Church.

54.
Indeed, it would not be difficult to show that the miracles of
Scripture are a far greater innovation upon the economy of nature than
the miracles of the Church upon the economy of Scripture. There is
nothing, for instance, in nature at all to parallel and mitigate the
wonderful history of the assemblage of all animals in the Ark,
or the multiplication of an artificially prepared substance, such as
bread. Walking on the sea, or the resurrection of the dead, is a plain
reversal of its laws. On the other hand, the narrative of the combats
of St. Antony with evil spirits is a development rather than a
contradiction of Revelation; viz., as illustrating such texts as speak
of Satan being cast out by prayer and fasting. To be shocked then at
the miracles of Ecclesiastical History, or to ridicule them for their
strangeness, is no part of a scriptural philosophy.

55.
Nor can the argument from \emph{à priori} ideas of propriety be made
available against Ecclesiastical Miracles with more safety than the
argument from experience. This method of refutation, as well as the
other, (to use the common phrase,) proves too much. Those who have
condemned the miracles of the Church by such a rule, have before now
included in their condemnation the very notion of a miracle
altogether, as the creation of barbarous and unphilosophical minds,
who knew nothing of the beautiful order of nature, and as unworthy to
be introduced into our contemplation of the providences of Divine
Wisdom. A miracle has been considered to argue a defect in the system
of moral governance, as if it were a correction or improvement of what
is in itself imperfect or faulty, like a patch of new cloth upon an
old garment. The Platonists of old were influenced by something like
this feeling, as if none but low and sordid persons would either
attempt or credit miracles truly such, and none but quacks and
impostors would profess them. The only true miracles, in the
conception of such a school, are miracles of knowledge;—words or
deeds which are the result of a greater insight into, or foresight of,
the course of nature, and are proofs of a liberal education and a
cultivated and reflective mind \footnote{Hence the charge against the Christians of
magic, or [\emph{go\={e}teia}]. Tertull. Apol. 23. Origen in Cels.
i. 38, ii. 9. Arnob. contr. Gent. i. Euseb. Dem. Ev. iii. 5 and 6, pp.
112, 130. August. Serm. xliii. 4, contr. Faust. xii. 45, Ep. cxxxviii.
fin. Julian calls St. Paul the greatest of rogues and conjurors, [\emph{ton
pantas pantachou tous p\={o}pote go\={e}tas kai apate\={o}nas
hyperballomenon Paulon}]. Ap. Cyr. iii. p. 100. Apollonius
professed a knowledge of nature as the secret of his miracles. Vid.
Philostr. Vit. Ap. v. 12. Also Quæst. ad Orthod. 24, where Apollonius
is said to have done his miracles [\emph{kata t\={e}n epist\={e}m\={e}n
t\={o}n physik\={o}n duname\={o}n}], not [\emph{kata t\={e}n
theian authentian}]. Philostratus illustrates this when he seems to
doubt whether the young woman was really dead, whom Apollonius raised,
iv. 45. [Vid. Kortholt. de Vit. et Mor. Christ. c. 3, 4.]}. It is easy to see how a
habit of this sort may grow upon scientific men, especially at this
day, unless they are on their guard against it. There is so much
beauty, majesty, and harmony in the order of nature, so much to fill,
satisfy, and tranquillize the mind, that by those who are accustomed
to the contemplation, the notion of an infringement of it will at
length be viewed as a sort of profanation, and as even shocking, as
the mere dream of ignorance, the wild and atrocious absurdity of
superstition and enthusiasm, (if it is right to use such language even
in order to describe the thoughts of others,) and as if analogous, to
take another and less serious subject, to some gross solecism, or
indecorum, or wanton violation of social usages or feelings. We should
be very sure, if we resolve on rejecting the Ecclesiastical Miracles,
that our reasons are better than that false zeal for our Master's
honour, which such philosophers express for the honour of the Creator,
and which reminds us of the exclamation, ``Be it far from Thee, Lord:
this shall not be unto Thee!'' as uttered by one who heard for the
first time that doctrine which to the world is foolishness.

8.

56.
The question has hitherto been argued on the admission, that a
distinct line can be drawn in point of character and circumstances
between the Miracles of Scripture and those of Church History; but
this is by no means the case. It is true, indeed, that the Miracles of
Scripture, viewed as a whole, recommend themselves to our reason, and
claim our veneration, beyond all others, by a peculiar dignity and
beauty; but still it is only as a whole that they make this impression
upon us. Some of them, on the contrary, fall short of the attributes
which attach to them in general, nay, are inferior in these respects
to certain ecclesiastical miracles, and are received only on the credit of the system, of which they form part. Again, specimens are
not wanting in the history of the Church, of miracles as awful in
their character and as momentous in their effects, as those which are
recorded in Scripture. The fire interrupting the rebuilding of the
Jewish temple, and the death of Arius, are instances, in
Ecclesiastical History, of such solemn events. On the other hand,
difficult facts in the Scripture history are such as these:—the
serpent in Eden, the Ark, Jacob's vision for the multiplication of his
sheep, the speaking of Balaam's ass, the axe swimming at Elisha's
word, the miracle on the swine, and various instances of prayers or
prophecies, in which, as in that of Noah's blessing and curse, words
which seem the result of private feeling are expressly or virtually
ascribed to a Divine suggestion.

57.
And thus, it would seem, there exists in matter of fact that very
connection and intermixture between ecclesiastical and Scripture
miracles, which, according to the analogy suggested in a former page,
the richness and variety of physical nature rendered probable.
Scripture history, far from being broadly separated from
ecclesiastical, does in part countenance what is strange in the
miraculous narratives of the latter, by affording patterns and
precedents for them itself. It begins a series which has, indeed, its
higher specimens and its lower, but which still proceeds in the way of
a series, with a progress and continuation, without any sudden
breaks and changes, or even any exact law of variation according to
the succession of periods. As in the natural world the animal and
vegetable kingdoms imperceptibly melt into each other, so are there
mutual affinities and correspondences between the two families of
miracles as found in inspired and uninspired history, which show that,
whatever may be their separate peculiarities, yet as far as concerns
their internal characteristics, they admit of being parts of one
system.

58.
For instance, there is not a more startling, yet a more ordinary gift
in the history of the first ages of the Church than the power of
exorcism; while at the same time it is open to much suspicion, both
from the comparative facility of imposture and the intrinsic
strangeness of the doctrine it inculcates. Yet, here Scripture has
anticipated the Church in all respects, even going the length of
testifying to the diabolical possession of brutes, which appears so
extravagant when introduced, as instanced above, into the life of
Hilarion by St. Jerome. Again, that relics should be the instruments
of exorcism is an aggravation of a doctrine already difficult; yet we
read in Scripture, ``And God wrought special miracles by the hands of
Paul, so that from his body were brought unto the sick handkerchiefs
or aprons, and the diseases departed from them, and the evil spirits
went out of them.'' [Acts xix. 11, 12.]

Similar precedents for a supernatural presence in
things inanimate are found in the miracles wrought by the touch of our
Saviour's garments, not to insist on what is told us about St. Peter's
shadow. Again, we have to take into account the resurrection of the
corpse which touched Elisha's bones, a work of Divine Power, which,
whether considered in its appalling greatness, the absence of apparent
object, and the means through which it was accomplished, we should
think incredible, with the now prevailing notions of miraculous
agency, were we not familiar with it. Elijah's mantle is another
instance of a relic endued with miraculous power. Again, the
multiplication of the wood of the Cross (the \emph{fact} of which is not here
determined, but must depend on the testimony and other evidence
producible) is but parallel to Elisha's multiplication of the oil, and
of the bread and barley, and to our Lord's multiplication of the
loaves and fishes. Again, the account of the consecrated host becoming
a cinder in unworthy hands is not so strange as the very first miracle
wrought by Moses, the first evidential miracle recorded in Scripture,
when his rod became a serpent, and then a rod again; nor stranger than
our Lord's first miracle, when water was turned into wine. When the
tree was falling upon St. Martin, he is said to have caused it to
whirl round and fall elsewhere by the sign of the Cross; is this more
startling than Elisha's causing the iron axe-head to swim by
throwing a stick into the water?

59.
It is objected by Middleton, that after the decree of the Council of
Laodicea, restricting exorcism to such as were licensed by the Bishop,
the practice died away \footnote{Inquiry, pp. 95, 96.}; this, indeed, implies a very
remarkable committal or almost abandonment of a Divine gift, supposing
it such, to the discretion of its human instruments; but how does it
imply more than we read of in the Apostolic history of the Corinthian
Christians, who had so absolute a possession of their supernatural
powers that they could use them disorderly, or pervert them to
personal ends? The miracles in Ecclesiastical History are often
wrought without human instruments, or by instruments but partially
apprehensive that they are such; but did not the rushing mighty wind,
at Pentecost, come down ``suddenly'' and unexpectedly? and were not the
Apostles forthwith carried away by it, not in any true sense using the
gift, but compelled to speak as the Spirit gave them utterance? It is
objected that ecclesiastical miracles are not so distinct and
unequivocal as to have a claim to be accounted true, but admit of
being plausibly attributed to fraud, collusion, or misstatement in
narrators; yet, in like manner, St. Matthew tells us that the Jews
persisted in maintaining that the disciples had stolen away our
Lord's Body, and He did not show Himself, when risen, to the Jews; and
various other objections, to which it is painful to do more than
allude, have been made to other parts of the sacred narrative. It is
objected that St. Gregory's, St. Martin's, or St. Hilarion's miracles
were not believed when first formally published to the world by Nyssen,
Sulpicius, and St. Jerome; but it must be recollected that Gibbon
observes scoffingly, that ``the contemporaries of Moses and Joshua
beheld with careless indifference the most amazing miracles,'' that
even an Apostle, who had attended our Lord through His ministry, did
not believe his brethren's report of His resurrection, and that St.
Paul's supernatural power of punishing offenders was doubted at
Corinth by the very parties who had seen his miracles and had been his
converts. That alleged miracles, then, should admit of doubt, or be
what is called ``suspicious,'' is not at all inconsistent with their
title to be considered the immediate operation of Divine Power.

60.
It is observable also, that this intercommunion of miracles, if the
expression may be used, which exists between the respective
supernatural agencies contained in Scripture and in Church history, is
seen also in the separate portions of Scripture history. The miracles
of Scripture may be distributed into the Mosaic, the Prophetical, and
the Evangelical; of which the first are mainly of a judicial and
retributive character, and wrought on a large field; the last are
miracles of mercy; and the intermediate are more or less of a romantic
or poetical cast. Yet, among the Mosaic we find the changing of the
rod into a serpent, and the sweetening of the water by a branch, which
belong rather to the second period; and among the Christian are the
deaths of Ananias and Sapphira, which resemble the awful acts of the
first; while Philip's transportation by the Spirit, and the ship's
sudden arrival at the shore, might be ranked among those of the
second.

9.

61.
And moreover this circumstance is worth considering, that a sort of
analogy exists between the Ecclesiastical and Evangelical histories,
and the Prophetical and Mosaic. The Prophetical and Ecclesiastical
are, each in its place, a sort of supplement to the supernatural
manifestations with which the respective Dispensations open, and
present to us a similar internal character. And, whereas there was an
interval between the age of Moses and the revival of miraculous power
in the Prophets, though extraordinary providences were never wholly
suspended, so the Ecclesiastical gift is restricted in its operation
in the first centuries compared with the exuberant exercise recorded
of it in the fourth and fifth; and as the Prophetical miracles
in a great measure belong to the schools of Elijah and Elisha, so the
Ecclesiastical have a special connection with the ascetics and
solitaries, and the orders of families of which they were patriarchs,
with St. Antony, St. Martin, and St. Benedict, and other great
confessors or reformers, who are the antitypes of the Prophets.
Moreover, much might be said concerning the romantic character of the
Prophetical miracles. Those of Elisha in particular are related, not
as if parts of the history, but rather as his ``Acta;'' as illustrations
indeed of that double portion of power gained for him by Elijah's
prayer, and perhaps with some typical reference to the times of the
Gospel, but still with a profusion and variety very like the
luxuriance which offends us in the miraculous narratives of
ecclesiastical authors. Elisha begins by parting Jordan with Elijah's
mantle; then he curses the children, and bears destroy forty-two of
them; then he supplies the kings of Judah, Israel, and Edom with water
in the wilderness, and gives them victory over Moab; then he
multiplies the oil; then he raises the Shunammite's son; then he
renders the poisonous pottage harmless by casting meal into it; then
he multiplies the bread and barley; then he directs Naaman to a cure
of his leprosy; then he reads Gehazi's heart, follows him throughout
his act of covetousness, and inflicts on him Naaman's leprosy; then he
makes the iron swim; then he reveals to the king of Israel the
counsels of Syria, and casts an illusion before the eyes of his army;
then he prophesies plenty in the siege; then he foretells Hazael's
future course. These wonderful acts are strung together as the direct
and formal subjects of the chapters in which they occur: they have no
continuity; they carry on no action or course of Providence. At length
Elisha falls sick, and, on the king's visiting him, promises him a
series of victories over the Syrians; then he dies and is buried, and
by accident a corpse is thrown into his grave; and ``when the man was
let down, and touched the bones of Elisha, he revived, and stood up on
his feet.'' [2 Kings xiii. 21.] Surely it is not too much to say, that
after this inspired precedent there is little in ecclesiastical
legends of a nature to offend as regards their matter; their
credibility turning first on whether they are to be expected at all,
and next whether they are avouched on sufficient evidence.

62.
Or take again the history of Samson; what a mysterious wildness and
eccentricity is impressed upon it, upon the miracles which occur in
it, and upon its highly favoured though wayward subject! ``At this
juncture,'' says a recent writer, speaking of the low estate of the
chosen people when Samson was born, ``the most extraordinary of the
Jewish heroes appeared; a man of prodigious physical power, which
he displayed, not in any vigorous and consistent plan of defence
against the enemy, but in the wildest feats of personal daring. It was
his amusement to plunge headlong into peril, from which he extricated
himself by his individual strength. Samson never appears at the head
of an army, his campaigns are conducted in his own single person. As
in those of the Grecian Hercules and the Arabian Antar, a kind of
comic vein runs through the early adventures of the stout-hearted
warrior, in which love of women, of riddles, and of slaying
Philistines out of mere wantonness, vie for the mastery. Yet his life
began with marvel, and ended in the deepest tragedy.'' \footnote{Milman's History of the Jews, vol. i. p. 204.} The
tone of this extract cannot be defended; yet what else has the writer
done towards the inspired narrative, but invest it in those showy
human colours which legendary writers from infirmity, and enemies from
malice, have thrown over the miracles of the Church? There is
certainly an aspect of romance in which Samson may be viewed, though
he was withal the instrument of a Divine Presence; and so again there
may have been a divinity in the acts and fortunes, and a spiritual
perfection in the lives, of the ancient Catholic hermits and
missionaries, in spite of whatever is wild, uncouth, and extravagant
in their personal demeanour and conduct, or rather in the record of
them. Once more; the books of Daniel and Esther are very different in
composition and style from the earlier portions of the sacred
volume, and present a view of the miraculous dealings of the Almighty
with His Church, very much resembling what we disparage in
ecclesiastical legends, or again in the historical portions of the
so-called Apocrypha, as if poetical or dramatic.

63.
The two Economies then, the Prophetical and the Ecclesiastical, thus
resembling each other in their character as well as their position in
their two Covenants respectively, should any one urge, as was stated
in a former place \footnote{Supra, p. 115.}, that the Ecclesiastical Miracles
virtually form a new dispensation, we need not deny it \emph{in the sense}
in which the Prophetical Miracles are distinct from the Mosaic; that
is, not as if the Law was in any respect or in any part repealed by
the Prophetical Schools, but that they, as well as other works of God,
had a character of their own, and, as in other things, so in their
miracles, were a new exhibition of that Supernatural Presence which
over-shadowed Israel from first to last. And it may be added, that, as
a gradual revelation of Gospel truth accompanied the miracles of the
Prophets, so to those who admit the Catholic doctrines as enunciated
in the Creed, and commented on by the Fathers, the subsequent
expansion and variation of supernatural agency in the Church, instead
of suggesting difficulties, will seem to be in correspondence, as they
are contemporaneous, with the developments and additions in
dogmatic statements which have occurred between the Apostolic and the
present age, and which are but a result and an evidence of spiritual
life.

10.

64.
Nor, lastly, is it any real argument against admitting the
Ecclesiastical Miracles on the whole, or against admitting certain of
them, that certain others are rejected on all hands as fictitious or
pretended. It happens as a matter of course, on many accounts, that
where miracles are really wrought, miracles will also be attempted, or
simulated, or imitated, or fabled; and such counterfeits become, not a
disproof, but a proof of the existence of their prototypes, just as
hypocrisy and extravagant profession are an argument for, and not
against, the reality of virtue \footnote{Douglas' Crit. p. 19.}. It is doubtless the tendency
of religious minds to imagine mysteries and wonders where there are
none; and much more, where causes of awe really exist, will they
unintentionally misstate, exaggerate, and embellish, when they set
themselves to relate what they have witnessed or have heard \footnote{Camp. Miracl. p. 122. Jenkins' Christ. Rel.
vol. ii. p. 455.}.
A fact is not disproved because the testimony is confused or
insufficient; it is only unproved. And further, the imagination, as is
well known, is a fruitful cause of apparent miracles \footnote{Le Moyne Miracl. pp. 486, 502. Douglas' Crit.
p. 93, etc.};
and hence, wherever there are works wrought which absolutely surpass
the powers of nature, there are likely to be others which surpass its
ordinary action. It would be no cause for surprise if, as the
destruction of Sodom is said to have arisen from volcanic influence,
so in the multitude of cures which the Apostles effected some were
solely attributable to natural, but unusual, effects of faith. And if
Providence sometimes makes use of natural principles even when
miracles seem intended as evidence of His immediate presence, much
more is He likely to intermingle the ordinary and the extraordinary,
when His object is not to prove a revelation, to accredit a messenger,
or to certify a doctrine, but to confirm or encourage the faithful, or
to rouse the attention of unbelievers. And it will be impossible to
draw the line between the two; and the possibility of explaining some
of them on natural principles will unjustly prejudice the mind against
accounts of those which cannot be so explained.

65.
Moreover, as Scripture expressly shows us, wherever there is
miraculous power, there will be curious and interested bystanders who
would fain ``purchase the gift of God'' for their own aggrandisement,
and ``cast out devils in the Name of Jesus,'' and who counterfeit what
they have not really to exhibit, and gain credit and followers among
the ignorant and perverse. The impostures, then, of various
kinds which from the first hour abounded in the Church \footnote{Vid. Acts viii. 9; xvi. 17; xix. 13. Vid.
Lucian. Peregr. etc. ap. Middlet. Inqu. p. 23.} prove
as little against the truth of her miracles as against the canonicity
of her Scriptures. Yet here too pretensions on the part of worthless
men will be sure to scandalize inquirers, and the more so if, as is
not unlikely, such pretenders manage to ally themselves with the
Saints, and have an historical position in the fight which is made for
the integrity or purity of the faith; yet, St. Paul was not less an
Apostle, nor have Confessors and Doctors been less his successors,
because ``as they have gone to prayer'' a spirit of Python has borne
witness to them as ``the servants of the most high God,'' and the
teachers of ``the way of salvation.''

66.
Nor is it any fair argument against Ecclesiastical miracles, that for
the most part they have a legendary air, while the miracles contained
in Scripture are on the contrary so soberly, so gravely, so exactly
stated; unless indeed it is an absurdity to contemplate a gift of
miracles without an attendant gift of inspiration to record them. Were
it not that the Evangelists were divinely guided, doubtless we should
have in Scripture that confused mass of truth and fiction together,
which the Apocryphal Gospels exhibit, and to which St. Luke seems to
allude. I repeat, the character of facts is not changed because they
are incorrectly reported; distance of time and place only does
injury to the record of them. The Scripture miracles were in
themselves what they are to us now, at the very time that the world
was associating them with the prodigies of Jewish strollers, heathen
magicians and astrologers, and idolatrous rites; they would have been
thus associated to this day, had not inspiration interposed; yet, in
spite of this, they would have been deserving our serious attention as
now, so far as we were able to separate the truth from the falsehood.
And such is the state in which Ecclesiastical miracles actually do
come to us, because inspiration was not continued; they are dimly seen
in twilight and amid shadows;—let us not, then, quarrel with them on
account of a characteristic which is but the necessary consequence of
external circumstances.

\xchapter{Chapter 4. On the State of the Argument in behalf of the Ecclesiastical Miracles}

67.
VARIOUS able writers, Leslie,
Paley, and Douglas, have laid down certain tests or criteria of
matters of fact, which may serve as guarantees that the miracles
really took place which are recorded in Scripture. They consider these
criteria to be of so rigid a nature that an alleged event which
satisfies them must necessarily have occurred, and that, as their
argument seems to imply, however great its antecedent improbability.
Thus they reply to objections such as Hume's drawn from the uniformity
of nature; not meeting them directly, but rather superseding the
necessity of considering them; for what is proved to be true, need not
be proved to be possible. Hume scruples not to use ``miracle'' and ``impossibility''
as convertible terms \footnote{``What have we to oppose to such a cloud of
witnesses but the absolute \emph{impossibility} or \emph{miraculous
nature} of the events which they relate?'' (Essay on Miracles.)};
Leslie before him, Douglas after him, seem to answer, ``Would you
believe a miracle if you \emph{saw} it? Now we are prepared to offer
evidence, if not as strong, still as convincing, as ocular
demonstration.'' Thus they escape from the abstract argument by a
controversial method of a singularly practical, and as it may be
called, English character.

68.
It would be well if such writers stopped here, but it was hardly to be
expected. Disputants are always exposed to the temptation of being
over-candid towards objections which they think they have outrun; they
admit as facts or truths what they have shown to be irrelevant as
arguments. Thus, even were there nothing of a kindred tone of mind in
Hume, who has assailed the Scripture miracles, and in some of our
friends who have defended them, it might have been anticipated that
the consciousness of possessing an irresistible weapon in the contest
would have led us to treat the arguments of our opponents with a
dangerous generosity. But, unhappily, there is much in Protestant
habits of thought actually to dispose our writers to defer to a
rationalistic principle of reasoning, the force of which they have
managed to evade in the particular case. Hence, though they are
earnest in their protest against Hume's summary rejection of all
miraculous histories whatever, they make admissions, which only do not
directly tell against the principal Scripture miracles, and do tell
against all others. They tacitly grant that the antecedent
improbability of miracles is at least so great that it can only be
overcome by the strongest and most overpowering evidence; that
second-best evidence does not even tend to prove them; that they are
absolutely incredible up to the very moment that all doubt is
decisively set at rest; that there can be no degrees of proof, no
incipient and accumulating arguments to recommend them; that no
relentings of mind or suspense of judgment is justifiable, as various
fainter evidences are found to conspire in their favour; that they may
be scorned as fictions, if they are not to be venerated as truths.

69.
It looks like a mere truism to say that a fact is not disproved,
because it is not proved; ten thousand occurrences are ever passing,
which leave no record behind them, and do not cease to have been
because they are forgotten. Yet Bishop Douglas, in his defence of the
New Testament Miracles in answer to Hume, certainly assumes that no
miracle is true which has not been proved so, or that it is safe to
treat all miracles as false which are not recommended by evidence as
strong as that which is adducible for the Miracles of Scripture.

70.
In estimating statements of fact, it is usual to allow that various
occurrences may be all true, which rest upon very different degrees of
evidence. It does not prove that this passage of history is false and
the fabrication of impostors, because that passage is attested
more distinctly and fully. Writers, however, like Douglas, are
constantly reminding us that we \emph{need} not receive the Ecclesiastical
miracles, \emph{though} we receive those of the New Testament. But the
question is not whether we \emph{need} not, but whether we \emph{ought} not to
receive the former, as well as the latter; and if it really is the
case that we ought not, surely this must be in consequence of some
positive reasons, not of a mere inferiority in the evidence. It is
plain, then, that such reasoners, though they deny that an \emph{à
priori} ground can be maintained in fact against the miracles of
Scripture, still at least agree with Hume in thinking such a ground
does exist, and that it is conclusive against ecclesiastical miracles
even antecedent to the evidence.

71.
In the title of his Dissertation, Douglas promises us ``a criterion by
which the \emph{true} miracles recorded in the New Testament are
distinguished from the \emph{spurious} miracles of the Pagans and
Papists;'' yet, when he proceeds to state in the body of the work the
real object to which he addresses himself, we find that it relates
quite as much to the \emph{evidence} for either class of miracles as
to the fact itself of their occurrence. He says, that whereas ``the
accounts which have been published to the world of miracles in
general,'' are concerned with events which are supernatural either in
themselves or under their circumstances, while the latter class
can be explained on natural principles, the former ``\emph{may}, from
the \emph{insufficiency} of the evidence produced in support of them,
be \emph{justly suspected} to have never happened.'' \footnote{Page 25.} But how does \emph{insufficiency} in the evidence create a positive
prejudice against an alleged fact? How can things depend on our
knowledge of them? This writer must mean that evidence of an inferior
kind is insufficient to overcome a certain \emph{pre-existing objection}
which attaches to the very notion of these miracles; otherwise even
slight evidence is sufficient to influence our minds, as Bishop Butler
would tell us, so far as it is positive, and evidence of this
defective kind may constitute the very trial of our obedience.

72.
Douglas continues: ``I flatter myself, that the evidence produced in
their support,''—in support of the miracles of ``Pagans and Papists,''—``will
appear to be so very defective and insufficient, as \emph{justly to
warrant} our rejecting them as idle tales that never happened, and \emph{the
inventions of bold and interested deceivers}.'' \footnote{Page 26.} There are many reasons to warrant our disbelieving alleged
facts, and ascribing them to imposture; for instance, if the evidence
is contradictory, or attended by suspicious circumstances; if the
witnesses are of bad character, or strong inducements to fraud exist;
but it is difficult to see how its mere \emph{insufficiency} or \emph{defectiveness}
is a justification of so decided a step. The direct effect of evidence
is to create a presumption, according to its strength, in favour of
the fact attested; it does not appear how it can create a presumption
the other way. The real explanation of this mode of writing certainly
must be, that the writer takes it for granted that all miraculous
accounts are already in a manner self-condemned, as being miraculous, \emph{till}
they are proved; and that evidence offered for them, which does not
amount to a proof, is but involved in that existing prejudice. There
is no medium then; the testimony must either prevail or be scouted; it
is certainly a fraud, if it is not an overpowering demonstration.

73.
But the author in question scarcely leaves us in doubt of his meaning,
when he avails himself of the following maxim of Dr. Middleton's: ``I
have already observed,'' he says, ``that the testimony supporting
[miracles] must be free from every suspicion of fraud and imposture.
And the reason is this: the history of miracles (to make use of the
words of an author whose authority you will think of some weight) is
of a kind totally different from that of common events; the one \emph{to
be suspected always of course without the strongest evidence} to
confirm it; the other to be admitted of course without as strong
reason to suspect it. So that, wherever the evidence urged for
miracles \emph{leaves grounds} for a suspicion of fraud and
imposition, \emph{the very suspicion} furnishes \emph{sufficient
reasons} for disbelieving them. And what I shall offer under this
head will make it evident, that those miracles which the Protestant
Christian thinks himself at liberty to reject have always been so \emph{insufficiently}
attested as to \emph{leave full scope} for fraud and imposition.'' \footnote{How much more cautious is Jortin! ``Though miracles,'' he says, ``\emph{may}
be wrought in secret, and cannot be disproved only because they were
seen by few, yet they \emph{often} afford motives for suspicion, and a
wise inquirer would perhaps \emph{suspend} his assent in such cases,
and pass \emph{no judgment} about them.'' (Eccl. History, Works, vol.
ii. p. 3, ed. 1810.) Again, ``As far as the subsequent miracles
mentioned by Christian writers fall short of the distinguishing
characters belonging to the works of Christ and His Apostles, \emph{so
far} they must fail of giving us \emph{the same full persuasion and
satisfaction}.'' (P. 20.)} That is, we may ascribe a story to fraud, whenever it is not
absolutely impossible so to ascribe it; we may summarily reject and
vilify all evidence \emph{up to} such evidence as is a moral
demonstration, though to such we must immediately yield, because we
cannot help it; and this as a matter ``of course.'' All this surely
implies the existence of some deep latent prejudice in the writer's
mind against miraculous occurrences, considered in themselves; else it
is not a reasonable mode of arguing.

74.
The Bishop continues in the same strain to ``lay down a few general
rules by which we may try those pretended miracles, one and all,
wherever they occur, and which may set forth \emph{the grounds on which
we suspect} \emph{them false}.'' \footnote{Page 27.} And then, ``by way of illustration,'' he selects three, telling
us that we suspect them false, or ``we \emph{may} suspect them false,''
when the existing accounts of miracles were not published till long
after the time when, or not at the place where, they are said to have
occurred; or, at least, if it seems probable that they were suffered
to get into circulation without examination at the time and place.
Here of course he does but act up to Middleton's bold principle which
he has adopted; he considers himself at liberty to bid defiance and
offer resistance to all evidence, till he is fairly subdued by it,
till it is impossible to doubt, and no merit to believe; while he
would never reject or impute fraud to a record of ordinary events,
merely because it was published in a foreign country, or a hundred
years after the events in question, however he might justly consider
such circumstances to weaken the force of the evidence.

75.
In a subsequent page of his work he speaks still more pointedly: ``When
the reporters of miracles,'' he says, ``content themselves with general
assertions and vague claims to a miraculous power, without ever
attempting to corroborate them by descending to particular facts, and
leave us strangely in the dark as to the persons by whom, the
witnesses before whom, and the objects upon whom these miraculous
powers are said to be exercised, omitting every circumstance necessary to be related by them before any inquiry can be made into
the truth of the pretension; when miracles, I say, are reported in
this \emph{unsatisfactory} manner, (and instances of miracles reported
on the spot by contemporary writers, in such a manner, might be
mentioned,) in this case it would be \emph{the height of credulity} to
pay \emph{any} regard to them in a distant age, because no regard
could possibly be paid to them in their own.'' \footnote{Page 50.} Yet it does not appear how this ``unsatisfactory manner'' in the
report can touch the events reported; if they took place, they were
before and quite independent of the evidence at present existing for
them, be it greater or less; our knowledge or ignorance does not
create or annihilate facts.

76.
Now these passages from Bishop Douglas have been drawn out, not simply
with a view of criticising him, but in order to direct attention to
the fact which he illustrates, viz., that our feeling towards the
Ecclesiastical Miracles turns much less on the evidence producible for
them, than on our view concerning their antecedent probability. If we
think such interpositions of Providence likely or not unlikely, there
is quite enough evidence existing to convince us that they really do
occur; if we think them as unlikely as they appear to Douglas,
Middleton, and others, then even evidence as great as that which is
producible for the miracles of Scripture would not be too much, nay, perhaps not enough, to conquer an inveterate, deep-rooted, and
(as it may be called) ethical incredulity.

77.
It shall here be assumed that this incredulity is a fault; and it is
the result of a state of mind which has been prevalent among us for
some generations, and from which we are now but slowly extricating
ourselves. We have been accustomed to believe that Christianity is
little more than a creed or doctrine, introduced into the world once
for all, and then left to itself, after the manner of human
institutions, and under the same ordinary governance with them, stored
indeed with hopes and fears for the future, and containing certain
general promises of aid for this life, but unattended by any special
Divine Presence or any immediately supernatural gift. To minds
habituated to such a view of Revealed Religion, the miracles of
ecclesiastical history must needs be a shock, and almost an outrage,
disturbing their feelings and unsettling their most elementary notions
and thoroughly received opinions. They are eager to find defects in
the evidence, or appearances of fraud in the witnesses, as a relief to
their perplexity, and as an excuse for rejecting, as if on the score
of reason, what their heart and imagination have rejected already. Or
they are too firmly persuaded of the absurdity, as they consider it,
which such pretensions on the part of the Church involve, to be moved
by them at all; and they content themselves with coldly
demanding to know points which cannot now be known, or to be satisfied
about difficulties which never will be cleared up, before they can be
asked to take interest in statements which they consider so
unreasonable. And certainly they are both philosophical and religious
in thus acting, granting that the Lord of all is present with
Christians only in the way of nature, as with His creatures all over
the earth. On the other hand, if we believe that Christians are under
an extraordinary Dispensation, such as Judaism was, and that the
Church is a supernatural ordinance, we shall in mere consistency be
disposed to treat even the report of miraculous occurrences with
seriousness, from our faith in a Present Power adequate to their
production. Nay, if we only go so far as to realize what Christianity
is, when considered merely as a creed, and what stupendous
over-powering facts are involved in the doctrine of a Divine
Incarnation, we shall feel that no miracle can be great after it,
nothing strange or marvellous, nothing beyond expectation.

2.

78.
All this applies to the view we shall take of the nature of the facts
which are laid before us, as well as of the character of the evidence.
If we disbelieve the divinity of the Church, then we shall do our best
to deny that the facts attested are miraculous, even admitting
them to be true. ``Though our not knowing on whom, or by whom, or
before whom, the miracles recorded by the Fathers of the second and
third centuries were wrought,'' says Douglas, ``should be allowed not to
destroy their credit (though this is a concession which very few will
make ... ), yet the facts appealed to are of so \emph{ambiguous} a kind,
that, granting they did happen, it will remain to be decided, by a
consideration of the circumstances attending the performance of them,
whether there was any miracle in the case or no.'' \footnote{Page 228.} Certainly it is a rule of philosophy to refer effects, if
possible, to known causes, rather than to imagine a cause for the
occasion; and, on the other hand, to be suspicious of alleged facts
for which no cause can be assigned, or which are unaccountable. If,
then, there is nothing in the Church more than in any other society of
men, it is natural to attribute the miracles alleged to have been
wrought in it to natural causes, where that is possible, and to
disparage the evidence where it is not so. But if the Church be
possessed of supernatural powers, it is not unnatural to refer to
these the facts reported, and to feel the same disposition to heighten
their marvellousness as otherwise is felt to explain it away. Thus our
view of the evidence will practically be decided by our views of
theology. There are two providential systems in operation among us,
the visible and the invisible, intersecting, as it were, each
other, and having a certain territory in common; and in many cases we
do not know the exact boundaries of each, as again we do not know the
minute details of those facts which are ascribed by their reporters to
a miraculous agency. For instance, faith may sometimes be a natural
cause of recovery from sickness, sometimes a miraculous instrument;
the application of oil may be a mere expedient of medical art, or
parallel to the application of water in Baptism. The Martyrs have
before now found red-hot iron, on its second application, even
grateful to their seared limbs; on the other hand, cases of a similar
kind are said to have occurred where religion was not in question, and
where a divine interposition cannot be conjectured. Sudden storms and
as sudden calms on the lake of Gennesareth might be of common
occurrence; yet the particular circumstances under which the waters
were quieted at our Lord's word may have been sufficient to convince
beholders that it was a miracle. The Red Sea may have been ordinarily
exposed to the influence of the East Wind, and nevertheless the
separation of its waters, as described in the Book of Exodus, may have
required a supernatural influence. In these and numberless other
instances men will systematize facts in their own way, according to
their knowledge, opinions, and wishes, as they are used to do in all
matters which come before them; and they will refer them to causes which they see or believe, in spite of their being referable to
other causes about which they are ignorant or sceptical.

79.
When, then, controversialists go through the existing accounts of
ecclesiastical miracles, and explain one after another on the
hypothesis of natural causes; when they resolve a professed vision
into a dream, a possession into epilepsy or madness, a prophecy into a
sagacious conjecture, a recovery into the force of imagination, they
are but expressing their own disbelief in the Grace committed to the
Church; and of course they are consistent in denying its outward
triumphs, when they have no true apprehension of its inward power.
Those, on the other hand, who realize that the bodies of the Saints
were in their lifetime the Temples of the Holiest, and are hereafter
to rise again, will feel no offence at the report of miracles wrought
through them; nor ought those who believe in the existence of evil
spirits to have any difficulty at the notion of demoniacal possession
and exorcism. And it may be taken as a general truth, that where there
is an admission of Catholic doctrines, there no prejudice will exist
against the Ecclesiastical Miracles; while those who disbelieve the
existence among us of the hidden Power, will eagerly avail themselves
of every plea for explaining away its open manifestations. All that
can be objected here is, that miracles which admit of this double
reference to causes natural and supernatural, taken by
themselves and in the first instance, are not evidence of Revealed
Religion; but I have nowhere maintained that they are. Yet, though not
part of the philosophical basis of Christianity, they may be evidence
still to those who admit the Divine Presence in the Church, and in
proportion as they realize it; they may be evidence in combination
with more explicit miracles, or when viewed all together in their
cumulative force; they may confirm or remind of the Apostolic
miracles; they may startle, they may spread an indefinite awe over
certain transactions or doctrines; they may in various ways subserve
the probation of individuals to whom they are addressed, more fully
than occurrences of a more marked character. The mere circumstance
that they do not carry their own explanation with them is no argument
against them, unless we would surrender the most sacred and awful
events of our religion to the unbeliever \footnote{[\emph{Meta tauta pros\={o}popoiei} \emph{Ioudaion aut\={o}i dialegomenon
t\={o}i 'I\={e}sou kai elenchonta auton} … \emph{h\={o}s
plasamenon autou t\={e}n ek parthenou genesin} ... \emph{ph\={e}si
de aut\={e}n kai hypo tou g\={e}mantos, tektonos t\={e}n
techn\={e}n ontos, exe\={o}sthai, elenchtheisan h\={o}s
memoicheum\={e}n\={e}n eita legei, h\={o}s ekbl\={e}theisa
hypo
tou andros, kai plan\={o}men\={e} atim\={o}s skotion egenn\={e}se
ton I\={e}soun}]. Orig. contr. Cels. i. 28.}. As the admission of a Creator is necessary for the
argumentative force of the miracles of Moses or St. Paul, so does the
doctrine of a Divine Presence in the Church supply what is ambiguous
in the miracles of St. Gregory Thaumaturgus or St. Martin.

3.

80.
The course of these remarks has now sufficiently shown that in drawing
out the argument in behalf of ecclesiastical miracles, the main point
to which attention must be paid is the proof of their antecedent
probability \footnote{``Men will be inclined to determine this controverted question
according to their preconceived notions and their accustomed way of
thinking; for there appears to be a sort of fatality in opinions of
this kind, which, when once taken up, are seldom laid down.'' (Jortin,
ibid. p. 24.) Yet he says elsewhere of Theophilus, an Arian
missionary, ``I blame not Tillemont for rejecting all these miracles,
which seem to have been rumours raised and spread to serve a party;
but the true reason of his disbelief is, that they were Arian
miracles; and if they had been reported concerning Athanasius, all
difficulties would have been smoothed over and accounted of small
moment.'' (P. 219.) As if a miracle wrought by Athanasius was not more
likely than miracles wrought by an Arian, though a missionary.}. If that is
established, the task is nearly accomplished. If the miracles alleged
are in harmony with the course of Divine Providence in the world, and
with the analogy of faith as contained in Scripture, if it is possible
to account for them, if they are referable to a known cause or system,
and especially if it can be shown that they are recognized, promised,
or predicted in Scripture, very little positive evidence is necessary
to induce us to listen to them or even accept them, if not one by one,
at least viewed as a collective body. In that case they are but the
natural effects of supernatural agency, and Middleton's canon, which
Douglas, as above quoted, adopts to their disadvantage, becomes
their protection. Then ``the history of miracles,'' instead of being ``suspected
always of course, without the strongest evidence to confirm it,'' is at
first sight almost ``to be admitted of course, without a strong reason
to suspect it;'' such suspicions as attach to it arising from our
actual experience of fraud, not from difficulties in its subject
matter. If ``the tabernacle of God is with men, and He will dwell with
them;'' if the Church is ``the kingdom of heaven;'' if our Lord is with
His disciples ``alway, even unto the end of the world;'' if He promised
His Holy Spirit to be to them what He Himself was when visibly
present, and if miracles were one special token of His Presence when
on earth; if moreover miracles are expressly mentioned as tokens of
the promised Comforter; if St. Paul speaks of ``mighty signs and
wonders by the power of the Spirit of God,'' and of his ``speech and
preaching'' being ``in demonstration of the Spirit and of power,'' and of
``diversities of gifts but the same Spirit,'' and of ``healing,'' ``working
of miracles,'' and ``prophecy,'' as among His gifts; surely we have no
cause to be surprised at hearing supernatural events reported in any
age, and though we may freely exercise our best powers of inquiry and
judgment on such and such reports, as they come before us, yet this is
very different from hearing them with prejudice, and examining them
with contempt or insult.

81.
This Essay, indeed, is not the place for doctrinal discussions: there
is one text, however, to which attention may be drawn, without
deviating into theology, in consequence of what may be called its
historical character, which on other accounts also makes it more to
our purpose,—our Lord's charge to His disciples at the end of St.
Mark's Gospel. It might in truth have been anticipated that, among the
promises with which He animated His desponding disciples when He was
leaving them, some mention would be made of those supernatural powers
which had been the most ready proof of His own divinity, and the most
awful of the endowments with which during His ministry He had invested
them. Nor does He disappoint the expectation; for in the passage
alluded to He distinctly announces a continuation of these pledges of
His favour, and that without fixing the term of it. At the very time
apparently when He said to them, ``Lo, I am with you alway, even unto
the end of the world,'' He also made two announcements, one for this
life, the other for the life to come. ``He that believeth and is
baptized shall be saved,'' was for the future; and the present promise,
which concerns us here, ran thus: ``These signs shall follow them that
believe; In My Name shall they cast out devils; they shall speak with
new tongues; they shall take up serpents; and if they drink any deadly
thing, it shall not hurt them; they shall lay hands on the sick, and
they shall recover.'' Now let us see what presumption is created
or suggested by this passage in behalf of the miraculous passages of
Ecclesiastical History as we have received them.

82.
First, let it be observed, five gifts are here mentioned as specimens
of our Lord's bequest to His disciples on His departure: exorcism,
speaking with new tongues, handling serpents, drinking poison without
harm, and healing the sick. When our Lord first sent out the Apostles
to preach during His ministry, He had specified four: ``Heal the sick,
cleanse the lepers, raise the dead, cast out devils.'' Comparing these
two passages together, we find that two gifts are common to both of
them, and thereby stand out as the most characteristic and prominent
constituents of the supernatural endowment. It is observable, again,
that these two gifts, of which there is this repeated mention, are not
so wonderful or so decisively miraculous as those of which mention
only occurs in one of the two texts. The power of exorcism and of
healing is committed by our Lord to the Apostles, both when He first
calls them, and when He is leaving them; but they are promised the
gift of tongues only on their second mission, and that of raising the
dead only on the first. This does not prove that they could not raise
the dead when our Lord had left them; indeed, we know in matter of
fact that they had, and that they exercised, the power; but it is
natural to suppose that a stress is laid on what is mentioned
twice, and to form thence some idea, in consequence, of the
predominant character of their miraculous endowment, when it was
actually brought into exercise. In accordance with this anticipation,
whatever it is worth in itself, St. Matthew heads his report of our
Lord's charge to His Apostles on their first mission with mention of
these very two gifts, and these only: ``And when He had called unto Him
His twelve Disciples, He gave them power \emph{against unclean spirits,
to cast them out, and to heal all manner of sickness and all manner of
disease}.'' And in like manner when the Seventy are sent, these two
gifts, and these only, are specified by St. Luke as imparted to them;
our Lord saying to them, ``Heal the sick,'' and they answering, ``Lord,
even the devils are subject unto us through Thy Name.''

83.
Further, when we turn to the history of the Book of Acts we find the
general tenor of the Apostles' miracles to be just such as these
passages in the Gospels would lead us to expect; that is, were a Jew
or heathen of the day, who had a fair opportunity of witnessing their
miracles, to be asked what those miracles consisted in, the general
impression left by them on his mind, and the best account which he
could give of them would be, that they were the healing the sick and
casting out devils. We have indeed instances recorded of their raising
the dead, but only two in the whole book, those of Tabitha and
Eutychus; and of these the latter was almost a private act, and
wrought expressly for the comfort of the brethren, not for the
conviction of unbelievers; and though the former was the means of
converting many in the neighbourhood, yet it was wrought at Joppa,
among a number of ``widows'' and ``saints,'' not in Jerusalem, where the
jealous eyes of enemies would have been directed upon it. In the same
book there are three instances of the gift of tongues, at Pentecost,
in Cornelius's house, and at Ephesus on the confirmation of St. John's
disciples. There is one instance of protection from the bite of
serpents, that of St. Paul at Melita. There is no instance of
cleansing leprosy, or of drinking poison without harm. With this
frugality in the display of their highest gifts is singularly
contrasted the bountifulness of the Apostles in exercising their
powers of healing and exorcising. ``They brought forth the \emph{sick}
into the streets, and laid them on beds and couches, that \emph{at the
least the shadow of Peter} passing by might overshadow some of
them. There came also a \emph{multitude} out of the cities \emph{round
about} unto Jerusalem, bringing \emph{sick} folks, and \emph{them that
were vexed with unclean spirits}; and they were \emph{healed every one}.''
Again, when St. Philip went down to Samaria, and ``the people with one
accord gave heed unto those things which Philip spake, hearing and
seeing the miracles which he did,'' what were the particular gifts
which he exercised? the inspired writer continues, ``For \emph{unclean
spirits}, crying with loud voice, came out of \emph{many}
that were possessed with them; and many taken with \emph{palsies}, and
that were \emph{lame}, were healed. And there was great joy in that
city.'' Again, we read of St. Paul, in a later part of the same book,
as has been already quoted in another connection, that ``from his body
were brought unto the sick handkerchiefs or aprons, and the \emph{diseases}
departed from them, and the \emph{evil spirits} went out of them.''
[Acts v. 15, 16; viii. 6-8; xix. 12.]

84.
If there is one other characteristic gift in the Book of Acts in
addition to these, it is the gift of visions and divine intimations.
And, as if to make up for our Lord's silence concerning it in the
Gospel of St. Mark, St. Peter opens the sacred history of the Acts
with a reference to the Prophet Joel's promise of the time, when ``their
sons and their daughters should prophesy, and their young men should
see visions, and their old men should dream dreams;'' an announcement
of which the narrative which follows abundantly records the fulfilment.
St. Stephen sees our Lord before his martyrdom; the Angel directs St.
Philip to go towards Gaza, and the Holy Spirit Himself bids him join
himself to the Ethiopian's chariot; St. Paul is converted by a vision
of our Lord; St. Peter has the vision of the clean and unclean beasts,
and Cornelius is addressed by an Angel; Angels release first the
Apostles, then St. Peter from prison; ``a vision appeared to Paul in
the night, there stood a man of Macedonia;'' at Corinth Christ ``spake
to Paul in the night by a vision, Be not afraid;'' Agabus and St.
Philip's four daughters prophesy; in prison ``the Lord stood by Paul,
and said, Be of good cheer;'' on board ship an Angel stood by him,
saying, ``Fear not, Paul, thou must be brought before Cæsar.'' [Acts
vii. 56; viii. 26, 29; ix. 3-6; x. 3, 10, etc.]

85.
Such is the general character of the miracles of the Book of Acts; and
next let it be observed, such is the character of our Lord's miracles
also, as they would strike the bulk of spectators. He raises indeed
the dead three times, He feeds the multitude in the desert, He
cleanses the leprosy, He gives sight to the blind, on various but
still definite occasions; but how different is the language used by
the Evangelists when His powers of healing and exorcising are spoken
of! We read of ``a great \emph{multitude} of people out of all Judea
and Jerusalem, and from the sea coast of Tyre and Sidon, which came to
hear Him and to be healed of their \emph{diseases}; and they that were
\emph{vexed with unclean spirits}; and they were healed. And the \emph{whole}
multitude sought to \emph{touch} Him; for there went \emph{virtue out of
Him}, and healed them \emph{all}.'' Again, ``\emph{Whithersoever He
entered}, into villages, or cities, or country, they \emph{laid the
sick in the streets}, and besought Him that they might touch if it
were but the corner of His garment; and \emph{as many as touched Him were
made whole}.'' Again, ``They brought unto Him \emph{all sick people}
that were taken with divers \emph{diseases} and \emph{torments}, and
\emph{those that were possessed with devils}, and those that
were \emph{lunatic}, and those that had the \emph{palsy}; and He
healed them.'' [Luke vi. 17-19; Mark vi. 56; Matt. iv. 24.] It may be
added that of other miraculous occurrences in the Gospels none are
more frequent than visions and voices, from the Angel which appeared
to Zacharias to the vision of Angels seen by the women after our Lord's
resurrection; as is obvious without proof.

86.
It appears, then, that the two special powers which were
characteristic, as of our Lord's miraculous working, so also of His
Apostles after Him, were exorcism and healing; and moreover that these
were in matter of fact the two gifts especially promised to the
Apostles above other gifts. It appears, also, that if one other gift
must be selected from the Gospels and Book of Acts as of greater
prominence than the rest it will be the gift of visions; so that
cures, exorcisms, and visions are on the whole the three
distinguishing specimens of Divine Power, by which our Lord
authenticated to the world the Religion He bestowed upon it. Now it
has already been observed \footnote{[N. 30, \emph{supra}.]}
that these are the very three especially claimed by the Primitive
Church; while, as to the more stupendous miracles of raising the dead,
giving sight to the blind, cleansing lepers, and the like, of these
certainly she affords instances also, but very rarely, as if after the
manner of Scripture. This surely is a remarkable coincidence;
and is the rather to be dwelt upon, because those who consider the
vagueness of the language with which the ecclesiastical miracles are
attested, as a proof that they were merely the fabrication of fraud or
credulity, have to explain how it was that, while the parties accused
were exercising their powers of imagination or imposture, they did not
embellish their pages with similar vague statements of miracles of a
more awful character, even from the mere love of variety, instead of
confining themselves to those which in appearance at least were shared
with them by Jews and heathen.

87.
Nor can it reasonably be urged that their acquaintance with Scripture
suggested to them in this matter an imitation of the Divine procedure
as there recorded; because Scripture does not on the face of it
impress upon the reader the fact which has been here pointed out. The
actual course of the \emph{events} which Scripture relates is one thing, and
the course of the \emph{narrative} is another; for the sacred writers do not
state events with that relative prominence in which they severally
occurred in fact. Inspiration has interfered to select and bring into
the foreground the most cogent instances of Divine interposition, and
has identified them by a number of distinct details; on the other
hand, it has covered up from us the ``many other signs'' which ``Jesus
did in the presence of His disciples,'' ``the which, if they should be
written every one, even the world itself,'' as St. John speaks, ``could
not contain the books that should be written.'' And doubtless there are
doctrinal reasons also for this circumstance, if we had means of
ascertaining them. But so it is, that the \emph{primâ facie}
appearance of the Gospel Miracles does not so correspond to that of
the Ecclesiastical Miracles, as probably it would have corresponded,
had St. John, for instance, given us a description of the second and
third centuries, instead of St. Justin and Origen, or had Sulpicius
described the Miracles of the Apostles at Jerusalem or Ephesus.

4.

88.
And now, if this representation has any truth in it, if our Lord, in
the passage of St. Mark in question, promised five gifts to His
disciples, two of which were those of exorcism and healing; if these
same two, distinguished in other places of the Gospels above the rest,
are the prominent external signs of power in the history both of our
Lord and of His Apostles; if these particular Miracles are the special
instruments of the conversion of whole multitudes; if on account of
the cures and exorcisms wrought by the twelve Apostles ``believers were
the more added to the Lord, multitudes both of men and women;'' if on
St. Philip's casting out devils, and curing palsy and lameness, ``the
people with one accord gave heed,'' and ``there was great joy in that
city;'' if when an evil spirit had confessed, ``Jesus I know, and
Paul I know, but who are ye?'' ``fear fell on them all,'' and ``the Name
of the Lord Jesus was magnified,'' and ``the word of God grew mightily
and prevailed;'' what is to be said of those modern Apologists for
Christianity who do their best to prove that these phenomena have
nothing necessarily miraculous in them at all? So much is evident at
once, that had they been the persons encountered by these miracles of
the Apostles, had they been the Samaritans to whom St. Philip came, or
the Ephesians who were addressed by St. Paul, they would have thought
it their duty to have felt neither ``much joy'' with the one, nor ``fear''
with the other; and that, if Samaritans and Ephesians had acted on the
modern view of what is rational and what is evidence, what sound
judgment and what credulity, Christianity would not have made way and
prospered, but we all should have been heathen at this day.

89.
Bishop Douglas, for instance, observes, that the circumstance that the
Fathers allow that ``cures of diseases, particularly of demoniacs by
exorcising them,'' ``were exercised by pagans with the assistance of
their demons and gods,'' and admit that ``there were exorcists among the
Jews and Gentiles, who by the use of certain forms of words, used as
charms, and by the practice of certain rites, cast out devils, as well
as the Christian exorcists,'' that this circumstance ``some may think
puts these feats of jugglers and impostors upon the same footing
of credibility with the works ascribed to Christians'' \footnote{Page 233-236.}:—why not with the works ascribed to Apostles? Again he
urges, that ``the cures ascribed to the prayers of Christians, to the
imposition of their hands, etc., in those early times, \emph{might, for
aught we know}, be really brought about in a \emph{natural} way,
and be accounted for in the same way in which we have accounted for
those ascribed to the Abbé Paris, and those attributed by the
superstitious Papists to the intercession of the Saints'':—perhaps
the acute unbelievers of Corinth or Ephesus by a parallel argument
justified their rejection of St. Paul. At Ephesus, when the demoniac
leapt on the Jewish exorcists, ``and overcame them, and prevailed
against them, so that they fled out of that house naked and wounded,'' ``fear''
in consequence ``fell on all the Jews and Greeks also dwelling at
Ephesus;'' but the Bishop would have taught them that ``a few grimaces,
wild gestures, disordered agitations, and blasphemous exclamations,
suited to the character of the supposed infernal inhabitants,
constitute all we know of their disease; and consequently, as \emph{all}
these symptoms are \emph{ambiguous}, and may be assumed at pleasure by
an impostor, a collusion between the exorcist and the person exorcised
will account for the whole transaction, and every one who would avoid
the character of being superstitiously credulous will naturally
account for it in this manner, rather than by supposing that any
supernatural cause intervened.'' \footnote{Page 146. Douglas is speaking here primarily of the Church of Rome;
afterwards he apparently refers to the passage when speaking of the
Primitive Church, p. 236.} Such is this author's judgment of one of the two exhibitions
of miraculous power with which our Saviour specially and singularly
gifted His Apostles, and by which they, in matter of fact, converted
the world. The question is not, whether in \emph{particular cases} its
apparent exercise may not be suspicious and inconclusive, for Douglas
is speaking against the gift as such; so that a heathen of Ephesus
would have been justified on his principles in demanding of St. Paul
to see a man raised from the dead, before he believed in Christ. And
such was the nature of the demand made by Autolycus upon St.
Theophilus at the end of the second century, and Middleton and Gibbon
justify it, and seem moreover to consider the mere silence of
Theophilus to be a proof that such a miracle was utterly unknown in
his days, as if resurrections abounded in the Acts \footnote{Defecere etiam mortuorum excitationes. Certe Autolyco roganti ut vel
unum ostenderet qui fuisset è mortuis revocatus, ita respondit
Theophilus quasi vel unum demonstrare minime potuerit. Dodw. in Iren.
Dissert. ii. 44. Jortin is more cautious. ``It is probable,'' he says, ``from
his [Theophilus's] silence, that he had heard of no instance of such a
miracle in his days; probable, I say, but not certain; because, though
he had heard of it, he might possibly have thought it to no purpose to
tell his friend that there were Christians who \emph{affirmed} such
things, and he might suspect that Autolycus \emph{would not have admitted
the testimony} of persons with whom he had no acquaintance, and for
whom he had little regard.'' Eccl. Hist. (Works, vol. ii. p. 92, ed.
1810). Vid. the striking statement of Origen. contr. Cels. i. 46.
Greg. Nyss. tom. ii. p. 1009.}.

90.
Again, St. Peter cured Æneas of the palsy, ``and all that dwelt
at Lydda and Saron saw him, and turned to the Lord;'' but the Bishop
would have advised them to wait till they had seen Tabitha
raised; because ``palsies, it is well known, arise from obstructions of
the spirits that circulate in the nerves, so that their influx into
the muscles is impeded, or from obstructions of the arterious blood.
Nothing more, therefore, was required here than to remove that
obstruction.'' \footnote{Page 82.}

91.
We read in Scripture of the sudden cure of the dropsy; but the Bishop
observes, ``That enthusiasm should warm its votaries to a holy madness,
and excite the wildest transports and agitations throughout their
whole frame, is an effect which, in a country so fruitful of this
production as is ours (though enthusiasm be the product of \emph{every}
soil and of \emph{every} religion), must be consistent with the
experience of many.'' \footnote{Page 104.}
Then he adds, speaking of some particular cases: ``As one of the
curative indications of a dropsy is an evacuation of the water by
perspiration, and as the medicines administered by the physician aim
to produce this effect, ... what could be more likely to excite
such copious perspiration than the enthusiastic transport with which
they prayed, and the convulsive struggles which shook their whole
frame?'' \footnote{Page 107.}

92.
Peter's wife's mother was raised from her fever at once, so as even to
be able to ``minister'' to the holy company; but Bishop Douglas would
have suggested to the Pharisees that, had there been more raising of
the dead, more restoring of sight to the blind, such cures might have
been dispensed with, because, where minds are ``heated and inflamed,
and every faculty of their souls burning with the raptures of devout
joy and enthusiastic confidence,'' it is ``far from being impossible ...
that in some cases a change might be wrought on the habit of the
body;'' \footnote{Page 102.} for ``in
this case the nervous system is strongly acted upon, and fresh and
violent motions are communicated to the fluids;'' \footnote{Page 106.} and ``such agitations necessarily suppose that the velocity of
the fluids'' is ``greatly accelerated;'' \footnote{Ibid.} and ``gouts, palsies, \emph{fevers} of all kinds, and even
ruptures, have been thus cured.'' \footnote{Page 101.} It certainly does not appear why a class of miracles which
was, in matter of fact, the principal means of the conversion of the
world in the age of the Apostles, should, when professed in the second
and third centuries, be put aside by our Apologists on the
excuse that ``powers were not appealed to, less ambiguous in their
nature,'' nor ``other works performed, which admit of no solution from
natural causes, and were incapable of being the effects of fraud and
collusion.'' \footnote{Page 236.}

93.
This being the language of so respectable a writer as Bishop Douglas,
the following sentiments from Middleton cannot surprise us. Of
miracles of healing he says: ``In truth, this particular claim of
curing diseases miraculously affords great room for … delusion and a
wide field for the exercise of \emph{craft}. Every man's experience
has taught him that diseases thought fatal and desperate are oft
surprisingly healed of themselves, by some secret and sudden effort of
nature impenetrable to the skill of man; but to ascribe this presently
to a miracle, as weak and superstitious minds are apt to do, to the
prayers of the living or the intercessions of the dead, is what
neither sound reason nor true religion will justify.'' \footnote{Page 79.} Of exorcisms: that certain circumstances ``concerning the
speeches and confessions of the devils, their answering to all
questions, owning themselves to be wicked spirits, etc., … may not
improbably be accounted for, either by the disordered state of the
patient, answering wildly and at random to any questions proposed, or
by the \emph{arts of imposture and contrivance} between the parties
concerned in the act.'' \footnote{Page 82.} And of visions: ``To declare freely what I think, whatever
ground there might be in those primitive ages either to reject or to
allow the authority of those visions, yet, from all the accounts of
them that remain to us in these days, there seems to be the greatest
reason to suspect that they were \emph{all contrived}, \emph{or
authorized at least, by the leading men of the Church}.'' \footnote{Page 109.}

94.
Such, then, is the opinion of Christian Apologists concerning the
nature of those miracles to which our Lord mainly entrusted the cause
of His sacred truth; for, however great the differences may be between
the Scripture and Ecclesiastical miracles, viewed as a whole, so far
is certain, that the actual and immediate instruments by which the
world was convinced of the Gospel were those which these writers
distinctly discredit as of an ambiguous and suspicious character. And,
if it be asked whether, after all, such miracles are not suspicious,
whatever be the consequence of admitting it, I answer, that they are
suspicious to read of, but not to see. The particular circumstances of
an exorcism, which no narrative can convey, might bring home to the
mind a conviction that it was a divine work, quite sufficient for
conversion; and much more a number of such awful exhibitions.
Generalized statements and abstract arguments are poor representations
of fact; but, as they are used to serve the purpose of those who
would disparage Saints, it is necessary to show that they can be
turned by unbelievers as plausibly, though as sophistically, against
Apostles.

5.

95.
To proceed. The same words of our Saviour which have introduced these
remarks in defence of the nature of the ecclesiastical gifts will
suggest an explanation of certain difficulties in the mode of their
exercise. Christ says, first, ``He that believeth shall be saved;'' and
then, ``These signs shall follow them that believe.'' Here it is obvious
to remark, that the power of working miracles is not promised in these
words to the \emph{preachers} of the Gospel merely, but to the \emph{converts}
\footnote{``Nec enim prædicantes illa secutura signa pollicetur, sed credentes;
nec eos qui jam antea credidissent, sed qui essent postea deinde
credituri. Responditque eventus accuratissimè; conversis enim, non
conversoribus, gratias illas donatas esse constat de quibus legimus in
primis Ecclesiarum conversionibus.'' Dodw. in Iren. Dissert. ii. 28.
This is so fully taken for granted by St. Bernard, that he thinks it
necessary to answer the objection why ``credentes'' did not work
miracles in his day: ``Quis enim ea, quæ in præsenti loco scripta
sunt, signa videtur habere credulitatis, sine quâ nemo poterit
salvari? quoniam qui non crediderit condemnabitur, et sine fide
impossibile est placere Deo.'' Serm. i. de Ascens. 2. He answers to the
question as St. Gregory does in the passage quoted, \emph{supra}, n.
40, making the miracles now wrought by the faithful to be moral ones.
Kuinoel says: ``Per [\emph{tous pisteuontas}] \emph{non omnes} Christi
sectatores intelligendi sunt, nam non omnes Christiani ejusmodi
miracula patrabant, qualia hoc loco describuntur, sed agit Christus
hoc loco ut locis parallelis, Luc. 24, 48. John 20, 19, \emph{cum legatis
suis}, atque adeo significantur imprimis Apostoli, et præter eos
alii tunc temporis præsentes, qui haud dubiè è numero septuaginta
discipulorum erant. Vid. Luc. 24, 33, coll. Luc. 10, 1; 9, 17, Etiam
infra v. 20, disertè commentorantur [\emph{ekeinoi}], illi Christi
discipuli, quibus ea dixit, quæ hoc loco leguntur, et ad hos [\emph{s\={e}meia}]
referuntur. Monuit præterea Storrius articulum [\emph{tous}] sæpe \emph{certos},
\emph{quosdam}, non omnes universos significare. Vid. Luc. 18, 15.
Coll. Marc. 10, 13. Matt. 21, 34. 36. 27, 62. 28, 12. Insignivit autem,
ut opinor, Christus discipulos suos, futuros religionis suæ doctores,
tunc temporis præsentes, voce [\emph{tois pisteusasi}], quoniam paulo
ante eorum incredulitatem vituperarat:'' in loc. This is such strange
reasoning, that it is the best argument for showing how futile the
attempt is to wrest our Lord's words from their plain meaning. The
elder school of Protestants was more candid. ``Non omnibus omnia,'' says
Grotius, ``ita tamen ut cuilibet, ut oportet, credenti, aliqua tunc
data sit admirabilis facultas, quæ se non semper quidem, sed datâ
occasione, explicaret.''}. It is not said, ``Preach
the Gospel to every creature, and these signs shall follow your
preaching,'' but ``these signs shall follow them that believe,'' the same
persons to whom salvation is promised in the verse preceding \footnote{Sulpicius almost grounds his defence of St. Martin's miracles on the
antecedent force of this text. He says of those who deny them, ``Nec
Martino in hac parte detrahitur, sed fidei Evangelii derogatur. Nam
cum Dominus ipse testatus sit istiusmodi opera, quæ Martinus implevit,
ab omnibus fidelibus esse facienda, qui Martinum non credit ista
fecisse, non credit Christum ista dixisse.'' Dial. i. 18.}. And further, whereas final salvation is there represented as
a personal gift, the gift of miracles is not granted here to ``\emph{him}
that believeth,'' but to ``\emph{them} that believe.'' And the particular word used, which the Authorized Version translates
``follow,'' suggests or encourages the notion that the miracles
promised were to \emph{attend upon} or to be \emph{collateral with}
their faith, as general indications and tokens \footnote{[\emph{S\={e}meia tauta parakolouth\={e}sei}]. ``Stephanus in
Thes. hæc citat ex Dioscoride in præf. lib. 6. [\emph{ta
parakolouthounta semeia hekast\={o}i t\={o}n pharmak\={o}n}].''
Raphel. Annot. in loc. Vid. ibid. in Luc. i. 3. In the last words of
the Gospel, where the ``signs following'' are wrought by the \emph{Apostles},
and in \emph{confirmation} the word is [\emph{epakolouthount\={o}n}].}; not indeed that they were to be the result of every act of
faith and in every person, but that on the whole, where men were
united together by faith in the name of Christ, there miracles would
also be wrought by Him who was ``in the midst of them.'' Thus the gift
was rather in the Church than of the Church.

96.
An important text already quoted teaches us the same thing: ``I will
pour out My Spirit upon all flesh; and your sons and your daughters
shall prophesy, your old men shall dream dreams, your young men shall
see visions: and also upon the servants and upon the handmaids in
those days will I pour out My Spirit.'' The young, the old, the bond
and the free, all flesh, all conditions of men, were to be the
recipients of the miraculous illuminations of the Gospel. The event
exactly accomplished the prediction. In the very opening of the New
Dispensation, not only Zacharias the Priest, but Mary the young maiden, Elizabeth the matron, Anna the widow of fourscore and four
years, and just and aged Simeon, were inspired to bear witness to it.
Again, in the Book of Acts, while Peter was preaching to Cornelius, ``the
Holy Ghost fell on \emph{all} them which heard the word.'' At Ephesus,
when St. Paul had laid his hands on John's disciples, the Holy Ghost
came on them, ``and all the men were about twelve.'' Moreover, we hear
of St. Philip's ``four daughters, virgins, which did prophesy.'' And the
disorders of the Church of Corinth plainly show that the miraculous
gifts were not confined to one or two principal persons of high
station or spiritual attainments, but were ``dispersed abroad'' with a
bountiful hand over all the faithful. The same inference may be drawn
from St. Peter's direction, ``As every man hath received the gift, even
so minister the same one to another, as good stewards of the manifold
grace of God.'' Such, then, is the Scripture account of the bestowal of
the miraculous powers in the Apostolic age; and, I repeat, it serves
to remove certain misapprehensions and objections which have been made
to their exhibition as instanced in the times that follow.

97.
For instance, there seems a fallacy in the mode in which a phrase is
used, which often occurs in the controversy. It has been contended
that there is no ``standing gift of miracles'' in the Church; and
then it is concluded that \emph{therefore} no manifestation of Divine Power takes place in it, but those rare and solemn
interpositions which we have reason to think actually occur even in
heathen countries. ``The position which I affirm,'' says Middleton, ``is
that, after the days of the Apostles, no \emph{standing power} of
miracles was continued \emph{to} the Church, to which they might
perpetually appeal for the conviction of unbelievers. Yet all my
antagonists treat my argument as if it absolutely rejected \emph{everything}
of a miraculous kind, whether wrought within the Church by the agency
of men, or on any other occasion by the immediate hand of God.'' \footnote{Vindic. p. 32, as quoted by Douglas, p. 224.} Now, there is an ambiguity in the words ``standing power,''
according as we take it to mean a capacity committed \emph{to}
particular persons and exercised \emph{by} them, or a Divine Agency
generally operating \emph{in} the Church and \emph{among} Christians,
as its Almighty Author wills. Middleton denies the standing power in
its former sense; but in our Lord's promise, as well as in St. Paul's
description of the presence of the Holy Spirit in the Church, the
latter is the prominent idea. Middleton speaks, just after the passage
above quoted, of ``the Church \emph{having} no standing power of \emph{working}''
miracles, and elsewhere of a ``standing power of \emph{working}
miracles, as \emph{exerted openly} in the Church, \emph{for} the
conviction of unbelievers.'' \footnote{Inquiry, p. 9.}
Again, he speaks of the ``opinion that after the days of the
Apostles there resided still in the primitive Church, through several
successive ages, a divine and extraordinary \emph{power of working
miracles}, which was frequently and openly \emph{exerted}, in
confirmation of the truth of the Gospel, and for the conviction of
unbelievers.'' \footnote{Introd. Disc. init.; but in Pref. p. xxxii. he speaks more to the
purpose.} In like
manner Douglas says of Middleton that ``his Free Inquiry is not,
whether any miracles were performed after the times of the Apostles,
but whether, after that period, \emph{miraculous powers subsisted} in
the Church; not whether God interposed at all, but whether He
interposed by making use of men as His instruments.'' \footnote{Page 224.} Here he makes ``the subsistence of miraculous powers''
equivalent with ``the instrumentality of men in their operation;''
meaning by the latter the conscious exercise of them by inspired
persons in proof of a divine mission, as a former passage of his work
shows \footnote{Page 216.}. The present
Bishop of Lincoln (Kaye) takes the same view of the controversy,
observing that Middleton's object ``was to prove that, after the
Apostolic age, no \emph{standing power} of working miracles existed in
the Church, that there was \emph{no regular succession} of favoured \emph{individuals}
upon whom God conferred supernatural powers, \emph{which they could
exercise} for the benefit of the Church of Christ, whenever their \emph{judgment},
guided by the influence of the Holy Spirit, told them that it
was expedient so to do.'' \footnote{Kaye's Tertull. p. 104.}
Certainly, if this is what Middleton set about to do, he had not a
difficult task before him.

98.
Yet Lord Barrington, before Middleton, had implied that the question
lay between the same two issues. ``There cannot be much doubt,'' he
says, ``of these gifts lasting as much longer as the oldest of those
lived to whom St. John imparted them ... Irenæus, speaking of the
prophetic gifts, mentions the gift of tongues and the discernment of
spirits. And that these did not last longer seems to have been the
case in fact, since Irenæus, who died about the year 190, in a very
old age, speaks of his having \emph{seen} these gifts, but says
nothing of \emph{his own} having them.'' \footnote{Vol. i. pp. 221, 222, ed. 1828.} That is, Barrington makes no medium between a definite
transmission of the gift from Christian to Christian by imposition of
hands or similar formal act, such as would involve Irenæus's own
possession of it, and on the other hand its having utterly failed.
Irenæus saw the gift, he had it not, therefore it was failing in his
time; else he would have had it.

99.
What ecclesiastical history rather inculcates is the doctrine of an
abiding presence of Divinity such as dwelt upon the Ark, showing
itself as it would, and when it would, and without fixed rules; which
was seated primarily in the body of Christians, and manifested
itself sometimes in persons, sometimes in places, as the case might
be, in saintly men, or in ``babes and sucklings,'' or in the very stones
of the Temple; which for a while was latent, and then became manifest
again; which to some persons, places, or generations was an evidence,
and to others was not \footnote{Dodwell has a theory (which agrees with what is said in the text,
except that he applies it only to the first ages) that miracles
abounded or became scarce according to the need, the conversion of the
nations being the chief object. ``Promisit Dominus majora editurum, qui
in illum postea crediderit, miracula quam quæ ipse Dominus ediderit.
Quod ego facile moderandum esse concessero, ut et de certis Evangelii
propagandi temporibus promissio illa fuerit intelligenda ... Sed nec
ita adimpleta est quin superesset adhuc satis amplus locus futuris
postea conversionibus, futurisque adeo miraculis … Trajano Imperante
novas Evangelii propagandi causâ susceptas expeditiones memorat
Eusebius, et quidem id novâ Dei comitante gratiâ atque [\emph{synergeiai}]
... Ortis jam sub Hadriano Hæreticis, ... factum est ut miracula
infidelium hæreticorum causâ præstanda fuerint etiam et ipsa
frequentiora ... A Marci temporibus deficere cœperunt, … cum nullas
aut raras admodum per ea sæcula expeditiones obirent Christiani ad
gentes ex professo convertendas; ... satis tamen liberalem adhuc
fuisse Deum multa ostendunt,'' etc., etc. Dissert. in Iren. ii. 28-45,
etc.}.
The ideas of ``regular succession,'' conscious ``exercise'' of power,
objects deliberately contemplated, discretionary use of a gift, and
the like, are quite foreign to a theory of miraculous agency of this
kind; yet, at the same time, it cannot surely be denied that in one
sense such an appointment may rightly be called a ``standing
power,'' and that it is very much more than such rare ``interpositions
of Providence,'' and such ``miracles of invisible agency,'' as the above
writers seem to consider the only alternative to the admission of a
discretionary, and conscious, and transmitted gift.

100.
The Ark was a standing instrument of miraculous operation, yet it did
not send forth its virtue at all times, nor at the will of man. What
was the nature of its mysterious powers we learn from the beginning of
the First Book of Samuel; where we read of it first as stationed in
the tabernacle, and of the Almighty speaking from it to the child
Samuel; next it is captured in battle by the Philistines; but next,
when it is set up in the house of Dagon, the idol, without visible
cause, falls down before it, and its worshippers are smitten. Next,
the cattle which are yoked to it are constrained against their natural
instinct to carry it back to Israel. And then the men of Bethshemeth
are smitten for looking into it. Was there, or was there not, then, a
standing power of miracles in the Jewish Church? There was not, in the
sense in which Middleton understands the phrase; there was no ``regular
succession'' of ``individuals'' who exercised supernatural gifts with a
divinely enlightened discretion; even the Prophets were not such a
body; yet the Divine Presence consisted in much more than an
occasional and extraordinary visitation or intervention in the
course of events. That such too should be the nature of the Presence
in the Christian Church is at least quite consistent with the tenor of
the new Testament; and is almost implied when, in the text which has
given rise to these remarks, our Lord bestows its miraculous
manifestations upon the body at large. The supernatural glory might
abide, and yet be manifold, variable, uncertain, inscrutable,
uncontrollable, like the natural atmosphere; dispensing gleams,
shadows, traces of Almighty Power, but giving no such clear and
perfect vision of it as one might gaze upon and record distinctly in
its details for controversial purposes. Thus we are told, ``The wind
bloweth where it listeth;'' ``a little while, and ye shall not see Me:
and again, a little while, and ye shall see Me;'' ``their eyes were
holden,'' and ``they knew Him, and He vanished;'' ``suddenly there came a
sound from heaven;'' when they had prayed, the place was shaken where
they were assembled together; ``all these worketh that One and the
selfsame Spirit, dividing to every man severally as He will.'' At one
time our Lord connects the gift with special holiness, as when He says
that certain exorcisms require ``prayer and fasting;'' at another He
allows it to the reprobate, as when He says that those whom He never
knew will in the last day appeal to the wonderful works they did in
His Name. At one time St. Paul, in evidence of his divine mission,
says, ``Truly the signs of an Apostle were wrought among you;'' at
another he seems to ascribe the power to an imposture: ``Though an
Angel from heaven preach any other Gospel unto you, let him be
accursed.''

6.

101.
Another difficulty which the text in question enables us to meet is
the indiscriminate bestowal of the miraculous gift, as we read of it
in ecclesiastical history. Its being in the Church, not of the Church,
implies this apparent disorder and want of method in its
manifestations, as has been already observed. Yet Middleton objects,
speaking of the Fathers, ``None of these venerable Saints have anywhere
affirmed, that either they themselves, or the Apostolic Fathers before
them, were endued with any power of working miracles, but declare only
\emph{in general}, that such powers \emph{were actually subsisting} in
their days and openly existed in the Church; that they had often seen
the wonderful effects of them; and that everybody else might see the
same, whenever they pleased; but as to the persons who wrought them,
they leave us strangely in the dark; for instead of specifying names, conditions,
and characters, their general style is, Such and such works are done
among us or by us; by our people; by a few; by many; by our exorcists;
by ignorant laymen, women, boys, and any simple Christian whatsoever.''
\footnote{Page 22.} That is, his objection
is against the very idea of a gift, committed to the body of the
Church, or abiding in the Church. Objectors are hard to please;
sometimes they imply dislike of the notion of the gift as delegated to
a ministerial succession, and formally transmitted from individual to
individual, and then, on the contrary, of its belonging to the Church
itself without the intervention of rites of appropriation or definite
recipients: what is this but saying that they will not entertain the
notion of a continuance of miracles at all? As to Middleton's
objection, it seems directed against the prophetic anticipation of the
times of the Gospel made to the Jews, as quoted already, that ``their
sons and daughters should prophesy, their young men see visions, and
their old men dream dreams,'' quite as much as against any seeming
incongruities and anomalies which are found in the early Church.

102.
Middleton's complaint, that the Fathers do not themselves profess a
miraculous gift, is echoed by Gibbon. ``It may seem somewhat
remarkable,'' he says, ``that Bernard of Clairvaux, who records so many
miracles of his friend St. Malachi, never takes any notice of his own,
which, in their turn, however, are carefully related by his companions
and disciples. In the long series of ecclesiastical history, does
there exist a single instance of a Saint asserting that he himself
possessed the gift of miracles?'' \footnote{Ch. xv. note s.} The concluding question concerns our present subject,
though St. Bernard himself is far removed from the period of history
on which we are engaged. I observe then, first, that it is not often
that the \emph{gift} of miracles is even ascribed to a Saint \footnote{``Hoc intercedit discrimen inter sanctos antiqui et Novi Testamenti, quòd
Deus, intercessione Sanctorum V. T., miracula operari dignatus est sæpius
in vitâ, et rarius post obitum corum; et quoad Sanctos N. T. sæpius
post obitum et rarius in vitâ ipsorum; cùm Sancti V. T. utpote a Deo
ipso canonizati, miraculis post obitum non indigerent; sancti autem N.
T. ab Ecclesiâ canonizandi, miraculis post obitum indigeant ... Cùm
nulla [S. Joannes B.] in vitâ miracula fecisset, putavit Herodes eum
post suam in Christo resurrectionem miracula fuisse editurum, 'Ait
pueris suis, Hic est Joannes Baptista, etc., et ideo virtutes
operantur in eo.''' Bened. xiv de Canon. Sanct. iv. i. § 26.}. In many cases miracles are only ascribed to their tombs or
relics; or when miracles are ascribed to them when living, these are
but single and occasional, not parts of a series. Moreover, they are
commonly what Paley calls \emph{tentative} miracles, or some out of
many which have been attempted, and have been done accordingly without
any previous confidence in their power to effect them. Moses and
Elijah could predict the result; but the miracles in question were
scarcely more than experiments and trials, even though success had
been granted them many times before \footnote{The present Bishop of London argues from Origen's expression, [\emph{ous
ho theos bouletai}], (Contr. Cels. ii. 33), ``that the attempts,
which no doubt were made to effect miraculous cures, were not always
successful;'' vid. Athan. Vit, Ant. 56, where this very thing is
confessed: then he continues: ``and \emph{if so}, we may \emph{safely infer}
that where they \emph{did} succeed, they were to be ascribed to the \emph{ordinary}
means of healing under the Divine blessing.'' Bishop Blomfield's
Sermons, p. 434. I cannot follow his Lordship in calling this
inference a safe one.}. Under these circumstances, how could the individual men
who wrought them appeal to them themselves? It was not till
afterwards, when their friends and disciples could calmly look back
upon their life, and review the various actions and providences which
occurred in the course of it, that they would be able to put together
the scattered tokens of Divine favour, none or few of which might in
themselves be a certain evidence of a miraculous power. As well might
we expect men in their lifetime to be called Saints, as workers of
miracles. But this is not all; the objection serves to suggest a very
observable distinction, which holds good between the conduct of those
whose miracles are designed to be evidence of the truth of religion,
and that of others though similarly gifted. The Apostles, for
instance, did their miracles openly, because these were intended to be
instruments of conversion; but when the supernatural Power took up its
abode in the Church, and manifested itself as it would, and not for
definite objects which it signified at the time of its manifestation,
it could not but seem to imply some personal privilege, when operating
in an individual, who would in consequence be as little inclined to
proclaim it aloud as to make a boast of his graces.

7.

103.
The same peculiarity in the gift will also account for that deficiency
in the evidence, and other unsatisfactory circumstances of a like
nature, which have already been spoken of. Since the Divine
manifestation was arbitrary, the testimony would necessarily be
casual. What else could be expected in the case of occurrences of
which there was no notice beforehand, and often no trace after, and
where we are obliged to be contented with such witnesses as happened
to be present, or, if they cannot be found, with the mere report which
has circulated from them? and when perhaps, as was noticed in the last
paragraph, the principal parties felt it to be wrong to court
publicity, after our Lord's pattern, and perhaps shrank from
examination? ``There is no man,'' said His brethren to Him, ``that doeth
anything in secret, and he himself seeketh to be known openly; if Thou
do these things, show Thyself to the world.'' In our Lord's own case
there was a time for concealment and a time for display; and, as it
was a time for evidence when miracles were wrought by the Apostles, so
afterwards there was a time for other objects and other uses, when
miracles were wrought through the Church; and as our Lord's miracles
were true, though the Jews complained that He ``made them so long to
doubt,'' so it is no disproof of the miracles of the Church, that those
who do not wish them true have room to criticise the character
or the matter of the testimony which at this day is offered in their
behalf.

8.

104.
One more remark is in point. Middleton, in the extract above quoted,
finds fault with the Fathers for ``declaring only in \emph{general}''
that miracles continued, that they had seen them themselves, and that
any one else might see them who would, while they made no attempt to
specify the names, conditions, and characters of the persons working
them. Yet surely this is but natural, if such miracles were as
frequent as ecclesiastical history represents. Instead of its being an
objection to them, it is just the state or things which must
necessarily follow, supposing they were such and so wrought as is
described. When we are speaking of what is obvious, and allowed on all
hands, we do not go about to prove it. We only argue when there is
doubt; we only consult documents, and weigh evidence, and draw out
proofs, when we are not eye-witnesses. If the Fathers had seen
miracles of healing or exorcisms not unfrequently, and were writing to
others who had seen the like, they would use the confident yet vague
language which we actually find in their accounts. The state of the
testimony is but in keeping with the alleged facts.

105. For instance,
St. Justin speaks of the Incarnation as having taken place ``for
the sake of believers, and for the overthrow of evil spirits;'' and ``\emph{you
may know this now},'' he continues, ``\emph{from what passes before your
eyes}; for many demoniacs all over the world, and in your own
metropolis, whom none other exorcists, conjurers, or sorcerers have
cured, these have many of our Christians cured, adjuring by the Name
of Christ, and still do cure.'' Again: ``With us even hitherto are
prophetical gifts, \emph{from which} you Jews ought to gather that
what formerly belonged to your race is transferred to us;'' and soon
after, quoting the passage from the prophet Joel, he adds, ``and with
us \emph{may be seen} females and males with gifts from the Spirit of
God.'' And St. Irenæus: ``In His Name His true disciples, receiving the
grace from Himself, work for the benefit of other men, as each has
received the gift from Him. For some cast out devils certainly and
truly, so that oftentimes the cleansed persons themselves become
believers, and join the Church. Others have foreknowledge of things
future, visions, and prophetical announcements. Others by imposition
of hands heal the sick, and restore them to health. Moreover, as I
have said, before now even the dead have been restored to life, and
have continued with us for many years. \emph{Indeed, it is not possible
to tell the number} of gifts which the Church throughout the world
has received from God in the Name of Christ Jesus, who was crucified
under Pontius Pilate, and exercises day by day for the benefit
of the nations, neither seducing nor taking money of any.'' Shortly
before he observes, that the heretics could not raise the dead, ``as
our Lord did, and the Apostles by prayer; and in the brotherhood \emph{frequently}
for some necessary object, (the whole Church in the place asking it
with much fasting and supplication,) the spirit of the dead has
returned, and the man has been granted to the prayers of the Saints.''
And again, he speaks of his ``\emph{hearing} many brothers in the
Church who had prophetical gifts, and spoke by the Spirit in all
tongues, and brought to light the hidden things of men for a
profitable purpose, and related the mysteries of God.'' And in like
manner Tertullian: ``\emph{Place some possessed person before your
tribunals}; any Christian shall command that spirit to speak, who
shall as surely confess himself to be a devil with truth, as elsewhere
he will call himself a god with falsehood ... What work can be
clearer? ... there will be no room for suspicion; you would say that
it is magic, or some other deceit, if your eyes and ears allowed you,
for what is there to urge against that which is proved by its naked
sincerity?'' Again Origen speaks of persons healing, ``with no
invocation over those who need a cure, but that of the God of all and
the Name of Jesus, with some narrative concerning Him. By these,'' he
adds, ``we, too, \emph{have seen many} set free from severe complaints,
and loss of mind, and madness, and numberless other such evils,
which neither men nor devils had cured.'' \footnote{Justin, Apol. ii. 6. Tryph. 82, 88. Iren. Hær. ii. 32, § 4, 31, §
2, v. 6, § 1. Tertull. Apol. 23. Origen, contr. Cels. iii. 24. Vid.
also Justin, Apol. 1, 40. Tryph. 30, 39, 76, and 85. Tertull. Apol.
37, 43. Scorp. 1. Test. Anim. 3 Ad. Scap. 4. Minuc. F. 27. Theoph. ad
Autol. ii. 8. Origen, contr. Cels. i. 46, 67, ii. 33, iii. 36.
Cyprian, Ep. 76, fin. ad. Magn. circa fin. vid. supr. n. 32. [Vid.
also note and passages in Murdoch's Mosheim, t. i. p. 128.]}

106.
This is the very language which we are accustomed to use, when facts
are so notorious that the \emph{onus dubitandi} may fairly be thrown
upon those who question them. All that can be said is, that the facts
are not notorious \emph{to us}; certainly not, but the Fathers wrote
for contemporaries, not for the eighteenth or nineteenth century, not
for modern notions and theories, for distant countries, for a
degenerate people and a disunited Church. They did not foresee that
evidence would become a science, that doubt would be thought a merit,
and disbelief a privilege; that it would be in favour and
condescension to them if they were credited, and in charity that they
were accounted honest. They did not feel that man was so
self-sufficient, and so happy in his prospects for the future, that he
might reasonably sit at home closing his ears to all reports of Divine
interpositions till they were actually brought before his eyes, and
faith was superseded by sense; they did not so disparage the Spouse of
Christ as to imagine that she could be accounted by professing
Christians a school of error, and a workshop of fraud and imposture.
They wrote with the confidence that they were Christians, and that
those to whom they transmitted the Gospel would not call them the
ministers of Antichrist.

\xchapter{Chapter 5. On the Evidence for Particular Alleged Miracles}

107.
IT does not strictly fall within
the scope of this Essay to pronounce upon the truth or falsehood of
this or that miraculous narrative as it occurs in Ecclesiastical
History; but only to furnish such general considerations as may be
useful in forming a decision in particular cases. Yet considering the
painful perplexity which many feel when left entirely to their own
judgments in important matters, it may be allowable to go a step
further, and without ruling open questions this way or that, to throw
off the abstract and unreal character which attends a course of
reasoning, by setting down the evidence for and against certain
miracles as we meet with them. Moreover, so much has been said in the
foregoing pages in behalf of the Ecclesiastical Miracles, antecedently
considered, that it may be hastily inferred that all miraculous
relations and reports should be admitted unhesitatingly and
indiscriminately, without any attempt at separating truth from
falsehood, or suspense of judgment, or variation in the reliance
placed in them one with another, or reserve or measure in the open
acknowledgment of them. And such an examination of particular
instances, as is proposed, may give opportunity to one or two
additional remarks of a general character for which no place has
hitherto been found.

108.
An inquirer, then, should not enter upon the subject of the miracles
reported or alleged in ecclesiastical history, without being prepared
for fiction and exaggeration in the narrative, to an indefinite
extent. This cannot be insisted on too often; nothing but the gift of
inspiration could have hindered it. Nay, he must not expect that more than a few can be exhibited with
evidence of so cogent and complete a character as to demand his
acceptance; while a great number of them, as far as the evidence goes,
are neither certainly true nor certainly false, but have very various
degrees of probability viewed one with another; all of them
recommended to his devout attention by the circumstance that others of
the same family have been proved to be true, and all prejudiced by his
knowledge that so many others on the contrary are certainly not true.
It will be his wisdom, then, not to reject or scorn accounts of
miracles, where there is a fair chance of their being true; but to
allow himself to be in suspense, to raise his mind to Him of whom they
may possibly be telling to ``stand in awe, and sin not,'' and to
ask for light,—yet to do no more; not boldly to put forward what, if
it be from God, yet has not been put forward by Him. What He does in
secret, we must think over in secret; what He has ``openly showed in
the sight of the heathen,'' we must publish abroad, ``crying aloud, and
sparing not.'' An alleged miracle is not untrue because it is unproved;
nor is it excluded from our faith because it is not admitted into our
controversy. Some are for our conviction, and these we are to ``confess
with the mouth'' as well as ``believe with the heart;'' others are for
our comfort and encouragement, and these we are to ``keep, and ponder
them in our heart,'' without urging them upon unwilling ears.

109.
No one should be surprised at the admission that few of the
Ecclesiastical Miracles are attended with an evidence sufficient to
subdue our reason, because few of the Scripture Miracles are furnished
with such an evidence. When a fact comes recommended to us by
arguments which do not admit of an answer, when plain and great
difficulties are in the way of denying it, and none, or none of
comparative importance, in the way of admitting it, it may be said to
subdue our reason. Thus Apologists for Christianity challenge
unbelievers to produce an hypothesis sufficient to account for its
doctrines, its rise, and its success, short of its truth; thus Lord
Lyttelton analyses the possible motives and principles of the human
mind, in order to show that St. Paul's conversion admits of but
one explanation, viz., that it was supernatural; thus writers on
Prophecy appeal to its fulfilment, which they say can be accounted for
by referring it to a Divine inspiration, and in no other way. Leslie,
Paley, and others have employed themselves on similar arguments in
defence of Revealed Religion. I am not saying how far arguments of a
bold, decisive, and apparently demonstrative character, however great
their value, are always the deepest and most satisfactory; but they
are those which in this day are the most popular; they are those, the
absence of which is made an objection to the Ecclesiastical Miracles.
It is right then to remind those who consider this objection as fatal
to these miracles, that the Miracles of Scripture are for the most
part exposed to the same. If the miracles of Church History cannot be
defended by the arguments of Leslie, Lyttelton, Paley, or Douglas, how
many of the Scripture miracles satisfy their conditions? Some infidel
authors advise us to accept no miracles which would not have a verdict
in their favour in a court of justice; that is, they employ against
Scripture a weapon which Protestants would confine to attacks upon the
Church; as if moral and religious questions required legal proofs, and
evidence were the test of truth.

110.
It is true that the Scripture miracles were, for the most part,
evidence of a Divine Revelation at the time when they were
wrought; but they are not so at this day. Only a few of them fulfil
this purpose now; and the rest are sustained and authenticated by
these few \footnote{Vid. \emph{supr}. Essay i. pp. 9, 55, 91, 92;
also pp. 187, 207.}. The many never
have been evidence except to those who saw them, and have but held the
place of doctrine ever since; like the truths revealed to us about the
unseen world, which are matters of faith, not means of conviction.
They have no existence, as it were, out of the record in which they
are found; they are not found as facts in the world, influencing its
course, and proving their reality by their power, but as sacred truths
taught us by inspiration. Such are the greater number of our Lord's
miracles viewed individually; we believe His restoration of the widow's
son, or His changing water into wine, as we believe His
transfiguration, on the word of His Evangelists. We believe the
miracles of Elisha, because our Lord has Himself recognised the book
containing the record of them. The great arguments by which
unbelievers are silenced do not reach as far as these particular
instances. As was just now noticed, one of the most cogent proofs of
the miracles of Christ and His Apostles is drawn from their effects;
it being inconceivable that a rival power to Cæsar should have
started out of so obscure and ignorant a spot as Galilee, and have
prevailed, without some such extraordinary and divine gifts; yet this
argument, it will be observed, proves nothing about the miracles
\emph{one by one} as reported in the Gospels, but only that the
Christian \emph{story} was miraculous, or that miracles attended it.
Paley's argument goes little beyond proving the fact of the
Resurrection, or, at most, that there were certain sensible miracles
wrought by our Lord, such as cures, to which St. Peter alludes in his
speech to Cornelius, yet without specifying what. Again, Douglas
considers that ``we may suspect miracles to be false,'' the account of
which was not published at the time or place of their alleged
occurrence, or if so published, yet without careful attention being
called to them; yet St. Mark is said to have written at Rome, St. Luke
in Rome or Greece, and St. John at Ephesus; and the earliest of the
Evangelists wrote some years after the events recorded, while the
latest did not write for sixty years; and, moreover, true though it be
that attention was called to Christianity from the first, yet it is
true also that it did not succeed at the spot where it arose, but
principally at a distance from it. Once more, Leslie almost confines
his tests to the Mosaic miracles, or rather to certain of them; and
though he is unwilling to exclude those of the Gospel from the benefit
of his argument, yet it is not easy to see how he brings them under it
at all.

111.
On the whole, then, it will be found that the greater part of the
Miracles of Revelation are as little evidence for Revelation at this
day, as the Miracles of the Church are now evidence for the
Church. In both cases the number of those which carry with them their
own proof now, and are believed for their own sake, is small; and
these furnish the grounds on which we receive the rest. The difference
between the two cases is this:—that, since an authentic document has
been provided for the miracles by which Revealed Religion was
introduced, which are thus connected together into one whole, we know
here exactly what miracles are to be received on warrant of those
which are already proved; but since the Church has never catalogued
her miracles, those which are known to be such do but create an
indefinite presumption in favour of others, but cannot be taken in
proof of any in particular.

112.
On the other hand, that fables have ever been in circulation, some
vague and isolated, others attached to particular spots or to
particular persons, is too notorious to need dwelling on: it is more
to the purpose to observe that the fact of such pretences has ever
been acknowledged even by those who have been believers or reporters
of miraculous occurrences. We have seen above \footnote{N. 28.} that one of St. Martin's first miracles in his episcopate, as
recorded by Sulpicius, was the detection of a pretended Saint and
Martyr, whose tomb had been an object of veneration to the ignorant
people. And in the very beginning of Christianity St. Luke, in
speaking of the ``many'' who had ``taken in hand to set forth in order a declaration of those things which are most surely believed
among us,'' seems to allude to the Apocryphal Gospels \footnote{Jones, On the Canon, part i. ch. 2, has collected the ancient and
modern authorities in proof that St. Luke was alluding to the
Apocryphal writings. Wolf denies it, Cur. Phil. in loc.}, which ascribe a number of trifling as well as fictitious
miracles to our Lord. And when St. Paul cautions the Thessalonians
against being ``soon shaken in mind or troubled, by spirit or by
letter, as from himself, as that the day of Christ was at hand,'' he
testifies both to the fact that spurious writings were then ascribed
to him, and that they contained professedly supernatural matter.

113.
What is confessed by Apostles and Evangelists in the first century,
and by Martyrologists in the fourth, would naturally happen both in
the interval and afterwards. Hence Pope Gelasius, while warning the
faithful against several Apocryphal works, mentions among them the
Acts of St. George, the martyr under Dioclesian, which had been so
interpolated by the Arians, that to this day, though he is the patron
of England, and in Chapters of the Garter is commemorated with honours
which even Apostles do not gain from us, nothing whatever is known for
certain of his life, sufferings, or miracles \footnote{Baron, Annal. 290. 35: Martyrol. Apr. 23.}. Again, we are told by St. John Damascene, and in the
Revelations of St. Bridget and St. Mathildis, that the Emperor Trajan was delivered from the place of punishment at the prayers of
St. Gregory the First; but Baronius says, concerning this and similar
stories, ``Away with idle tales; silence once for all on empty fables;
be they buried in eternal silence. We excuse those who, accounting
true what they received as fact, committed it to writing; praise to
their zeal, who, when they found it asserted, discussed in scholastic
fashion how it might be; but more praise to them who, scenting the
falsehood, detected the error.'' \footnote{Emunctis naribus odorati. Annal. 604. 49.} Melchior Canus, again, a Dominican and a Divine of Trent, uses
the same language even of St. Gregory's Dialogues and the
Ecclesiastical History of Bede. ``They are most eminent persons,'' he
says, ``but still men; they relate certain miracles as commonly
reported and believed, which critics, especially of this age, will
consider uncertain. Indeed, I should like those histories better, if
their authors had joined more care in selection to severity in
judgment;'' \footnote{Loc. Theol. xi. 6.} though he adds
that far more was to be retained in their works than was to be
rejected. He does not, however, speak even in these measured terms of
the Speculum Exemplorum, and the Aurea Legenda of Jacobus de Voragine;
the former of which, he says, contains ``monsters of miracles rather
than truths;'' and the latter is the production of ``an iron mouth, a
leaden heart, and an intellect without exactness or discretion.''
Avowals such as these from the first century to the sixteenth, from
inspired writers to the schools of St. Dominic and the Oratory, may
serve to prepare us for fictitious miracles in ecclesiastical history
in no small measure, and to show us at the same time that such
fictions are no fair prejudice to others which possess the characters
of truth \footnote{The illustration of this subject might be pursued without limit.
Tillemont quotes from a writer of the thirteenth century the broad
maxim: ``Quand la raison se trouve contraire à l'usage, il faut que l'usage
cede à la raison;'' and proceeds to quote Papebrok as saying that we
cannot too often repeat this excellent rule, ``à ceux qui trouvent
mauvais qu'on accuse de fausseté diverses choses qui se sont
introduites dans l'Eglise par l'ignorance de l'histoire.'' vol. vii. p.
640. The Bollandists say, ``Nimiâ profecto simplicitate peccant qui
scandalizantur quoties audiunt aliquid ex jam olim creditis, et juxta
Breviarii præscriptum hodiedum recitandis, in disputationem adduci.''
Dissert. Bolland. tom. ii. p. 140. Vid. also Alban Butler's Saints,
Introd. Disc. p. xlvii., etc., ed 1833. Bauer's Theolog. tom. i. art.
ii. p. 487, and works there referred to. Benedict. XIV. de Canon.
Sanct. iv. p. 1. c. 5, etc. Farmer, On Miracles, p. 320; also the
passages from various authors quoted in Geddes' Tracts, vol. iii. pp.
115-118, ed. 1730; who also furnishes, though not in a good spirit,
a number of specimens of the sort of miracles which such authors
condemn.}. And in like
manner, if it be necessary, exceptions might be taken to certain of
the miracles recorded by Palladius in his Lausiac, and by Theodoret in
his Religious History, and by the unknown collector of the miracles of
St. Stephen, which a late writer has brought forward with the hope of
thereby involving all the supernatural histories of antiquity in
a general suspicion and contempt. That Palladius has put in writing a
report of a hyena's asking pardon of a solitary for killing sheep, and
of a female turned by magic into a mare, or that one of the Clergy of
Uzalis speaks of a serpent that was seen in the sky, will appear no
reason, except to vexed and heated minds, for accusing the holy
Ambrose of imposture, or the keen, practised, and experienced
intellect of Augustine of abject credulity \footnote{``Ambrose occupies a high position among the Fathers; and there
was a vigour and dignity in his character, as well as a vivid
intelligence, which must command respect; but in proportion as we
assign praise to the man individually, we condemn the system which
could so far vitiate a noble mind, and impel one so lofty in temper to
act a part which heathen philosophers would utterly have abhorred …
Under the Nicene system, Bishops in the great cities could stand up in
crowded churches, without shame, and with uplifted hands appeal to
Almighty God in attestation of that, as a miracle, which themselves
had brought about by trickery, bribes, and secret instructions.''
Ancient Christ. part vii. pp. 270, 271. ``He [Augustine] was the
dupe of his own credulity, not the machinator of fraud.'' P. 318.}.

114.
Nor is there anything strange or startling in this mixture of fable
with truth, as appeared from what was said on the subject in a former
page. It as little derogates from the supernatural gift residing in
the Church that miracles should have been fabricated or exaggerated,
as it prejudices her holiness that within her pale good men are mixed
with bad. Fiction and pretence follow truth as its shadows; the
Church is at all times in the midst of corruption, because she is in
the midst of the world, and is framed out of human hearts; and as the
elect are fewer than the reprobate, and hard to find amid the chaff,
so false miracles at once exceed and conceal and prejudice those which
are genuine. Nor would the difficulty be overcome, even if we took on
ourselves to reject all the Ecclesiastical Miracles altogether; for
the fictions which startle us must in fairness be viewed as connected,
not only with the Church and her more authentic histories, but with
Christianity, as such. Superstition is a corruption of Christianity,
not merely of the Church; and if it discredits the divine origin of
the Church, it discredits the divine origin of Christianity also.
Those who talk even most loudly of the corruptions of the fourth and
fifth centuries, seem, when closely questioned, still to admit that
Christianity was not extinct, but overlaid by corruptions. If, then,
the Church herself, and her miracles \emph{in toto}, are to be
included in that corruption, then of course the corruption was only
deeper and broader, than if she is to be accounted as in herself a
portion of Apostolic Christianity; and if such greater corruption does
not compromise the divinity of Christianity, so the lesser surely does
not compromise the real power and gifts of the Church. On both sides
fanaticism, imposture, and superstition are admitted as existing in the history of miracles; and on neither side must these evil agents
be held to throw suspicion on particular miracles which have no direct
or probable connection with them.

And now, after these preliminary considerations,
let us proceed to inquire into the evidence and character of several
of the miracles in particular, which we meet with in the first
centuries of Christianity.

\xsection{Section 1. The Thundering Legion}

115.
CLAUDIUS APOLLINARIS,
Bishop of Hierapolis, addressed an Apology for Christianity to the
Emperor Marcus, about A.D. 176.
It is lost, but reference to it, as it would appear, or at least to
one of his works, is made by Eusebius \footnote{Hist. v. 5.}, in which Apollinaris bore witness to a remarkable answer to
prayer received a year or two before by the Christian soldiers of that
very Emperor's army in the celebrated war with the Quadri. Tertullian,
writing about A.D. 200, and also
in a public Apology, urges the same fact upon the Proconsul of Africa
whom he is addressing.

116.
The words of Eusebius, introductory of the evidence of Apollinaris and
Tertullian, are these: ``It is said that when Marcus Aurelius Cæsar
was forming his troops in order of battle against the Germans and
Sarmatians, he was reduced to extremities by a failure of water.
Meanwhile the soldiers in the so-called Melitene \footnote{On the question of this Melitine or Melitene legion, vid. Vales. in
loc. Euseb.} legion, which for its faith remains to this day, knelt down
upon the ground, as we are accustomed to do in prayer, and betook
themselves to supplication. And whereas this sight was strange to the
enemy, another still more strange happened
immediately,—thunderbolts, which caused the enemy's flight and
overthrow; and upon the army to which the men were attached, who had
called upon God, a rain, which restored it entirely when it was all
but perishing by thirst.'' He adds, that this account was given by
heathens as well as by Christians, though they did not allow that the
prayers of Christians were concerned in the event. Then he quotes
Apollinaris for the fact that in consequence the legion received from
the Emperor the name of ``Thundering.'' Again, Tertullian speaks of ``the
letters of Marcus Aurelius, an Emperor of great character, in which he
testifies to the quenching of that German thirst by the shower gained
by the prayers of soldiers who happened to be Christians.'' \footnote{Apol. c. 5.} He adds that, ``while the Emperor did not openly remove the
legal punishment from persons of that description, yet he did in fact
dispense with it by placing a penalty, and that a more fearful one, on
their accusers.'' And in his Ad Scapulam: ``Marcus Aurelius in the
German expedition obtained showers in that thirst by the prayers
offered up to God by Christian soldiers.'' \footnote{Ad Scap. c. 4.} The statement, then, as given by two writers, one writing at
the very time, the other about twenty years later, is this: that the
soldiers in or of one of the Roman legions, gained by their prayers a
seasonable storm of rain and thunder and lightning, when the army was
perishing by thirst, and was surrounded by an enemy; and they add two
evidences of it—Apollinaris, that the legion in which these soldiers
were found was thenceforth called the Thundering Legion; and
Tertullian, that the Emperor in consequence passed an edict in favour
of the Christians.

117.
Here we are only concerned with the \emph{fact}, not with its alleged \emph{evidences};
and this is worth noticing, for it so happens that the fact is true,
but the evidences, \emph{as} evidences, are not true; that is, there
is just enough incorrectness in the statement to hinder their availing
as evidences. This, I say, is worth noticing, because it may serve in
other cases to make us cautious of rejecting facts stated by the
Fathers because we discredit (rightly or wrongly is not the question)
the grounds on which they rest them. Did we know no other evidence
than what Apollinaris and Tertullian allege for the sudden relief of
the Roman legions in Germany, we should have rejected the fact when we
had invalidated the evidence; but this, as the event shows,
would have been a hasty proceeding. Sometimes facts are so notorious
that proof is \emph{ex abundanti}; and sometimes writers like those in
question hurt a good cause by not leaving it to itself.

118.
Now, as to the corroborative statement made by Apollinaris, writers of
great authority assume that he, or other early writers, speak as if a
legion in the Roman army was composed \emph{wholly} of Christians \footnote{Vales. In Euseb. Hist. v. 5. Moyle's Posthumous Works, Vol. ii. p. 82.
Jablonski's Opusc. Tom. iv. p. 9.}. Yet even Eusebius does but speak of ``the soldiers in the
Melitene legion,'' which is an ambiguous form of expression; while
Tertullian uses the phrase, ``Christianorum \emph{forte} militum
precationibus,'' ``Christianorum militum orationibus,'' no mention being
made of a legion at all, and the word ``forte'' strongly opposing
the idea that the Christians formed an entire body of troops. As to
Apollinaris, he, it is true, stated in his lost work, that in
consequence of the miracle a legion was called ``Thundering''; but we
may not assume that he said more than that the Christians who prayed
were \emph{in} the legion, since there is nothing strange in the idea
of a whole body obtaining a name from the good deed of some of them,
nor strange, again, considering that bodies of troops were drawn, then
as now, from particular places, and were open to various local or
other influences, that Christians should have been numerous enough in
one particular legion to give a character to it. This
difficulty, however, being disposed of, a more important objection
remains; there was indeed a Thundering Legion, as Apollinaris says,
but then it was as old as the time of Trajan, nay, of Augustus \footnote{Moyle's Posthumous Works, Vol. ii. p. 90, and Scaliger and Valesius
before him. Baronius accounts for the fact by supposing that the
Christian soldiers were in all parts of the army, and after this were
incorporated into the existing Thundering Legion. ``Par est credere,
ipsum eosdem ob tam egregium atque mirandum facinus Fulminantium
nomine nobilitasse, ac eosdem simul ejusdem nominis legioni pariter
aggregasse.'' Ann. 176, 20; vid. also Witsius, Diatrib. 46. Mr. King,
too, observes that Xiphilin is the only author who ``absolutely affirms
the soldiers of the Melitenian Legion to be all Christians.'' Ap. Moyle,
p. 116; vid. also Milman, Christ. Vol. ii. p. 190. Moyle answers that
King is the first person who has interpreted Eusebius, etc.,
otherwise, p. 212. Lardner, Testim. Vol. ii. Ch. 15; and Mosheim, ant.
Constant. Sec 2. Ch. 17, side with Moyle. Mosheim connects ``forte''
with ``precationibus impetrato.'' [Lumper, t. 7, p. 510 note, says that ``forte''
is African Latin for ``fortuito;'' he seems to agree with Mosheim
in the construction. He gives a list of authors who have treated of
the occurrence, p. 515.]}. This circumstance, of course, is fatal to his argument. Moyle,
upon this, observes that ``Apollinaris, the first broacher of the
miracle, was grossly mistaken, to say no worse;'' \footnote{He retracts and throws the blame on Eusebius, p. 221, almost denying
that Apollinaris made the statement imputed to him. So does Neander,
Church Hist., Vol. i. 1, 2.} but, though it was a mistake, it surely is not grosser than if
a country clergyman at this day were to commit a blunder in speaking
of the Queen's regiments serving in India or Canada. In spite of
our advantages from the present diffusion of knowledge, certainly our
parish priests do not know much more of the constitution or history of
the British army than the Bishop of Hierapolis of the military
establishments of Rome.

119. Tertullian, on the other hand, tells us that the Emperor, in a formal
document, acknowledged the miracle as obtained by the prayers of the
Christians, and favoured the whole body in consequence; not, indeed,
repealing the laws against them, but putting a heavier punishment on
informers against them than on themselves. And it would appear that
the Emperor did issue a rescript in their favour in an earlier period
of his reign, which Eusebius has preserved \footnote{Moyle denies the genuineness of this Rescript, and Dodwell suspects
it. Dissert. Cypr. xi. 34, fin. Moyle adds, p. 337, that G. Vossius
wrote a Dissertation to prove it a forgery. Pagi and Valesius maintain
it; so does Jablonski, 1. c., assigning it with Pagi to the ninth year
of Antoninus, while Valesius assigns it to the first.}, to the effect that ``the parties accused of Christianity shall
be pardoned, though it be proved against them, and the informer shall
undergo the penalty instead;'' and in the reign of Commodus, the son of
Marcus, a Pagan actually had his legs broken, and was put to death,
for bringing an accusation against a Christian \footnote{Jablonski, ibid. p. 18. Moyle suspects the story, yet without strong
grounds, p. 249. It is found in Eusebius.}. And, further, that the Emperor, about the time of the
German war, showed a leaning towards ``foreign rites,'' which might
easily be mistaken by the Christians to include or even to imply
Christianity, is made clear by one of the authors to whom reference
has just been made at the foot of the page \footnote{Jablonski, ibid. Moyle, with a different purpose, gives instances of
the Emperor's leaning towards Chaldeans, magicians, etc., p. 235; vid.
also p. 356.}. Moreover, that the Emperor recognized the miracle is very
certain, as will appear directly; but, all this being undeniable,
still there is no evidence for the very point on which the force of
Tertullian's proof depends, viz., that his act of grace towards the
Christians was in \emph{consequence} of his belief in the miracle, and
his belief that they were the cause of it \footnote{Moyle maintains, p. 244, that Tertullian does not assert this
connection of Antoninus's acknowledgment of the miracle with his
edict, nor any other ancient writer.}. So far from it, he was in a course of persecution against the
Church, both before and after its date. How severely that persecution
raged a few years afterwards, the well-known epistle of the Churches
of Gaul informs us \footnote{Witsius, to evade the difficulty, maintains that the persecution was
the consequence of a riot, and the hostility of local governors,
Diatrib. c. 66. King maintains the same, ap. Moyle, p. 309. Eusebius
certainly speaks of it as [\emph{ex epithise\={o}s t\={o}n d\={e}m\={o}n}].
Hist. v. proœm.};
though its force must at least have been suspended as regards Asia
Minor, otherwise Apollinaris, writing at the time, could not have fancied that the Emperor had recognized the miracle as the result
of Christian prayer.

120.
Dismissing, however, these two statements, which, though they cannot
be maintained as they stand, still are not necessary conditions of the
alleged miracle, and which admit, as we have seen, of a very ready
explanation, we have, nevertheless, the following decisive evidence in
proof of the occurrence of some extraordinary and providential storm,
when the Roman army was in very critical circumstances in the course
of the German war.

121.
Eusebius observes that even the Pagans confessed the miracle, though
they did not allow that it was attributable to the prayers of the
Christians; and what is left of antiquity sufficiently confirms his
statement. Indeed, so certain was the fact, that nothing was left to
the Pagans but to record it and to account for it. They accounted for
it by referring it to their own divinities; they recorded it on medals
and on monuments. Dio Cassius calls it a ``wonderful and providential''
preservation, and attributes it to an Egyptian magician, of the name
of Arnuphis, who invoked ``Mercury, who is in the air, and other
spirits.'' Julius Capitolinus attributes it to the Emperor's
prayers. Themistius, who says the same, adds that he had seen a
picture, ``in the middle of which the Emperor was praying in the line
of battle, and his soldiers were catching the rain in their helmets,
and quenching their thirst with the draught thus providentially
granted.'' Moreover, the memorial of it is sculptured on the celebrated
Antonine column at Rome, where is a figure of Jupiter Pluvius
scattering lightning and rain, the enemy and their horses lying
prostrate, and the Romans, sword in hand, rushing on them. A medal,
too, is or was extant, of the very year of the occurrence, with the
head of Antoninus crowned with laurel on one side, and a figure of
Mercury on the reverse.

122.
The very fact of this event being recorded with such formality on the
column of Antoninus, is of itself a sufficient proof of its
importance; but perhaps the reader will be more impressed by the pagan
Dio's description of it, which runs as follows: ``When the Barbarians
would not give them battle, in hopes of their perishing by heat and
thirst, since they had so surrounded them that they had no possible
means of getting water, and when they were in the utmost distress from
sickness, wounds, sun, and thirst, and could neither fight, nor
retreat, but remained in order of battle and at their posts in this
parched condition, suddenly clouds gathered, and a copious rain fell,
not without the mercy of God. And when it first began to fall, the
Romans, raising their mouths towards heaven, received it upon them;
next, turning up their shields and helmets, they drank largely out of
them, and gave to their horses. And when the Barbarians charged
them, they drank as they fought; and numbers of them were wounded, and
drank out of their helmets water and blood mixed. And while they were
thus incurring heavy loss from the assault of the enemy, because most
of them were engaged in drinking, a violent hail-storm and much
lightning were discharged upon the enemy. And thus water and fire
might be seen in the same place falling from heaven, that some might
drink refreshment, and others be burned to death; for the fire did not
touch the Romans, or if so, it was at once extinguished; nor did the
wet help the Barbarians, but burned like oil; so that, drenched with
rain, they still needed moisture, and they wounded their own selves,
that blood might put out the fire.'' \footnote{This is translated from Baronius; but it agrees with the original in
all important points, though not always literal. Dion. Hist. lxxi. p.
805. Vid. also Themist. Orat. 15.} This of course is rhetorically written, but men do not write
rhetorically without a cause, and the effort of the composition shows
the marvellousness of the occurrence.

123.
We are sure, then, of the providential deliverance of the army, as
Eusebius and the others state it. And that there were Christians in
the army we may be quite sure, from what we gather from the general
history of the times \footnote{Moyle indeed contends ``that there were few or none at all in the
army,'' and observes, ``Considering the passive principles of the
age, I would as soon believe my Lord Marlborough had a whole regiment
of Quakers in his army as that Antoninus had a whole legion of
Christians in his.'' pp. 84, 85. He argues from the testimonies of the
early Fathers, of Celsus, etc., and from the oaths and other
idolatrous acts to which soldiers were obliged to submit, adding, ``that
it was impossible for a Christian to serve in them unless it were by
the help of Occasional Conformity. At least in such a case the prayers
of such mock Christians would hardly work wonders.'' p. 87. This is an
objection which, if valid, strikes deeper than any of those which I
have noticed in the text. Mr. Milman observes of the alleged
apparition of the Cross to Constantine, ``This irreconcilable
incongruity between the symbol of universal peace and the horrors of
war, in my judgment, is conclusive against the miraculous or
supernatural character of the transaction.'' Hist. of Christ, Vol. ii.
p. 354. He adds, ``This was the first advance to the military
Christianity of the middle ages.'' He refers in a note to Mosheim for
similar sentiments, ``for which,'' he says, ``I will readily
encounter the charge of Quakerism.'' He then refers to the Empress
Helena's turning the nails of the Cross into a helmet and bits for
Constantine's war-horse. ``True or false,'' he observes, ``this story is
characteristic of the Christian sentiment then prevalent.'' [Vid. also
Lupus, Opp. t. xi. p. 94, etc. Pusey on Tertullian, p. 184. Gibbon,
Miscellan. Works, p. 759, ed. 1837.]},
even independently of what these writers state. And further, we
may be sure also, even before we have definite authority for the fact,
that they offered up prayers for deliverance.

124.
Under these circumstances I do not see what remains to be proved. Here
is an army in extreme jeopardy, with Christians in it; the enemy is
destroyed and they are delivered. And Apollinaris, Tertullian, and
Eusebius, attest \footnote{Moyle indeed maintains that the Christians in general did not believe
it to be a miracle; he argues from the silence of St. Theophilus, St.
Clement, Origen, St. Cyprian, Arnobius, and Lactantius, p. 277. W.
Lowth, however, refers to a passage in St. Cyprian, ad Demetrian.
Routh, t. i. p. 153. It really seems unreasonable to demand that every
Father should write about everything.} that
these Christians in the army prayed, and that the deliverance
was felt at the time to be an answer to their prayers; what remains
but to accept their statement? We, who are Christians as well as they,
can feel no hesitation on the score that pagan writers attribute the
occurrence to another cause, to magic or to false gods. Surely we may
accept the evidence of the latter to the fact, without taking their
hypothetical explanation of it. And we may give our own explanation to
it for our own edification, in accordance with what we believe to be
divine truth, without being obliged to go on to use it in argument for
the conversion of unbelievers. It may be a miracle, though not one of
evidence, but of confirmation, encouragement, mercy, for the sake of
Christians.

125.
Nor does it concern us much to answer the objection that there is
nothing strictly miraculous in such an occurrence, because sudden
thunder-clouds after drought are not unfrequent; for in addition to
other answers which have been made to such a remark in other parts of
this Essay, I would answer, Grant me such miracles ordinarily in the
early Church, and I will ask no other; grant that upon prayer benefits
are vouchsafed, deliverances are effected, unhoped-for success
obtained, sickness cured, tempests laid, pestilences put to flight,
famines remedied, judgments inflicted, and there will be no need of
inquiring into the causes, whether supernatural or natural, to which
they are to be referred \footnote{Moyle is obliged to allow so much as this, saying of the defeat of the
Philistines by a storm on Samuel's prayers, ``This fact, though it
cannot properly, in the strict and genuine sense of the word, be
called a miracle, yet well deserves a place in the lower form of
miracles, because it was \emph{preternatural}, and not performed by
the ordinary concurrence of second causes, but by the immediate hand
of God.'' p. 286. Vid Benedict. xiv. de Can. Sanct. iv. part i. 11, who
instances the hail-stones in Joshua's battle as ``præter naturam.'' Vid
infr. n. 143, 193.}.
They may or they may not, in this or that case, follow or surpass the
laws of nature, and they may surpass them plainly or doubtfully, but
the common sense of mankind will call them miraculous; for by a
miracle, whatever be its formal definition, is popularly meant an
event which impresses upon the mind the immediate presence of the
Moral Governor of the world. He may sometimes act through nature,
sometimes beyond or against it, but those who admit the fact of such
interferences will have little difficulty in admitting also their
strictly miraculous character, if the circumstances of the case
require it, and those who deny miracles to the early Church will be
equally strenuous against allowing her the grace of such intimate
influence (if we may so speak upon the course of Divine Providence,) as
is here in question, even though it be not miraculous.

126.
On the whole then we may conclude that the facts of this memorable
occurrence are as the early Christian writers state them; that
Christian soldiers did ask, and did receive, in a great distress, rain
for their own supply, and lightning against their enemies; whether
through miracle or not we cannot say for certain, but more probably
not through miracle in the philosophical sense of the word. All we
know, and all we need know is, that ``He made darkness His secret
place, His pavilion round about Him, with dark water and thick clouds
to cover Him; the Lord thundered out of heaven, and the Highest gave
His thunder; hail-stones and coals of fire. He sent out His arrows,
and scattered them; He sent forth lightnings against them.''

\xsection{Section 2. The Change of Water into Oil at the Prayer of St. Narcissus of Jerusalem}

127.
NARCISSUS, Bishop of Jerusalem,
when oil failed for the lamps on the vigil of Easter, sent the persons
who had the care of them to the neighbouring well for water. When they
brought it, he prayed over it, and it was changed into oil \footnote{Euseb. Hist. vi. 9.}. Narcissus was made Bishop about A.D.
180, at the age of eighty-four; he was at a Council on the question of
Easter 195, and lived through some years of the third century, dying
at the extraordinary age of a hundred and sixteen, or more.

128.
It is favourable to the truth of this account, that the instrument of
the miracle was an aged, and, as also was the case, a very holy man.
It may be added that he was born in the first century, before St. John's
death, and was in some sense an Apostolical Father, as Jortin
observes.

129.
But there are certain remarkable circumstances connected with him,
which, as persons regard them, will be viewed in contrary
lights, as making the miracle more or as less probable. Eusebius
informs us that Narcissus was for some years the victim of a malignant
calumny. Three men, disliking his strictness and the discipline he
exercised, accused him of some great crime, with an imprecation on
themselves if they spoke falsely; the first that he might perish by
fire, the second that he might be smitten with disease, and the third
that he might lose his eyesight. Narcissus fled from his Church, and
lived many years in the wild parts of the country, as a solitary. At
length the first of his three accusers was burned in his house, with
all his family; the second was covered from head to foot with the
disease which he had named; and the third confessed his crime, but,
overcome with shame and remorse, lost his eyes by weeping. Narcissus
was restored, and died in possession of his see.

130.
Now it may be said that the extraordinary nature of this history only
increases the improbability of the miracle. It reads like a made
story; there is a completeness about it; and there is an extravagance
in the notion of the loss of sight by weeping. Yet the same thing
happened to St. Francis. ``His eyes,'' says Butler, ``seemed two
fountains of tears, which were almost continually falling from them,
insomuch that at length he almost lost his sight.'' He was seared with
red-hot iron from the ear to the eye-brow, with the hope of
saving it. In his last illness ``he scarce allowed himself any
intermission from prayer, and would not check his tears, though the
physician thought it necessary for the preservation of his sight;
which he entirely lost upon his death-bed.'' \footnote{Lives of the Saints, Oct. 4.} However, even though we allow that the history in question is
embellished, still the general outline may remain, that Narcissus was
unjustly accused and by a wonderful providence vindicated. In this
point of view it surely adds to the probability of the miracle before
us, that it is attributed to a man, not only so close upon Apostolic
times and persons, so holy, so aged, but in addition so strangely
tried, so strangely righted. It removes the abruptness and
marvellousness of what at first sight looks like ``naked history,'' as
Paley calls it, or what we commonly understand by a legend. Such a man
may well be accounted ``worthy for whom Christ should do this.'' And if
the foregoing circumstances are true, not only in outline, but in
detail, then still greater probability is added to the miracle.

131.
Jortin objects that ``the change of water into oil to supply the church
lamps has the air of a miracle performed upon an occasion rather too
slender.'' \footnote{Dissert. in Iren. ii. 49.} But Dodwell \footnote{Vol. ii. p. 103.} had already observed that the mystical idea connected with the
sacred lights gives a meaning to it, and particularly at that
season; and Eusebius tells us that the people were much troubled \footnote{[\emph{Dein\={e}s athumias dialabous\={e}s to pan pl\={e}thos}].} at their failure.

132.
Jortin also observes that ``in the time of Augustus a fountain of oil
burst out at Rome, and flowed for a whole day. In natural history
there are accounts of greasy and bituminous springs, when something
like oil has floated on the water. Pliny, and Hardouin in his notes,
mention many such fountains, 'qui explent olei vicem,' and 'quorum aquâ
lucernæ ardeant.''' This circumstance perhaps adds probability to the
miracle, both as lessening its violence, (if the word may be used,) as
the accompanying history of the Bishop's trials lessens it in another
way, and because in matter of fact Almighty Wisdom seems, as appears
from Scripture, not unfrequently to work miracles beyond, rather than
against nature.

133.
Eusebius notices pointedly that it was the \emph{tradition} of the
Church of Jerusalem \footnote{[\emph{H\={o}s ek paradose\={o}s t\={o}n kata diadoch\={e}n
adelph\={o}n}].}. It
should be recollected, however, that the tradition had but a narrow
interval to pass from Narcissus to Eusebius,—not above fifty or
sixty years, as the latter was born about A.D.
264.

134.
On the whole then there seems sufficient ground to justify us in
accepting this narrative as in truth an instance of our Lord's
gracious presence with His Church, though the evidence is not so
definite or minute as to enable us to \emph{realize} the miracle. This
is a remark which is often in point; belief, in any true sense of the
word, requires a certain familiarity or intimacy of the mind with the
thing believed. Till it is in some way brought home to us and made our
own, we cannot properly say we believe it, even when our reason
receives it. This occurs constantly as regards matters of opinion and
doctrine. Take any characteristic point of detail in the religious
views of a person whom we revere and follow on the whole; do we
believe this particular doctrine or opinion of his, or do we not? We
do not like to pledge ourselves to it, yet we shrink from saying that
it is not true, and we defend it when we hear it attacked. We have \emph{no
doubt} about it, yet we cannot bring ourselves to say positively
that we \emph{believe} it, because belief implies an habitual presence
and abidance of the matter believed in our thoughts, and a familiar
acquaintance with the ideas it involves, which we cannot profess in
the instance in question. Here we see the use of reading and studying
the Gospels in order to true belief in our Lord; and, again, of acting
upon His words, in order to true belief in them \footnote{[Vid. Essay towards a Grammar of Assent, Ch.
iv. § 2.]}. This being considered, I do not see that we can be said
actually to \emph{believe} in a miracle like that now in question, of
which so little is known in detail, and which is so little
personally interesting to us; but we cannot be said to disbelieve it,
there being sufficient grounds for conviction in the sense in which we
believe the greater part of the accounts of general history.

\xsection{Section 3. The Miracle wrought on the course of the River Lycus by St. Gregory Thaumaturgus}

135.
DOUGLAS, in his great
earnestness to prove that no real miracles were wrought by the Fathers
and Saints of the second and third centuries, tells us that the
miracles of St. Gregory Thaumaturgus, some of which have been detailed
above, ``are justly rejected as inventions of a later age, and can be
believed by those only who can admit the miracles ascribed to
Apollonius, or those reported so long after his death to Ignatius.
Gregory of Nyssa,'' the biographer of Thaumaturgus, ``according to Dr.
Cave's character of him, was apt to be too credulous. No wonder,
therefore, he gave too much credit to old women's tales, as the
anecdotes of the Wonder-worker must be allowed to be, when related, as
we learn from St. Basil, by his aged grandmother Macrina.'' \footnote{Page 327, note.} This is not respectful either to St. Macrina or to St. Gregory
Thaumaturgus, to say nothing of his treatment of St. Gregory
Nyssen; plainly, it can mean nothing else but that St. Gregory did no
miracles, and that it is weak, nay, even heathenish, to believe he
did. Otherwise thinks a very careful and learned writer, not a member
of our Church, and his statement may fitly be placed in contrast with
the opinion of one who was a Bishop in it. ``His history,'' says Lardner,
speaking of Thaumaturgus, ``as delivered by authors of the fourth and
following centuries, particularly by Nyssen, it is to be feared, has
in it somewhat of fiction; but there can be no reasonable doubt made
but he was very successful in making converts to Christianity in the
country of Pontus about the middle of the third century; and that,
beside his natural and acquired abilities, he was favoured with
extraordinary gifts of the Spirit, and wrought miracles of surprising
power. The plain and express testimonies of Basil and others, at no
great distance of time or place from Gregory, must be reckoned
sufficient grounds of credit with regard to these things. Theodoret,
mentioning Gregory, and his brother, and Firmilian, and Helenus, all
together, ascribes miracles to none but him alone. They were all
Bishops of the first rank; nevertheless Gregory had a distinction even
among them. It is the same thing in Jerome's letter to Magnus; there
are mentioned Hippolytus, Julius Africanus, Dionysius of Alexandria,
and many others, of great note and eminence for learning and piety.
But Theodore, afterwards called Gregory, is the only one who is
called a man of Apostolical signs and wonders.'' \footnote{Credib. ii. 42, § 5.}

135.
These remarks of Lardner should be kept in mind by those who would
examine the miracles attributed to St. Gregory. For it is obvious to
reflect, that if we once believe that he did work miracles, it is the
height of improbability that in the course of a century all of these
should be forgotten, and a set of pretended miracles substituted in
their place, and that among a people who are noted for a particular
attachment to their old customs, and especially to the rites and
usages introduced by St. Gregory. ``The people of Neocæsarea,'' says
Lardner, ``retained for a long while remarkable impressions of
religion; and they had an affection for the primitive simplicity, very
rare and uncommon, almost singular, at that time, when innovations
came into the Church apace.'' \footnote{Vid. also above, n. 20.}
And if reasons can be given for believing one of the miracles, a
favourable hearing will be gained for the rest, which belong to one
family with it, and are conveyed to us through the same channels. All
are of the romantic kind, all come to us on tradition committed to
writing by St. Gregory Nyssen. That is, we shall have reason for
believing his narrative in its substance, for still there is nothing
to prevent misstatement in its detail. Against this, indeed,
inspiration alone could secure us.

137.
This absence of a perfection which only attaches to inspired
documents, has often been made an objection to receiving the miracles
which ecclesiastical history records \footnote{n. 23.}. But there is another peculiarity about its existing materials,
which applies in particular to Nyssen's Life of Thaumaturgus. That
Life does not answer the purpose for which critics and
controversialists require it at this day; it is very unsatisfactory as
an attestation of miracles, and would not read well in a process of
canonization. For in truth the author did not set himself to attest
them at all; he wrote a sacred panegyrical discourse, and from the
nature of the composition he left out names, dates, places,
particulars, all of them necessary indeed for critical proof, had he
been engaged in furnishing evidence. But, on the contrary, he was only
an encomiast of a departed Saint; and why not, if he found it fitting
for his people? He even omits particulars which he certainly knew
well, and every one else; as the name of Gregory's see, and of the
Emperor under whose persecution he fled from it; why then need
we be suspicious of other omissions, as if they necessarily reflected
on the general authenticity of his narrative? why may he not put off a
secular style and manner, when he is treating a religious subject? Why
is he to be compelled to turn the Church into a court of law, and to
introduce a prosaic phraseology into the hymns, the anthems, and the
lessons of the Euchology? He wrote for the faithful and devout, and he
had a right to do so.

138.
Dr. Lardner, a calm and impartial critic, as we have seen, here loses
himself. He seems to forget that a Bishop may, if he pleases, write
homilies and panegyrics, and read martyrologies, and that inspired
Scripture itself is not over-careful in dating, locating, and naming
the sacred persons and sacred things which it introduces. Would not
then St. Gregory's simple answer to such criticism as the following
be, that he did not write for Dr. Lardner? That candid writer seems to
forget this, when he makes the following additional remarks on his
Oration:—``It is plain it is a panegyric, not a history. Nyssen is so
intent upon the marvellous, that he has scarce any regard to common
things; he relates distinctly the mysterious faith which Gregory
received one night from John the Evangelist, but he despatches in a
very few words the instructions which Gregory received from Origen,
though he was five years under his tuition, and had before him
excellent materials to enlarge upon concerning that part of our
Bishop's history. Then he takes little or no notice of circumstances
of time and place, or the names of persons; these he omits as things
of no moment. Indeed, he has been so good as to inform us of Gregory's
native city and country, and that he studied some time, as he says, at
Alexandria; but he does not let us know where Gregory was acquainted
with Origen, whether at Alexandria or at Cæsarea. He does not inform
us of the Temple where Gregory lodged and silenced the demon; neither
where it stood, nor to what god it was dedicated. He has not so much
as once mentioned the name of the Priest who was converted in so
extraordinary a manner; nor has he mentioned the name of any one of
the many persons, subjects of Gregory's miraculous works.'' ``Possibly
it will be said,'' he continues, ``that it was contrary to the rules of
rhetoric to be more particular in an oration. If that be so, and all
that Nyssen aimed at was to entertain his hearers or readers with a
fine piece of oratory, we must consider it as such; but then, though
it may afford us some good entertainment, it will hardly be a ground
for much faith; for a story to be amusing is one thing, to be credible
is another.'' But is there no refuge from the rostrum on the one hand,
and the witness-box on the other? Must a style be either rhetorical or
controversial? May it not be ecclesiastical? However, Lardner grants
that the miracle wrought upon the Lycus is at least
particularized by the name of the river; and as some may think that it
even approaches to fulfil Leslie's celebrated criterion of a miracle,
a few words shall be given to it here.

139.
Leslie's tests of the truth of a miracle are these: that it should be
sensible; public; verified by some monument and observance; and that
set up at the very time when it was wrought. Now St. Gregory is said
by his biographer, as we have seen above \footnote{n. 23.}, to have restrained a mountain-stream within its mounds, which
had been accustomed to flood the plain country into which it
descended. He tells us the place, as well as the river, and the
mountains from which it flowed; he describes its impetuosity as
recurring, according as it was swollen by the waters from the
mountains, and its ravages as very serious. But he adds, that, after
Gregory had visited the spot and prayed, the calamity was stopped once
for all, though the stream descended with fury as before, and came up
to the very place which St. Gregory had marked as its limit. He
specifies this place by referring his readers to a monument standing
upon it, and that from the time of the miracle, and moreover, a
monument which in its history involves an additional miracle. He says
that the Saint took his staff, and fixed it at the opening of the
mound which the current had forced; that the staff grew into a tree,
and that the waters swept it, but never passed it, and that it
remained to his own day, being known by the people of the country as ``the
staff.'' Moreover, it may be said there was an observance, as well as a
monument, which dates from the time of the miracle; for surely the
conversion of the people benefited, upon their receipt of the benefit,
is, in its results, of the nature of a standing observance, and well
fitted to preserve and continue the knowledge of the supernatural act.
And further, as some immediate extraordinary occurrence is necessary
to enable us to account for so extraordinary an event as the
conversion of a whole people, but the success of St. Gregory's
restraint upon the stream could not be known till after an interval,
or rather only in a course of years, some probability is thereby added
to the idea that in the manner or circumstances of his action itself
there was something impressive and convincing, and such the miracle
wrought upon the staff would have been in an eminent way.

140.
Further, Nyssen not only lived too near the times to allow of a
spurious tradition fastening itself on the history, whether of the
tree or of the people, but he was a native and inhabitant of that part
of Asia Minor, and his family before him. His grandmother, Macrina,
was brought up at Neocæsarea, Gregory's see, by his immediate
disciples. Should the account be false, it will be somewhat of a
parallel to suppose a person, at this day, in high ecclesiastical
station, born and educated and writing in the Isle of Man, and
assuring us that Bishop Wilson once laid a storm in behalf of the
fishermen, and that a lighthouse was built at the time, and still
remains, in commemoration of the event; and writers, moreover, of this
day, in England, Scotland, and Ireland, confirming the testimony, by
incidentally observing, without allusion to the particular story, that
Bishop Wilson had the gift of miracles. We should say it was
impossible that such evidence could be offered in behalf of a fiction
now; and why not say the same of a similar case then? ``But a fiction
was possible then,'' it may be argued, ``because the age was more
superstitious than now.'' I answer, ``And so was a miracle possible,
because Christendom was more Catholic and Apostolic.''

141.
Of course, an objection may be raised, on the score of the miracle not
being of such a kind as to preclude the possibility of our referring
it to physical causes, in a country where earthquakes were not
uncommon. But miracles of degree, which admit abstractedly and
hypothetically of being explained by the joint operation of nature and
of exaggeration in the informant, are among the most common in
Scripture, and may nevertheless be cogent and convincing in the
particular case, as has been already observed. No east wind could
raise the waters of the Red Sea, as Scripture describes them to be
raised at the time of the Exodus; no supposition of earthquake or other physical disturbance will suffice to deprive Nyssen's
narrative, as it stands, of its miraculous character. Nor may we take
on ourselves to mutilate or deface either it or the Book of Exodus, or
any other professed statement of fact, without first assigning reasons
for our proceeding.

\xsection{Section 4. Appearance of the Cross in the sky to Constantine}

142.
WHEN Constantine was on his
march to Rome to attack Maxentius, at a time when he was as yet
undecided about the truth of Christianity, a luminous Cross is said to
have appeared in the sky at mid-day, in sight of himself and his army,
with the inscription, ``In this conquer.'' \footnote{Philostorgius, Nicephorus, and Zonaras say
that the inscription was in Latin. Eusebius gives the impression that
it was in Greek. So says the Emperor Leo expressly; vid. Grets. de
Cruc. tom. ii. p. 37, who, mentioning this difference of statement, \emph{ibid}.,
also determines against Brentius that the apparition was that of the
Cross with the monogram of χρ, and not merely of the
monogram. Rivetus too contends for the monogram only. Cath. Orth. ii.
19, p. 168.} His victory and his conversion followed. The date of these
transactions is A.D. 311 and
312.

143.
Now, here the fact reported is plainly miraculous. No known physical
cause could have formed a sentence of Greek or Latin in the air. It
has sometimes been supposed, indeed, that letters were not
really exhibited, but only some emblem, such as a crown, which denoted
conquest \footnote{Fabric, Dissert. de Cruce 10.}; and that then
what remains of the phenomenon may be resolved into meteoric effects.
But since any extraordinary appearance at such a juncture, whatever be
its physical cause, or whether it have one or no, is undeniably the
result of an immediate Divine superintendence, it is not easy to see
what is gained by an hypothesis of this nature. If in matter of fact
our Lord was then really addressing Constantine, it seems trifling to
make it a grave point to prove that He did so in this way, and not in
that. In such a case nature either would be made to minister, or would
be no impediment, to His Will; and His Will to address Constantine is
sufficient surely by itself to account for a contravention or
suspension of the laws of nature, and to overcome the presumption
which \emph{primâ facie} lies against the miracle. That He \emph{should}
address Constantine intelligibly is a miracle already. And surely to
sway and overrule the physical system towards a moral object, is a
miracle only different in degree from an interference with it for such
an object. For this is to impose on nature a constraint \emph{beyond and
above itself}, i.e. a \emph{supernatural} constraint; and if it is
subordinate to moral laws, why should it not sometimes give way to
them? In short, does the case ever stand thus, if it may be
reverently said that the Almighty \emph{would} address man, \emph{unless}
nature stood in the way? Does He fetter Himself with its laws, who,
even in the days of His flesh, did but submit to them, in order in the
event to dispense with them? Such explanations, then, either imply
that the inviolability of creation is more sacred than a Purpose of
the Creator, or they tamper with historical evidence for an
insufficient end. To mutilate the evidence is to incur all the
difficulty of denying it, with none of the gain. So this question may
be passed over.

144.
In the next place the \emph{à priori} aspect of the reported miracle,
if it is so to be called \footnote{``Confudit Danæus apparitiones Crucis cum miraculis Crucis; etsi enim
apparitiones Crucis possunt miracula appellari, Bellarminus tamen
apparitiones a miraculis distinxit, miracula vocans ea, quæ per
Crucem supra naturæ vim et ordinem patrata sunt.'' Gretzer de Cruce,
iv. 12, p. 253, ed. 1734.},
is in its favour. The approaching conversion of the Roman empire, in
the person of its head, was as great an event as any in Christian
history. Constantine's submission of his power to the Church has been
a pattern for all Christian monarchs since, and the commencement of
her state establishment to this day; and, on the other hand, the
fortunes of the Roman empire are in prophecy apparently connected with
her in a very intimate manner, which we are not yet able fully to
comprehend. If any event might be said to call for a miracle, it
was this; whether to signalize it or to bring it about. Thus it was
that the fate of Babylon was written on the wall of the banqueting
hall; also portents in the sky preceded the final destruction of
Jerusalem, and are predicted in Scripture as forerunners of the last
day. Moreover our Lord's prophecy of ``the Sign of the Son of Man in
heaven'' \footnote{On the sign of the Son of Man being understood of the Cross by the
Fathers, vid. G. Voss. Thes. Theolog. xvi. 9. Cornel. a Lapid. in loc.
Matt. and Maldonat. in loc.} was
anciently understood of the Cross. And, further, the sign of the cross
was at the time, and had been from the beginning, a received symbol
and instrument of Christian devotion, and cannot be ascribed to a then
rising superstition. Tertullian speaks of it as an ordinary rite for
sanctifying all the ordinary events of the day; it was used in
exorcisms; and, what is still more to the point, it is regarded by St.
Justin, Tertullian, and Minucius as impressed with a providential
meaning upon natural forms and human works, as well as introduced by
divine authority into the types of the Old Testament.

One would be inclined then to receive the
wonderful event in question on very slight evidence, if that evidence
were good as far as it went; and now let us see what, and of what
kind, is producible in its behalf. It is on the whole sufficient, yet
not without its difficulties.

145.
In the panegyrical oration delivered immediately upon the victory, the
speaker, who is a Pagan, asks, ``What God, what Divine Presence
encouraged thee, that when nearly all thy companions in arms and
commanders not only had secret misgivings but had open fears of \emph{the
omen}, yet against the counsels of men, against the warnings of the
diviners, thou didst by thyself perceive that the time of delivering
the city was come?'' \footnote{Baron. Ann. 312, 14.} Now
here an omen is mentioned of a public nature, which dismayed the
heathen priests and soldiers; it is remarkable too that \emph{what} it
was is not mentioned. All this would be sufficiently accounted for, if
it was the sign of the Cross which they had seen; a spectacle of all
others of bad augury with the hierarchy of the pagan city \footnote{Julian is said to have found a cross upon the entrails of a victim he
was offering in sacrifice; the sight [\emph{phrik\={e}n paresche kai
ag\={o}nian}]. Naz. Orat. iv. 54. He upbraids the Christians
with their worship of the wood of the Cross, and signing it upon their
foreheads and sculpturing it upon their dwellings. Cyril contr.
Julian, p. 194.}. And in corroboration of this interpretation, Eusebius, in his
own account of the miracle, tells us that on sight of the apparition
Constantine, who was still fluctuating between Christianity and
Paganism, was at first much distressed from a doubt what it portended.

146.
Next, about the year 314 or 315, that is, three years after the event,
Constantine erected his triumphal arch at Rome, which still
remains, with an inscription testifying that he had gained the victory
``\emph{instinctu divinitatis}, mentis magnitudine.'' \footnote{Burton, however, tells us that ``the words \emph{instinctu divinitatis}
are supposed to have been added afterwards, as the marble is there
rather sunk in, and the holes for the bronze letters are confused.''
Rome, p. 215. Yet the inscription reads oddly without them.}

147.
Further, before \footnote{\emph{i.e}. Before the first breach between Constantine and Licinius.
Vid. Gibbon, Ch. xx. note 40.} 314,
Lactantius or Cæcilius, as we determine the author, published his De
Mortibus Persecutorum; in which he asserts, not in a rhetorical tone
or in the form of panegyric, but in the grave style of history, that
Constantine, in consequence of a dream, caused the initial letter of
the word Christ to be inscribed on the shields of his soldiers, and
that he thereby gained the victory. ``Constantine,'' he says, ``was
admonished in sleep to mark the heavenly sign of God on the shields,
and so to engage the enemy. He did as he was bidden, and marks the
name of Christ on the shields, by the letter \emph{χ} drawn
across them, with the top circumflexed. Armed with this sign, his
troops take up arms. The enemy marches to meet them without their
imperial Commander, and passes over the bridge,'' etc. \footnote{De M. P. § 44.} Here is no mention of an \emph{apparition}, but still the
author speaks of the ``\emph{heavenly} sign.''

148.
On the 1st of March, 321, Nazarius, a pagan orator of celebrity,
pronounced, apparently at Rome, and not in Constantine's presence, a
panegyrical oration upon the Emperor. In this he speaks of the
assistance which the latter had received against Maxentius in the
following terms:—``Thou didst fight, O Emperor, by compulsion; but it
was thy best claim upon victory, that thou didst not seek it. Peace
was denied to him for whom victory was destined ... In short, it is
the talk of all the Gallic provinces, that hosts were seen, who bore
on them the character of divine messengers. And though heavenly things
use not to come to sight of man, in that the simple and uncompounded
substance of their subtle nature escapes his heavy and dim perception,
yet those, thy auxiliaries, bore to be seen and to be heard; and when
they had testified to thy high merit, they fled from the contagion of
mortal eyes. And what accounts are given of that vision, of the vigour
of their frames, the size of their limbs, the eagerness of their zeal!
Their flashing bosses shot an awful radiance, and their heavenly arms
burned with a fearful light; such did they come, that they might be
understood to be thine. And thus they spoke, thus they were heard to
say, 'We seek Constantine; we go to aid Constantine.' Even divine
natures have their boastings, and heavenly natures are touched by
ambition. Warriors who had glided down from heaven, warriors who were
divinely sent, even they did glory that they were marching with
thee. Their leader, I suppose, was thy father Constantius,'' etc. \footnote{Ap. Baron. Ann. 312. 11.}

149.
It is impossible to doubt from these contemporaneous witnesses,
witnesses more exactly contemporaneous than are commonly producible,
that some remarkable portent appeared, or was generally believed in,
when Constantine was in anticipation of his engagement with Maxentius,
and about the time that he first professed Christianity. After all
allowances for the rhetoric of Nazarius, his story surely must have
had some foundation; by it he is virtually doing homage to a religion
which he disowns, though he adroitly converts it to the service of
Paganism, by recurring to the old heathen prodigies, such as the
appearance of Castor and Pollux, and seeking to authenticate them by
the recent apparition. Even if the Cross appeared, he could not be
expected to mention it; he could not have done more than he has done.
The same may be said for the still earlier orator, who is obliged to
allude to the Emperor's Christianity, while he is complimenting him on
having rightly interpreted what his friends thought an omen of evil.
Lactantius, though he adds nothing to the evidence of the apparition
in the sky \footnote{Socrates, Philostorgius, Gelasius, Nicephorus, say that the Cross was
in the sky. Sozomen first speaks of it as seen in a dream, and then,
on the authority of Eusebius, describes the apparition in the sky.
Rufinus also gives both accounts.}, testifies
to the general idea of some wonderful occurrence having attended
the conversion of the Emperor. He testifies also to a fact which from
its boldness requires accounting for, Constantine's marking the symbol
of the Cross upon the arms of his soldiers.

150.
Nor is this the only indication of some extraordinary influence then
exerted upon the Emperor's mind. Not to dwell on the words already
quoted from his arch, which make no express mention of the Cross, we
find him even going so far as to form a new military standard, and
that is the Labarum, or Standard of the Cross. And on his entering
Rome in triumph, he forthwith erected a statue of himself with a Cross
in his hand, and an inscription to the effect that ``with that saving
sign'' he had delivered the city from a tyrant. But the most remarkable
evidence in point is a medal, extant in the last century, which bears
the figure of the Labarum with the very words, ``In this sign thou
shalt conquer.'' \footnote{So says Gibbon, referring to the Abbé du Voisin and a Jesuit, the Père
de Grainville. Such a medal is not described in Baronius, Gretser, or
Lipsius. Fabricius says, ``Nullus extat nummus, nullum vetus
monumentum, quo crux, in cœlo a Constantino visa, diserte confirmatur.''
Script. Græc. lib. v. c. 40. (t. 6. p. 706, ed. Harles.)} Thus
his assaults upon Paganism and the supernatural explanation of them go
together; one and the same auspicious omen is repeated, whether in
ensigns, medals, or monuments. And indeed, if we may dare to judge of
the course of Providence in this instance by its general laws,
it is scarcely possible to think that no divine direction was given to
such an instrument of its purposes on so great an occasion. In
junctures of such awful moment, nay, in far inferior ones, men are not
left alone, but strange impressions come over them, without which they
would not have nerve for bold deeds. It was an act surely of no
ordinary courage in Constantine to introduce the Labarum into the
Roman armies to the virtual disparagement of those standards which had
carried them to victory through so many fights, whether we regard the
feelings of his soldiers or the misgivings of his own mind.

151.
From this strictly contemporaneous testimony little or no part of
which can be called ecclesiastical, we seem to gather thus much, that
an omen happened to Constantine and his army, which most men thought
bad, but which he trusted;—there was some appearance in the heavens
visible to all;—some vision granted to himself;—and a Cross,—but
where seen does not appear, whether in his dream, or as part of the
visible appearance, and in connection with the omen spoken of; we are
but able to discern it in its reflection,—upon the shields, helmets,
and standards of his forces, and in his public commemorations of his
victory.

152.
Thus rests the evidence of the miracle in Constantine's lifetime;
after his death Eusebius gives the Emperor's own account of it, which
certainly does in a remarkable way explain those acts of his
which we have been recounting, and combine the scattered rumours which
accompanied them. Eusebius declares on the word of Constantine, who
confirmed it with an oath, that Constantine on his march saw, together
with his whole army, a luminous Cross in the sky above the mid-day
sun, with the inscription, ``In this conquer;'' and that in the ensuing
night he had a dream, in which our Lord appeared with the Cross, and
directed him to frame a standard like it as a means of victory in his
contest with Maxentius. Such is the statement ascribed by Eusebius to
Constantine; and it must be added that the historian had no leaning
towards over-easiness of belief, as many passages of his history show
\footnote{\emph{E.g}. He omits mention of the dove in the martyrdom of St.
Polycarp, of the miracles of St. Gregory Thaumaturgus, etc. In such
miracles as he does record, he is careful not to commit himself to an
absolute assent to them, but commonly introduces qualifying phrases.
And his answer to Hierocles is written in a very sober tone. Vid.
Kestner. de Euseb. Auct. et Fid. § 56, 57.}.

153.
This then is the state of the argument in behalf of the miracle; on
the other hand, there are these two difficulties in the way of
receiving it. First, Constantine's testimony, which alone is direct
and trustworthy, is not given till many years after the event;
moreover, it is given with an oath and in private, though it concerns
an occurrence of public notoriety; and it is not published in
his lifetime, nor till twenty-six years after the time to which it
refers \footnote{This objection is urged by Gibbon, Ch. 20. Lardner, Credib. ii. 70. §
3. Hoornebeek ap. Noris. Hist. Donat. App. 8.}. And next, it
is supported by no independent and by no ecclesiastical testimony. ``The
advocates for the vision,'' says Gibbon, ``are unable to produce a
single testimony from the Fathers of the fourth and fifth centuries,
who in their voluminous writings repeatedly celebrate the triumph of
the Church and of Constantine.'' \footnote{Ch. xx. note 52.} It is remarkable too that even Eusebius does not mention it in
his history, but in his Life of Constantine, as if, instead of its
being a public event, it were but a visitation or providence personal
to the Emperor.

154.
This, however, may be said in reply: It has already been shown that
rumours of some or other extraordinary occurrence abounded from almost
the time of the Gallic march \footnote{It is remarkable, however, as is observed by Gothofred. Diss. in
Philostorg. i. 6, that Optatianus Porphyrius, in his Panegyric. ad
Constant. written in the year 326, does not mention the apparition,
except that he calls the Cross ``cœleste signum.'' He wrote, however,
from banishment, though the place is not known.}; Nazarius says that it was the talk of the whole of Gaul; and
we see from his own account of it that it was mixed up with fiction,
as such popular reports are sure to be. An army is not like a
neighbourhood, or a class of society; it is cut off from the world, it
has no home, it acts as one man, it is of an incommunicative nature,
or at least does not admit of questioning. The troops of Constantine
saw the vision, and marched on; they left behind them a vague
testimony, which would fall misshapen and distorted on the very ears
that heard it, which would soon be filled out with fictitious details
because the true were not forthcoming, and which took a pagan form in
a country of Pagans. It was not unnatural that under such
circumstances Constantine should have been led formally to impart to
Eusebius the fact as it really took place; nor, considering the
misstatements that abounded, and the apparent unbelief of intelligent
Pagans \footnote{Vid. Gelas. Conc. Nic. i. 4.}, that he should
have confirmed his account of it with an oath. Nor is it wonderful
that Eusebius should not appeal to living witnesses of it, an omission
which Gibbon urges, as if an army, or the constituent parts of an
army, had a residence and an address, and that at the distance of
twenty-six years; or as if an ecclesiastic, a native of Palestine,
must have had many acquaintances among the veterans of Gaul \footnote{Eusebius says that he \emph{once} or \emph{sometimes} happened to see
the Labarum. [\emph{o d\={e} kai \={e}mas ophthalmois pote suneb\={e}
paralabein}]. V. Const. i. 30.}. Nor is it any great difficulty that, in a work professedly
panegyrical, and not historical, and written with much oratory
of phrase and circumlocution, and continual vagueness and
indeterminateness in statement \footnote{Vid. infr. n. 157. note t.}, the writer should not have mentioned the time and place of
the miraculous occurrence.

155.
It is a more serious difficulty that Eusebius's statement is not
supported by other Fathers of his own and the following century; yet
this is not so great as at first sight appears. It is not pretended
that any of them contradicts or interferes with his account of the
matter; and at the very time, there were no great ecclesiastical
writers to speak one way or the other. The miracle is said to have
taken place in 311 or 312; the only writer of note extant during the
first fifty years of the century, besides Eusebius is Athanasius; and
his writings are taken up with later transactions and a far different
subject. Nor does there seem any special reason why later writers
should mention it \footnote{Columbus, however, in Lactant. de Mort. Pers. 44, refers to St.
Gregory Nazianzen's second invective against Julian, where he speaks
of the Cross seen in the air when the works at the Temple were
miraculously stopped, and observes that he certainly would have
alluded to Constantine's Cross, had he known of it. Yet surely he
would have been going out of his way to do so, considering it was but
one portent out of many which he was recounting, and another Cross had
been seen over Jerusalem in St. Cyril's time.}. The
real miracle was encompassed with even heathen fables; the classical
or philosophical description contained in the panegyric of Nazarius had been almost coeval with its occurrence, and was not
likely to prejudice the Church in its favour, and yet was, as far as
we know, the only testimony by which it was conveyed to the Fathers of
the fourth century. At least Gibbon himself grants that they were not
acquainted with Eusebius's statement, and grants it for the very
reason that they did not avail themselves of it. He confirms this
opinion by the fact of St. Jerome's ignorance of the Life of
Constantine, in which Eusebius reports the miracle, a work which he
considers ``was recovered by the diligence of those who translated or
continued his Ecclesiastical History.'' \footnote{Ibid. Ch. xx. note 52. [Vid. Danz. de Euseb. Cæsar. p. 71.]} Nor does it appear why the Fathers of the Church should have
mentioned the miracle, even had they known it. It was not a miracle
especially addressed to them, or wrought for the uses of the Church at
large. It was, first, a fitting rite of inauguration when Christianity
was about to take its place among the powers to whom God has given
rule over the earth; next, it was an encouragement and direction to
Constantine himself and to the Christians who marched with him; but it
neither seems to have been intended, nor to have operated, as a
display of Divine power to the confusion of infidelity or error. In
like manner, while the Fathers appeal to the fiery eruption at the
Jewish Temple, because it was the means of a signal triumph over an
enemy, on the other hand, they refer to the destruction of
Jerusalem without mentioning the prodigies which attended it. The
distinction is clear. First, the taking of Jerusalem and the
conversion of Constantine were events of a past day; Julian's
antichristian attempt was of their own day. Again, the portents in the
sky and the luminous Cross did but concur with and, as it were,
illustrate the march of events, which was evident to all men without
them; but the fire which burst forth when Julian would rebuild the
Temple, was in opposition to the apparent course of things, and
arrested them and defeated them. It did a deed, whereas the luminous
Cross did but herald one.

156.
It may be added, that there is a beautiful harmony and contrast in the
omens by which the overthrow of Judaism and Paganism were respectively
preceded. The omens in the former instance were only evil, for the
chosen people was falling away; but since the nations were to be
brought into the Church who had hitherto been outcasts, the sign in
the heavens in the latter case was the Cross itself, a terror indeed
and dismay at first sight to the ignorant Pagan beholders, but their
redemption and salvation under the awful compulsion of Him who
suffered on it.

\xsection{Section 5. The Discovery of the Holy Cross}

157.
IN the year after the Nicene
Council, A.D. 326, St. Helena,
mother of Constantine, and then nearly eighty years of age, went on a
pilgrimage, as it was afterwards called, to the Holy Land, and
especially to Jerusalem. Her purpose was to visit the scene of the
wonderful events recorded in Scripture, and the spots consecrated by
the presence of our Lord. Among other objects of her pious search was
the Cross upon which He suffered. It was the custom of the Jews to
bury the instruments of death with the corpses of the malefactors \footnote{``Accedit consuetudo Judæorum quibus solemne
instrumenta suppliciorum juxta cadavera sontium obruere.'' Gretser de
S. Cruc. tom. 1. i. 37, he refers to Baronius and Velser. S. Basnage
agrees, Annal. 326. g. [Vid. also Aringhi Rom. Subterr. p. 98, ed.
1659.]}; and, considering their eagerness to remove the bodies both of
Christ and of His two companions before the approaching Feast, there
seemed no reason to doubt that, after Joseph had begged His body of
Pilate, and placed it in the neighbouring tomb, His Cross, and those of the two thieves, as well as their corpses, had hastily been
thrown into the ground on the very place of crucifixion. But where
that place was, at first sight, was not so easy to determine. The city
had been destroyed, and its soil (it is said) ploughed up, in
punishment of the very deeds, of which Helena was seeking to recover
the memorial. Our Lord had suffered outside the walls; but the
population, driven beyond Mounts Sion and Acra, which it had hitherto
occupied, had overflowed toward the north, and, without as yet
covering Calvary itself \footnote{St. Cyril says of Calvary, that ``it was a garden before, and the
tokens and traces thereof remain.'' Catech. xiv. 5.},
had obliterated the features of the immediate neighbourhood. And
Hadrian, by erecting statues of the Pagan divinities over the sacred
spots which were in question \footnote{St. Jerome mentions Hadrian by name. Eusebius, in the vague way which
he adopts on other occasions (as not writing a history but a
panegyric, V. Const. i. 11), says, ``Ungodly men formerly, or rather
the whole race of demons by means of them.'' V. Const. iii. 26. In like
manner he speaks of ``tyrants of our days who essayed to fight against
the God of all, and oppressed His Church,'' i. 12, meaning Dioclesian;
``the Emperor who had first rank,'' i. 14, meaning the same;
``tyrannical slavery,'' i. 26, \emph{i.e}., the sovereignty of
Maxentius; vid. also i. 33, etc., ``news came that some dreadful beast
was attacking,'' etc. viz. Licinius. i. 49, ``news came of no small
disturbance having possessed the Churches,'' ii. 61, meaning the Arian
controversy; ``a man well approved by Constantine for the sobriety of
faith,'' etc., meaning Hosius, ii. 63; ``the ruling city of Bithynia,''
meaning Nicomedia, iii. 50, etc., etc.},
had driven away worshippers, and at length effaced all general
recollections of their respective localities. But what had destroyed
the tradition about them with the many might reasonably be expected to
be the means of preserving it with the few; nor did it seem difficult,
even without such accidental advantage, to recover, with proper pains,
at least the general position of the spot where so great and memorable
a deed had been done. The Empress availed herself of the assistance of
the most learned both of Christians and of Jews \footnote{This consideration answers, as far as the present question is
concerned, Professor Robinson's remark, that ``the Fathers of the
Church in Palestine, and their imitators the Monks, were themselves
for the most part not natives of the country.'' Palestine, Vol. i. p.
373.}, and she seems to have been animated by a hope, surely not
presumptuous, that she was under a guidance greater than human \footnote{Calvin, however, considers that St. Helena was urged by ``stulta
curiositas,'' or ``ineptus religionis zelus,'' De Reliqu. p. 276.}. At length there is said to have been a general agreement as to
the place; it was covered, first with a vast quantity of earth, next
with the Pagan edifices; the place of Crucifixion and Burial lay
beneath. Helena gave the word, and the soldiers who attended her began
to clear away both buildings and soil.

158.
Hitherto the main outlines of the history are confirmed by Eusebius,
though he speaks of Constantine, his Imperial Patron, instead of St.
Helena, and only of the Holy Sepulchre, not of the Holy Cross.
And though Constantine seems, during the years 326, 327, to have
remained in the parts of his Empire between Thessalonica, Sirmium, and
Rome, yet under his direction or authority Helena doubtless acted.
Eusebius attests the intention of Constantine to build a church over
the Holy Sepulchre; its desecrated state; the huge mound and
stone-work which covered it; the shrine of Venus which had been raised
at the top; and then the demolition of the whole mass of heathenism at
the Emperor's command, statues, altars, buildings, mound, and the
earth which lay under it. He then continues thus: ``And when
another level appeared instead of the former, viz., the ground which
lay below, then at length the solemn and all-holy memorial of the
Saviour's Resurrection appeared beyond all hope; and thus the cave, a
holy of holies, imaged the Saviour's revival, and, after being sunk in
darkness, came to light again, and to those who came to the sight
presented a manifest history of the wonders which had there been done,
witnessing by facts more eloquently than by any voice the Resurrection
of the Saviour.'' \footnote{V. Const. iii. 28.} Here
Eusebius ends his narrative; he proceeds, indeed, to speak of the
church which Constantine built upon the spot, but he says nothing of
any discovery besides that of the Sepulchre itself. As to the Cross on
which our Lord suffered, judging from the course of his narrative,
we should conclude, not only that it was not found, but that it
was not even sought after, nay, according to his literal statements,
that St. Helena did not come to Jerusalem, only to Mount Olivet and
Bethlehem \footnote{Dallæus objects that St. Cyril ``nihil addit de Helenâ.'' Rel. Cult.
Obj. v. 1. p. 709. Yet we shall see Professor Robinson's reluctant
admission presently, note d \textbf{[}Note
12\textbf{]}. As to Eusebius, his object
was to praise Constantine. In V. Const. iii. 41, he first speaks as if
Constantine founded the Churches at Bethlehem and Mount Olivet, and
then corrects himself. [It is remarkable, moreover, that the Bourdeaux
Pilgrim, A.D. 333, vid.
Wesseling Itinerarium, pp. 589-596, whose silence about the Cross is
sometimes brought in corroboration, (vid. Gibbon, Hist. Ch. xxiii.,
note 63), is also silent about St. Helena's Church on Olivet, which no
one doubts about; vid. Euseb. V. Const. iii. 42, 43.]}.

159.
Yet, though he is silent himself, he has preserved Constantine's
Letter to Macarius, Bishop of that see, on occasion of the proposed
Martyry or Church of the Resurrection. This letter does not contain
any express mention of the Cross; and yet, did we read it without
knowing the fact of the historian's silence when writing in his own
person, we certainly should have the impression that it is of the
Cross that Constantine was speaking. He says that ``the \emph{token} of the
Saviour's most holy \emph{passion}'' had been ``buried under the earth
for many years;'' and he speaks of it as a discovery ``surpassing all
human \emph{calculation} and all \emph{amazement},'' and again and
again of the \emph{miracle} which it involved or had effected \footnote{Const. iii. 30. Mr. Taylor says, that ``the phrases he [Eusebius]
employs [\emph{i.e}. in Constantine's Letter] clearly imply the
invention of the Cross, although apart from other evidence they would
leave us in the dark as to the facts.'' Anc. Christ. Part vii. p. 296.
I should say the same; but it is not granted by the Centuriators (who
say that St. Ambrose is the first to mention the discovery), nor by
Dallæus, S. Basnage (who speaks of the ``intoleranda Bellarmini
sive inscitia sive audacia'' in maintaining it, Annal. 326, 9),
Hospinian, etc. The elder Protestants wish to put the ``cultus
reliquiarum'' as late as possible; Mr. Taylor as early. Tillemont
understands Constantine to speak of the Cross, but Zaccaria, Dissert.
t. 1. v. 4. § 5, is disposed, at least for argument's sake, to give
up the point.}.

160.
It is remarkable, too, that Eusebius also, though silent about the
Cross, makes mention of miracles as attending the discovery of the
Sepulchre, in a passage of his Commentary upon the Psalms. Treating of
the words, ``Dost Thou show wonders among the dead?'' he says, ``If any
one will give his attention to the marvels which in our time have been
performed at the Sepulchre and the Martyry of our Saviour, truly he
will perceive how the prediction has been fulfilled in the event.'' \footnote{Psalm lxxxvii. 13. Thus Montfaucon, but Zaccaria strangely denies the
allusion.} Yet, commenting on the 108th (109th) Psalm, he mentions the
honours paid to ``the Sepulchre of Him who was delivered over to
the Cross and death,'' without saying a word of honours paid to the
Cross itself. Eusebius died about A.D.
338, \emph{i.e}. eleven years after St. Helena's visit to Jerusalem;
and this is all the evidence which we have on the subject,
whatever is its value, for about the first twenty years.

161.
St. Cyril of Jerusalem is our next informant concerning the discovery
of the Cross. He was one of the Clergy of the Church of Jerusalem, and
delivered his Catechetical Lectures about A.D.
347, in the very Church of the Resurrection which by that time
Constantine had built; and in the first year of his Episcopate, A.D.
351, he wrote his letter to Constantius concerning a luminous Cross
which had just then appeared in the air over Jerusalem. As he died A.D.
386, and was a Priest and (as St. Jerome says) a young man in 347, he
must have been in his boyhood at the time of St. Helena's visit;
whether in Jerusalem is not known. In his Catechetical Lectures he
speaks of the Holy Cross as discovered, though he does not mention the
circumstances or the time of its discovery. Speaking of our Lord's
crucifixion, he says, ``Shouldest thou be disposed to deny it, the very
place which all can see refutes thee, even this blessed Golgotha, in
which, on account of Him who was crucified on it, we are now
assembled; and further, the whole world is filled with the portions of
the Wood of the Cross.'' Again, speaking of the witness borne in such
manifold ways to Christ, he says, ``The Holy Wood of the Cross is His
witness, which is seen among us to this day, and by means of those who
have in faith taken thereof, has from this place now almost
filled the whole world.'' Once more, speaking against the heretics who
denied the reality of our Lord's passion, ``Though I should now deny
it, this Golgotha confutes me, near which we are now assembled; the
wood of the Cross confutes me, which has from hence been distributed
piecemeal to all the world.'' \footnote{Cat. iv. 10; x. 19; xiii. 4.} Considering then that we hear nothing of the Wood of the Cross
in Ecclesiastical history before this date, and that this date follows
close upon the discovery of the Holy Sepulchre, it does not need
further proof, though St. Cyril said nothing else, that there is some
connection between this alleged discovery of the Cross, and that of
the Sepulchre. It does not need St. Cyril's express statement to that
effect in his Letter to Constantius a few years later, as we there
read it; and it does not matter, even though that letter be spurious,
as some Protestant critics, though without strong reason, contend \footnote{The authenticity of this Epistle is denied from its omission in St.
Jerome's list of his works, and its mention of the [\emph{homoousion}];
\emph{e.g}., by Dallæus, Rel. Cult. Obj. v. 1. p. 707. J. Basnage
Hist, de l'Eglise, p. 3. xviii. 13. § 2. Mr. Taylor seems to grant
it, Anc. Christ. Part vii. p. 292. The mention of the [\emph{homoousion}]
would be decisive against it, did it not occur at the end, in a sort
of doxology which will admit of being considered an addition. The
question is discussed at length by Zaccaria, ibid. in answer to
Oudinus.}. His words are these: ``In the time of thy father, the divinely
favoured Constantine of blessed memory, the salutary Wood of the
Cross was found in Jerusalem, divine grace granting the discovery of
the hidden holy places to one who laudably pursued religious objects.''

162.
From the evidence of St. Cyril and the passages of Eusebius, we gain,
then, as much as this; that the discovery of the Holy Cross was a
received fact twenty years after St. Helena's search for the Holy
Sepulchre; that it was by that time notorious throughout the world,
because portions of the Cross had been sent in all directions; hence,
that the professed discovery must have taken place some years before
the end of the twenty years when St. Cyril mentions it, to allow of
such a general publication and dispersion of it; and that it must have
been well known to Eusebius, who wrote his Life of Constantine A.D.
337, only ten years before St. Cyril's Lectures, whether he believed
it or not; further, that his silence about it did not necessarily
proceed from disbelief, because he is silent about St. Helena's search
after it, nay, as I have said above, even about her visiting
Jerusalem, an historical fact which cannot be gainsaid \footnote{``Such is the account which Eusebius, the contemporary and eye-witness,
gives of the Churches erected in Palestine by Helena and her son
Constantine. Not a \emph{word}, not a \emph{hint}, by which the reader
would be led to suppose that the mother of the Emperor had anything to
do with the discovery of the Holy Sepulchre, or the building of a
Church upon the spot. But ... all the writers of the \emph{following}
century relate as with one voice that the mother of Constantine,'' etc.
Robinson, Palestine, Vol. ii. p. 14. Yet this same writer says in the
very next page, ``Leaving out of view the obviously legendary portions
of this story, it would seem \emph{not improbable}, that Helena \emph{was}
the prime mover in searching for and discovering the sacred Sepulchre!''}, and again, because, in his Comment on the Psalms, he
speaks also of miracles wrought at the Sepulchre, of which
nevertheless he says nothing at a later date, in his Life of
Constantine; lastly, that Constantine recognized the discovery of the
Cross at the very time, because, while the terms \footnote{[\emph{Gn\={o}risma tou pathous}], V. C. iii. 30; whereas Eusebius
calls the Sepulchre [\emph{t\={e}s athanasias mn\={e}ma}], 26; [\emph{t\={e}s
anastase\={o}s martyrion}], 28.} he employs in his Letter to Macarius are more suitable to
denote the Cross than the Sepulchre, the strong expressions of his
amazement are also more suitable to the discovery of the former than
of the latter, a discovery which, as we have seen, was certainly
reported and generally believed a few years later.

163.
I conceive, then, that the evidence already brought together is
conclusive of the fact of the alleged discovery of the Cross about the
time of St. Helena's visit to Jerusalem, and in connection with that
visit. Eusebius's silence is of course a difficulty, and, as it would
appear, cannot satisfactorily be accounted for \footnote{``Notwithstanding the silence of Eusebius, there would seem to be
hardly any fact of history better accredited than this alleged
discovery of the True Cross.'' Robinson's Palest. Vol. ii. pp. 15, 16;
vid. also, ``However difficult,'' etc., p. 76.}. Yet he is silent at other times about facts which he
cannot be said to disbelieve \footnote{He does not mention St. Antony, or Methodius of Tyre, or the
Martyrdoms of Perpetua and Felicitas, etc., etc.}. We should also ask ourselves \emph{what} it is that his
silence is to be taken to prove; not that he had not \emph{heard} of
the alleged discovery, for that it was alleged is undeniable; it can
only be taken to show that he did not \emph{believe} in it \footnote{Dallæus contends Eusebium \emph{nescivisse} quod tacet. Rel. Cult.
Obj. v. 1. p. 706. The Centuriators are vague, so is J. Basnage, Hist.
de l'Egl. p. 3. xviii. 13. § 2. S. Basnage implies Eusebius's \emph{knowledge}
but \emph{disbelief} of the story. Annal. 326. 9. Jortin says that
Eusebius ``\emph{either} knew nothing \emph{or} believed nothing of it.''
Eccles. Hist. Vol. ii. p. 223.}. Yet his statement elsewhere, that certain miracles occurred
at the Sepulchre, while it suggests some further story which he does
not relate, is favourable, as far as it goes, to his belief in the
received one. Moreover, if the discovery was not really made, there
was \emph{imposture} in the proceeding \footnote{Basnage considers it a pious fraud of St. Cyril's. Annal, 326. 9. Mr.
Taylor prefers to impute it to Macarius to imagining ``Cyril and his
colleagues to have hatched the fraud coolly and at leisure twenty
years afterwards.'' Anc. Christ. Part vii. p. 297. ``Cyril of Jerusalem
and Augustine are the two Fathers who may be believed to have been the
\emph{dupes} of, rather than the actors in, the frauds of their times.''
Ibid. p. 292. ``It would perhaps not \emph{be doing injustice} to the
Bishop Macarius and his clergy,'' says Professor Robinson, ``if we
regard the whole as a \emph{well-laid} and \emph{successful} plan for
restoring to Jerusalem its former consideration, and elevating his see
to a higher degree of influence and dignity.'' p. 80.}; an imputation upon the Church of Jerusalem, nay, in the event
on the whole Christian world, so heavy, as to lead us to weigh
well which is the more probable hypothesis of the two, so systematic
and sustained a fraud, or the discovery of a relic, or in human
language an antiquity, three hundred years old.

164.
Now let it be observed that hitherto this passage of history has had
nothing definitely miraculous in it. It does but relate to the
discovery, by ordinary methods of inquiry, of an instrument of death
used by Roman executioners three centuries before. And perhaps it is
right to draw a line between the above testimony and the evidence
which follows at a later date, and which is next to be considered,
except so far as the later evidence happens to be confirmatory of the
earlier. It would seem impossible that the original story should not
receive a colour or an exaggeration \footnote{We have an instance of such exaggeration in the report of the
Samaritan woman, ``Come, see a man which told me \emph{all} things that ever I
did.'' If the men of her city had been instructed in Protestant
divinity, they would have cross-examined her as to what she meant by \emph{all},
and then said that there was evident inaccuracy, and grounds for
suspicion, that they were not called upon to stir, that they were not
obliged to believe her, etc., etc.} when taken up as a matter of popular belief, and that in
countries far removed from the scene to which it belongs. While, then,
we may be prepared for additions which will not compromise the
original evidence, those additions, when we find them in subsequent
writers, whether true or false, are exposed \emph{primâ facie} to a
suspicion which does not attach to the particulars which we have
hitherto been reviewing.

165.
Now St. Ambrose, in his discourse upon the death of Theodosius, (A.D.
395,) and St. Chrysostom, in his Homilies upon St. John, (about A.D.
394,) speak of three Crosses, not one; and say that the True Cross was
known by the title which Pilate fixed on it.

166.
St. Paulinus, St. Sulpicius, and Theodoret, agreeing in the main in
the additional circumstances related by St. Ambrose and St. Chrysostom,
differ from them in assigning a miracle as the test by which our Lord's
Cross was ascertained. Paulinus writes to Sulpicius, and the latter
reports in his history, that it was distinguished from the other two
by its restoration of a corpse to life. These two authors write about A.D.
400, Paulinus in Italy, Sulpicius in Gaul. Theodoret, who wrote his
Church History, about A.D. 440,
in Syria, speaks, not of a corpse restored to life, but of a sick
woman restored to health.

167.
Again, Rufinus, (about A.D.
400,) and Socrates and Sozomen, (both about A.D.
440,) say that the inscription was detached from the Cross, and that
the female, who was the subject of the miracle, was only on the point
of death. Moreover St. Ambrose, Rufinus, Theodoret, Socrates, and
Sozomen speak of the nails as found at the same time. A further
miracle is spoken of by Paulinus; that the portion of the Cross kept at Jerusalem gave off fragments of itself without diminishing \footnote{Jortin, Eccles. Hist. (Works, Vol. ii. p. 222), translates Tillemont
as believing this miracle, and as saying that ``St. Paulinus relates a
very \emph{singular} thing,'' putting the words in italics, and
prefacing his extract with observing that ``the words of Tillemont are
full of what the French call \emph{unction}, and the English \emph{canting};''
whereas in fact Tillemont at the least doubts Paulinus's account. Mem.
Eccles. Vol. vii. p. 8.}; and he adds, that ``it has imbibed this undecaying virtue and
this unwasting solidity from the blood of that Flesh, which underwent
death, yet saw not corruption.'' \footnote{Ep. 31. fin.} This is mentioned here, as being one of the alleged miracles
which followed upon the discovery of the Cross, though it has no
connection with the discovery itself, which is our proper subject.

168.
Such is the evidence arranged in order of time, in behalf of this most
solemn and arresting occurrence \footnote{St. Jerome too says of St. Paula, A.D. 386, ``Prostrataque
ante Crucem, quasi pendentem Dominum cerneret, adorabat.'' Ep. 108. n.
9.}, which is kept in memory, even to this distant generation, in
the Greek, Latin, and English Calendars on the 3rd of May and the 14th
of September. It seems hardly safe absolutely to deny what is thus
affirmed by the whole Church \footnote{``This history of the discovery of the Holy Cross,'' says W. Lowth
in Socr. i. 17. ed. Read., ``is not found in Eusebius. But Cyril,
Bishop of Jerusalem, who lived in the same age, openly witnesses that
the Wood of the Holy Cross was divinely shown to the Emperor
Constantine; also in his Catechetical Lectures he speaks of its
discovery, as of a thing known to all. Wherefore of the faith of this
history we cannot doubt.'' Upon this Jortin asks, ``What did this
Protestant Divine of ours mean? Could he believe that the True Cross
was found? or would he only say that a pretended one was discovered?''
Ibid.}; whether however miracles accompanied the discovery,
must ever remain uncertain. That a sick, dying, or dead person was
restored by means of the Cross, rests on the authority of Latins
writing at the distance of seventy years after the discovery, and of
Greek authors of forty years later still; not on any testimony given
with particularity or at the time. Moreover, such an occurrence is
inconsistent with the account, taken in the letter, of St. Ambrose and
St. Chrysostom, who say that the True Cross was recognized by its
title. On the other hand, whether there was one Cross or three \footnote{S. Basnage urges, however, ``vero non esse proximum, latronum cruces
cum illa Christi, uno eodemque loco fuisse conditas;'' Annal. 326. 9. \emph{because}
crosses were always buried with the bodies, \emph{but} no bones were
found with the Cross,—an assumption. There were too many bones
surely in ``the place of a skull,'' to discriminate, or to mention the
fact.}, some \emph{mode} of recognizing it is implied in the very
idea of recognition; and a miraculous recognition is perhaps the most
natural and obvious hypothesis. Nay, the very fact that a beam of wood
should be found undecayed after so long a continuance in the earth
would in some cases be a miracle \footnote{Quæri et illud potest, num annis pænè trecentis in solo absque
putredine cessante miraculo,'' etc. S. Basnag. Ann. 326, 9.}. And perhaps there are few imaginations, which are once
able to surmount the shock of hearing that the very Cross on which our
Lord suffered was really recovered, but will be little sensitive of
difficulty in the additional statement, that miracles were wrought by
means of it. It must not be forgotten too, that Eusebius himself,
though silent about the Cross, alludes to the occurrence of miracles
at the Sepulchre; and these of course become more credible, if we
suppose some great object, such as the recognition of the Cross, to
account for them.

169.
An objection, however, has from time to time been urged with much
earnestness by several writers, which, if substantiated, would
altogether overthrow the history of the discovery of the Cross; viz.,
that Helena chose a wrong site for the Holy Sepulchre. This was Dr.
Clarke's opinion, whose reasons were discussed and answered by a
writer, it is believed Bishop Heber, in the \emph{Quarterly Review}
for March, 1813. It has lately been revived with some additional
considerations by one or two controversialists, one of them at least
with the view, not simply of disproving the fact, which is a point of
secondary importance, but of fixing upon the Fathers and Church of the
fourth century the imputation of deliberate imposture, and that for
selfish ends. Indeed, a drift of this kind is the only intelligible
explanation of the earnestness which such writers manifest \footnote{The Quarterly Reviewer of December, 1841, professes to consider it
only a question of ``poetic statement,'' ``fond reminiscences,'' ``reverential
feelings,'' ``pleasing visions,'' and the like; and contrasts with them ``the
stern voice of truth,'' etc., etc., whereas the simple question is
whether we shall consider the Church of the fourth century very
credulous or very profligate. Mr. Taylor is far more perspicacious.}. It might not seem to be worth any great exertion to construct
a proof that the Holy Sepulchre was not found by St. Helena, if no
conclusion is to follow but that we need not pay attention to the
festivals of the Cross, or to the claims of particular pieces of wood
professing to be fragments of it. Even admitting the True Cross was
discovered, it would be still open to Protestants to treat it with
neglect or disrespect, and they would doubtless exercise their right.
The Cross on which Christ suffered would be in their eyes but a piece
of wood; or again, as they sometimes speak both of it and of the sign
of it, it would be a something loathsome and hateful, bringing our
Lord under the curse rather than sanctified by Him; and that the more,
because, like the Brazen Serpent, it had been the occasion of
superstition and idolatry. When then writers set themselves to oppose
passages of history such as that now before us, it is for a far bolder
purpose than is directly implied in their opposition; it is of course
in order to depreciate or destroy the authority of the Church. It is
an attempt to transfer the quarrel between her religion and
their own, from the province of opinion to the ground of matter of
fact; nor can there be in itself a fairer procedure towards the
Church, or one of which her children have less cause to complain,
however they may be pained at the spirit in which it originates.

170.
Nay, perhaps such controversialists are fairer to the Church than to
themselves; and though undoubtedly, if they once prove their point,
they will but gain the greater credit and the more decisive victory,
by so frank a procedure, yet it is plain there is at first sight a
very strong probability \emph{against} their proving it. The chance
is, that they have undertaken more than they can accomplish. For it
stands to reason, which of two parties is the more likely to be right
in a question of topographical fact, men who lived three hundred years
after it and on the spot, or those who live at a distance of eighteen
hundred, and at the Antipodes? Granting that the fourth century had
very poor means of information, it does not appear why the nineteenth
should have more ample. There are indeed branches of knowledge in
which we have decidedly the advantage of the early Church. If it were
a point of philology which was under review or a question about the
critical interpretation of the Hebrew text, or the etymological force
or derivation of a Greek term; or a problem in physics, such as
whether or not such and such an occurrence were beyond, or beside, or
according to the laws of nature, and properly miraculous; or if
it required some subtle analysis, and could only be wrought out by a
mathematical formula; or if, though a question of history, or
chronology, or topography, it was disputed by writers of the fourth
century one with another, so that we must oppose this great name or
that, and choose a side; or if it was advanced by some one Father, and
by him unsupported, and in himself of no great authority, then the
attempt to contradict him would be plausible; but in such a matter
and under such circumstances as the present, when Calvary is the spot
and Eusebius the informant, when a very learned and not over-credulous
writer, whose silence about the Cross is thought so ominous of his
disbelief, reports and assents to the unanimous decision of the local
Church in favour of the discovery of the Sepulchre, and is supported,
not merely by the silent assent, but by the concurrence of the
literature of his own century \footnote{The main \emph{authority} for the present site of the Holy Sepulchre
is Eusebius; and the warrant for its preservation or recovery is the
Pagan Temple raised over it upon the destruction of the city by
Hadrian, which became a lasting record of the spot. What is to be
urged against Eusebius I know not; but it is urged against the
argument from the Pagan Temple, first, that only St. Jerome, and not
Eusebius, attributes its erection to Hadrian. But why is not St.
Jerome, a learned Father, well acquainted with Palestine, and no
friend at all (as Mr. Taylor allows) to the superstitions and
pollutions of which in his time Jerusalem was the scene, why is he not
a sufficient authority? whereas Eusebius (\emph{after his manner})
does but say ``ungodly men;'' vid. supr. n. 157, note t \textbf{[}Note
3\textbf{]}. Next, it is objected that
there was no \emph{recognized tradition} of the spot, because St.
Helena had to \emph{search} for it, and to summon learned Jews and
Christians to her assistance: but it does not follow, because there
was no popular tradition, that therefore there was no historical and
antiquarian knowledge of the fact, or, again, no means of recovering
it, though forgotten. Further, it is urged that it was unlike Hadrian's
character to insult the Christians, when he was but punishing the
Jews. But, granting his general leniency towards the former, what
Sulpicius says, Hist. Sacr. ii. 45, suggests the conjecture, that,
from the circumstance of the Jewish Bishops not only being natives and
inhabitants of the place, but practising \emph{circumcision}, he
confused them and their flocks and the objects of their veneration
with the Jews. From these three considerations, (1) that St. Jerome is
the first informant that Hadrian placed a Temple over the Sepulchre;
(2) that there was no continuous public local tradition to that
effect; and (3) that it is a deed unlike Hadrian, it is proposed to
infer that a place was pitched upon at random as the site of the
Sepulchre, and that, among all places in Jerusalem, a heathen temple.
As to the actual Sepulchre found under the mound, that of course is
the work of \emph{fraud}.}, the presumption is very great, before going into the
case, that such acute and ingenious persons, as now for the first and
last time in their lives traverse Jerusalem with their measuring tape,
are wrong, and those who were natives of the place fifteen hundred
years ago are right. Of course such presumption constitutes no plea at
all for declining to examine their arguments; it weighs nothing
against an overwhelming proof, when they have brought it; but, to use
the language of their school, when speaking of the miracles of the
Church, unless the proof is overwhelming we are not ``obliged'' to
accept it; when there is but a balance of arguments, ``we may suspect
it'' to be fallacious, and may pronounce it ``unsatisfactory.'' It
may be entered upon with a just prejudice, listened to with suspicion,
criticized with fastidious precision, and rejected on the ground of
counter-arguments in themselves of inferior cogency.

171.
Let it carefully be observed that a point of evidence like this has
nothing to do with the question of honesty or dishonesty in the
parties who give it. Were Macarius or St. Cyril and the Clergy of
Jerusalem the most covetous and unprincipled of hypocrites, why should
this lead them to fix on a false site for our Lord's crucifixion and
burial? Why should they not do their best to fix on the right one? why
should they subject themselves to an additional chance of detection,
and give to persons like their present impugners a gratuitous
advantage, as if it were not enough to fabricate a Cross, but they
must hazard a superfluous mistake in fabricating a Sepulchre? Were
they then knaves and impostors of the deepest dye, this would be no
reason for their omitting, nay, the strongest reason for their taking,
all possible pains to find the very and true spot where our Lord
suffered. And therefore the question returns to the issue on which it
has already been put—which is the more \emph{likely}, that
inhabitants of Jerusalem in the fourth century, or of New York in the
nineteenth, should be able rightly to determine Calvary and the
Holy Sepulchre \footnote{As if to meet this presumption, Dr. Robinson and Mr. Taylor set
themselves to prove incontrovertibly that St. Helena did fix the site
of the Ascension on Mount Olivet wrongly; and if she was wrong in one
case, she might be in another. And their proof is as follows: 1st, St.
Luke in his Gospel says that our Lord led out His disciples as far as \emph{Bethany},
therefore He ascended from Bethany (in spite of his saying in the
Acts, that they returned from \emph{Olivet}); 2nd, St. Luke says that
from Mount Olivet to Jerusalem is a Sabbath-day's journey, but the
Church of the Ascension is only half a mile from Jerusalem. It has
been usual to say in answer (1) that Bethany was not only a village,
but a district which extended over a portion of Olivet. Thus Bethany
is considered by Lightfoot, Chorogr. in Matt. 37; ibid. 4; Chorogr. in
Marc. 4; Hor. Hebr. in Luc. xxiv. 50, and in Act. Ap. i. 12. In a
previous work he had thought otherwise, vid. Comment, in Act. Ap. on
the ground that names of towns and names compounded with ``Beth''
never were extended to a district. He gets over this difficulty, in
his later work, by saying that the town was called from the district,
not the district from the town. The same explanation of ``Bethany'' is
given by Beza, Grotius, Sanctius, and De Dieu in Poole's Synopsis.
Again, Spanheim calls it ``tractus montis Oliveti.'' Geogr. Sacr. part i.
fin. (2) As to the alleged difficulty of the Sabbatical distance, it
is not really such, till critics are agreed \emph{what} that distance
is. ``Iter Sabbaticum \emph{octo stadia} excepit aut totum milliare.''
Lightfoot, Chorogr. in Matt. 40. Elsewhere he says, ``Iter Sabbaticum
ex septem et dimidio,'' in Luc. xxiv. 50, adding, that it was ``bis
mille cubitorum.'' Yet, Comment. in Act. i. 12, he says that, while the
Sabbatical distance is nine stadia if the cubit is three feet, it is
but four and a half if the cubit be a foot and a half; and that the
latter is the true calculation. ``What space is a Sabbatical journey?''
says Drusius in Poole's Synopsis; ``in the number two thousand most
agree; but some say cubits, others paces, Jerome feet, Origen ells,
which Origen's translator calls cubits.'' De Dieu (ibid.) with
Lightfoot in Act. makes the cubit a foot and a half, or the Sabbatical
journey about five stadia, which is the distance of Mount Olivet,
according to Josephus, and the actual distance of the Church of the
Ascension, from Jerusalem. Grotius considers it eight stadia. Reland
quotes Origen for its being eight (Palæst. i. 52. fin.), but thinks
this too much; and quotes Epiphanius for its being six, which
according to him (vid. Wolf in Act. Ap. i. 12.) would make the
Sabbatical journey a quarter of an hour's walk. It must be added that,
if the Church of the Ascension is short of the Sabbatical distance
(which, as we see, it is not proved to be), at all events Bethany is
in excess of it.}. I mean
no disrespect to the traveller to whom I allude, which is not due to
one who accuses the Church of Jerusalem of the fourth century of
deliberate fraud. And I make a great distinction between a learned
person like himself, who writes with gravity and temper, and the
English writer who has made use of his statements in a work to which
reference has been several times made at the foot of the page. Yet I
do not see why weak arguments should be treated with indulgence
because they are directed against sacred persons and better times. In
order to form a due estimate of them, we must now consider with some
attention the site of Jerusalem.

172.
Jerusalem in our Lord's time occupied four hills: Sion, on which was
the city of David, or the upper city, on the south (A); on the north
of it, Acra, the site of Jerusalem proper, or the lower city (B); to
the east of Acra, Moriah, on which was the Temple (C); and to
the north of Moriah, Bezetha (D), on which lay the new city,
containing the overflowings of the population, which were at that time
very considerable. Denoting them by the four letters, A, B, C, and D,
we shall have the nearest idea of their relative position by
considering B and C on a line running from west to east; A under B,
and D above C. The Church of the Holy Sepulchre is in the space
between B and D, and the question, roughly stated, is whether the city
wall which went across from B to D fell within or without its site.

173.
The first and most ancient wall, which included A and C, Mounts Sion
and Moriah, ran nearly in a straight line, on its north side, which
alone it concerns us here to consider, from the north-west corner of
Sion, where Herod afterwards built a tower called Hippicus, to the
western portico of the Temple on Moriah. Accordingly the lower city
(B) was exterior to it, or, if any part was included, it was whatever
lay in the angle between Sion and Moriah, A and C. This wall is
supposed to be as early as David's time; and in its northern line,
which thus divided Sion and Acra (A and B), stood the gate Gennath,
from which in process of time was drawn a second wall, across or
around Acra, terminating in the tower Antonia, which stood at the
north-west corner of the Temple, opposite Bezetha (D). After our Lord's
death a third wall was drawn by Agrippa, which inclosed Bezetha also;
but with this we are not concerned. At the time of His
crucifixion the second wall was the limit of the city; and the
question in controversy is whether that second wall went across Acra
or outside it. Our Lord suffered and was buried outside the wall: now
St. Helena's traditionary site is nearly upon the descent of Acra on
the north \footnote{Clarke and Maundrell make it on the edge of Mount Moriah; but Dr.
Robinson ``directly on the ridge of Acra.'' Vol. i. p. 391.}. If, then,
the wall traversed the hill, that site falls without the wall; if it
inclosed the hill, within it.

174.
The argument advanced by the learned writer in question for the latter
of these alternatives is of the following kind: he admits that a
straight line drawn from Gennath to Antonia would fall short of the
Church of the Holy Sepulchre, but he observes truly that Josephus
expressly says that the wall had a curvature in the interval between
its extreme points. Next he draws attention to a certain pool called
the pool of Hezekiah, which is situated only a little to the
south-west of the Church of the Sepulchre, and which he remarks
Hezekiah formed inside the city. If, then, the pool was within the
wall, the Church must have been within the wall too; for the wall
could not include the former and exclude the latter without making a
short turn at their point of contact, which we have no warrant in
supposing. Further, unless we suffer the wall to extend beyond
the alleged site of the Holy Sepulchre, we shall not allow room enough
for the lower city, which in our Saviour's time was extensive and
populous. And, lastly, there are certain circumstances in the ground
which interfere with the supposition of a narrower circuit; and there
are some large hewn stones far to the north, near the gate of
Damascus, which, being masonry probably of an age long anterior to
that of our Lord, cannot belong to the third, and therefore are
probably a part of the second wall.

175.
Now, of these arguments it is obvious that not much stress can be laid
upon the last. The ground on which Jerusalem stands has gone through
many alterations at various times: valleys have been filled up,
summits have been levelled. If the surface was so much changed that
Helena could not at once find the Holy Sepulchre, surely its changes
are great enough to hinder a modern traveller from determining, by its
present appearance, the course of the second wall. Nor does it follow,
though the ruins at the Damascus gate are as early as the date of the
Jewish Kings, as this writer implies, that therefore they belong to
the second wall, for he does not prove that the second wall existed so
early as that time, as will be noticed presently; and they may be
remains of some other ancient work, even if it did.

176.
Nor surely is it any great objection that the lower city will be
much straitened if we draw its boundary short of the present Holy
Sepulchre; for Josephus expressly speaks of the scantiness of the
limits of the city, and of the population exceeding them in
consequence. The population covered the north side of the Temple
Mount, and then crossing the deep trench which bounded it, overflowed
upon the opposite hill Bezetha, where it formed two large suburbs, one
or both of which were called the ``New City,'' \footnote{Josephus in one place speaks of ``Bezetha \emph{and} the new city;''
Bell. Jud. ii. 19. § 4; at another of ``[\emph{h\={e} kat\={o}ter\={o}
kainopolis}],'' v. 12. § 2. If this second suburb was to the west
of Bezetha, it must occupy the north of the present Sepulchre; which
would almost be a proof that the present Sepulchre was without the
second wall.} and both of which were external to the then walls. Mount Sion
itself also contained ample room for a population, the original city
of David being but a citadel upon it. And, for what appears, the
eminence or ridge Ophel to the south of Moriah afforded additional
accommodation \footnote{Sion seems to have been covered with streets and private houses, in
spite of its public buildings. Josephus says that ``the houses'' on Sion
and on Acra ``\emph{ended} in the ravine between them:'' v. 4 ; and says
that Sion was called ``the Upper Agora,'' which implies a population. De
Bell. Jud. v. 4. In the sacking of Sion, he speaks of small houses on
it ([\emph{d\={o}matia}]), and lanes or alleys, [\emph{sten\={o}poi}],
ibid. vi. 8. § 5. Indeed, we might infer a population from the length
of the hill, which was far beyond the needs of a citadel, palace, and
public buildings. Manasseh, too, seems to have taken in a space beyond
the city of David to the south; 2 Chron. xxxiii. 14. Bezetha seems to
have been very thickly inhabited; Josephus speaks of the shops of
mercers and braziers, the clothes-market, and the alleys running upon
the city wall. De Bell. Jud. v. 7. fin. The lower city, too, was full
of alleys or narrow lanes, as appears by the following chapter. The
Tyropœum, or deep valley between Acra and Sion [and Acra and Moriah],
was ``ædificiis densa.'' Spanheim's Geograph. p. 49. The height of the
houses, too, in such localities should be considered. At Rome a law
was passed by Augustus, that houses should not be above seventy feet
high. The poor inhabited them in floors. Vid. Gibbon, Ch. xxxi. [vid.
also Merivale, Rom. Empire, Vol. 5. Ch. 40.] It is a question, too,
whether a portion of the inhabitants did not live in the excavations
under Sion and Moriah. The deeper caves were used for the purposes of
concealment in the sack of the city by Titus. Lightfoot tells us,
Chorogr. in Luc. 1. § 6, that both Iturea and Idumæa were remarkable
for their caverns, and he even derives the name of the former from
this circumstance. Strabo speaks of two caverned mountains, one of
which would hold four thousand men; Lib. xvi. p. 1074. The cave of
Zedekiah, according to a Rabbinical authority, whom Lightfoot quotes,
held eighteen thousand. And according to William of Tyre there was a
cave on the other side of Jordan, sixteen miles from Tiberias, with
different stories in it. Vid. also Joseph. Antiqu. xv. 10. § 1. It is
the Ecclesiastical tradition that a cave was the place of the
Nativity; S. Justin Martyr notices it, and Origen says that in his day
it was visited by pilgrims. However, Dr. Robinson brings this
tradition specially as a sample of the spuriousness of traditions
about sacred history in general, because a cave or grotto is
introduced. Nothing, he says, is done without grottoes. As if, too,
some traditions might not be true and some false; the latter
imitations of the earlier.}. But
however this be, the simple question for us to consider is
whether the \emph{deduction} from the supposed larger area of the city
which adherence to the present site of the Sepulchre requires, will
materially lessen that area. I conceive not \footnote{The Quarterly Reviewer for December, 1841, says: ``One argument appears
to us absolutely insuperable. To exclude the Church of the Holy
Sepulchre, the ancient city, that is, \emph{the part between the western
wall and} the hill of the Temple, must be narrowed to less than a
quarter of a mile.'' This is an inexplicable statement. It assumes that
the second wall always continued \emph{at the same distance} from the
Temple Mount which it had over against the Sepulchre.}. If the area be too scanty for the population with this
reservation, will it be sufficient without it? Sion, the greater part
of Acra, and Moriah, within the walls, and Bezetha outside of them,
remain; and if we suppose the wall which is in question, on starting
from Gennath, first to run north, and then to curve round, when it
came over against the site of the present Latin convent, very little
of Acra will be lost. Dr. Robinson refers to a passage in Josephus, in
which that historian speaks of a northern and a southern portion of
the second wall, a mode of expression which requires some such change
of direction to account for it.

177.
Nor is it easy to see how the New City can be altogether excluded from
the second wall, as we know it to have been, if the second wall is
extended any great way beyond the present Holy Sepulchre; and if it is
not extended, how any great increase of room will be obtained merely
by including the Holy Sepulchre. Nor is it natural in Josephus to
speak of the population overflowing across the trench of Moriah upon
Bezetha, if it lay all along the west of the latter hill
already, and had thence extended itself upon Bezetha eastward. In
short, if there be a difficulty in accommodating the population, it
lies in this, that the hill of Acra, from Hippicus, on the north-west
corner of Sion, to the east side of the Temple, is little more than
the third of a mile across, as Dr. Robinson measures it \footnote{Dr. Robinson says that ``the \emph{breadth} of the city is the same now
as anciently,'' Vol. ii. p. 67; \emph{i.e}., to show that it could
spare nothing in length; now he says elsewhere that the breadth from
the brow of the valley of Hinnom near the Yaffa gate to the brink of
the valley of Jehoshaphat is 1,020 yards; while the length, measured
on his map, from Herod's Gate to the limit of the ancient city on the
south, is 1,700 yards or short of a mile. Therefore an area of a mile
by 5/8 of a mile is \emph{greater} than the site of the old city,
and makes no allowance for the Temple, fort, etc., etc., yet even this
is little more than half a square mile. Here then is a fixed limit
agreed on by all who do not adopt the random hypothesis of Dr. Clarke,
that the Hill of Evil Counsel is Sion. Might not an objection be made
to the smallness of even such an area by those who do not consider how
the population of fortified cities \emph{packs}? Nothing seems known
for certain about the \emph{ordinary} population of Jerusalem. Mr.
Greswell makes several calculations, Dissert. xxiii., which exceed
what at first sight the space could seem possibly to admit.}. No theory about the north wall of the city can dispose of
this fact \footnote{Such difficulties are of frequent occurrence in history; \emph{e.g}.,
Oxford in the middle ages is said to have had 30,000 students.}.

178.
Putting aside, then, considerations such as these, which might be
useful to corroborate a proof, but have very little intrinsic force to
create one, we come to the \emph{main circumstance} on which the
author's argument depends, and which certainly deserves a
careful consideration,—viz., the position of the pool of Hezekiah.
To judge from his plan, this pool nearly joins the Holy Sepulchre on
the south-west; and was once even considered as attached to it, and
was called from it. Now Hezekiah formed his pool or reservoir within
the city; either, then, the Holy Sepulchre lay within the city also,
or the wall ran between the pool and the Sepulchre.

179.
Now, first I would observe that there is no absurdity in the latter
supposition. Let us allow that it would involve a sharp bend in the
second wall \footnote{Mr. [Dr.] Milman has no difficulty in such a supposition; ``the second
wall,'' he says, ``began at a gate in the old or inner one, called
Gennath, the gate of the gardens; it intersected the lower city, and
having struck northward for some distance, \emph{turned} to the east,
and joined the north-west corner of the town of Antonia.'' Hist. of
Jews, Vol. iii. p. 16. And he even represents it on his plan of the
city as turning at an acute angle. Dr. Robinson himself, as is said
over-leaf cannot escape a bend. When he has brought his supposed
second wall near Bezetha, he speaks of its ``\emph{bending} southward
to the corner of Antonia.'' Vol. i . p. 468.}, which is
our author's objection to it; yet Josephus, as we have seen, expressly
speaks of a northern and southern portion of the wall, which implies a
change of direction somewhere; and even though a range be supposed for
the wall beyond the present Sepulchre, it could not materially change
its direction without considerable abruptness. Dr. Robinson observes
that the wall could not exclude the Holy Sepulchre, unless it ``made
an angle expressly in order to exclude'' it; but let it be observed,
the angle \emph{must} be made anyhow in order to arrive at Antonia;
nay, and such an angle he himself makes in his own conjectural
description of it.

180.
Again, it is obvious to remark that, supposing Calvary was a place
used for the execution and burial of criminals, as is not unnatural to
suppose, and as its name may be taken to mean, there was a reason why
the second wall, whenever drawn, should avoid it. And we know that,
wherever it was, it was close upon the wall, both from the Apostle's
saying that it was ``without the \emph{gate},'' and from the custom
of the Jews fixing their places of execution outside their cities \footnote{Deut. xvii. 5; Luke iv. 29; also 2 Kings x. 8; vid. also Lev. xxiv.
14; Numb. xv. 35. Zorn. Opusc. Sacr. Vol. ii. p. 193, upon Heb. xiii.
12, refers to 1 Kings xxi. 13; Acts vii. 59. And for the like custom
among the Romans, to Plaut. Mil. Act. ii. sc. 4; Tac. Ann. ii. 32; Hor.
Epod. 5. 99. On the Jewish cemeteries as without the cities, vid.
Lightfoot, Chorograph. in Matt. 100. However, they were far enough to
be out of sight of the inhabitants. The cemeteries of the Levitical
cities were two thousand cubits off. Ibid. Chorograph. in Marc. 8. §
8.}.

181.
But, next, dismissing this question, we come to this most important
and remarkable circumstance, which will strike most readers even at
first sight; viz., that the author under review, whose learning none
can question, and whose zeal for Scripture all must honour, has fixed
the site of Hezekiah's pool by \emph{tradition}, and \emph{tradition
alone}. He says that Hezekiah ``built within the city a pool,
apparently the same which now \emph{exists under his name},'' and upon
this traditionary determination of the pool of Hezekiah he proceeds to
deny the faithfulness of the tradition concerning the site of the Holy
Sepulchre. Yet it does not at all appear why the latter tradition is
not as good as the former, especially since far greater pains have
been taken to ascertain the site of the Sepulchre than that of the
pool. Nor can it here be urged that springs of water are not of a
nature to be formed at will; that they have a perpetuity and a
possession of the soil which mounds, or walls, or sepulchres have not;
that they are not of common occurrence in Jerusalem, and that there is
no great choice of pools between which the tradition might err. This,
indeed, would be an argument if the pool were any more than a
reservoir; and that too, as Dr. Robinson himself observes, in part at
least of modern workmanship; but as the case stands, of course it is
quite inapplicable. Nor can he intend to make a distinction between
Christian tradition and Jewish, as if the Jews were deserving of more
consideration and credit than the interested clergy or the
superstitious laity of the Christian Church. For he candidly admits
that the destruction of Jerusalem by Hadrian involved the destruction
of all their local recollections. ``It may perhaps be asked,'' he says, ``whether
there does not exist a Jewish tradition, which would also be
trustworthy? \emph{not in respect to Jerusalem itself}; for the
Jews for centuries could approach the Holy City only to weep over it.''
\footnote{Vol. i. p. 376, note; vid. also p. 350; [vid. Lumper, P. H. t. 6, p.
660.]} By a law of Hadrian
they were forbidden to approach within some miles of the city, and
Constantine did but permit them to view it from the neighbouring
hills.

182.
It seems, then, that our author's argument against the alleged site of
the Holy Sepulchre depends on a definite and single fact, and for that
single fact he offers no proof whatever, except that very kind of
proof, and that not so good in its kind as that on which the site of
the Holy Sepulchre is at present received. He cannot tell how long the
reservoir has been called Hezekiah's pool, though he does tell us that
it used to be called by the Monks the Pool of the Sepulchre; while we
know, on the other hand, that the Sepulchre was fixed in its present
site as much as fifteen hundred years ago. He does not know under what
circumstances the pool was determined to be Hezekiah's; whereas we do
know that the site of the Sepulchre was settled after a public formal
examination, and, as it is reported, with the united aid of learned
Jews and Christians, and with a unanimous decision. Yet, if the real
pool was within the wall, and the real Sepulchre without it, and if
their professed sites are so close to each other that both must
have been without or both within (a point which itself, as we have
seen, is not at all clear), he asserts that the tradition concerning
the Sepulchre must be the false, and the tradition concerning the pool
must be the true \footnote{Professor Robinson, after speaking of Hippicus, Antonia, and Hezekiah's
pool, says: ``We have then three points for determining the probable
course of this wall'' (the second); ``we repaired \emph{personally} to \emph{each}
of these \emph{three points},'' etc. Vol. ii. p. 67. Now of the first
he does but say himself, ``it early occurred to us that [the tower of
David] was \emph{very probably} a remnant of the tower of Hippicus,''
Vol. i. p. 455; ``this \emph{impression} was strengthened,'' etc.; of
the second Lami says, ``I have set down several places in the map,
whose true situation is \emph{not known}; as, for instance, the castle
Antonia;'' App. Bibl. p. 76, ed. 1723, London; though Dr. Robinson
considers he has ascertained it. And what reliance is to be placed on
the site of the pool we have seen in the text. In like manner Dr.
Robinson can but say of Gennath ``\emph{apparently} near Hippicus,'' p.
411; ``\emph{doubtless} near Hippicus,'' p. 461. And of the second wall,
``Josephus's description of the second wall is very short and \emph{unsatisfactory},''
p. 461. And he locates the Tyropœum differently from other writers.
Yet on these private inferences from doubtful conjectures on probable
assumptions from unsatisfactory testimony, the Catholic Church is to
be convicted of fraud and folly.}.

183.
To proceed: it will be observed that Dr. Robinson takes for granted
another point, besides that of the existing pool being really Hezekiah's;
viz., that Hezekiah's pool was within the \emph{second} wall.
Certainly it was within the city wall, as it ran in Hezekiah's time;
but it is obvious to ask, why was it only within the \emph{second}
wall, and not within the \emph{first}, which was drawn short of
the second? Did the second wall exist as early as Hezekiah? But if
Hezekiah's pool was within the first wall, and the existing pool is Hezekiah's,
then Dr. Robinson will have proved too much; for he will have brought
up the city of David all across the valley of the Tyropœum to the
ridge of Acra on which the Holy Sepulchre stands, and within ``less
than a quarter of an English mile'' of the north-west corner of the
Temple.

184.
It is necessary then for his argument that he should clearly show, not
only that the pool really was Hezekiah's, but also that the second
wall was built and bounded the city, in Hezekiah's time. On this
point, however, he does but speak as follows; ``Of the date of this
erection,'' \emph{i.e}. the second wall, ``we are \emph{nowhere informed};
but it must \emph{probably} have been \emph{older} than the time of Hezekiah,
who built within the city a pool, \emph{apparently} the same which now
exists under his name.'' \footnote{Vol. ii. p. 67.}
That is, on the one hand, Hezekiah's pool was within the second wall,
not within the first, because the second wall, not the first, was in
Hezekiah's time the boundary of the city; and on the other hand, the
second wall, not the first, was in Hezekiah's time the boundary of the
city, because Hezekiah's pool is within, not the first wall, but the
second. Such is the author's proof of the second fact by which
he shows that the Church of the Sepulchre was built upon a pretended
site.

185.
But it may be asked whether Scripture throws no light upon the
position of the pool; for in this way perhaps the tradition respecting
it may gain an authority which it has not in itself. No tradition
certainly is tenable which contradicts Scripture; but many a tradition
deserves attention or commands assent about which Scripture is silent,
or to which it devotes but a few words or a passing allusion. Dr.
Robinson is more rigorous on this point than I should be myself;
``This is the point,'' he says, ``to which I would particularly
direct the reader's attention, that all Ecclesiastical tradition
respecting the ancient places in and around Jerusalem and throughout
Palestine \emph{is of no value}'' (and he prints the words in
capitals,) ``except so far as it is supported by circumstances known to
us from the Scriptures or from other contemporary testimony.'' \footnote{Vol. i. p. 374.} It would seem, then, as if according to his deliberate
principle, distinctly and formally avowed, some Scriptural argument
ought to be forthcoming in favour of the traditionary settlement of
the site of Hezekiah's pool;—what Scripture does say, may be told in
a very few words.

186.
In the Second Book of Chronicles we simply read as follows: ``This
same Hezekiah also stopped the upper water-course of Gihon, and
brought it straight down to the \emph{west side} of the city of
David.'' \footnote{2 Chron. xxxii. 30. If it is necessary to appeal to authority, Calmet
considers Hezekiah's pool to have been in the western quarter of the
city of David, in 2 Paralip. xxxii. 30, and fed by Gihon, in 2 Esdr.
ii. 14. So does Lightfoot, Chorograph. in Matt. 25, and in Joan. 5.
§§ 2, 3. Reland places the fount of Gihon, from which it was fed, at
the south-west. Palest. iii. p. 859.} Now, what
Gihon is, and where, is not here the question; Dr. Robinson has some
very interesting remarks on the subject, on its concealment by
Hezekiah, and on the subterraneous channels by which he fed the
reservoirs in Jerusalem. All that here concerns us to observe in this
passage are two distinct statements, each of them quite inconsistent
with the tradition that the supposed pool of Hezekiah is really the
work of that king. First, the inspired writer tells us above that
Hezekiah brought the water to \emph{the city of David}, and the
pretended pool is \emph{not} in that city; and next, that he brought
it to the \emph{west} side of the city, and the pool is on the \emph{north}
of it. What then can be said, but that this author's argument against
the truth of the alleged site of the Holy Sepulchre is based, not only
on a blind Jewish tradition, the like of which he elsewhere
reprobates, but on a disregard of the sacred text which it is the
special object of his work to exalt?

187.
In conclusion, I will but draw attention to the light which this
discussion has thrown upon the extreme improbability, which was
noticed before entering into it, that the parties who aided St. Helena
in her search should have placed the Sepulchre where we find it,
unless it were the true site. If facts are as clear as Dr. Robinson
would consider them, they were too clear for any one to miss them \footnote{Dr. Robinson begins by speaking of the ``difficulty arising from the \emph{present}
location'' of the Sepulchre ``in the heart of the city,'' which ``has been
felt by \emph{many pious minds}.'' Yet what so natural, as Maundrell
observes, as that the Sepulchre, when found, should attract the city
round it? Again, why is it not a difficulty that Sion is now so
deserted? Is not this extension, if not change, of site, what happens
to all cities of any standing? Was Dr. Robinson sceptical about St.
Giles's in the Fields when he came to London? Pope Gregory was
perfectly aware of the change of site of the city. ``Hoc quoque quod
additur,'' he says, ``Non relinquent in te lapidem super 'lapidem,'
etiam ipsa jam ejusdem civitatis \emph{transmigratio} testatur; quia
dum nunc in eo loco constructa est, ubi \emph{extra portam} fuerat
Dominus crucifixus, prior illa Jerusalem, ut dicitur, funditus est
eversa.'' Hom. in Evang. 39. init.}. If the present pool of Hezekiah was then acknowledged to be
such, close upon the present Sepulchre, is it credible that, with that
intimate knowledge of the letter of the inspired writings which no one
denies to their times, the clergy of Jerusalem should have fixed on a
site for the Sepulchre which, according to Dr. Robinson, they would be
confessing to be, not only within the lower city, but even within the
city of David? Did the pool escape their eyes, or its title their
ears, or the sacred text their memory, or the conclusion from these data their reason? Could it be that a pool, which Scripture says
was within the walls, should be situated upon a place of execution
which Scripture as surely places without them? And in like manner we
might ask, were it worth while, if the stones near the Damascus gate
wear an antique look now, were they not likely to tell their own story
better, if they were on the spot then? and have traces of the old wall
become fainter or stronger in the course of years? and had the
disposition of the ground undergone more alteration then than now, or
less, considering Hadrian rebuilt the city on the site on which he
found it \footnote{Those who deny that the Pagan Temple was built on the site of the
Sepulchre, have to account for the utter oblivion to which, on their
hypothesis, the place of our Lord's crucifixion was consigned; whereas
the circumstances attendant on that profanation which the Temple
occasioned will explain such partial ignorance concerning it as seems
to have obtained among the Christians of Jerusalem.}? But it is
needless to dwell on the improbability of an hypothesis which has been
shown to be altogether gratuitous.

188.
On the whole, then, I cannot doubt that the Holy Sepulchre was really
discovered as Eusebius declares it to have been; and I am as little
disposed to deny that the Cross was discovered also, as that the
relics of St. Cuthbert or the coffin of Bishop Coverdale have been
found here in England, in our own day.

\xsection{Section 6. The Death of Arius}

189.
CONSTANTINE, being gained over
by the Arian party, called Arius to Constantinople, with the intention
of obliging Alexander, the Bishop of that see, to restore him to the
communion of the Church. The old man, who was at that time
ninety-seven years of age, betook himself with his people to prayer
and fasting. He shut himself up in the church, and continued in
supplication for several days and nights. The coming Sunday was
appointed for the reception of Arius, and on the preceding day
Alexander was summoned before Constantine, and commanded to comply
with his wish. On his refusal the Emperor grew angry, and Alexander
withdrew in silence to urge the cause of Catholic truth with greater
earnestness in a more suitable Presence. He fell on his face before
the altar, and he conjured Christ, the Lord of all and King of kings,
to deliver the Church from the danger and disgrace which threatened
it. One of the persons attendant on him was Macarius, from whom St.
Athanasius relates it. Macarius followed his prayer as he spoke
it, and it ran thus: ``If Arius communicates tomorrow, then let
Thy servant depart, and destroy not the righteous with the wicked. But
if Thou sparest Thy Church, and I know Thou sparest it, have respect
unto the words of the Eusebians, and give not Thine heritage unto ruin
and reproach; and take Arius away, lest if he enter into the Church
his heresy seem to enter with him, and henceforth religion be counted
as irreligion.'' \footnote{De Mort. Ar.} This
prayer is said to have been offered about three p.m. on the Saturday;
that same evening Arius was in the great square of Constantine when he
was suddenly seized with indisposition. On retiring, he was overtaken
by what is commonly considered to be the fate of Judas, as described
in the Book of Acts. The building where this event took place became a
record of it to future times, and, as Socrates tells us,
``rendered the manner of his death ever memorable, all passers-by
pointing the finger at it.'' \footnote{Hist. i. 38.}

190.
Now of this occurrence it is obvious to remark, first of all, that it
is strictly of an historical character. It enters into the public
transactions of the times, and is one of a chain of events which are
linked together, and form a whole. It has a meaning, and gives a
meaning to the course of action in which it is found. It is in no
sense what Paley calls ``naked history,'' and in this respect differs
from certain other extraordinary occurrences, such, for
instance, as are recorded in the lives of the Monks; nay, from certain
miracles of Scripture, such as St. Paul's preservation from the viper,
of which nothing comes, and still more the resurrection wrought by
Elisha's bones. ``It has been said,'' says Paley, ``that if the prodigies
of the Jewish history had been found only in fragments of Manetho or
Berosus, we should have paid no regard to them; and I am willing to
admit this. If we knew nothing of the facts but from the fragment; if
we possessed no proof that these accounts have been credited and acted
upon from times, probably, as ancient as the accounts themselves; if
we had no visible effects connected with the history, no subsequent or
collateral testimony to confirm it; under these circumstances I think
that it would be undeserving of credit.'' He goes on to say that this
is not the case as regards the introduction of Christianity; nor, as
we may add, as regards the history of Arius.

191.
Again, it must be observed that this is more strictly a miracle of the
Church than many which occur within her pale and among her members;
that is, it is done by the Church as the Church. Though it bears a
tentative character, it is the result of a solemn intercession, a
solemn anathema, of the Church. Miracles happened in the kingdom of
Israel where there was no Church; but here is a contest between an
Emperor and a heresy on the one side, and the Church on the
other; the Church speaks through her constituted authorities, and the
judgment which is inflicted on her enemy is an attestation to her
divinity.

192.
Further, it was done in the presence of hostile power, which was awed
by it, and altered its line of action in consequence. Paley observes,
when arguing for the miracles of the Gospel, ``We lay out of the case
those which come merely in affirmance of opinions already formed. It
has long been observed that Popish miracles happen in Popish
countries, that they make no converts. In the moral as in the natural
world, it is change which requires a cause. Men are easily fortified
in their old opinions, driven from them with great difficulty.'' \footnote{Evidences, Part ii. Ch. i.} Now the event in question was a Catholic miracle in an Arian
city, before an Arian court, amid a prevalent Arianism extending
itself all through the East.

193.
``But after all, was it a miracle? for if not, we are labouring at a
proof of which nothing comes.'' The more immediate answer to this
question has already been suggested several times. When a Bishop with
his flock prays night and day against a heretic, and at length begs of
God to take him away, and when he is suddenly taken away almost at the
moment of his triumph, and that by a death awfully significant, from
its likeness to one recorded in Scripture, is it not trifling to
ask whether such an occurrence comes up to the definition of a
miracle? The question is not whether it is formally a miracle, but
whether it is an event, the like of which persons who deny that
miracles continue will consent that the Church should be considered
still able to perform. If they are willing to allow to the Church such
extraordinary protection, it is for them to draw the line, to the
satisfaction of people in general, between these and strictly
miraculous events; if, on the other hand, they deny their occurrence
in the times of the Church, then there is sufficient reason for our
appealing here to the history of Arius in proof of the affirmative.
This is what suggests itself at first sight; however, that it was
really miraculous, Gibbon surely is a sufficient voucher. ``Those,'' he
says, ``who press the literal narrative of the death of Arius, must
make their option between poison and miracle.'' Now, considering that
this awful occurrence took place in an Arian city and court, and in
the face of powerful and quick-sighted adversaries, who had every
means and every interest to detect an act of such dreadful wickedness
as Gibbon insinuates, surely, putting aside all higher considerations,
there are insuperable difficulties in the theory of poison; while
those who do not deny the moral governance of God, and the heretical
and ungodly character of Arianism, will have no difficulty in
referring the catastrophe to miracle.

194.
One other question may be asked, though it is of a doctrinal nature,
and therefore hardly needs to be considered here; whether so solemn a
denunciation as that adopted by Alexander, and so positive a reference
of the event which followed to that denunciation as a cause, are not
modes of acting and judging uncongenial to the Christian religion. One
passage there certainly is in the New Testament which at first sight
seems in opposition to it. When James and John wished to be allowed to
call down fire from heaven upon the Samaritans, as Elijah had done
upon Ahaziah's messengers, Christ answered, ``Ye know not what manner
of spirit ye are of: for the Son of man is not come to destroy men's
lives, but to save them.'' However, it is obvious to reply, first, that
Elijah, in the passage in question, called down a miraculous
punishment on the soldiers of Ahaziah mainly in his own defence; and
it is observable that the Apostles asked leave to do the same, when
the Samaritans had refused to receive their Lord and them; whereas the
great rule of the Gospel is to ``avenge not ourselves, but rather give
place unto wrath,'' as our Lord exemplified when ``they went to another
village.'' But whether there be any force in this distinction or not,
certain it is that in the Acts, in which we surely have the principles
of the Gospel drawn out into action, two precedents occur in
justification of the conduct of St. Alexander, one given us by St.
Peter and the other by St. Paul. St. Peter's denunciation of
Ananias and Sapphira was followed by their instantaneous deaths; St.
Paul's denunciation of Elymas, by his immediate blindness. These
instances, moreover, suggest that our Lord's earthly ministry might
probably be conducted on different laws from those which belonged to
His risen power, when the Spirit had descended, and light was spread
abroad; according to the text in which blasphemy against the Son of
man and blasphemy against the Spirit are contrasted. Hence St. Paul
calls Elymas, who was ``seeking to \emph{turn away} the deputy
from the faith,'' an ``\emph{enemy} of all righteousness,'' and a ``\emph{perverter}
of the right ways of the Lord;'' and St. Peter still more expressly
accuses Ananias and his wife of ``lying against the Holy Ghost,'' and ``tempting
the Spirit of the Lord.'' It is obvious also to refer to St. Paul's
imprecation on Alexander the copper-smith, that the Lord would reward
him according to his works. Here St. Paul, who had the gift of
inspiration, speaks of Alexander personally; but the Bishops of the
Church did not venture so much as this; they did but contemplate her
enemies \emph{in} their opposition, \emph{as} heretics or rebels, and dealt with
them accordingly, without any direct reference to their real and
absolute state in the sight of God.

\xsection{Section 7. The Fiery Eruption on Julian's attempt to rebuild the Jewish Temple}

195.
BISHOP WARBURTON,
as is\emph{}well known, has
written in defence of the miraculous character of the earthquake and
fiery eruption which defeated the attempt of the Emperor Julian to
rebuild the Jewish Temple. Though in many most important respects he
shows his dissent from the view of the Ecclesiastical Miracles taken
in these pages, yet the propositions which he lays down in the
commencement of his work are precisely those which it has been here
attempted to maintain; first, ``that not all the miracles recorded
in Church history are forgeries or delusions;'' next, ``that their
evidence doth not stand on the same foot of credit with the miracles
recorded in Gospel history.'' In drawing out the facts and the
evidence of the miracle in question, I shall avail myself of the work
of this learned and able writer, with which I agree in the main,
though of course there is room for difference of opinion, both as
regards the details of the one and the other, and as regards the view
to be taken of them.

196.
In the year 363, Julian, in the course of his systematic hostilities
against Christianity, determined to rebuild the Temple at Jerusalem.
The undertaking was conducted on a magnificent scale, large sums being
assigned out of the public revenue for its execution. Alypius, an
intimate friend of Julian, was set over the work; the Jews aided him
with a vast collection of materials and of workmen. Both sexes, all
ranks, took part in the labour, entering upon the ruins, clearing away
the rubbish, and laying bare the foundations \footnote{It was quite an enthusiastic movement. We are
told that the spades and pickaxes were of silver, and the rubbish was
removed in mantles of silk and purple. Vid. Gibbon, Ch. xxiii.}. What followed is attested by a number of authorities, who
agree with each other in all substantial respects, though, as was to
be expected, no single writer relates every one of the particulars.
First, we have the contemporary testimony of the Pagan historian
Ammianus Marcellinus, and we may add of Julian himself; then of St.
Gregory Nazianzen \footnote{Orat. v. 4-7. The Oration was composed the very year of the miracle. [Vid.
Fabric. Salutaris Lux. p. 124, Gothofr. in Philostorg. p. 296.]}, St.
Ambrose, and St. Chrysostom, who were more or less contemporaries; and
of Rufinus, Socrates, Sozomen, and Theodoret, of the century
following. They declare as follows. The work was interrupted by a
violent whirlwind, says Theodoret, which scattered about vast
quantities of lime, sand, and other loose materials collected for the
building. A storm of thunder and lightning followed; fire fell,
says Socrates; and the workmen's tools, the spades, the axes, and the
saws, were melted down. Then came an earthquake, which threw up the
stones of the old foundations of the Temple, says Socrates; filled up
the excavation, says Theodoret, which had been made for the new
foundations; and, as Rufinus adds, threw down the buildings in the
neighbourhood, and especially the public porticoes, in which were
numbers of the Jews who had been aiding the undertaking, and who were
buried in the ruins. The workmen returned to their work; but from the
recesses laid open by the earthquake, balls of fire burst out, says
Ammianus; and that again and again, so often as they renewed the
attempt. The fiery mass, says Rufinus, ranged up and down the street
for hours; and St. Gregory, that when some fled to a neighbouring
Church for safety, the fire met them at the door, and forced them back
with the loss either of life or of their extremities. At length the
commotion ceased; a calm succeeded; and, as St. Gregory adds, in the
sky appeared a luminous Cross surrounded by a circle. Nay, upon the
garments and upon the bodies of the persons present Crosses were
impressed, says St. Gregory; which were luminous by night, says
Rufinus; and at other times of a dark colour, says Theodoret; and
would not wash out, adds Socrates. In consequence, the attempt was
abandoned.

197.
There is no reason for doubting any part of this narrative; however,
enough will remain if we accept only the account given us by Ammianus,
who, to use the words of Warburton, was ``a contemporary writer, of
noble extraction, a friend and admirer of Julian, and his companion in
arms, a man of affairs, learned, candid, and impartial, a lover of
truth, and the best historian of his times,'' and ``a Pagan professed
and declared.'' ``Though Julian,'' says this writer, ``with anxious
anticipation of contingencies of every kind, was keenly engaged in the
prosecution of the numberless arrangements incident to his [Persian]
expedition, yet that no place might be without its share in his
energy, and that the memory of his reign might continue in the
greatness of his works, he thought of rebuilding at an extravagant
expense the proud Temple once at Jerusalem, which after many conflicts
and much bloodshed, in the siege under Vespasian first, and then
Titus, was with difficulty taken; and he committed the accomplishment
of this task to Alypius of Antioch, who had before that been
Lieutenant of Britain. Alypius therefore set himself vigorously to the
work, and was seconded by the governor of the province; when fearful
balls of fire, breaking out near the foundations, continued their
attacks, till the workmen, after repeated scorchings, could approach
no more; and thus, the fierce element obstinately repelling them, he
gave over his attempt.''

198.
Julian, too, seems awkwardly to allude to it in a fragment of a letter
or oration, which Warburton has pointed out, and which is so curious
an evidence of his defeat and its extraordinary circumstances that it
may be fitly introduced in this place. He is encouraging the zeal of
the Pagans for the honour of their divinities, and he says: ``Let no one
disbelieve the gods, from seeing and hearing that their statues and
their temples have been insulted in some quarters. Let no one beguile
us by his speeches, or unsettle us on the score of providence; for
those who reproach us on this head, I mean the Prophets of the Jews,
what will they say about their own Temple, which has been thrice
overthrown, and is not even now rising \footnote{Fabricius and De la Bleterie consider the ``three times'' to include
Julian's own attempt to rebuild; yet it is harsh, as Warburton
observes, to call a hindrance in rebuilding an actual destruction of
the building, though the hindrance was a destruction as far as it
went. But Lardner and Warburton seem to mistake when they argue
against Fabricius that [\emph{egeiromenou de oude nun}] means ``not \emph{raised
again} to this day,'' whereas it must rather be construed ``not \emph{rising}''
or ``in \emph{course of building}.'' Warburton reckons the alterations
and additions under Herod as by implication a destruction of the
second Temple; and as another hypothesis he suggests the profanation
under Antiochus. Lardner thinks Julian spoke vaguely or rhetorically,
or that he referred to the calamities which came upon Jerusalem in the
time of Adrian. ``Julien loin de conclure de ce qui étoit arrivé à
Jerusalem la vérité de la religion Chrétienne, en inferoit que la
revelation judaïque étoit fausse.'' De la Bleterie. Julian, v.
p. 399.}? This I have said with no wish to reproach them, inasmuch as I myself, at so late a day, had in purpose to rebuild it for the
honour of Him who was worshipped there. Here I have alluded to it,
with the purpose of shewing that of human things nothing is
imperishable, and that the Prophets who wrote as I have mentioned,
raved, and were but the gossips of canting old women. Nothing, indeed,
contradicts the notion of that God being great, but He is unfortunate
in His Prophets and interpreters; I say that they did not take care to
purify their souls by a course of education, nor to open their
fast-closed eyes, nor to dissipate the darkness which lay on them.
And, like men who see a great light through a mist, not clearly nor
distinctly, and take it not for pure light, but for \emph{fire}, and
are blind to all things around about it, they cry out loudly, 'Shudder
and fear; fire, flame, death, sword, lance,' expressing by many words
that one destructive property of \emph{fire}.'' \footnote{Page 295. Ed. Spanh. Lardner contends that this letter from its tone
must have been written \emph{before} any attempt to rebuild the
Temple; which indeed he considers Julian never to have put into
execution. This is a paradox more in the style of Warburton, whom he
is opposing, than of so sensible and sober a writer.} When it is considered that Julian was, as it were, defeated by
the prophets of that very people he was aiding; that he desired to
rebuild the Jewish Temple, and the Christians declared that he could
not, for the Jewish Prophets themselves had made it impossible; we
surely may believe, that in the foregoing passage this was the
thought which was passing in his mind, while the prophetic emblem of
fire haunted him, which had been so recently exhibited in the
catastrophe by which he had been baffled.

199.
The fact then cannot be doubted \footnote{It is objected by Lardner that St. Jerome, Prudentius, and Orosius are
silent about the miracle. Others have alleged the silence of St. Cyril
of Jerusalem. But if, as a matter of course, good testimony is to be
overborne because other good testimony is wanting, there will be few
facts of history certain. Why should Ammianus be untrue because Jerome
is silent? Sometimes the notoriety of a fact leads to its being passed
over. Moyle is ``unwilling to reject all [miracles since the days of
the Apostles] without reserve, for the sake of a very remarkable one
which happened at the rebuilding of the Temple,'' etc. Posth. Works,
Vol. i. p. 101. He professes to be influenced by the testimony and the
antecedent probability. Douglas speaks of Warburton's defence of it as
``a work written with a solidity of argument which might always have
been expected from the author, and with a spirit of candour which his
enemies thought him incapable of.'' These admissions are very strong,
considering the authors. Mosheim takes the same side. J. Basnage,
Lardner, Hey, etc., take the contrary.}; it may be asked, however, whether the perpetual ruin of the
Temple was actually predicted in the Prophets; and if not, what was
the drift of this miracle, and how it was connected with the Church.
It is connected with the Church and the Prophets by one circumstance,
if by no other, and that a remarkable one; that before the actual
attempt to rebuild, a Bishop of the Church had denounced it,
prophesied its failure, and that from the light thrown upon the
subject by the Prophets of the Old Testament. ``Cyril, Bishop of
Jerusalem,'' says Socrates, ``bearing in mind the words of the Prophet
Daniel, which Christ had confirmed in the Holy Gospels, declared to
many beforehand, that now the time was come, when stone should not
remain upon stone for that Temple, but the Saviour's prophecy should
be fulfilled.'' \footnote{Hist. iii. 20. Lardner (Testimonies, Ch. 46. 3) says, that ``it is very
absurd for any Christians to talk in this manner. Christ's words had
been fulfilled almost 300 years before;'' and refers to Rufinus as
giving the true account of St. Cyril's words, viz., that ``it could not
be that the Jews should be \emph{able} to lay them stone upon stone;''
but St. Cyril himself expressly says what Socrates reports of him,
Catech. xv. 15: ``Antichrist shall come \emph{at a time when} there
shall not be left one stone upon another in the Temple.'' This was
written before Julian's attempt; and St. Chrysostom, after it,
pronounces the prophecy of ``not one stone upon another'' not fulfilled
even then. Hom. 75. in Matth.} St. Cyril
seems to have argued that since our Lord prophesied the utter
destruction of the Temple, and since that destruction was not yet
fully accomplished, but only in course of accomplishment, for the old
foundations at that time still remained, therefore Julian was
reversing the Divine order of things, and building up when God was
engaged in casting down, and in consequence was sure to fail. And as
Julian probably understood Daniel's and our Lord's words in the same
way, and did set himself deliberately and professedly to contravene
them, viewed as fulfilled in the fortunes of the Temple, he was
evidently placing himself in open hostility to Christ and His Prophet,
and challenging Him to the encounter. No circumstances then could be
more fitting for the interposition of a miracle in frustration of his
undertaking.

200.
The same conclusion may be argued from our Lord's words to the
Samaritan woman. He does not indeed mention the Temple by name, but he
must be considered to allude to it, when He says that men should not ``\emph{worship}
at Jerusalem.'' They were indeed to worship there, as everywhere, but
to worship without the Temple; and that because they were to worship ``in
spirit and in truth.'' A spiritual worship was incompatible with the
Judaic services; so that when Christianity appeared the Temple was
destroyed. Julian then, in building again the Temple, was doing what
he could to falsify Christianity.

201.
But, again, the Jewish Temple was confessedly the \emph{centre} of the
Jewish worship and polity; to rebuild the Temple, then, was to
establish the Jews, \emph{as Jews}, in their own land, an event which,
if prophecy is sure, never is to be. ``The building of [the Temple],''
says Mr. Davison, ``was directed for this reason, that God had given '\emph{rest} to His people,' and henceforth would not suffer them to \emph{wander}
or be disturbed; so long as they enjoyed the privilege of being His
people at all. 'Moreover, I will appoint a \emph{place} for my people
Israel, and will plant them, that they may \emph{dwell} in a place
of their own, and \emph{move no more}.' This promise of rest was
connected with the Temple, for it was spoken when God confirmed and
commanded the design of building it.'' He continues presently, ``Their
national estate was henceforth attached to this Temple. It fell with
them; when they returned and became a people again, it rose also ...
Excepting around this Temple, they have never been able to settle
themselves, as a people, nor find a public home for their nation or
their religion ... So that the long desolation of their Temple, and
their lasting removal from the seat of it, are no inconsiderable
proofs that their \emph{polity and peculiar law} are come to an end in
the purposes of Providence, and according to the intention of the
Temple-appointment, as well as in the fact.'' \footnote{Discourses on Prophecy, v. 2. § 2.} Julian then, in proposing to rebuild the Jewish Temple, aimed
at the re-establishment of Judaism,—of that ceremonial religion
which in its day indeed had been the instrument of Divine Providence
towards higher blessings in store, and those for all men, but which,
when those blessings were come, forthwith was disannulled in the
Divine counsels ``for the weakness and unprofitableness thereof.''

202.
And next the question may be asked whether there was after all any
miracle in the case, as in the instance of most of the other
extraordinary occurrences which have passed under review. The luminous
Crosses upon the garments and bodies of the persons present were
apparently of a phosphoric nature; the Cross in the air resembled
meteoric phenomena; the earthquake and balls of fire had a volcanic
origin; and other marvellous circumstances are referable to
electricity. This all may be very true, and yet it may be true also
that the immediate cause, which set all these various agents in
motion, and combined them for one work, was supernatural; just as the
agency of mind on matter, in speaking, walking, writing, eating, and
the like, is not subject to physical laws, though manifesting itself
through them \footnote{``The mineral and metallic substances which, by their accidental
fermentation, are wont to take fire and burst out in flame, were the
native contents of the place from which they issued; but in all
likelihood \emph{they would have there slept}, and still continued in
the quiet innoxious state in which they had so long remained, \emph{had
not} the breath of the Lord \emph{awoke} and \emph{kindled} them.
But when the Divine Power had thus miraculously interposed to \emph{stir
up} the rage of these fiery elements, and yet to restrain their
fury to the objects of His vengeance, He then again suffered them to
do \emph{their ordinary office}; because nature, thus directed, would,
by the exertion of its own laws, answer all the ends of the moral
designation.'' Warburt. Julian. p. 246. Again, ``We see why \emph{fire}
was the scourge employed; as we may be sure \emph{water} would have
been, were the region of Judea naturally subject to inundations. For
miracles, not being an ostentatious, but a necessary instrument of God's
moral government, we cannot conceive it probable that He would \emph{create}
the elements for this purpose, but \emph{use} those which already lay
stored up against the day of visitation.'' Ibid. p. 250.}. Again, even
supposing that these phenomena were not in themselves miraculous, yet
surely their concurrence with the moral system of things, their
happening at that time and place and in that subserviency to the
declamations of ancient prophecy, is in itself of the nature of a
miracle. It is observable too, that though the Cross in the air be
attributable to meteoric causes, yet such an occurrence is after all
very unusual; now we read of three such occurrences in the course of
the fifty years between Constantine's accession to power and Julian,
during which period Christianity was effecting its visible triumph and
establishment in the world;—viz., the Cross at the conversion of
Constantine, that which hung over Jerusalem in the reign of
Constantius, and the Cross which forms part of the awful events now in
question; and while any accumulation of extraordinary phenomena
creates a difficulty in finding a cause in nature adequate to their
production, the recurrence of the same phenomena argues design, or the
interference of agency beyond nature. It must be added, too, that the
occurrence of a whirlwind, an earthquake, and a fire, especially
reminds us of Elijah's vision in Horeb, and again of the manifestation
of the Divine Presence in the first and fourth of the Acts, yet it
does not appear as if the writers to whom we have referred had these
events in their mind; rather it is only by the union of their separate
testimonies, each incomplete in itself, that the parallel is formed \footnote{It should be observed that the \emph{order} in which the miraculous
phenomena have been arranged above is not found in the original
authorities; Warburton has been followed except in one instance.}.

203.
Moreover the events in question did the work of a miracle; they
defeated powerful enemies, who would not have been unwilling to detect
imposture, and who would not have been deterred from their purpose by
interruptions which are extraordinary only in a relation. If the
purpose of the Scripture miracles be to enforce on the minds of men an
impression of the present agency and of the will of God, His approval
of one man or doctrine, and His disapproval of another, not even the
clearest of those recorded in the Gospel could have secured this
object more effectually than did the wonderful occurrence in question.
And did we see at this day a great attempt made to reinstate the Jews
as Jews in their own land, to build their Temple, and to recommence
their sacrifices, did the enemies of the Catholic Church forward it,
did heretical bodies and their officials on the spot take part in it,
and did some catastrophe, as sudden and unexpected as the fiery
eruption, befall the attempt, I conceive, whatever became of abstract
definitions, we should feel it to be a Divine interference, bringing
with it its own evidence, and needing no interpretation. It must be
recollected, too, that certain of the miracles of Scripture, such as
the destruction of Sodom, may be plausibly attributed to physical
causes, yet without disparagement of their Divine character. And
lastly, as to the extravagance of some writers who have considered the
miracle an artifice of the Christian body, the same scepticism which
has wantonly ascribed it to combustibles of the nature of gunpowder,
has at other times suggested a like explanation of the thunders and
lightnings when the Law was given, and of the deaths of Korah, Dathan,
and Abiram.

\xsection{Section 8. Recovery of the Blind Man by the Relics of St. Gervasius and St. Protasius at Milan}

204.
THE broad facts connected with
this memorable interposition of Divine Power are these: St. Ambrose,
with a large portion of the population of Milan, was resisting the
Empress Justina in her attempt to seize on one of the churches of the
city for Arian worship. In the course of the contest he had occasion
to seek for the relics of Martyrs, to be used in the dedication of a
new church, and he found two skeletons, with a quantity of fresh
blood, the miraculous token of martyrdom. Miracles followed, both
cures and exorcisms; and at length, as he was moving the relics to a
neighbouring church, a blind man touched the cloth which covered them,
and regained his sight. The Empress in consequence relinquished the
contest; and the subject of the miracle dedicated himself to religious
service in the Church of the Martyrs, where he seems to have remained
till his death. These facts are attested by St. Ambrose himself,
several times by St. Augustine, and by Paulinus, secretary to
St. Ambrose, in his Life of the Saint addressed to St. Augustine.

205.
This miracle, it is to be presumed, will satisfy the tests which
Douglas provides for verifying events of that nature. That author lays
down, as we have already seen, that miracles are to be suspected, when
the accounts of them were first published long after the \emph{time}\emph{}or far from the \emph{place}
of their alleged occurrence; or, if not, yet at least were not
then and there subjected to examination. Now in the instance before us
we have the direct testimony of three contemporaries, St. Ambrose, St.
Augustine, and Paulinus; two of whom at the least were present at the
very time and place, while one of those two wrote his account
immediately upon or during the events, as they proceeded. These three
witnesses agree together in all substantial matters; and the third,
who writes twenty-six years after the miracle, when St. Ambrose was
dead, unlike many reporters of miracles, adds nothing to the
narrative, as St. Ambrose and St. Augustine left it. Douglas observes
in explanation of the third of his conditions, that we may suspect
miracles of having ``been admitted without examination, first, if they
coincided with the favourite opinions and superstitious prejudices of
those to whom they were reported, and who on that account might be
eager to receive them without evidence; secondly, if they were set on
foot, or at least were encouraged and supported, by those who
alone had the power of detecting the fraud, and who could prevent any
examination which might tend to undeceive the world.'' \footnote{Pages 28, 52.} Now here all the power was on the side of those against whom
the miracle was wrought; and, though the popular feeling was with St.
Ambrose, yet the whole city had had an Arian clergy for nearly twenty
years, and could not but be in a measure under Arian influence. But
however this might be, at least Ambrose had to cope with Arian princes
armed with despotic power, an Arian court, an Arian communion lately
dominant and still organised, with a bishop at its head. His enemies
had already made attempts to assassinate him; and again, to seize his
person, and to carry him off from the city. They had hitherto been the
assailants, and he had remained passive. Now, however, he had at last
ventured on what in its effects was an aggressive act. As I have said,
he has to dedicate a Church, and he searches for relics of Martyrs. He
is said to find them; miracles follow; the sick and possessed are
cured; at length in the public street, in broad day, while the relics
are passing, a blind man, well known in the place by name, by trade,
by person, and by his calamity, professes to recover his sight by
means of them.

206.
Here surely is a plain challenge made to the enemies of the Church,
almost as direct as Elijah's to the idolatrous court and false
prophets of Israel. St. Ambrose supplies them with materials, nor do
they want the good-will to detect a fraud, if fraud there be. Yet they
are utterly unable to cope with him. They denied the miracle indeed,
and they could not do otherwise, if they were to remain Arians; as
Protestant writers deny it now, that they may not be forced to be
Catholics. They denied the miracle, and St. Ambrose, in a sermon
preached at the time, plainly tells us that they did; but they did not
hazard any counter statement or distinct explanation of the facts of
the case. They did not so much as the Jews, who, on the Resurrection,
at least said that our Lord's Body was stolen away by night. They did
nothing but deny,—except indeed we let their actions speak for them.
One thing then they did; they gave over the contest. The Miracle was
successful.

207.
This miracle answers to Leslie's criteria also. It was sensible; it
was public; and the subject of it became a monument of it, and that
with a profession that he was so. He remained on the spot, and
dedicated himself to God's service in the Church of the Martyrs who
had been the means of his cure; thus by his mode of life proclaiming
the mercy which had been displayed in his behalf, and by his presence
challenging examination.

208.
An attempt has lately been made to resolve this miracle into a mere
trick of priestcraft; but doubtless the Arians would have been
beforehand with the present objector, could a case have been made out
with any plausibility. This anticipation is confirmed by an inspection
of the inferences or conjectures of which he makes the historical
facts the subject. The blood, he says, was furnished by the blind
Severus, who had been a butcher, and might still have relations in the
trade. And since St. Ambrose translated the relics at once, instead of
waiting for the next Sunday, this is supposed to argue that he was
afraid, had the ceremony been postponed, of the fraud being detected
by the natural consequence of the delay.

209.
But all facts admit of two interpretations; there is not the
transaction or occurrence, consisting of many parts, but some of them
may be fixed upon as means of forcing upon it a meaning contrary to
the true one, as is shewn by the ingenuity exercised in defence of
clients in the courts of law. What has been attempted by the writer to
whom I allude, as regards St. Ambrose, has been done better, though
more wickedly, by the infidel author of the New Trial of the Witnesses
as regards the History of the Resurrection. In such cases inquirers
will decide according to their prepossessions \footnote{This has been dwelt on at length, supr. n. 71 to n. 80. Gibbon gives
us a curious illustration of it in his remark on the miracle of the
Confessors, which is presently to be related. He says: ``This
supernatural gift of the African Confessors, who spoke without
tongues, will command the assent of those, and of those only, who \emph{already}
believe that their language was pure and orthodox. But the \emph{stubborn
mind} of the infidel is \emph{guarded by secret incurable suspicion};
and the Arian or Socinian, who has seriously rejected the doctrine of
the Trinity, \emph{will not be shaken by the most plausible evidence}
of an Athanasian miracle.'' Ch. xxxvii.}; if they are prepared to believe that the Fathers and
Doctors of the Church would introduce the blood of the shambles into a
grave, and pretend that it was the blood of God's saints, and hire men
first to feign themselves demoniacs and then to profess themselves
dispossessed on approaching the counterfeit relics, they will be
convinced in the particular case by very slight evidence, and will
catch at any circumstances which may be taken as indications of what
they think antecedently probable; but if they think such proceedings
to be too blasphemous, too frightful, too provocative even of an
immediate judgment, for any but the most callous hearts and the most
reckless consciences to conceive and carry out, they would not believe
even plausible evidence in their behalf. If it appears to them not
unlikely that miracles continue in the Church, they will find that it
is easier to admit than to reject what comes to them on such weighty
testimony; but if they think miracles as improbable after a revelation
is given, as they appeared to Hume before it, then they will judge
with him that ``a religionist may know his narration to be false, and
yet persevere in it, with the best intentions in the world, for
the sake of promoting so holy a cause; or even where this delusion has
no place, vanity, excited by so strong a temptation, operates on him
more powerfully than on the rest of mankind in any other
circumstances, and self-interest with equal force.'' \footnote{Essay on Miracles.}

210.
There are circumstances, however, in this miracle, which may be felt
as difficulties by those who neither deny the continuance of a Divine
Presence in the Church, nor accuse her Pastors and Teachers of impious
imposture. Yet it is difficult to treat of them, without entering upon
doctrinal questions which are not in place in the present Essay. One
or two of them, which extend to the case of other alleged miracles of
the early Church, besides the one immediately before us, shall here be
briefly considered, and that in the light which the analogy or the
pattern of Scripture throws upon them, which is the main view I have
taken of objections all along.

211.
Now, first it may be urged that the discovery of the blood of the
Martyrs is not after the precedent of anything we meet with in
Scripture, which says very little of relics, and nothing of relics of
such a character as this, involving as it does a miracle. What is the
true doctrine about relics, how they are to be regarded, what is their
use and their abuse, is no question before us. If it could be shown
that the doctrine involved in the discovery of the Martyrs, is
on Scriptural grounds such as plainly to prove either that it did not
take place, or that it cannot be referred to Divine Agency, this of
course would supersede all other considerations. Meanwhile I will but
observe, as far as the \emph{silence} of Scripture is concerned, that
Scripture could not afford a pattern of the alleged miracle, from the
nature of the case. The resurrection of the body is only a Christian
dogma; and martyrdom, that is, dying for a creed, is a peculiarity of
the Gospel, and was instanced among the Jews, only in proportion as
the Gospel was anticipated. The blood was the relic of those whose
bodies had been the temple of the Spirit, and who were believed to be
in the presence of Christ. Miracles were not to be expected by such
instruments, till Christ came; nor afterwards, till a sufficient time
had elapsed for Saints to be matured and offered up, and for pious
offices and assiduous attentions to be paid by others towards the
tabernacles which they left behind them. Precedents then to our
purpose, whether in Old or in New Testament, are as little to be
expected, as precedents to guide us in determining the relations of
the Church to the State, or the question of infant baptism, or the
duty of having buildings for worship. Time alone could determine what
the Divine purpose was concerning the earthly shrines in which a
Divine Presence had dwelt: whether, as in the case of Moses and
Elijah, they were to be withdrawn from the Church, or, as in the
case of Elisha, to fulfil some purpose, even though the soul had
departed; and if the latter, whether their bones were to be
employed,—or whether their bodies would be preserved incorrupt, as
St. Jerome reports of Hilarion,—or whether the Levitical sacrifices,
which as types were once for all fulfilled when our Lord's blood was
shed, were nevertheless to furnish part of the analogy existing
between the Christian and the Mosaic Dispensations. Nor is there
anything that ought to shock us in the idea that blood, which had
become coagulate, should miraculously be made to flow. A very
remarkable prototype of such an event seems to be granted to us in
Scripture, in our Lord's own history. The last act of His humiliation
was, after His death, to be pierced in His side, when blood and water
issued from it. A stream of blood from a corpse can hardly be
considered to be other than supernatural. And it so happens that St.
Ambrose is the writer to remark upon this solemn occurrence in his
comment on St. Luke, assigning at the same time its typical meaning.
``Blood,'' he says, ``undoubtedly congeals after death in our
bodies; but in that Body, though incorrupt, yet dead, the life of all
welled forth. There issued water and blood; water to wash, blood to
redeem. Let us drink then what is our price, that by drinking we may
be redeemed.'' \footnote{In Luc. lib. x. § 135. Eutheymius Zigab. says of the same in loc.
Joan. Theophylact. in loc. says that in order to place the miracle
beyond doubt the water issued also. That the flowing of the blood was
miraculous would appear from the description St. John gives of it, ``\emph{forthwith
came there out};'' which implies a stream, and not a few drops.
Calov. in Joan. xix. 35. Nor is there any reason to suppose that the
water came forth by drops; yet the words just quoted are common to the
blood and to the water. Again, the water was miraculous (for ``medical
men tell us that the fluid of the pericardium is yellow in colour,
bitter in taste, and therefore different from what we mean by water,''
S. Basnag. Ann. 33. § 126; and the wound was most probably on the
right side, as St. Augustine and the most ancient pictures and coins
represent it, and the Arabic or Ethiopic version, vid. Greta. de Cruc.
t. 1. i. 35. Lamp. in loc. Joan.), and therefore there is no reason
for a strained interpretation only to escape believing that the blood
was miraculous. Further, St. John's solemn asseverations, ``He who saw
it bare record,'' etc., which seems to intimate something miraculous,
applies to the blood as well as the water. And moreover, in 1 John v.
6, the blood is insisted on even more than the water; ``not by water
only,'' etc. Another parallel to this miracle is to be found in the
reported instances of blood flowing from a corpse at the approach of
the murderer; vid. an instance introduced into a Scotch court as late
as 1688, in the notes to the Waverley Novels, vol. xliii. p. 127. It
is scarcely necessary to say that, whatever truth there may be in any
such stories, or in certain others of which the blood of Martyrs is
the subject, they are so encompassed by fictions and superstitions,
that it seems hopeless at this day to trace the Divine Agency, as and
when It really wrought, though we may believe in Its presence
generally.}

212.
Another objection which has been made to the miracles ascribed by St.
Ambrose to the relics which he discovered, is the encouragement which
they are supposed to give to a kind of creature-worship, unknown to
Scripture. This is strongly urged by the objector whom I just now had
occasion to notice. He observes that miracles can be of no avail
against the great principles of religious truth, such as the Being and
Attributes of Almighty God; that no miracles can sanction and justify
idolatry; if then the Nicene Miracles (so he calls them), ``when
regarded in the calmest and most comprehensive manner,'' ``have
constantly operated to debauch the religious sentiments of mankind, if
they have confirmed idolatrous practices, if they have enhanced that
infatuation which has hurried men into the degrading worship of
subordinate divinities, we then boldly say that, whether natural or
preternatural, such miracles are not from God, but from 'the
enemy.' \footnote{Anc. Christ. Part vii. p. 361.} ``Do you choose,''
he continues, ``to affirm the supernatural reality of the Nicene
Miracles? you then mark the Nicene Church as the slave and agent of
the Father of Lies;'' and then he proceeds to quote the charge of Moses
to his people: ``If there arise among you a prophet or a dreamer of
dreams, and giveth thee a sign or a wonder, and the sign or the wonder
come to pass, whereof he spoke unto thee, saying, 'Let us go after
other gods which thou hast not known, and let us serve them,' thou
shalt not hearken unto the words of that prophet or that dreamer of
dreams; for the Lord your God proveth you, to know whether ye love the
Lord your God with all your heart and with all your soul.''

213.
But the objection, which of course demands a careful consideration,
admits of being met, perhaps of being overcome, by reference to an
analogy contained in the Old Testament, to which the appeal is made.
It is well known that the Divine revelations concerning Angels
received a great development in the course of the Jewish Dispensation.
When the people had lately come out of Egypt, with all the forms of
idolatry familiar to their imaginations, and impressed upon their
hearts, it did not seem safe, if we may dare to trace the Divine
dealings in this matter, to do much more than to set before them the
great doctrine of the Unity and Sovereignty of God. To have disclosed
to them truths concerning angelic natures, except in the strictest
subserviency to this fundamental Verity, might have been the occasion
of their withdrawing their heart from Him who claimed it whole and
undivided \footnote{[Eusebius says this, contrasting Genesis with Daniel, Eccl. Theol. ii.
p. 20, and vid. Suicer de Symb. Nic. pp. 89-91.]}. Hence, though
St. Stephen tells us that they ``received the Law by the disposition of
Angels,'' and St. Paul that ``it was ordained by Angels,'' in the Old
Testament we do but read of ``the voice of the trumpet exceeding loud,''
and its ``waxing louder and louder,'' and ``Moses speaking, and God
answering him by a voice,'' and of ``the Lord talking with them face to
face in the mount.'' In like manner, when Angels appeared, it was for
the most part in the shape of men; or if their heavenly nature
was disclosed, still they are called ``wind'' or ``flame,'' or
represented as a glory of the Lord, and so intimately and mysteriously
connected with His Presence that it was impossible that God should be
forgotten, and a creature worshipped. Thus it is said of the Angel who
went before the Israelites, ``Obey his voice, \emph{for My Name is in him};''
and it was the belief of the early Church that the Second Person of
the Holy Trinity did really condescend to manifest Himself in such
angelic natures. Again, the title of ``the Lord of hosts'' does not
occur till the times of Samuel, who uses it when he sends Saul against
the Amalekites, whereas it is the ordinary designation of Almighty God
in the Prophets who lived after the captivity \footnote{Vid. e.g. Hagg. ii. 4-9; Zech. viii.}. And so again, in the Book of Daniel, Angels are made the
ordinary instruments of Divine illumination to the Prophet, and are
represented as the guardians of the kingdoms of the world, and that
without any mention of the Divine Presence at all, which, on the
contrary, had been awfully signified in the vision of Isaiah, when the
Seraph touched his lips.

214.
Still more striking is the difference of language in different parts
of the inspired volume as to the doctrine of an Evil Spirit, whom even
to name might have been to create a rival to the All-Holy Creator in
carnal minds which had just left the house of spiritual as well
as temporal bondage. The contrast between the earlier and later books
of the Old Testament in this point has often been observed. Satan is
described in the Book of Job and in Micaiah's vision as appearing
before God, and acting under his direction. Again, while in the Second
Book of Samuel we are merely told that ``the anger of the Lord was
kindled against Israel, and He moved David against them to say, Go
number Israel and Judah;'' in the First Book of Chronicles we read that
``Satan stood up against Israel, and provoked David to number Israel.''

215.
Yet, in spite of this merciful provision on the part of Almighty God,
it would appear that the revelation of Angels, when made, did lead
many of the Jews into an idolatrous dependence upon them. It is the
very remark of Theodoret upon St. Paul's mention of Angel-worship in
his Epistle to the Colossians, that ``the advocates of the Law induced
men to worship Angels, because the Law was given by them, and for
humility-sake, and because the God of all is invisible and
inaccessible and incomprehensible, so that it was fitting to procure
the Divine favour through the Angels.'' \footnote{In Col. ii. 18. Vid. also the passage from the Prædic. Petr. in Clem.
Strom. vi. 5. and Origen, in Joann. tom. xiii. 17; also contr. Cels.
v. 6, etc., Hieron. ad. Algas. Ep. 121. § 10: vid. references to
Rabbinical and other writings, Calmet. Dissert. 2. in Luc.} The Essenes, too, are said to have paid to the Angels an
excessive honour, and several of the early heresies, which did the
same, sprang from the Jews. What place afterwards the invocation of
Angels for magical purposes held in the practical Cabbala, as Brucker
calls it, is well known.

216.
Such is the history of the revelation of the doctrine of Angels among
the Jews; and it is scarcely necessary to draw out at length its
correspondence with the history of the introduction and abuse of
several of the tenets and usages which characterize the Christian
Church. In its origin, the Jewish as well as the Pagan institutions
with which the Apostles were surrounded, suggested to them a cautious
economy in the mode in which they set Divine truth before their
disciples, lest a resemblance of external rites and offices, or of
phraseology, between Christianity and the prevailing religions, should
be the means of introducing into their minds views less holy and
divine than those which they were inspired to reveal. It is on this
supposition that some English divines even account for the omission in
the New Testament of the words ``priest,'' ``sacrifice,'' and the like, in
their plain Christian sense; as if the Jewish associations which
attached to them would not cease till the Jewish worship had come to
an end. The remark may obviously be extended to the miracle under
review, so far as no parallel is found for it in the New
Testament. As the doctrine of priesthood might be almost necessarily
Judaic in the minds of the Jewish converts, so that of piety towards
Saints and Martyrs was in the minds of Pagans necessarily idolatrous;
and it may be for this, as well as other reasons, that so little
explicit mention is made in the New Testament of the honours due to
Saints, as also of the Christian Priesthood, after the pattern of that
silence, which has been above noticed, about the offices of good
Angels and about the Author of Evil in the earlier books of the Old;
and it may be as rash to say that a miracle was not from God because
it was wrought by a Martyr's relics, or because such relics have, in
other instances been idolatrously regarded, as to say that the Prophet
Daniel was not divinely inspired, because we hear nothing of Michael
or Gabriel in the Books of Moses, or because the names of those Angels
were afterwards superstitiously used in the charms of the Cabbalists.
The holy Daniel's profound obeisance and prostrations before the Angel
are a greater innovation, if it must so be called, on the simplicity
of the Mosaic ritual, than the treasuring the blood of the Martyrs
upon the ecclesiastical observances of the Apostles; and as no one
would say that Daniel's conduct incurred the condemnation pronounced
by Moses on those who introduced the worship of other gods, so much
less was the reverence paid by St. Ambrose and other Saints to the
relics of the Martyrs inconsistent with precepts which in their
direct force belong to an earlier Dispensation.

217.
There is a third difficulty, which may be raised upon the passage of
history before us, not arising, however, out of the miracle, but out
of the circumstances under which it took place. It may be represented
as giving a sanction to a subject's playing the part of a demagogue,
and heading a mob (as we may speak) against his lawful sovereign. The
crowds which attended Ambrose, whether in the church which the Empress
had seized, or on occasion of the translation of the relics, would
have been dispersed at this day among ourselves by the officers of the
peace; and with our present notions of law and of municipal and
national order, not to say of the subserviency of the Church to the
State, and our interpretations of the Scripture precepts concerning
civil obedience, there is something strange and painful to us in the
sight of a Christian Bishop placed in opposition to the powers that
be. But it must be recollected, according to a former remark, that
everything that happens has two aspects; and the outside or political
aspect is often the reverse of its inward or true meaning. We are used
to put together the particulars which meet our eye, to parallel them
with other transactions which bear a similar appearance, to suggest
for them such motives of action as our own principles or disposition
suggests, and thus to form what seems to us a philosophical view
of the whole case. And if our own habitual feelings and opinions, and
the parallels to which we betake ourselves, are not of a very exalted
nature, as may easily happen, while the subject contemplated, be it a
person, or an act, or a work, is of such a nature, then we produce a
theory as shallow, and as far from the truth, as a naturalist, who,
judging of men by their anatomical peculiarities, should rank them
among the brute creation. Every day brings evidence in great things or
little, how incapable the run of men are of doing justice to minds of
even ordinary refinement and sincerity, and how, rather than ascribe
to them the honesty and purity of purpose which is the most natural
and straightforward account of their actions, they will even go out of
their way, and distort facts, thereby to be at liberty to impute petty
motives; and much more will they catch at any circumstances which
admit of being plausibly perverted into an evidence of such motives.
Indeed, of such continual occurrence are instances of this sort, that
in tales of fiction nothing is more commonly taken as a plot of the
story than the troubles in which an innocent person is involved by an
ingenious but perverse selection and collocation of his actions or of
circumstances connected with him, to the detriment of his character.

218.
As to the case immediately before us, it is enough to observe that an
imputation of disloyalty, if preferred against St. Ambrose, is
only what the notorious Paine, I believe, throws out against the
Jewish Prophets; and it is obvious what plausible materials are
afforded by the history of Elijah and Elisha, in the hands of
irreligious persons, for such a charge. Nor is it to be doubted that a
secular historian, who heard the Prophet Jeremiah's public declaration
on Nebuchadnezzar's invasion, ``He that abideth in this city shall die,
but he that goeth out, and falleth to the Chaldeans, shall live,''
would have decided that he was in the pay of the King of Babylon, and
justified the Jews in their treatment of him. It must be recollected,
too, that one charge against our Lord was that He ``stirred up the
people.'' We indeed have learned from the Gospel that He withdrew
Himself from the multitude ``when He perceived that they would come and
take Him by force to make Him a king;'' but a secular historian either
would not know the fact, or might not believe the sincerity of His
withdrawal, if He did. A more exact instance in point is afforded us
in the history of St. John Baptist. No man surely has less of a
political character upon him than this holy ascetic, as described in
the Gospel; but it seems, according to Josephus, that Herod was of
another mind, and the view he took of him as a popular leader is so
curious that I will quote the words of a recent writer on the subject.
``Herod,'' says Mr. Milman, ``having formed an incestuous connection
with the wife of his brother Herod Philip, his Arabian queen
indignantly fled to her father, who took up arms to revenge her wrongs
against her guilty husband. How far Herod could depend in this contest
on the loyalty of his subjects was extremely doubtful. It is possible
he might entertain hopes that the repudiation of a foreign alliance,
ever hateful to the Jews, and the union with a branch of the Asmonean
line (for Herodias was the daughter of Herod the Great by Mariamne),
might counterbalance in the popular estimation the injustice and
criminality of his marriage with his brother's wife. The \emph{influence}
of John, according to Josephus, was almost \emph{unlimited}. The
subjects, and even the soldiery, of the tetrarch \emph{crowded with
devout submission} around the Prophet. \emph{On his decision} might
depend the \emph{wavering loyalty} of the whole province. But John
denounced with open indignation the royal incest, and declared the
marriage with a brother's wife to be a flagrant violation of the law.
Herod, before long, \emph{ordered him to be seized and imprisoned} in
the \emph{strong fortress} of Machærus, on the \emph{remote border} of his
trans-Jordanic territory.'' \footnote{Hist. Christ. Vol. i. p. 176.}

219.
Such was the light thrown upon the Holy Baptist by the secular events
in which he was encompassed, in the opinion of one who nevertheless,
as we know, ``feared him, knowing he was a just man.'' And as St.
John seemed to be a demagogue and a mere organ of the popular voice,
yet spoke from heaven, so in like manner it need not take from the
sanctity of St. Ambrose, or the truth of his cause, that the people
sided with him, even tumultuously, and the Imperial Court accused him
of insubordination.

\xsection{Section 9. The Power of Speech continued to the African Confessors deprived of their Tongues}

220.
ARIANISM, though speedily
exterminated from the Roman Empire, had taken refuge among the
Barbarians of the North, who were then hanging over it, and soon to
overwhelm it. Among these nations were the Vandals, who in the early
part of the fifth century took possession of the Roman provinces on
the African coast. Genseric forthwith commenced, and his successors
continued, a terrible persecution of the Catholic Church, which they
found there. Hunneric his son, to whose reign the miracle which is to
be related belongs, began his series of cruelties by stationing
officers violently to assault and drag off all Vandals whom they found
attending the Churches, and by sending off the dependents of his court
who were Catholics to work in the country as agricultural labourers.
Others he deprived of their civil functions, stripped of their
property, and banished to Sicily and Sardinia. Next he summoned
the nuns out of their convents, accused them of the vilest crimes, and
submitted them to the most miserable indignities. Further, he caused
them to be hung up without clothes, with weights to their feet, and to
be tortured with red-hot irons in various parts of the body, in order
to make them admit the charges he brought against them. His next
measure was the wholesale cruelty of banishing a number of bishops,
priests, deacons, and others, as many as four thousand nine hundred
and sixty-one \footnote{The number is given differently; Gibbon says
four thousand and ninety-six; Fleury four thousand nine hundred and
sixty-six; that in the text is as it stands in the Bibl. Patr. Par.
1624.}, to the
desert. He began by assembling them in the two towns of Sicca and
Laribus; and in one or other of these places Victor, who has preserved
the history of the transaction, saw them. His account is too horrible
to be translated. They had been shut up, how long does not appear, in
a small prison, and when Victor entered he sank up to his knees in the
filth of the place. At length they set forth for the desert, with
their faces and clothes in this defiled condition, chaunting the
words, ``Such glory have all His saints.'' They journeyed chiefly by
night, on account of the heat of the days; when they flagged, their
conductors goaded them and pelted them, or if this did not quicken
them, they tied them by the feet, and dragged them after them along
the rocky roads. Those who survived the journey found themselves
in places abounding in venomous reptiles, and the food given them was
the barley provided for the beasts of burden.

221.
In the beginning of 484 Hunneric convened four hundred and sixty-six
Catholic Bishops at Carthage, for the purpose of holding a disputation
on the faith of Nicæa; and, to intimidate them, he began by burning Lætus
alive, who was one of their most learned members. This not succeeding,
he dismissed them again to their homes, allowing them neither the
beasts of burden on which they had come, nor their servants, nor their
clothes, and forbidding all persons to lodge or feed them; when they
remonstrated, he set his cavalry to charge them. Jealous of their
orthodoxy as a bond of union with the Catholic world, he next proposed
to them to swear allegiance to his son and successor, and abstention
from all ecclesiastical correspondence beyond sea. Forty-six refused
it on the plea of our Lord's prohibition in the Sermon on the Mount;
three hundred and two, on the stipulation that their flocks and
themselves should be restored to their churches, took it. The latter
he distributed as serfs up and down the country, as having broken the
Gospel precept against swearing; the former he transported to Corsica
to cut timber for his navy. Of the rest, twenty-eight had succeeded in
escaping from Carthage, and eighty-eight conformed. A general persecution followed, in which neither sex nor age was pitied, nor
torture, mutilation, nor death was spared.

222.
These particulars, which form but a portion of the atrocities which
this savage was permitted to perpetrate, have here been mentioned,
because they form a suitable antecedent, and (if the word may be used)
a justification of the miracle which followed. It was no common
occasion that called forth what was no common manifestation of the
wonderful power of God. The facts, as stated by one who in such a case
cannot be called a too favourable witness, were as follows: ``Tipasa,''
says Gibbon, ``a maritime colony of Mauritania, sixteen miles to the
east of Cæsarea, had been distinguished in every age by the orthodox
zeal of its inhabitants. They had braved the fury of the Donatists;
they resisted, or eluded, the tyranny of the Arians. The town was
deserted on the approach of an heretical Bishop; most of the
inhabitants who could procure ships passed over to the coast of Spain;
and the unhappy remnant, refusing all communion with the usurper,
still presumed to hold their pious, but illegal, assemblies. Their
disobedience exasperated the cruelty of Hunneric. A military Count was
despatched from Carthage to Tipasa; he collected the Catholics in the
Forum, and, in the presence of the whole province, deprived the guilty
of their right hands and their tongues. But the holy Confessors
continued to speak without tongues.'' \footnote{Hist. Ch. xxxvii.} ``The gift continued through their lives. Their number is not
mentioned by any of the original witnesses; but is fixed by an old
Menology at sixty.'' \footnote{Ibid.} Such
was the miracle; the evidence on which it rests shall next be stated.

223.
Victor, Bishop of Vite, who has been already mentioned, published in
Africa his history of the persecution only two years after it took
place. He says, ``The King in wrath sent a certain Count with
directions to hold a meeting in the forum of the whole province, and
there to cut out their tongues by the roots, and their right hands.
When this was done, by the gift of the Holy Ghost, they so spoke and
speak, as they used to speak before. If however any one will be
incredulous, let him now go to Constantinople, and there he will find
one of them, a sub-deacon, by name Reparatus, speaking like an
educated man without any impediment. On which account he is regarded
with exceeding veneration in the court of the Emperor Zeno, and
especially by the Empress.'' \footnote{Hist. Pers. Vand. iii. p. 613.}
It has been asked why Victor refers his readers to Constantinople,
instead of pointing out instances of the miracle in the country in
which it is said to have taken place \footnote{This is suggested in the article on Miracles in the Encycl. Metrop. [\emph{i.e}.
Essay i. supr. p. 87, \textbf{[}Note
26\textbf{]}], in which I could wish
some correction of opinion, but more of tone, in my treatment of the
primitive miracles. The Essay aims, indeed, at bringing out the
characteristics of the evidence for the Scripture Miracles, in
contrast with all others so considered; but Middleton and Douglas are
unsafe guides, and it is no exaltation of Christ to lower His Saints.}. But persecution scattered the Catholics far and wide, as
St. Gregory observes in a passage which is to follow; many fled the
country; others concealed themselves. Under such circumstances, a
writer would not know even where his nearest friends were to be found;
and in this case Victor specified one of the Confessors who had been
welcomed by an orthodox capital and court, and had the opportunity of
exhibiting in security the miraculous gift wrought in him.

224.
Æneas of Gaza was the contemporary of Victor. When a Gentile, he had
been a philosopher and rhetorician, and did not altogether throw off
his profession of Platonism when he became a Christian. He wrote a
Dialogue on the Immortality of the Soul and the Resurrection of the
Body; and in it, after giving various instances of miracles, he
proceeds, in the character of Axitheus, to speak of the miracle of the
African Confessors: ``Other such things have been and will be; but what
took place the other day, I suppose you have seen yourself. A bitter
tyranny is oppressing the greater Africa; and humanity and orthodoxy
have no influence over tyranny. Accordingly this tyrant takes offence
at the piety of his subjects, and commands the priests to deny their
glorious dogma. When they refuse, O the impiety! he cuts out
that religious tongue, as Tereus in the fable. But the damsel wove the
deed upon the robe, and divulged it by her skill, when nature no
longer gave her power to speak; they, on the other hand, needing
neither robe nor skill, call upon nature's Maker, who vouchsafes to
them a new nature on the third day, not giving them another tongue,
but the faculty to discourse without a tongue more plainly than
before. I had thought it was impossible for a piper to show his skill
without his pipes, or harper to play his music without his harp; but
now this novel sight forces me to change my mind, and to account
nothing fixed that is seen, if it be God's will to alter it. I myself
saw the men, and heard them speak; and wondering at the articulateness
of the sound, I began to inquire what its organ was; and distrusting
my ears, I committed the decision to my eyes. And opening their mouth,
I perceived the tongue entirely gone from the roots. And astounded I
fell to wonder, not how they could talk, but how they had not died.''
He saw them at Constantinople.

225.
Procopius of Cæsarea was secretary to Belisarius, whom he accompanied
into Africa, Sicily, and Italy, and to Constantinople, in the years
between 527 and 542. By Belisarius he was employed in various
political matters of great moment, and was at one time at the head of
the commissariat and the fleet. He seems to have conformed to
Christianity, but Cave observes, from his tone of writing, that
he was no real believer in it, nay, preferred the old Paganism, though
he despised its rites and fables \footnote{Cave, Hist. Liter. Procop.}. He wrote the History of the Persian, Vandalic, and Gothic War,
of which Gibbon speaks in the following terms: ``His facts are
collected from the personal experience and free conversation of a
soldier, a statesman, and a traveller; his style continually aspires,
and often attains, to the merit of strength and elegance; his
reflections, more especially in the speeches which he too frequently
inserts, contain a rich fund of political knowledge; and the
historian, excited by the generous ambition of pleasing and
instructing posterity, appears to disdain the prejudices of the people
and the flattery of courts.'' \footnote{Hist. Ch. xl.}
Such is Procopius, and thus he speaks on the subject of this
stupendous miracle: ``Hunneric became the most savage and iniquitous of
men towards the African Christians. For, forcing them to Arianize,
whomever he found unwilling to comply, he burnt and otherwise put to
death. And of many he cut out the tongue as low down as the throat \footnote{[\emph{Ap' pharungos}].}, who, even as late as my time, were alive in Byzantium, and
talked without any impediment, feeling no effects whatever of the
punishment. But two of them, having allowed themselves to hold
converse with abandoned women, ceased to speak.'' \footnote{Bell. Vand. i. 10.}

226.
Our next witness, and of the same date, is the Emperor Justinian, who,
in an edict addressed to Archelaus, Prætorian Prefect of Africa, on
the subject of his office, after Belisarius had recovered the country
to the Roman Empire, writes as follows: ``The present mercy, which
Almighty God has deigned to manifest through us for His praise and His
Name's sake, exceeds all the wonderful works which have happened in
the world; viz., that Africa should through us recover in so short a
time its liberty, after being in captivity under the Vandals for
ninety-five years, those enemies alike of soul and body. For such
souls as could not sustain their various tortures and punishments, by
rebaptizing they translated into their own misbelief; and the bodies
of free men they subjected to the hardships of a barbaric yoke. Nay,
the very churches sacred to God did they defile with their deeds of
misbelief; some they turned into stables. We have seen the venerable
men, who, when their tongues had been cut off at the roots, yet
piteously recounted their pains. Others after diverse tortures were
dispersed through diverse provinces, and ended their days in exile.'' \footnote{Cod. Just. lib. i. tit. 30, ed. 1553.}

227.
Count Marcellinus, Chancellor to Justinian before he came to the
throne, is the fourth layman to whose testimony we are able to appeal.
He too, as two of the former, speaks as an eye-witness, and the additional circumstances, with which he commences seem to throw light
upon Æneas's singular account that the Confessors spoke ``more plainly
than before.'' ``Through the whole of Africa,'' he says in his
Chronicon, under the date 484, ``the cruel persecution of Hunneric,
King of the Vandals, was inflicted upon our Catholics. For after the
expulsion and dispersion of more than 334 Bishops of the orthodox, and
the shutting of their churches, the flocks of the faithful, afflicted
by various punishments, consummated their blessed conflict. Then it
was that the same King Hunneric ordered the tongue to be cut out of a
Catholic youth, who from his birth had lived without speech at all;
soon after he spoke, and gave glory to God with the first sounds of
his voice. In short, I myself have seen at Byzantium a few out of this
company of the faithful, religious men, with their tongues cut off,
and their hands amputated, speaking with perfect voice.''

228.
Victor, Bishop of Tonno in Africa Proconsularis, another contemporary,
and a strenuous defender of the Tria Capitula, which were condemned in
the Fifth Ecumenical Council, has left behind him a Chronicon also;
which at the same date runs as follows: ``Hunneric, King of the
Vandals, urging a furious persecution through the whole of Africa,
banishes to Tubunnæ, Macrinippi, and other parts of the desert, not
only Catholic Clerks of every order, but even Monks and laymen,
to the number of about four thousand; and makes Confessors and
Martyrs; and cuts off the tongues of the Confessors. As to which
Confessors, the royal city, where their bodies lie, attests that after
their tongues were cut out they spoke perfectly even to the end. Then
Lætus, Bishop of the Church of Nepte, is crowned with Martyrdom, etc.''
It is observable from this statement that the miracle was recorded for
the instruction of posterity at the place of their burial.

229.
Lastly, Pope Gregory the First thus speaks in his Dialogues: ``In the
time of Justinian Augustus \footnote{This date is a mistake of St. Gregory's; also he calls them Bishops.},
when the Arian persecution raised by the Vandals against the faith of
Catholics was raging violently in Africa, some Bishops, courageously
persisting in the defence of the truth, were brought under notice whom
the King of the Vandals, failing to persuade to his misbelief with
words and offers, thought he could break with torture. For when, in
the midst of their defence of the truth, he bade them be silent, but
they would not bear the misbelief quietly, lest it might be
interpreted as assent, breaking out into rage he had their tongues cut
off from the roots. A wonderful thing, and known to many senior
persons; for afterwards, even without tongue, they spoke for the
defence of the truth, just as they had been accustomed before to
speak by means of it ... These then, being fugitives at that time,
came to Constantinople. At the time, moreover, that I was myself sent
to the Emperor to conduct the business of the Church, I fell in with a
certain senior, a Bishop, who attested that he had seen their mouths
speaking, though without tongues, so that with open mouths they cried
out, 'Behold and see; for we have not tongues, and we speak.' And it
appeared to those who inspected, as it was said, as if, their tongues
being cut off from the roots, there was a sort of open depth in their
throat, and yet in that empty mouth the words were formed full and
perfect. Of whom one, having fallen into licentiousness, was soon
after deprived of the gift of miracle.'' \footnote{Dial. iii. 32.}

230.
Little observation is necessary on evidence such as this. What is
perhaps most striking in it is the variety of the witnesses, both in
their persons and the details of their testimony, together with the
consistency and unity of that testimony in all material points. Out of
the seven writers adduced, six are contemporaries; three, if not four,
are eye-witnesses of the miracle; one reports from an eye-witness, and
one testifies to a permanent record at the burial-place of the
subjects of it. All seven were living, or had been staying, at one or
other of the two places which are mentioned as their abode. One is a
Pope; a second a Catholic Bishop; a third a Bishop of a
schismatical party; a fourth an Emperor; a fifth a soldier, a
politician, and a suspected infidel; a sixth a statesman and courtier;
a seventh a rhetorician and philosopher. ``He cut out the tongues by
the roots,'' says Victor, Bishop of Vite; ``I perceived the tongue
entirely gone by the roots,'' says Æneas; ``as low down as the throat,''
says Procopius; ``at the roots,'' say Justinian and St. Gregory. ``He
spoke like an educated man, without impediment,'' says Victor of Vite;
with ``articulateness,'' says Æneas, ``better than before;'' ``they talked
without any impediment,'' says Procopius; ``speaking with perfect voice,''
says Marcellinus; ``they spoke perfectly even to the end,'' says the
second Victor; ``the words were formed full and perfect,'' says St.
Gregory.

231.
One of the striking points then in this miracle, as contained in the
foregoing evidence, is obviously its \emph{completeness}. We know that
even deaf and dumb persons can be made in some sense to utter words;
and there may be attempts far superior to theirs, yet wanting in that
ease and precision which characterize the ordinary gift of speech. But
the articulateness, nay, the educated accent of these Confessors is
especially insisted on in the testimony. ``A cure left thus imperfect,''
says Douglas, speaking of a Jansenist miracle, ``has but little
pretension to be looked upon as miraculous; because its being so
imperfect naturally points out a failure of power in the cause
which brought it about.'' \footnote{Page 111.}
Whatever be the truth of this position, it cannot be applied to the
miracle under review.

232.
The \emph{number} on which it was wrought is another most important
circumstance, distinguishing this history from others of a miraculous
character. It both increases opportunities for testimony, and it
prevents the interposition of what is commonly called chance, which
could not operate upon many persons at once in one and the same way.
This is the proper answer to Middleton's objection, that cases are on
record of speech without a tongue, when no special intervention of
Providence could be supposed. Not to say that a person \emph{born} without a
tongue, as in the instance to which he refers, may more easily be
supposed to have found a compensation for her defect by a natural
provision or guidance, than men who had ever spoken by the ordinary
organ till they came suddenly to lose it. ``If we should allow after
all,'' says he, ``that the tongues of these Confessors were cut away to
the very roots, what will the learned Doctor [Berriman] say if this
boasted miracle, which he so strenuously defends, should be found at
last to be no miracle at all? The tongue, indeed, has generally been
considered as absolutely necessary to the use of speech; so that to
hear men talk without it might easily pass for a miracle in that
credulous age.'' \footnote{Page 184.}
And then he mentions the case of a girl born without a tongue, who yet
talked as distinctly and easily as if she had enjoyed the full benefit
of that organ, according to the report of a French physician who had
carefully examined her mouth and throat, and who refers at the same
time to another instance published about eighty years before, of a boy
who at the age of eight or nine years lost his tongue by an ulcer
after the small-pox, yet retained his speech,—whether as perfectly
as before, does not appear \footnote{[Vid. Note at the end.]}.

233.
Now, taking these instances at their greatest force, does he mean to
say that if a certain number of men lost their tongues at the command
of a tyrant for the sake of their religion, and then spoke as plainly
as before, nay, if only one person was so mutilated and so gifted, it
would not be a miracle? if not, why does he not believe the history of
these Confessors? At least he might believe that some of them had the
gift of speech continued to them, though the numbers be an
exaggeration. It is his canon, as Douglas assures us, that while the
history of miracles is ``to be suspected always of course, without the
strongest evidence to confirm it,'' the history of common events is ``to
be \emph{admitted of course}, without as strong reason to suspect it.'' \footnote{Page 26. vid. supr. n. 73.} Now here all the reason or evidence is on the side of
believing; yet he does not believe it; why? simply because, as
common sense tells us, and as he feels, it \emph{is} a miraculous
story. It is far more difficult to believe that a number of men were
forbidden to profess orthodoxy, did continue to profess it, were
brought into the forum, had their tongues cut out from the roots,
survived it, and spoke ever afterwards as they did before, \emph{without}
a miracle than \emph{with} it. But Middleton would secure two weapons
at once for his warfare against the claims of the Catholic
Church:—it \emph{is} a miracle, and therefore it is incredible as a
fact; it is \emph{not} a miracle, and therefore it is irrelevant as an
argument.

234.
Another remarkable peculiarity of this miracle is what may be called
its \emph{entireness}, by which I mean that it carried its whole case
with it to every beholder. When a blind man has been restored to
sight, there must be one witness to prove he \emph{has been} blind,
and another, that he \emph{now sees}; when a cure has been effected,
we need a third to assure us that no medicines were administered to
the subject of it; but here the miracle is condensed in the fact that
there is no tongue, and yet a voice. The function of witnessing is far
narrower and more definite, yet more perfect, than in other cases.

235.
A further characteristic of this miracle is its \emph{permanence}; and
in this respect it throws light upon a remark made in a former page to
account for the deficiency of evidence which generally attaches to the
Ecclesiastical Miracles. It was there observed that they
commonly took place without notice beforehand, and left no trace after
them; and we could not have better or fuller testimony than that which
happened to be found on the spot where they occurred \footnote{n. 103.}. The instance before us, however, being of a permanent
character, and carrying its miraculousness in the very sight of it,
admitted of being witnessed in a higher way, and so it is witnessed.
Supposing the miracles of St. Gregory Thaumaturgus or St. Martin to
have had advantage of similar publicity, at least they would have been
disengaged from the misstatements and exaggerations which at present
prejudice them;—are we sure they would not have gained, instead, a
body of testimony to their substantial truth?

236.
It may be thought a drawback on this miracle that it produced no
impression on the brutal Prince who was the occasion of it. He
continued the persecution. Yet it must be recollected that his death
followed in no long time; and that, under that horrible and loathsome
infliction with which it has in other cases pleased Almighty God to
visit those who have used their power, committed to them by Him in
cruelties towards His Church.

237.
And now, after considering this miracle, or that of the recovery of
the blind Severus by the relics, as described in its place, or the
death of Arius, how unreal real does the remark appear with
which Douglas concludes his review of the alleged miracles of the
first ages! ``I shall only add,'' he says, ``that \emph{if ever there were
any accounts} of miracles, which passed current \emph{without being
examined into} at the first publication, and which consequently
will not bear the test of the third rule which I laid down in this
treatise, \emph{this may be affirmed of the miracles} recorded by
writers \emph{of the fourth and fifth ages}, when Christianity, \emph{now
freed from the terrors of persecution}, and \emph{aided by civil
magistrates}, began to be corrupted by its credulous or
ill-designing professors, and the foundation was laid of those
inventions which have gathered like a snow-ball in every succeeding
age of superstitious ignorance, till at last the sunshine of the
Reformation began to melt the monstrous heap.'' \footnote{Page 239.} Surely, if there are miracles prominent above others in those
times, in that number are the three which I have just specified; they
are great in themselves and in their fame. What then is meant by
saying that in Arius's death the Church was ``aided by the civil
magistrate?'' or that she was ``freed from the terrors of persecution''
when Severus was restored to sight? or that the report of the power of
speaking given to Reparatus and his brethren ``passed current without
being examined into?'' But if these are true, why should not others be
true also, whether at this day they have evidence sufficient for
our conviction or not? That superstitions and imposture \emph{accompanied}
the civil establishment of Christianity, all will allow; but they
could but obscure,—they could not reverse or undo,—and why should
they prejudice?—that true work of God in his Church, of which they
were but the mockery.

\xchapter{Conclusion}

238.
MUCH stress has been laid
throughout this Essay on the \emph{differences}
existing between the Miracles recorded in Scripture, and those
which are found in Ecclesiastical history; but from what has come
before us in the course of it, it would seem that those differences
are for the most part merely such as \emph{necessarily} attend
the introduction of a religion to the world compared with its
subsequent course, the miraculous agency itself being for the most
part the same throughout. For instance, the miracles of Scripture are
wrought by persons \emph{conscious of
their power and of their exercise of it};
for these persons are the very heralds of Almighty God, whom He
has commissioned, whom He has instructed, and whom He has gifted for
their work. The Scripture miracles are wrought \emph{as
evidence of revealed truth},
because they are wrought before that truth had as yet been
received. They are \emph{grave and simple} in
their circumstances, because they are wrought by persons who know
their gift, and who, as being under immediate Divine direction, use it
without alloy of human infirmity or personal peculiarity. They are \emph{definite and certain}, drawn
out in an orderly form, and finished in their parts, because they are
found in that authoritative Document which was intended by God’s
Providence to be the pattern of His dealings and the rule of our
thoughts and actions. They are \emph{undeniably
of a supernatural character},
not only because it is natural that the most cogent miracles
should be wrought in the beginning of the Dispensation, but because
the sacred writers have been guided to put into the foreground those
works of power which are the clearest tokens of a Divine Presence, and
to throw the rest into the distance. They have \emph{no
marks of exaggeration} about
them, and are \emph{none of them false
or suspicious}, because
Inspiration had dispersed the mists of popular error, and the
colouring of individual feeling, and has enabled the writers to set
down what took place, and nothing else. But when once Inspiration was
withdrawn, whether as regards those who wrought or those who recorded,
then a Power which henceforth was mysterious and inscrutable in
operation, became doubly obscure in report; and fiction in the
testimony was made to compensate for incompleteness in the
manifestation.

239.
In conclusion I will but observe what, indeed, is very obvious,
but still may require a distinct acknowledgment; that the view here
taken of the primitive miracles is applicable in defence of those of
the medieval period also. If the occurrence of miraculous
interpositions depends upon the presence of the Catholic Church, and
if that Church is to remain on earth until the end of the world, it
follows, of course, that what will be vouchsafed to Christians at all
times, was vouchsafed to them in the middle ages inclusively. Whether
this or that alleged miracle be in fact what it professes to be, must
be determined, as in the instances already taken, by the particular
case; but it stands to reason, that, where the views and
representations drawn out in the foregoing pages are admitted, no
prejudice will attend the medieval miracles at first hearing, though
no distinct opinion can be formed about them before examination.

240.
On the other hand, I am quite prepared to find those views themselves
condemned by many readers as subtle and sophistical. This is ever the
language men use concerning the arguments of others, when they dissent
from their \emph{first principles},—which take them by surprise, and
which they have not mastered.

\xsection{NOTE ON PAGE 383}

March, 1870. Since the date of these Essays
facts have been published, bearing upon the apparent miracle wrought
in favour of the African confessors in the Vandal persecution, which
have led me, in my ``Apologia'' (p. 306, ed. 2), to write as
follows:

``Their tongues were cut out by the Arian
tyrant, and yet they spoke as before. I insisted \footnote{That is, in my Second Essay, as supr. Ch. v. Sec.
9, especially at p. 383.}
on this fact as being strictly miraculous. Among other remarks
(referring to the instances adduced by Middleton and others in
disparagement of the miracle, viz., of 'A girl born without a tongue,
who yet talked as distinctly and easily as if she had enjoyed the full
benefit of that organ,' and of a boy who lost his tongue at the age of
eight or nine, yet retained his speech, whether perfectly or not,) I
said, 'Does Middleton mean to say that, if a certain number of men
lost their tongues \emph{at the command of a tyrant}, for the \emph{sake
of their religion}, and then spoke as plainly as before, nay, \emph{if
only one person was so mutilated} and so gifted, it would not be a
miracle?'—p. 210. And I enlarged upon the minute details of the fact
as reported to us by eyewitnesses and contemporaries.

``However, a few years ago an article appeared
in \emph{Notes and Queries} (No. for May 22, 1858), in which various
evidence was adduced to show that the tongue is not necessary for
articulate speech.

``1. Colonel Churchill, in his 'Lebanon,'
speaking of the cruelties of Djezzar Pacha, in extracting to the root
the tongues of some Emirs, adds, 'It is a curious fact, however,
that the tongues grow again sufficiently for the purposes of speech.'

``2. Sir John Malcolm, in his 'Sketches of
Persia,' speaking of Zâb, Khan of Khisht, who was condemned to lose
his tongue, 'This mandate,' he says, 'was imperfectly executed, and
the loss of half this member deprived him of speech. Being afterwards
persuaded that its being cut close to the root would enable him to
speak so as to be understood, he submitted to the operation and the
effect has been that his voice, though indistinct and thick, is yet
intelligible to persons accustomed to converse with him ... I am not
an anatomist, and I cannot therefore give a reason why a man, who
could not articulate with half a tongue, should speak when he had none
at all; but the facts are as stated.'

``3. And Sir John McNeil says, 'In answer to
your inquiries about the powers of speech retained by persons who have
had their tongues cut out, I can state from personal observation that
several persons whom I knew in Persia, who had been subjected to that
punishment, spoke so intelligibly as to be able to transact important
business ... The conviction in Persia is universal, that the power of
speech is destroyed by merely cutting off the tip of the tongue; and
is to a useful extent restored by cutting off another portion as far
back as a perpendicular section can be made of the portion that is
free from attachment at the lower surface … I never had to meet with
a person who had suffered this punishment, who could not speak so as
to be quite intelligible to his familiar associates.' So far these
writers.

``I should not, however, be honest, if I
professed to be simply converted by their testimony, to the belief
that there was nothing miraculous in the case of the African
confessors. It is quite as fair to be sceptical on one side of the
question as on the other; and if Gibbon is considered worthy of
praise for his 'stubborn incredulity' in receiving the evidence for
the miracle, I do not see why I am to be blamed, if I wish to be quite
sure of the full appositeness of the recent evidence which is brought
to its disadvantage. Questions of fact cannot be disproved by
analogies or presumptions; the inquiry must be made into the
particular case in all its parts, as it comes before us. Meanwhile, I
fully allow that the points of evidence brought in disparagement of
the miracle are \emph{primâ facie} of such cogency, that, till they
are proved to be irrelevant, Catholics are prevented from appealing to
it for controversial purposes.''
